% Paper 34: GeoVac as a Two-Layer Framework
% Projection Taxonomy and the Empirical Matches Catalog
% GeoVac Project, May 2026
% v0.1 (skeleton) -- Standalone observation paper.  Names the
% framework's two-layer architecture (bare graph + projection
% family), gives a per-projection dictionary tagging each by
% physical variables introduced, dimensions introduced, and
% transcendental signature, and consolidates the running list of
% machine-precision matches as a tagged catalog.

\documentclass[11pt]{article}
\usepackage{amsmath,amssymb,amsthm}
\usepackage{booktabs}
\usepackage{array}
\usepackage{longtable}
\usepackage[margin=1in]{geometry}
\usepackage{hyperref}

\newtheorem{observation}{Observation}
\newtheorem{definition}{Definition}
\newtheorem{remark}{Remark}
\newtheorem{prediction}{Prediction}
\newtheorem{theorem}{Theorem}

\title{GeoVac as a Two-Layer Framework:\\
Projection Taxonomy and the Empirical Matches Catalog}
\author{GeoVac Project}
\date{May 2026}

\begin{document}
\maketitle

% ===================================================================
\begin{abstract}
% ===================================================================
The GeoVac framework decomposes into two structurally separate
layers.  \textbf{Layer~1} is a bare combinatorial graph: quantum-number
labels $(n,\ell,m,\ldots)$, integer eigenvalues, rational matrix
elements, no physics.  \textbf{Layer~2} is a finite family of named
projections from the bare graph onto physical settings.  Every
physical variable, every dimensional unit, and every transcendental
constant that appears in any GeoVac calculation enters through
exactly one projection, and is fully traceable to it.

We catalogue twenty-eight projection mechanisms and tag each one along
three independent axes: (a)~which physical variables it introduces
(charge $Z$, coupling $\alpha$, energy scale $\hbar\omega$, internuclear
distance $R$, gauge coupling $\beta$, cutoff $\Lambda$, $\ldots$);
(b)~which physical dimension or unit it attaches to graph data
(energy, length, dimensionless, solid angle); and (c)~its
transcendental signature in the sense of Paper~18 (rational, $\pi^{2k}$,
half-integer Hurwitz, odd~$\zeta$, Dirichlet~$L$, $\ldots$).  The
projection axis is a structural sibling of Paper~18's transcendental
axis: each transcendental tier corresponds to a specific projection
type, and reading the dictionary in either direction recovers the
other.

We then assemble the empirical matches catalogue: a tagged table of
GeoVac quantities that match standard known physics values to
machine precision (or stated uncertainty).  Each row is annotated
by projection(s), variable(s), dimension, transcendental class, and
match precision.  Reading the catalogue surfaces a falsifiable
structural prediction:
\begin{prediction}[Layer-2-presence bound.]
The residual error of a catalogue row is bounded above by the
larger of (i)~the framework's basis-quality residual at zero
Layer-2 inputs, and (ii)~the largest single Layer-2 input
consumed by the chain.  Equivalently:
$|\epsilon(\text{observable})|
\leq \max\bigl(\epsilon_\text{basis}(\text{framework}),
\max_i |L_2^{(i)}|\bigr)$.
The genuine axis ordering catalogue precision is observable-class
(zero-, one-, multi-Layer-2-input) crossed with the framework's
basis-quality limit; chained projection depth is a downstream
proxy that fails in both directions.
\end{prediction}
This replaces an earlier naive form that scaled residual linearly
with chain depth; the May~2026 synthesis pass falsified the depth
form on the catalogue itself (depth-3 chains exist at residual~$0$
and at $+286$~ppm; depth-4 chains exist at $+2$~ppm and at
$+0.534\%$).  The Lamb shift example (depth-4 chain, $-0.534\%$
residual at converged basis) is consistent with the framework's
one-loop QED basis-quality limit; the H~1S polarizability example
(depth-3 chain, residual $0$) is consistent with the
zero-Layer-2-input class for rational-analytic-answer
observables.

This paper introduces no new computation.  It is a structural
consolidation of the framework's self-presentation: language for
discussing where physics enters, what carries dimensions, and which
transcendentals are forced by which projection.
\end{abstract}

% ===================================================================
\section{Introduction: The Two-Layer Thesis}
\label{sec:intro}
% ===================================================================

The recurring difficulty in discussing GeoVac is that the framework
sits between two languages.  In one language, GeoVac is a sparse
graph Hamiltonian framework for quantum chemistry: nodes labeled by
quantum numbers, edges from selection rules, eigenvalues that
predict atomic and molecular spectra (Paper~14, Paper~17, Paper~20).
In the other language, GeoVac is a discrete spectral triple in the
Marcolli--van~Suijlekom lineage of graph-spectral-triple
constructions (Paper~32), with a heat-kernel expansion that
reproduces the QED $\beta$-function at one loop (Paper~28) and
recovers seven of eight QED selection rules with a single Hopf-base
calibration constant $1/(4\pi)$ per loop (Paper~33).

These two languages are not in conflict, and we have not been able
to choose between them, because they are descriptions of two
different layers of the same construction.  This paper names the
layers and the bridge between them.

\begin{definition}[Layer~1: the bare graph]
\label{def:layer1}
The Layer~1 object is the dimensionless combinatorial graph.  Nodes
are quantum-number labels $(n, \ell, m, \kappa, m_j, \ldots)$.
Eigenvalues of the bare graph Laplacian are integers ($n^2 - 1$ on
the Coulomb $S^3$ graph, $n + 3/2$ on Dirac $S^3$, $N$ on the
Bargmann--Segal $S^5$ graph).  Matrix elements of bare graph
operators (Casimir, ladder, adjacency) are rational, integer, or
algebraic over $\mathbb{Q}[\sqrt{2}, \sqrt{3}, \sqrt{6},\ldots]$
(from Wigner~$3j$ algebra); $\pi$ does not appear in any Layer~1
quantity (Paper~24's $\pi$-free certificate; the same applies on
the Coulomb side).  Layer~1 carries no physical units, no
coupling constants, no observable energies.
\end{definition}

\begin{definition}[Layer~2: the projection family]
\label{def:layer2}
The Layer~2 object is a finite family of named maps from the bare
graph to physical settings.  Each projection is a structural
construction with a name in the literature (Fock 1935 conformal
projection; Hopf bundle; Bargmann--Segal transform; Connes--Chamseddine
spectral action; Sturmian basis at exponent $\lambda = Z/n$;
Camporesi--Higuchi spinor lift; Wigner~$3j$ angular coupling;
Wigner~$D$-matrix rotation between molecular centers; Wilson
plaquette; molecule-frame hyperspherical separation; vector-photon
promotion; Mellin/heat-kernel regularization).  Each projection
adds (i)~specific physical variables, (ii)~specific physical
dimensions, and (iii)~specific transcendental content to the bare
graph data on which it acts.
\end{definition}

\begin{observation}[Single-projection origin of physical content]
\label{obs:single_origin}
Every coupling constant, energy scale, length scale, and
transcendental constant that appears in any GeoVac computation can
be traced to exactly one projection in Layer~2.  The bare graph
does not contain $\alpha$, $Z$, $\hbar\omega$, $R$, $\Lambda$, $\pi$,
or $\zeta(2k)$.  These quantities are introduced by, and only by,
the projection that the calculation invokes.
\end{observation}

The remainder of this paper is the dictionary.

% ===================================================================
\section{Layer 1: The Bare Graph}
\label{sec:layer1}
% ===================================================================

% TODO: This section should cite Papers 0, 7, 24, 22 in detail.
% Brief here; expand from existing material.

\subsection{Quantum-number labelling}
\label{sec:layer1_labels}
Nodes of the bare GeoVac graph are tuples of quantum numbers.  The
specific tuple depends on which natural geometry is in play:
\begin{itemize}
\setlength\itemsep{0pt}
\item Atomic Coulomb $S^3$: $(n, \ell, m_\ell)$ for the scalar sector,
$(n, \kappa, m_j)$ for the Dirac sector.
\item Bargmann--Segal $S^5$ (3D HO): $(N, \ell, m_\ell)$ in $\mathrm{SU}(3)$
$(N, 0)$ irreps.
\item Composed molecular: per-block tuples of the above, with
inter-block edges labeled by the multipole order $L$ of the
cross-center expansion.
\end{itemize}
No physical quantities are required to define these labels.  They
are pure representation-theoretic data of compact Lie groups.

\subsection{$\pi$-free certificate}
\label{sec:layer1_pi_free}
On both the Coulomb $S^3$ Fock graph and the Bargmann--Segal $S^5$
graph, every adjacency matrix entry, every diagonal eigenvalue, and
every CG-coupling matrix element is exact rational or algebraic in
$\mathbb{Q}[\sqrt{2}, \sqrt{3}, \sqrt{5}, \sqrt{6}, \sqrt{10},
\sqrt{15}, \ldots]$ (Paper~24~\cite{paper24}, Paper~28's
graph-native QED certificate).  No $\pi$ appears in Layer~1 by
exact computation.

\subsection{Layer-1 observables}
\label{sec:layer1_obs}
Some observables are computable entirely within Layer~1.  These are
the ``zero-projection'' matches in the catalogue
(Section~\ref{sec:catalog}).  Examples:
\begin{itemize}
\setlength\itemsep{0pt}
\item Pauli term count $N_\text{Pauli} = 11.10 \cdot Q$ (across 35
composed molecules; Paper~14).
\item Magic numbers $2, 8, 20, 28, 50, 82, 126$ (HO + spin-orbit
graph eigenvalue counting; Paper~23).
\item Angular sparsity density $1.44\%$ at $\ell_\text{max}=3$
(Paper~22, potential-independent).
\item All Gaunt selection rules (CG sparsity).
\end{itemize}
These exist on Layer~1 alone.  They have no $\pi$, no $\alpha$, no
energy scale, no transcendentals of any kind, and they match
experiment exactly because there is nothing to project---they are
consequences of compact-Lie-group representation theory acting on
discrete quantum-number labels.

% ===================================================================
\section{Layer 2: The Projection Family}
\label{sec:layer2}
% ===================================================================

We enumerate twenty-eight projection mechanisms.  For each: source and
target spaces, physical variables introduced, dimensions
introduced, transcendental signature, and existing GeoVac papers
that exercise the projection.  Projections~14 (rest-mass) and~15
(observation / temporal-window) were added in sprint
KG\footnote{Sprint/track designations are anchors into the project's
internal chronicle (\texttt{CHANGELOG.md} in the repository).} /
Paper~35 (May~2026); projection~16 (two-body Dirac / Breit retardation) was
added in the precision-catalogue sprint of 2026-05-08, motivated by
the positronium 1S--2S framework-native residual that surfaced as
the sixth (and structurally cleanest) instance of the
multi-focal-composition wall named in CLAUDE.md~\S~2.
Projections~17--19 (nuclear charge-density / Foldy--Friar,
magnetization-density / Zemach, and tensor multipole) were added in
the intra-particle hunting sprint of 2026-05-09, after the audit
memo \texttt{debug/paper34\_v\_base\_unit\_audit\_memo.md} flagged
magnetization-density as a single 17th-projection candidate and the
follow-up hunt established that the three mechanisms are
categorically distinct (different Hamiltonian operators, different
selection rules; documented in
\texttt{debug/intra\_particle\_L\_hunting\_memo.md}).  A boundary
subsection (\S~\ref{sec:layer2_boundary}) names the
internal-graph-derived vs.\ Layer-2-input distinction, with the
external-field hunting memo
(\texttt{debug/external\_field\_hunting\_memo.md}, same date)
verifying that uniform external classical fields (Zeeman, Stark) are
Layer-2 inputs rather than \S~\ref{sec:layer2} projections.

\subsection{Fock conformal projection: $\mathbb{R}^3 \to S^3$}
\label{sec:proj_fock}
\textbf{Source:} continuum Coulomb Hamiltonian $H = -\tfrac{1}{2}\nabla^2 - Z/r$
on $\mathbb{R}^3$.\\
\textbf{Target:} unit $S^3$ via stereographic-like map at focal
length $p_0^2 = -2E$.\\
\textbf{Variables introduced:} $Z$ (nuclear charge), $E$ (energy
eigenvalue).\\
\textbf{Dimension introduced:} energy (in atomic units, Hartree).
The dimensional content is carried by the choice
$m_e = \hbar = e = 1$.\\
\textbf{Transcendental signature:} rational prefactor
$\kappa = -1/16$ (the Rydberg-to-graph-eigenvalue conversion).
$\pi$ enters through $\mathrm{Vol}(S^3) = 2\pi^2$ when subsequent
spectral integrals are performed.\\
\textbf{Used in:} Paper~7 (S$^3$ proof), Papers 13/15/17
(hyperspherical hierarchy), Paper~18 \S~III (Peter--Weyl
realization), all atomic energy computations.

\subsection{Hopf bundle: $S^3 \to S^2 \times S^1$ (locally)}
\label{sec:proj_hopf}
\textbf{Source:} $S^3 \cong \mathrm{SU}(2)$.\\
\textbf{Target:} $S^2$ base $\times S^1$ fibre.\\
\textbf{Variables introduced:} $\alpha$ (the fine-structure constant
treated as a candidate observable in Paper~2's Observation).\\
\textbf{Dimension introduced:} dimensionless.\\
\textbf{Transcendental signature:} $\pi = \mathrm{Vol}(S^2)/4$
(Hopf measure factor; Sprint~A $\alpha$-PI identification).\\
\textbf{Used in:} Paper~2 (fine-structure Observation), Paper~25
(Wilson--Hodge gauge structure), Paper~30 (SU(2) Wilson sibling).

\subsection{Bargmann--Segal transform: $\mathbb{R}^3$ HO $\to S^5$ Hardy space}
\label{sec:proj_bargmann}
\textbf{Source:} 3D isotropic harmonic oscillator
$H_{\text{HO}} = \tfrac{1}{2}p^2 + \tfrac{1}{2}\omega^2 r^2$.\\
\textbf{Target:} holomorphic Hardy-space sector of $S^5$, restricted
to $\mathrm{SU}(3)$ $(N,0)$ symmetric irreps.\\
\textbf{Variables introduced:} $\hbar\omega$ (HO frequency).\\
\textbf{Dimension introduced:} energy.\\
\textbf{Transcendental signature:} \emph{none}.  The Bargmann--Segal
graph is bit-exactly $\pi$-free in rational arithmetic at every
finite $N_\text{max}$ (Paper~24).  $\pi$ appears only as
$\mathrm{Vol}(S^5) = \pi^3$ in continuum integration measures, never
in lattice data.\\
\textbf{Used in:} Paper~24 (HO discretization), Paper~31 (universal
vs.\ Coulomb partition).

\subsection{Stereographic / conformal coordinate change}
\label{sec:proj_stereo}
\textbf{Source:} compact sphere ($S^2$, $S^3$, $S^5$).\\
\textbf{Target:} flat-space coordinates ($\mathbb{R}^2, \mathbb{R}^3,
\mathbb{R}^5$).\\
\textbf{Variables introduced:} length scale (radial coordinate $r$).\\
\textbf{Dimension introduced:} length.\\
\textbf{Transcendental signature:} conformal factor producing the
$1/r$ Coulomb potential as coordinate distortion (Paper~7).\\
\textbf{Used in:} Paper~7 (Coulomb-as-coordinate), all calculations
that quote results in physical bohr/Hartree.

\subsection{Sturmian reparameterization at $\lambda = Z/n$}
\label{sec:proj_sturmian}
\textbf{Source:} bare Coulomb $S^3$ graph.\\
\textbf{Target:} same graph relabelled as a Coulomb Sturmian basis
spanning bound states plus a discretized continuum.\\
\textbf{Variables introduced:} $Z$, $n$ (re-used; no new variables).\\
\textbf{Dimension introduced:} preserves existing.\\
\textbf{Transcendental signature:} preserves rationality of
Layer~1.  No new $\pi$ or other transcendentals enter at this step;
calibration-tier transcendentals (e.g.\ Bethe logarithm) enter
through subsequent spectral mode sums.\\
\textbf{Used in:} Papers 8--9, Paper~19, sprint LS-2 (Bethe logarithm
from Sturmian sum at 2--3\%).

\subsection{Connes--Chamseddine spectral action / heat kernel}
\label{sec:proj_spectral_action}
\textbf{Source:} Dirac operator $D$ on bare $S^3$ graph (or
extension).\\
\textbf{Target:} heat-kernel expansion
$\mathrm{Tr}\, f(D^2/\Lambda^2)$ with cutoff function $f$ and scale
$\Lambda^2$.\\
\textbf{Variables introduced:} $\Lambda$ (UV cutoff), $\alpha$ (when
gauge sector included).\\
\textbf{Dimension introduced:} energy via $\Lambda$.\\
\textbf{Transcendental signature:} $\sqrt{\pi} \times \mathbb{Q}$
Seeley--DeWitt coefficients on unit $S^3$ ($a_0 = a_1 = \sqrt{\pi}$,
$a_2 = \sqrt{\pi}/8$); resulting observable coefficients are
$\pi^{2k} \times \mathbb{Q}$ at one loop (T9 theorem, Paper~28).\\
\textbf{Used in:} Paper~28 (QED on $S^3$), Paper~32 (spectral triple
synthesis), sprint LS-1 (Uehling shift $-27.13$ MHz from
$\Pi = 1/(48\pi^2)$).\\
\textbf{Structural reading (master Mellin engine framing).}  This
projection is the $k=2$ sub-case of the master Mellin engine of
Paper~32~\S~VIII and Paper~18~\S~III.7:\ the proper-time integration
$\int_0^\infty dt\, f(t)\,\mathrm{Tr}\,e^{-tD^2}$ that converts $D$-spectral
data to physical observables.  The same Mellin technology applied
at $k=0$ to $\mathrm{Tr}\,e^{-tD^2}$ recovers the M1 Hopf-base measure
content (\S~\ref{sec:proj_hopf} and \S~\ref{sec:proj_vector_photon}),
and applied at $k=1$ recovers the M3 vertex-parity Hurwitz content
(composed with \S~\ref{sec:proj_spinor} and vertex restriction;
Paper~28~\S~IV).  The Mellin transform itself is the
\emph{evaluator} (a meta-mechanism organizing the spectral-side
projection family), not a $\S$~III peer projection; consequently the
Schwinger proper-time formalism and the Connes--Chamseddine spectral
action are the same operation under two names, both absorbed into the
present entry (Sprint~3 Track~12 scoping verdict, 2026-05-15;
\texttt{debug/sprint3\_drafts/track\_12\_heat\_kernel\_regularization\_scoping\_memo.md}).
The $\alpha$-power index $n$ in iterated applications of this projection
counts loop order; per-iteration transcendental content is the
$\sqrt{\pi}\cdot\mathbb{Q}$ ring at the heat-kernel level (T9 theorem,
Paper~28) and the $1/(4\pi)$ vector-photon Hopf-base calibration per
loop (\S~\ref{sec:proj_vector_photon} in Coulomb gauge, \S~\ref{sec:proj_gauge_choice}).
The finite parts of $\alpha^n$ observables are framework-native at every
$n$; UV-divergent multi-loop integrands are reproduced structurally
(LS-8a verdict, Paper~36~\S~VII) but their counterterm-renormalized
finite extraction requires external $Z_2/\delta m$ input not
autonomously generated by the bare spectral action (Paper~35~\S~VII.3
Refined Prediction~1; Sprint~3 Track~11 absorption verdict).  The open
question about fractional-$s$ behaviour of the Mellin engine
(\S~\ref{sec:open} (a)) is structurally separate from the present
entry and remains as catalogued.

\subsection{Camporesi--Higuchi spinor lift}
\label{sec:proj_spinor}
\textbf{Source:} scalar bare $S^3$ graph.\\
\textbf{Target:} Dirac-on-$S^3$ graph with eigenvalues
$|\lambda_n| = n + 3/2$, degeneracies $g_n = 2(n{+}1)(n{+}2)$.\\
\textbf{Variables introduced:} $\alpha$ enters through $(Z\alpha)^2$
relativistic content.\\
\textbf{Dimension introduced:} dimensionless (the spinor lift itself
is unit-preserving).\\
\textbf{Transcendental signature:} half-integer Hurwitz zetas
$\zeta(s, n + 3/2)$, producing odd-$\zeta$ content $\zeta(3),
\zeta(5), \ldots$ at integer $s$ (Paper~28 \S~IV); Catalan~$G =
\beta(2)$ and Dirichlet~$\beta(4)$ at the vertex level via
quarter-integer shifts $\zeta(s, 3/4), \zeta(s, 5/4)$.  Lives in
the spinor-intrinsic algebraic ring $R_\text{sp} =
\mathbb{Q}(\alpha^2)[\gamma]/(\gamma^2 + (Z\alpha)^2 - 1)$
(Paper~18 \S~IV).\\
\textbf{Used in:} Paper~28 \S~IV, Paper~33, Tier-2 spin-ful composed
qubit encoding (Paper~14 \S~V).

\subsection{Wigner $3j$ angular coupling}
\label{sec:proj_3j}
\textbf{Source:} pair of angular-momentum bare-graph nodes.\\
\textbf{Target:} coupled angular-momentum eigenstate.\\
\textbf{Variables introduced:} none.\\
\textbf{Dimension introduced:} dimensionless.\\
\textbf{Transcendental signature:} pure rational
($\mathbb{Q}[\sqrt{2k+1}]_{k}$).\\
\textbf{Used in:} all Gaunt integral evaluations, Paper~22 angular
sparsity, every Slater $F^k$ computation.

\subsection{Wigner $D$-matrix rotation between molecular centers}
\label{sec:proj_wignerD}
\textbf{Source:} composed orbital block at center $A$.\\
\textbf{Target:} same block in molecular coordinate frame at center
$B$.\\
\textbf{Variables introduced:} nuclear position vector
$\vec{R}_{AB}$.\\
\textbf{Dimension introduced:} length.\\
\textbf{Transcendental signature:} preserves rationality with
$\mathbb{Q}[\sqrt{2}, \sqrt{3}, \sqrt{6}]$ algebraic content from
the Wigner $d$-matrix at non-collinear angles.\\
\textbf{Used in:} multi-center molecules in Paper~17 and Paper~14
(LiF, CO, N$_2$, NaCl, CH$_2$O, C$_2$H$_2$, C$_2$H$_6$, F$_2$).

\subsection{Wilson plaquette: edge $\to$ 2-cell}
\label{sec:proj_wilson}
\textbf{Source:} Hopf graph 1-cochains.\\
\textbf{Target:} 2-cells (plaquettes) carrying $\mathrm{SU}(2)$ link
variables.\\
\textbf{Variables introduced:} $\beta = 1/g^2$ (gauge coupling).\\
\textbf{Dimension introduced:} dimensionless.\\
\textbf{Transcendental signature:} non-abelian $\mathrm{SU}(2)$ Haar
content; in the maximal-torus reduction yields the abelian
$U(1)$ content of Paper~25.\\
\textbf{Used in:} Paper~30 (SU(2) Wilson lattice gauge), Paper~25
(Hodge-1 Laplacian as gauge propagator), transverse photon
construction (Paper~28 \S~transverse\_photon).

\subsection{Vector-photon promotion}
\label{sec:proj_vector_photon}
\textbf{Source:} scalar 1-cochain photon on bare graph.\\
\textbf{Target:} vector-harmonic photon with explicit $(q, m_q)$
quantum numbers and Wigner~$3j$ vertex coupling.\\
\textbf{Variables introduced:} none (combinatorial labels).\\
\textbf{Dimension introduced:} dimensionless; specifically a
solid-angle measure on the Hopf $S^2$ base.\\
\textbf{Transcendental signature:} $1/(4\pi)$ per loop, identified
as the $S^2$ Weyl exchange constant of the Hopf base (Paper~33).\\
\textbf{Used in:} Paper~33 (recovers six angular-momentum selection
rules), sprint VP-1 (vertex correction calibration).

\subsection{Molecule-frame hyperspherical separation}
\label{sec:proj_molframe}
\textbf{Source:} composed atomic blocks.\\
\textbf{Target:} molecule-frame hyperspherical coordinates
$(\rho, \alpha, \theta_{12}, \ldots)$ with internuclear distance $R$
appearing in the charge function.\\
\textbf{Variables introduced:} $R$ (internuclear distance); for
polyatomics, additional bond angles.\\
\textbf{Dimension introduced:} length.\\
\textbf{Transcendental signature:} piecewise-smooth in $R$ at the
adiabatic potential level (Track~S, Paper~15); Gaunt integrals
preserve rationality at the angular level.\\
\textbf{Used in:} Paper~15 (Level~4), Paper~17 (composed geometry),
all H$_2$/HeH$^+$/H$_2$O calculations.

\subsection{Drake--Swainson asymptotic subtraction}
\label{sec:proj_drake_swainson}
\textbf{Source:} divergent finite-basis spectral sum
$J_v(nl) / I_v(nl)$ for $\ell > 0$ Bethe-logarithm states, where the
velocity-form closure $I_v(nl) = (Z^4/n^3)\,\delta_{\ell,0}$ vanishes
exactly for $\ell > 0$ on the bare Sturmian-projected graph.\\
\textbf{Target:} K-independent regularized value
$\ln k_0(n,\ell) = J_v(n,\ell) / D_{\text{drake}}(n,\ell)$, with the
structural denominator
\[
  D_{\text{drake}}(n,\ell) = \frac{2(2\ell + 1)\,Z^4}{n^3}.
\]
This recovers $D_{\text{drake}}(nS) = I_v(nS) = 2Z^4/n^3$ for $\ell = 0$
(matching the closure value), and provides a non-vanishing structural
normalization for $\ell > 0$.\\
\textbf{Variables introduced:} subtraction scale $K$ (transient; the
final result is $K$-independent over a plateau spanning $\sim 3.6$
orders of magnitude in $K$, as verified at sprint LS-4).\\
\textbf{Dimension introduced:} energy via $K$ (transient,
cancelled in $\beta_{\text{low}}(K) + \beta_{\text{high}}(K) = J_v$);
the resulting $\ln k_0$ is dimensionless.\\
\textbf{Transcendental signature:} flow tier (Paper~18 \S~IV).
$D_{\text{drake}}$ itself is rational ($Z^4 \cdot \mathbb{Q}$); the
transcendental content sits in $J_v$, where the spectral-sum
logarithms $\ln |2 \Delta E_m / \mathrm{Ry}|$ are inherited from the
Sturmian projection and carry calibration $\alpha$-dependent pieces
through the matching to the standard Bethe formula.  The structural
denominator
$D_{\text{drake}}(n,\ell) = (\text{spin})\cdot(\text{angular degeneracy})\cdot Z^4/n^3$
is a purely combinatorial quantity (spin~$\cdot$~$2\ell+1$~$\cdot$~hydrogenic
density), making the projection's structural content
combinatorial-rational while its operational content is calibration.\\
\textbf{Used in:} sprint LS-4 (Bethe logs for 2P at $+2.4\%$, 3P at
$+6.3\%$, 3D at $-0.24\%$ versus Drake-Swainson 1990 tabulated
values).  The combined hydrogen Lamb shift via this projection
chained with Sturmian and spectral-action gives $-0.39\%$ residual at
the sweet-spot $N=40$ (transient cancellation; converges to LS-1
baseline $-3.10\%$ at $N \to \infty$, reflecting the residual
$\alpha^5$ multi-loop ceiling).  The projection is structurally a
sibling of the Schwartz~(1961) integral form, simplified to act on
the discrete Sturmian pseudostate basis directly.

\subsection{Rest-mass projection}
\label{sec:proj_restmass}
\textbf{Source:} bare relativistic spectrum on $S^3 \times \mathbb{R}$
(or any natural-geometry spectrum), in the massless limit.\\
\textbf{Target:} massive spectrum
$\omega_n^2 = \omega_n^{(0)\,2} + m^2$, where $\omega_n^{(0)}$ is the
massless eigenvalue from the ambient natural geometry.\\
\textbf{Variables introduced:} $m$ (rest mass; equivalently rest-energy
$mc^2$ at $c=1$).\\
\textbf{Dimension introduced:} mass $=$ energy at $c=1$.\\
\textbf{Transcendental signature:} \textit{trivial}.  The bare-graph
algebraic-extension ring is preserved when $m^2 \in \mathbb{Q}$ (or
any ring extension of $\mathbb{Q}$ already present in the ambient
spectrum).  Verified explicitly in sprint KG-1 for the Klein-Gordon
spectrum on $S^3 \times \mathbb{R}$: $\omega_n = c_n \sqrt{d_n}$ with
$c_n \in \mathbb{Q}$ and $d_n$ a positive square-free integer for all
$n \in [1, 50]$ across the rational $m^2$ panel
$\{0, 1, 1/4, 2\}$ (Paper~35).  The mass projection is parameter-only;
it does not introduce $\pi$ or any other transcendental.\\
\textbf{Used in:} muonic / electronic / positronium reduced-mass
corrections; Paper~23 nuclear-electronic composed Hamiltonian
($I \cdot S$ hyperfine coupling); any future Klein-Gordon or
Dirac-on-$S^3$ relativistic mass calculation; Paper~35
Proposition~1.A.

\subsection{Observation / temporal-window projection}
\label{sec:proj_observation}
\textbf{Source:} bare relativistic spectrum on a spatial natural
geometry crossed with open temporal direction $\mathbb{R}$.\\
\textbf{Target:} same spectrum on the spatial geometry crossed with
periodic Euclidean time $S^1_\beta$ of period $\beta$ (equivalently:
finite-time observation window of duration $\beta$ in real time;
finite-temperature partition function at temperature $T = 1/\beta$).
For bosons, temporal Matsubara modes $\omega^t_k = 2\pi k / \beta$,
$k \in \mathbb{Z}$; for fermions with antiperiodic time,
$\omega^t_k = (2k+1)\pi / \beta$.\\
\textbf{Variables introduced:} $\beta$ (or $T = 1/\beta$).\\
\textbf{Dimension introduced:} time $=$ inverse energy.\\
\textbf{Transcendental signature:} $2\pi \cdot \mathbb{Q}$ per
Matsubara mode in the spectrum directly; in spectral integrals over
the compactified direction, $\pi^{2k} \cdot \mathbb{Q}$ from
$\zeta_R(2k)$ values, including the Stefan-Boltzmann
constant $\pi^2/90$ in the high-$T$ limit of $S^3 \times S^1_\beta$.
The first $\pi$-bearing eigenvalue of the compactified Klein-Gordon
spectrum is the $(n{=}0, k{=}1)$ Matsubara mode at
$\omega^2 = 4\pi^2/\beta^2$, with no earlier step injecting $\pi$
(verified in sprint KG-2; Paper~35).\\
\textbf{Used in:} all finite-temperature observables (Casimir at
$T \neq 0$, free-energy density, Stefan-Boltzmann constants); all
QED loop integrals that involve continuous time integration (Schwinger
proper-time, Mellin-Barnes, etc.); spectral-action evaluation
$\mathrm{Tr}\, f(D^2/\Lambda^2)$ when the trace runs over a
compact-time direction; Paper~35 Proposition~1.B.  This projection is
structurally what an observer does when integrating an instantaneous
spectrum over a finite measurement window.\footnote{%
Bisognano--Wichmann reading.\ Under the standard Wick-rotation chain
(Hartle--Hawking 1976~\cite{hartle_hawking1976};
 Sewell 1982~\cite{sewell1982};
 Bisognano--Wichmann 1976~\cite{bisognano_wichmann1976}),
the temporal-window projection's $\beta$ IS the period of the
modular-flow / Lorentz-boost orbit on a wedge or static-exterior
algebra:\ Bisognano--Wichmann establishes that the modular automorphism
of the algebra of observables localized in a Rindler wedge, with
respect to the vacuum state, IS the unitary representation of the
Lorentz boost subgroup that maps the wedge to itself; Sewell 1982
generalizes this to globally hyperbolic Lorentzian manifolds with
bifurcate Killing horizons, including Schwarzschild at
$\beta_{\mathrm{cigar}} = 8\pi M$.  Sprint TD Track~4
(\texttt{debug/sprint\_td\_track4\_memo.md}, May 2026) reproduces this
$\beta = 8\pi M$ from the framework's M1 Hopf-base measure /
$\mathrm{Vol}(S^1_\tau)$ mechanism alone.  Track~D
(\texttt{debug/bisognano\_wichmann\_track\_d\_memo.md}, May 2026)
documents the structural correspondence between the framework-side
$2\pi$ and the Lorentzian-side modular-flow / boost-orbit period.
The algebraic structure underwriting this reading is the
Tomita--Takesaki modular flow on the wedge / static-exterior algebra
(modular operator $\Delta = e^{-K}$, modular conjugation $J$,
$\sigma_t(A) = \Delta^{it}A\Delta^{-it}$ closing with $\sigma_{i\cdot 2\pi}
= $ identity at the BW wedge); Sprint~3 Track~10 (2026-05-15;
\texttt{debug/sprint3\_drafts/track\_10\_tomita\_takesaki\_scoping\_memo.md})
absorbed Tomita--Takesaki into the present footnote and \S~\ref{sec:proj_wick_rotation},
with the operator-system-level framework-internal realization (modular
Hamiltonian $K$ on $\mathcal{T}_{n_{\max}}$) named as the open extension
in \S~\ref{sec:open}.  The
Hawking temperature $T_H = 1/(8\pi M)$ enters the catalogue as an
off-precision row in \S~\ref{sec:offprecision} (error class
C, $M$ external Einstein-equations input, machinery-witness only); the
Unruh pendant (\texttt{debug/unruh\_pendant\_memo.md}, May 2026,
a sprint-scale follow-up to Track~D) closes the Track~D~\S 5.4 flag with
$T_U = a/(2\pi)$ as the third Wick-rotation witness via the same M1
Hopf-base measure / $\mathrm{Vol}(S^1)$ mechanism (Hawking--Sewell--BW--Unruh
four-witness chain).  See Paper~32~§VIII Remark on the
Bisognano--Wichmann reading and Paper~35~§VIII for the framework-level
structural correspondence.  This is a structural correspondence, not a
literal identification: the operator-level statement
($\sigma_{i\cdot 2\pi} = $ identity on the truncated metric spectral
triple) is the principal falsifier flagged in Track~D \S 3.1, a
follow-up beyond sprint scale.
}

\subsection{Two-body Dirac / Breit retardation}
\label{sec:proj_breit_retardation}
\textbf{Source:} framework-native single-particle Eides~\S~3.2
self-energy / Lamb-shift bracket on the Sturmian-reparameterized
Fock graph, valid in the $\lambda_l \ll \lambda_n$ hierarchy where the
$1/M_\text{nucleus}$ recoil expansion is regular.\\
\textbf{Target:} full two-body Breit / Bethe--Salpeter Hamiltonian for
the bound state at order $\alpha^4$, retaining all retardation and
crossed-photon terms that the $1/M$ expansion drops as subleading.\\
\textbf{Variables introduced:} the rest-mass ratio $m_l/m_n$ (and not
the absolute mass~$m$ already introduced by the rest-mass projection
of \S~\ref{sec:proj_restmass}).  This ratio does \emph{not} vanish in
the equal-mass leptonic limit $m_l = m_n$ where the $1/M$ expansion
is invalid.  Distinct from \S~\ref{sec:proj_restmass}: rest-mass
projection rescales single-particle observables under
$m_e \to m_\text{red}$; this projection adds the genuine two-body
Breit kinematic content that has no single-particle analog.\\
\textbf{Dimension introduced:} energy.\\
\textbf{Transcendental signature:} $\alpha^4 \cdot \mathbb{Q}$.  Pure
rational coefficient $\times \alpha^4$, no transcendental injected
beyond $\alpha$ already carried by the spectral-action projection.
Structurally the cleanest tier: same calibration class as the leading
magnetic dipole-dipole structure, ring-preserving over
$\mathbb{Q}(\alpha)$.  Distinct from the magnetization-density /
Zemach projection class (Paper~34 catalogue Layer-2 inputs), which
couples to nuclear spatial structure; this projection is a relativistic
two-body kinematic effect at the leading $\alpha^4$ order.\\
\textbf{Used in:} positronium 1S--2S two-photon transition (sprint
precision-catalogue 2026-05-08, equal-mass leptonic limit
$\lambda_a = \lambda_b = 0.5$); positronium 1$^3$S$_1$ HFS (annihilation
channel separated as a distinct Layer-2 mechanism); future positronium
fine-structure tests; future tests of the equal-mass $\mu^+ e^-$
limit if measured.\\
\textbf{Empirical anchor:} Ps 1S-2S framework-native at $+64.75$ ppm
vs Fee~1993 $1{,}233{,}607{,}216.4(3.2)$ MHz; Layer-2 $\alpha^4$
Breit input from Penin--Pivovarov~1998 (PRL~80, 2101) at the
$\sim 80$ GHz scale closes the residual to sub-MHz.\\
\textbf{Honest scope.}  This projection has one empirical anchor as
of 2026-05-08 (Ps 1S-2S).  Its definition is structural --- the
construction (single-particle Eides bracket $\to$ full two-body
Breit Hamiltonian) is established physics (Bethe--Salpeter~1957,
Karplus--Klein~1952, Penin--Pivovarov~1998, Czarnecki--Melnikov--Yelkhovsky~1999,
Karshenboim~2005~\S~4) --- but the framework's catalogue
verification rests on this single observable until additional
equal-mass or near-equal-mass two-body precision data are tested.\\
\textbf{Operator-level verification (2026-05-09, Sprint Calc-Ps-1S2S
Track 4 autopsy, \S\ref{sec:autopsy_ps_1s2s}):} PARTIALLY VERIFIED.
The retardation kernel $r_<^k/r_>^{k+3}$ is implemented as production
module \texttt{geovac.breit\_integrals} (26 tests; exact
Fraction/sympy arithmetic).  Closed-form matrix elements at $Z=1$,
$k=0$ for the $(1s,1s)$, $(1s,2s)$, $(2s,2s)$ orbital pairs that
participate in the Ps 1S--2S level structure live in
$\mathbb{Q}[\log 2, \log 3]$ (Paper~18 embedding-log content) for
the diagonal orbital pairs and in $\mathbb{Q}$ for the
$(1s,2s;1s,2s)$ exchange,
\emph{e.g.\,} $R^0_\text{BP}(1s,1s;1s,1s) = -5 + 8\log 2$,
$R^0_\text{BP}(1s,2s;1s,2s) = 4/81$
(production-verified by \texttt{tests/test\_breit\_integrals.py}
and \texttt{tests/test\_paper34\_projection\_spot\_checks\_batch1.py};
the $(1s,2s;1s,2s)$ value is a pure rational, no $\log$ content).
Roothaan multipole termination at $L_\text{max} = 2 l_\text{max}$ is
preserved at the equal-mass limit $\lambda_a = \lambda_b = 0.5$
because it is angular content (Wigner-3j triangle inequality,
\S~\ref{sec:proj_3j}), not a small-parameter expansion.
\emph{Named structural extension:} state-specific bound-state matrix
element evaluation $\langle\text{Ps}, 1S | H_B | \text{Ps}, 1S\rangle
- \langle\text{Ps}, 2S | H_B | \text{Ps}, 2S\rangle$ in equal-mass
two-body Dirac normalization requires the Bethe--Salpeter quasi-energy
expansion machinery, structurally outside the framework's
single-particle Eides~\S~3.2 SE bracket; this is the W2a-class
infrastructure target.  The $\sim 0.55\%$ match between the
calibration-by-construction Layer-2 input ($-79{,}861.9$ MHz $=$
$-0.2279 \times m_e c^2 \alpha^4/h$) and the canonical
Karshenboim~2005~\S~4 closed-form $-(11/48) = -0.2292$ is consistency
evidence at the magnitude level.\\
\textbf{Structural reading.}  This is the \emph{structurally cleanest
known instance} of the multi-focal-composition wall named in
CLAUDE.md~\S~2.  The wall's other five instances (Sprint H1
Yukawa, LS-8a $Z_2/\delta m$, HF-3 recoil, HF-4 Zemach, HF-5
multi-loop $a_e$) all surface in the multi-loop / renormalization
sector at $\alpha^5$ or higher, where two-body-projection failure
is entangled with renormalization failure.  The Breit retardation
content surfaces at $\alpha^4$ --- one full order \emph{before}
the LS-8a renormalization wall --- so the two-body-projection
failure is isolated and visible without renormalization confounds.
This is also what places the projection in the Layer-2 input class
(structurally a sibling of magnetization-density / Zemach in
\S~\ref{sec:proj_magnetization_density}): both take the framework's
single-particle output and add a structural correction the bare
spectral action does not autonomously generate.

\subsection{Nuclear charge-density correction (Foldy/Friar)}
\label{sec:proj_charge_density}
\textbf{Source:} bare Coulomb $V_\text{ne}(\vec{r}) = -Z e^2 / |\vec{r}|$
on the Fock-projected $S^3$ graph (point-source nucleus).\\
\textbf{Target:} convolved Coulomb $V_\text{ne}(\vec{r}) \to \int
d^3r' \, \rho_E^\text{nuc}(\vec{r}')\,(-Ze^2/|\vec{r}-\vec{r}'|)$,
with $\rho_E^\text{nuc}$ the nuclear electric charge distribution; at
leading order in the small-$r$ expansion the correction reduces to
the Foldy / Friar contact term
$\Delta V = +\tfrac{2\pi}{3} Z\alpha\,\langle r^2\rangle_E\,\delta^3(\vec{r})$.\\
\textbf{Variables introduced:} $\langle r^2\rangle_E$ (or
equivalently the charge radius $r_E$ for a given species).\\
\textbf{Dimension introduced:} length (one new $[L]$ scale per
nuclear species: proton $r_p$, deuteron $r_d$, etc.; structurally
distinct from atomic / internuclear length scales).\\
\textbf{Transcendental signature:} ring-preserving over
$\mathbb{Q}(\alpha)$.  The Foldy / Friar correction is rational in
the charge form-factor moments at every order of the small-$r$
expansion; the only transcendental content is the foundational
$\alpha$ already present in the spectral-action chain.\\
\textbf{Used in:} hydrogen Lamb shift $r_p$ contribution (currently
consumed externally via Eides Tab.~7.3 / Friar 1979); $\mu$H Lamb
shift Friar moment leading-order (Sprint MH Track A; $4.3\%$ at
leading); future framework-native operator-level Foldy module
(sibling of \texttt{geovac/magnetization\_density.py}, not yet
implemented).\\
\textbf{Empirical anchor:} proton charge radius
$r_p = 0.84075(64)$ fm (CODATA 2022); deuteron charge radius
$r_d = 2.1415(45)$ fm (Pohl et al.~2017).  Hydrogen 2$S$ Lamb shift
contribution from $r_p$ at the $\sim +1.18$ MHz level (Sprint LS-7
reframing; Eides Tab.~7.3).\\
\textbf{Honest scope.}  The framework currently consumes $r_p, r_d$
as external scalars (\texttt{R\_PROTON\_BOHR} and analogs) at the
point-source level.  Operator-level construction of $V_\text{ne}
\to V_\text{ne} \star \rho_E$ is the natural next step but has not
been implemented; the catalogue rows that depend on this projection
treat $r_p$ as a Layer-2 input.\\
\textbf{Structural reading.}  Categorically distinct from
\S~\ref{sec:proj_magnetization_density}: the Hamiltonian operator
modified is $V_\text{ne}$ (Coulomb), not $A_\text{hf}^\text{contact}$
(Fermi contact).  Empirical separation: hydrogen Lamb shift depends
on $r_p$ but only trivially on $r_Z$, while 21~cm hyperfine depends
on $r_Z$ but only trivially on $r_p$.

\subsection{Nuclear magnetization-density correction (Zemach)}
\label{sec:proj_magnetization_density}
\textbf{Source:} bare Fermi-contact hyperfine operator
$A_\text{hf}^\text{contact} = (8\pi/3)\,\mu_e \cdot \mu_N\,\delta^3(\vec{r})$
on the Fock-projected $S^3$ graph (point-source nuclear magnetic
moment).\\
\textbf{Target:} extended-magnetization Fermi-contact operator
$A_\text{hf}^\text{contact}(1 - 2 Z\alpha m_e r_Z + O(r_Z^2))$,
where $r_Z = \int d^3r\, d^3r'\,\rho_E^\text{nuc}(\vec{r}')\,
\rho_M^\text{nuc}(|\vec{r}-\vec{r}'|)\,|\vec{r}|$ is the Zemach
radius defined as the convolution moment of the nuclear charge and
magnetization profiles.\\
\textbf{Variables introduced:} $r_Z$ per nuclear species
(sub-features: nuclear magnetic radius $r_M$ and Friar moment
$\langle r^3\rangle_{(2)}$ enter as profile-shape parameters).\\
\textbf{Dimension introduced:} length (a new $[L]$ scale,
structurally distinct from \S~\ref{sec:proj_charge_density}).\\
\textbf{Transcendental signature:} ring-preserving over
$\mathbb{Q}(\alpha)$ at leading order; profile-dependence at
sub-leading order admits Gaussian, exponential, and dipole forms
with structurally identical leading $-2 Z\alpha m_e r_Z$ behaviour
(profile independence at leading order verified in Sprint HF
Track 4 / Sprint MH Track C).\\
\textbf{Used in:} hydrogen 21~cm hyperfine HF-4 row (Sprint HF;
$r_Z = 1.045$ fm yields $-39.5$ ppm framework-native shift); $\mu$H
1S Zemach mass-enhancement (Sprint MH Track B/C; factor $185.94
\times$ ep enhancement at leading order); deuteron 1S hyperfine
(Sprint precision-catalogue; $r_Z(D) = 2.593$ fm yields $-98.0$ ppm
framework-native shift).\\
\textbf{Empirical anchors:} proton Zemach radius
$r_Z = 1.045(20)$ fm (Eides 2024 atomic compilation) or $r_Z =
1.013(10)(12)$ fm (lattice QCD 2024); deuteron Zemach radius
$r_Z(D) = 2.593(16)$ fm (Friar--Payne 2005).\\
\textbf{Honest scope.}  Operator-level construction is mature:
\texttt{geovac/magnetization\_density.py} ($\sim 600$ lines, 57
tests; Sprint MH Track C + W1b recoil-mixing extension, May 2026)
implements the convolved Fermi-contact operator with the
\texttt{lepton\_mass} parameter eliminating manual mass-scaling for
muonic systems and an opt-in \texttt{include\_recoil\_mixing} flag
extending the kernel to next-to-leading order with the recoil-mixing
prefactor $m_l/(m_l + m_n)$ (per arXiv:2604.06930 eq.~95) and the
Friar moment correction $(1/2)(Z\alpha m_\text{red})^2 \langle r^2
\rangle_{(2)}$ (per Friar~1979 / Eides~\S6.2).  At leading order the
kernel matches Eides analytical to floating-point precision; with
recoil-mixing on, the framework muH 1S HFS prediction matches
Krauth~2017's full-theory Zemach line to $\sim 1\%$ (vs $\sim 3\%$
overshoot at LO-only).  $r_Z$ values themselves are Layer-2 inputs
(cannot be predicted from framework structure; consumed from QCD
lattice or atomic spectroscopy).  The W1b recoil-mixing extension
sprint (May 2026, see \texttt{debug/w1b\_recoil\_mixing\_extension\_memo.md})
showed that this kernel improvement does NOT close the $\sim 0.22$~fm
muH-alone vs Eides offset diagnosed by Sprint Calc-rZG-extended;
that offset is now attributable to Layer-2 itemization-convention
mismatch between Krauth~2017, Eides~2024, and Antognini-Krauth-Pohl~2015,
not to kernel approximation.  Operator-level verified at I=1 in
\S\ref{sec:autopsy_d_hfs} (D 1S HFS autopsy, May 9 2026):
substituting $r_Z(D) = 2.593$~fm reproduces the Eides analytic LO
scalar at $\sim 1.45 \times 10^{-14}\%$ of the LO shift (machine
precision); profile (Gaussian vs exponential) independence preserved
at I=1; the operator does not depend on nuclear spin (it operates
on the spatial register only, with $r_Z$ as a scalar moment).  Pauli
encoding remains 4 terms ($II, Z_e, Z_p, Z_e Z_p$) at I=1, identical
minimal sparse encoding as I=1/2.\\
\textbf{Structural reading.}  Sibling of
\S~\ref{sec:proj_charge_density} (different Hamiltonian operator)
and of \S~\ref{sec:proj_breit_retardation} (Breit retardation:
both take the framework's single-particle output and add a
structural correction the bare spectral action does not
autonomously generate).  Categorical distinction from
charge-density is empirical: 21~cm hyperfine depends on $r_Z$ but
trivially on $r_p$, electronic Lamb shift depends on $r_p$ but
trivially on $r_Z$.

\subsection{Nuclear tensor multipole coupling}
\label{sec:proj_tensor_multipole}
\textbf{Source:} multipole expansion of the nuclear charge potential
$V_\text{ne}(\vec{r}) = -Ze^2/|\vec{r}|$, truncated at the
monopole on the bare Fock-projected $S^3$ graph (spherical-symmetric
nucleus).\\
\textbf{Target:} same potential with $L \geq 2$ multipole
corrections retained.  At leading rank-2: quadrupole coupling
$H_Q = -e\,Q_N\,T^{(2)}_{ij}\,(\partial_i E_j)/6$, with $Q_N$ the
nuclear electric quadrupole moment and $T^{(2)}_{ij}$ a rank-2
spin tensor on the nuclear-spin Hilbert space.  Higher ranks
(octupole, hexadecapole) follow the same pattern at higher angular
momentum.\\
\textbf{Variables introduced:} $Q_N$ per nuclear species (rank-2
leading); higher-rank moments enter at higher $L$.\\
\textbf{Dimension introduced:} length-squared (the rank-2
quadrupole has dimension $[L]^2$; rank-$L$ multipoles have dimension
$[L]^L$).  This is the only projection in \S~\ref{sec:layer2} whose
dimensional injection is not strictly $[L]^1$ or $[E]$; the
squared-length structure reflects the tensor character of the
operator.\\
\textbf{Transcendental signature:} ring-preserving over $\mathbb{Q}$
at the angular level (multipole expansion uses Wigner $3j$ at every
step).  Sub-leading corrections involve the same algebraic ring as
\S~\ref{sec:proj_charge_density}.\\
\textbf{Used in:} deuteron 1S hyperfine s-state contribution
(sub-ppm; \emph{not} load-bearing in current catalogue); future
molecular HD / D$_2$ rotational hyperfine where electronic
gradients at the nucleus are nonzero and $Q_d$ couples at the
percent level (Komasa--Pachucki 2020 sub-percent precision regime);
future heavy-nucleus hyperfine spectra where $L \geq 1$ valence
orbitals couple to $Q_N$ at leading order.\\
\textbf{Empirical anchors:} deuteron quadrupole moment
$Q_d = 0.285699(15)(18)$ fm$^2$ (Komasa--Pachucki 2020); $^{17}$O
and $^{14}$N quadrupole moments for molecular spectroscopy
benchmarks; sub-percent precision spectroscopy in HD / D$_2$
rotational HFS as the load-bearing test.\\
\textbf{Honest scope.}  Forward-looking projection slot.  No
catalogue row currently uses this projection at leading order
(deuteron 1S HFS s-state sees only sub-ppm contribution; the
load-bearing tests are in molecular rather than atomic
spectroscopy).  Operator-level construction has not been
implemented; this projection occupies a structural slot, with the
empirical anchor in molecular HD / D$_2$ HFS not currently
catalogued.\\
\textbf{Operator-level test status.}  Track~4 of the multi-track
Roothaan autopsy sprint (2026-05-18) produced an order-of-magnitude
scoping autopsy on HD molecule $J=1$ deuteron quadrupole coupling
$\chi_d \approx 224.54$~kHz (Code-Ramsey 1971; Komasa-Pachucki
2020 ab~initio 224.51(5)~kHz).  Using only quantities the
framework currently exposes analytically (multipole expansion of
the partner nucleus's bare Coulomb potential at $R_\text{eq} =
1.4011$~bohr, plus a free-atom 1s electron screening estimate,
plus the $J=1, m_J=0$ rigid-rotor Wigner-3j angular reduction
$\langle V_{zz}\rangle = -(2/5)\,V_{zz}^\text{mol}$), the
framework-side OoM estimate gives $\chi_d^\text{OoM} = 135$~kHz
vs experiment $224.54$~kHz, residual $-89.5$~kHz $= -39.9\%$.
The bare-nuclear-only contribution alone reaches $195$~kHz, 87\%
of experiment; the 40\% gap is the absence of molecular bond-charge
redistribution from the free-atom picture.  The architecture is
confirmed structurally consistent (the framework's existing
multipole machinery captures the dominant LO contribution), but
sub-percent reproduction requires a position-space electronic-density
operator at the deuteron coordinate.  This is the named follow-on
sprint \texttt{HD-EFG-position-space} (beyond sprint scale; analog of the
W1a-D Roothaan recoil and \S\ref{sec:proj_magnetization_density}
operator-level extensions).  Source: \texttt{debug/HD\_rotational\_HFS\_scoping\_track4\_memo.md}.\\
\textbf{Structural reading.}  Categorically distinct from
\S~\ref{sec:proj_charge_density} and
\S~\ref{sec:proj_magnetization_density}: tensor coupling to a
gradient is structurally different from scalar coupling to a contact
density.  Different selection rules: $Q_N$ requires $l\geq 1$ on the
electronic side (does not contribute to s-states at leading order),
while charge-density and magnetization-density Foldy / Friar / Zemach
contributions are pure $l=0$ contact corrections.  The
Wigner--Eckart theorem ensures the tensor and scalar projections
are mutually orthogonal: they cannot be unified by a single
profile-broadening prescription.

\subsection{Phillips--Kleinman / core--valence orthogonality}
\label{sec:proj_phillips_kleinman}
\textbf{Source:} a valence-only Hamiltonian on the Fock-projected $S^3$
graph in which the core orbitals $\{\phi_c\}$ of a heavy atom are
\emph{frozen out} of the active Hilbert space, leaving the valence
electrons without any structural memory of the Pauli-exclusion
constraint they must satisfy with respect to the core.\\
\textbf{Target:} same valence-only Hamiltonian with the Phillips--Kleinman
projection added:
\begin{equation*}
  H_\text{val} \;\longrightarrow\; H_\text{val}
  \;+\; \sum_{c \in \text{core}} (E_v - E_c)\,
  |\phi_c\rangle\langle\phi_c|\,,
\end{equation*}
or equivalently at the matrix-element level
$\Delta H_{pq}^\text{PK} = \sum_c (E_v - E_c)\,S_{pc}\,S_{cq}$ for
valence indices $p,q$ and cross-overlap $S_{pc} = \langle\phi_p\,
|\,\phi_c\rangle$, with $\{\phi_c, E_c\}$ the (external) frozen-core
orbitals and binding energies and $E_v$ a valence reference energy
(typically $E_v = 0$ for the standard absolute-PK barrier, which is
purely repulsive whenever $E_c < 0$).\\
\textbf{Variables introduced:} the indexed set
$\{\phi_c, E_c\}_{c \in \text{core}}$ of core orbital wavefunctions
and binding energies.  Structurally one bookkeeping integer (the
core orbital count) plus an external Layer-2 input for each $E_c$ and
each radial profile $\phi_c$.  These are not generated by the
framework's bare graph or projection chain; they are consumed from
Clementi--Roetti HF tabulations or from the framework's own core-shell
solution (Paper~17~\S~IV's \texttt{CoreScreening} object,
which solves the same-center core-electron problem independently and
hands its eigenvectors to the valence solver).\\
\textbf{Dimension introduced:} none.  Orthogonality is dimensionless;
the projector $\sum_c |\phi_c\rangle\langle\phi_c|$ takes the same
dimension as its argument.  The valence reference $E_v$ and the
core eigenvalues $E_c$ carry energy, but they were already part of the
spectrum-of-states ledger before the projection was applied.  In this
sense Phillips--Kleinman is in the same dimensionless class as
Drake--Swainson asymptotic subtraction (\S\ref{sec:proj_drake_swainson})
--- both are bookkeeping projections that re-organize an already-
dimensioned quantity rather than introduce a new physical scale.\\
\textbf{Transcendental signature:} ring-preserving.  No transcendental
is injected beyond what the source spectrum already carries.  In the
finite hydrogenic / Slater basis used by the composed and balanced
architectures, the cross-overlaps $S_{pc}$ are evaluated in closed
form (radial integrals via the hypergeometric Slater machinery
\texttt{geovac/hypergeometric\_slater.py}) and live in $\mathbb{Q}$ at
the angular level (Wigner~$3j$, \S\ref{sec:proj_3j}) and in the same
algebraic-extension ring as the underlying orbitals at the radial
level.  The PK weights $(E_v - E_c)$ are external Layer-2 numerical
inputs whose transcendental content is whatever the source compilation
delivers.  No new $\pi$, no new $\zeta$ value, no new Catalan-class
transcendental enters from the projection itself.\\
\textbf{Used in:} the same-center PK barrier in the composed
architecture (Paper~17~\S~IV, module
\texttt{geovac/ab\_initio\_pk.py}) for LiH and BeH$^+$ at
$l_\text{max} \leq 2$; the $l$-dependent variant
(Paper~17~\S~IV, channel-blind $\to$ $l = 0$ projected) which
reduces LiH $R_\text{eq}$ error from $6.9\%$ to $5.3\%$; the
cross-center PK barrier
(Paper~19~\S~6, module
\texttt{geovac/phillips\_kleinman\_cross\_center.py}) for second-row
chemistry where the heavy-atom core is frozen out via a
\texttt{FrozenCore} potential rather than carried as an explicit
electron block.\\
\textbf{Empirical anchor:} LiH balanced-coupled $R_\text{eq}$ at
$5.3\%$ error with the channel-blind $l$-dependent PK at
$l_\text{max} = 2$ (Paper~17~\S~IV, Track~CD baseline; the
PK contribution is necessary --- removing it produces an unbound PES,
see Track~CB record in CLAUDE.md~\S~3); core-valence
orthogonality satisfied numerically in the explicit-core architectures
(LiH balanced-coupled at $n_\text{max}=2$ has the Li-core block
diagonalized simultaneously with the LiH-bond block; Pauli is
enforced by Slater-determinant antisymmetry in the FCI sector and the
PK projector is redundant for the explicit-core slice but
load-bearing for the frozen-core slice).\\
\textbf{Honest scope.}  Three layers of scope-honesty are required
here, all flagged for CLAUDE.md~\S~1.5 rhetoric compliance.

First, in the composed / balanced architecture as implemented, the PK
projection is \emph{essential} --- it is not redundant in second
quantization.  The Track~CB sprint (CLAUDE.md~\S~3 record, April~2026)
tested whether explicit cross-block ERIs could substitute for the PK
barrier by carrying core-valence repulsion explicitly; the result was
that adding cross-block $V_{ee}$ without simultaneously balancing
cross-center $V_{ne}$ produces a $29\%$ FCI energy error at LiH
$n_\text{max}=2$ (worst of three tested configurations, CLAUDE.md~\S~3
record).  The composed orbital basis uses different effective charges
$Z$ on different blocks, and PK is the infrastructure that handles
the corresponding orthogonality without requiring two-center one-body
integrals in a single-center hydrogenic basis.  PK is \emph{not} a
calibration choice; it is the projection that makes the composed
architecture's Hilbert-space factorization consistent with Pauli
exclusion.

Second, the same-center PK in
Paper~17~\S~IV \cite{paper17} closes the
core-valence-orthogonality problem cleanly at the $l_\text{max} = 2$
operating point.  Six prior PK modification attempts
(channel-blind, projector, spectral, self-consistent, $R$-dependent,
$l$-dependent on the 2D solver; see the
six-row PK modifications digest in the CLAUDE.md~\S~3
failed-approaches table) all carry the same lesson:
PK provides coordinate-space exclusion that angular projectors cannot
replicate.  The $l_\text{max}$ divergence (Paper~17~\S~V
and outlook section) at higher $l_\text{max}$ is a structural feature
of the composed architecture's interaction between PK and angular
basis expansion, not a deficiency of the PK projection itself; the
2D variational solver reproduces the same drift, confirming that PK
is not the bottleneck (Paper~17 Track~BQ record).

Third, the cross-center PK barrier
(Paper~19~\S~6, May~2026) extends the same
projection to the heavy-atom-frozen-core regime
($Z \geq 11$, [Ne] / [Ar] / [Kr] / [Xe] cores) and was tested as the
named engineering closure for the W1c-residual orthogonality wall
identified in the Phase~C chemistry sprint.  The honest verdict was
NEGATIVE on binding outcome but POSITIVE on architectural completeness:
PK cross-center reduces the W1c-residual NaH overattraction by
$14.6\%$ ($0.357$ Ha $\to$ $0.305$ Ha at standard absolute-PK
$E_v = 0$) but does \emph{not} close the wall.  The NaH PES remains
monotonically descending with no internal equilibrium minimum.  The
remaining residual ($\sim 290$~mHa, $4\times$ experimental
$D_e \approx 75$~mHa) lives in mechanisms beyond W1c + PK orthogonality:
the hydrogenic $Z_\text{orb} = 1$ basis on the Na valence side is
qualitatively wrong (it is not the actual Na 3s orbital but a generic
$Z_\text{eff} = 1$ hydrogenic function), the cross-block ERI between
H~1s (proper) and Na ``valence'' (wrong shape) inherits this error,
and a strong-PK probe at unphysical $E_v = 100$~Ha shows a
minimum-depth signature ($0.148$~Ha descent) that represents parameter
tuning rather than derivation.  The CLAUDE.md~\S~3 failed-approaches
table carries the explicit row ``Phillips-Kleinman cross-center
barrier for second-row chemistry binding (post-Track-2 sprint,
2026-05-08)'' documenting this honest negative.  The PK projection
remains a clean architectural addition; the wall is structurally
deeper than PK orthogonality alone.\\
\textbf{Structural reading.}  Phillips--Kleinman sits in the
Layer-2~$\to$~Layer-2 input class, structurally a sibling of the three
nuclear-structure projections (\S\ref{sec:proj_charge_density},
\S\ref{sec:proj_magnetization_density}, \S\ref{sec:proj_tensor_multipole}):
each takes the framework's bare single-particle output and adds a
structural correction whose load-bearing data is external to the
framework's graph and projection chain.  For the nuclear-structure
trio the external data is a QCD-internal scalar
($\langle r^2\rangle_E$, $r_Z$, $Q_N$); for Phillips--Kleinman the
external data is the set of frozen-core orbitals
$\{\phi_c, E_c\}$ derivable either from an independent atomic
Hartree--Fock solve (Clementi--Roetti) or from the framework's own
core-shell Level~3 solution.  In both cases the projection mechanism
is structurally well-defined and ring-preserving; in both cases the
input scalars / orbitals are calibration content the framework does
not autonomously derive.

Two further structural distinctions deserve naming.  First,
Phillips--Kleinman differs categorically from the three
nuclear-structure projections in its targeted Hilbert-space sector:
the nuclear-structure projections modify the
electron-nucleus or hyperfine operators by convolving with a nuclear
density profile, while Phillips--Kleinman modifies the kinetic-energy
\emph{plus} potential-energy effective single-particle Hamiltonian
on the valence electrons via an orthogonality projector.  Different
mechanism, same input class.

Second, the same-center and cross-center variants of PK are
structurally one projection in two architectures.  In the same-center
(explicit core block) architecture, the core orbitals live inside the
FCI sector and Pauli is enforced by Slater-determinant antisymmetry;
the PK barrier in Paper~17~\S~IV is the residual content
that handles the kinetic-energy cost of orthogonalizing the angular
component of the valence basis against the radial core.  In the
cross-center (frozen-core) architecture, the core orbitals are
external to the FCI Hilbert space (they live in a one-body local
$Z_\text{eff}(r)$ potential), and the PK barrier (Paper~19~\S~6)
is the non-local one-body operator that
re-injects core-valence orthogonality the frozen-core approximation
removed.  Same operator-form projection, two architectures, two
load-bearing roles, the same load-bearing limitation: PK gets the
projection structure right but cannot independently fix a wrong
valence basis.

\subsection{Multipole expansion / Gaunt termination}
\label{sec:proj_multipole_gaunt}
\textbf{Source:} bare two-body Coulomb kernel $1/|\vec{r}_1 - \vec{r}_2|$
(for $V_{ee}$) or $1/|\vec{r} - \vec{R}_B|$ (for cross-center $V_{ne}$)
acting on a product of angular-momentum eigenstates on the Fock-projected
$S^3$ graph.\\
\textbf{Target:} the same kernel re-expressed as a finite sum over
multipoles
\begin{equation*}
  \frac{1}{|\vec{r}_1 - \vec{r}_2|}
  \;=\; \sum_{L=0}^{L_\text{max}}
        \frac{r_<^L}{r_>^{L+1}}\, P_L(\cos\theta_{12}),
\end{equation*}
with the angular content $P_L(\cos\theta_{12})$ contracted against
spherical-harmonic vertices via Gaunt integrals
$G(l_1, L, l_2) = 2 \begin{pmatrix} l_1 & L & l_2 \\ 0 & 0 & 0
\end{pmatrix}^{\!2}$, i.e.\ Wigner-$3j$ angular coupling
(\S~\ref{sec:proj_3j}) at each vertex.\\
\textbf{Variables introduced:} the multipole order $L$ (integer, runs
$0, 1, \dots, L_\text{max} = 2 l_\text{max}$).  No physical parameter
is introduced; $L$ is an internal summation index whose upper bound is
fixed by the angular content of the basis.\\
\textbf{Dimension introduced:} dimension-preserving.  Each multipole
term $r_<^L / r_>^{L+1}$ carries the same $[E]$ scaling as the parent
Coulomb matrix element; the decomposition reorganizes the angular
content of an existing operator without injecting a new $[L]$, $[E]$,
$[M]$, or $[T]$ scale.\\
\textbf{Transcendental signature:} ring-preserving over the Layer~1
algebraic ring.  The Gaunt $3j$ coefficients live in
$\mathbb{Q}[\sqrt{2k{+}1}]_k$ (the Wigner-$3j$ ring of
\S~\ref{sec:proj_3j}); the radial $R^L$ matrix elements are algebraic
in the hydrogenic / $Z_\text{eff}$-rational basis (Slater-type
integrals evaluate to rationals at integer screening and to
incomplete-gamma combinations at fractional screening, per
\S~\ref{sec:proj_3j} sibling derivations and Paper~19 \S~III.C).  No
transcendental is introduced at any step.\\
\textbf{Used in:} essentially every multi-electron matrix element in
the codebase.  Specific instances: Level~4 mol-frame hyperspherical
$V_{ee}$ and $V_{ne}$ coupling (Paper~15 \S~V, split-region Legendre
expansion); cross-center $V_{ne}$ in the balanced coupled architecture
(Paper~19 \S~III.C; \texttt{geovac/shibuya\_wulfman.py}); single-center
Slater $F^k$ and $G^k$ integrals at every level (Paper~7 \S~VI.B,
Paper~12); the Roothaan cross-register $V_{eN}$ kernel at multi-focal
$\lambda_a \ne \lambda_b$ (\texttt{geovac/cross\_register\_vne.py});
the multipole structure underlying \S~\ref{sec:proj_tensor_multipole}
(nuclear rank-$L$ couplings).\\
\textbf{Empirical anchor:} the multipole expansion underwrites the
angular sparsity of the Pauli-string Hamiltonian across the catalogue
(Paper~22 angular sparsity theorem; ERI density $1.44\%$ at
$l_\text{max} = 3$ depends on the exact $L \leq l_1 + l_2$
truncation).  Direct numerical verification in Sprint CD (April 2026,
Paper~19 \S~III.C, Track CD v2.0.39): cross-center $V_{ne}$ for LiH at
$n_\text{max} = 2$ contributes exactly 33 nonzero one-body matrix
elements (multipole sum runs $L = 0, 1, 2$, terminates at $L_\text{max}
= 2$); at $n_\text{max} = 3$, 168 nonzero elements ($L_\text{max} =
4$).  Machine-precision evaluation via the regularized incomplete
gamma function (\texttt{scipy.special.gammainc} / \texttt{gammaincc})
replaces grid-based trapezoid quadrature with zero truncation error at
the angular level.\\
\textbf{Honest scope.}  The termination $L \leq L_\text{max} =
2 l_\text{max}$ is \emph{exact}, not asymptotic.  It is forced by the
Wigner-$3j$ triangle inequality $|l_1 - l_2| \leq L \leq l_1 + l_2$
applied to the angular Gaunt integral at each vertex
(\S~\ref{sec:proj_3j}): for $L > l_1 + l_2$ the $3j$ symbol vanishes
identically, so the multipole sum has no nonzero contribution past
$L_\text{max}$.  There is no truncation parameter to tune, no UV
cutoff to remove, no convergence series to evaluate.  This is
categorically different from the slowly convergent Shibuya--Wulfman
\emph{basis} expansion (Paper~19 \S~III.B), which expands a complete
orbital basis on one center in the Sturmian basis of another center
and does \emph{not} terminate; the multipole expansion of the
$1/|\vec{r}_1 - \vec{r}_2|$ \emph{kernel} terminates trivially while
the Shibuya--Wulfman expansion of the \emph{basis} does not.
Computational note: where the radial integrand at fractional screening
is non-polynomial (the smooth $[Z_A - Z_\text{eff}(r)]/r$ correction
in composed geometries), the algebraic expansion does not apply and
Gauss--Legendre quadrature is used for the radial integral only; the
angular content remains exactly truncated regardless
(Paper~15 \S~V, closing paragraph).\\
\textbf{Structural reading.}  Layer~1 $\to$ Layer~1.  This projection
reorganizes the angular content of the Coulomb kernel \emph{prior to}
any focal-length, mass, or observation projection; it does not move
the chain across the Layer-1/Layer-2 boundary of
\S~\ref{sec:layer2_boundary}.  Categorically, it is the coordinate-side
counterpart of \S~\ref{sec:proj_3j}: the multipole expansion of
$1/|\vec{r}_1 - \vec{r}_2|$ produces the angular vertices at which
Wigner-$3j$ coupling is then applied, so the two projections are
inverse images of each other across the kernel-versus-vertex divide
($P_L(\cos\theta_{12}) \leftrightarrow G(l_1, L, l_2)$).  Operationally
they almost always compose: in the framework's matrix-element pipeline
the multipole expansion is the kernel-side step that produces the
$3j$-symbol calls at each vertex.  Structurally distinct from
\S~\ref{sec:proj_tensor_multipole} (nuclear tensor multipole coupling),
which uses the same Wigner-$3j$ angular machinery but applies it to
nuclear-structure tensor moments (one new $[L]^L$ length scale per
rank); the present projection introduces no physical scale at all and
sits cleanly inside Layer~1.  This is the workhorse projection that
underwrites the entire angular-sparsity story of Paper~22 and the
$O(Q^{2.5})$ Pauli scaling of Paper~14: every claim about ERI density,
selection-rule survival, or composed-architecture sparsity ultimately
rests on the exact $L \leq 2 l_\text{max}$ termination of this
projection.

\subsection{Bipolar harmonic / Drake combining}
\label{sec:proj_bipolar_drake}
\textbf{Source:} pair of single-electron angular tensor operators
acting on distinct electrons of a multi-electron state on the
Fock-projected $S^3$ graph, in the LS coupling scheme.\\
\textbf{Target:} a bipolar tensor harmonic
$\{Y^{(k_1)}(\hat r_1) \otimes Y^{(k_2)}(\hat r_2)\}^{(K)}_M$
coupling the two angular factors into a single rank-$K$ object
on the joint two-electron angular space, with $K$-dependent
matrix elements in the coupled $|LSJ\rangle$ basis given by
$(-1)^{L+S+J}\,6j\{L\,S\,J;\,S\,L\,k\}$ (rank-$k$ Racah J-pattern).\\
\textbf{Variables introduced:} the bipolar coupling indices
$(k_1, k_2)$ and the total angular rank $K$ (integer triple with
triangle constraint $|k_1 - k_2| \le K \le k_1 + k_2$).  These are
internal-graph-derived combinatorial labels in the sense of
\S~\ref{sec:layer2_boundary}; they introduce no new dimensioned
scale and no continuous parameter, only an enumeration of
Gaunt-allowed angular channels.\\
\textbf{Dimension introduced:} preserved (the projection is a
purely angular reorganization on the multi-electron Hilbert
space; no $[L]$, $[E]$, or $[T]$ scale enters).\\
\textbf{Transcendental signature:} pure rational.  The bipolar
harmonic expansion of a two-electron tensor operator (Breit
retardation, spin-spin, spin-other-orbit, orbit-orbit) decomposes
on the angular side into Wigner $3j$ Gaunt coefficients (one per
single-electron factor) times rank-$k$ $6j$ symbols (for the
J-pattern in $|LSJ\rangle$) times rank-$(k_1, k_2, K)$ $9j$
symbols (for the bipolar recoupling to the coupled
two-electron space).  Each ingredient is exactly rational over
$\mathbb{Q}[\sqrt{2k+1}]_{k}$ --- the same Wigner ring as
\S~\ref{sec:proj_3j}.  The Drake (1971) combining coefficients
$(\tfrac{3}{50}, -\tfrac{2}{5}, \tfrac{3}{2}, -1)$ for spin-spin
and spin-other-orbit on the helium triplet are pure rationals
(Paper~18 intrinsic tier).  No transcendental injection at any
step of the angular reduction.\\
\textbf{Used in:} two-body Breit-Pauli spin-spin (rank-$K{=}2$
spatial, rank-2 coupled spin tensor $[s_1 \otimes s_2]^{(2)}$) and
spin-other-orbit (rank-$K{=}1$ spatial, rank-1 sum tensor
$\vec s_1 + 2\vec s_2$) interactions on the relativistic composed
qubit encoding (Paper~14~\S~V); helium-class fine-structure
predictions through the He~$2{}^3\!P$ Roothaan autopsy
(\S~\ref{sec:autopsy_he_2_3p}); structurally compatible with all
two-body retardation kernels in
\texttt{geovac/breit\_integrals.py}; angular machinery underlying
Drake's $M^K_{\rm dir} / M^K_{\rm exch}$ direct/exchange
decomposition of the helium triplet, with Sprint~5 DV
characterizing the bipolar-to-Drake basis change for the
He~$(1s)(2p)\,{}^3\!P$ configuration as a
convention-dependent piecewise recombination of Region-I and
Region-II radial pieces.\\
\textbf{Empirical anchor:} He $2{}^3\!P$ fine structure dominant
intervals at NIST precision.  Framework-native operator-level
decomposition (Sprint multi-track Roothaan autopsy 2026-05-09,
\S~\ref{sec:autopsy_he_2_3p}) reproduces
$E({}^3\!P_0) - E({}^3\!P_1) = +29{,}612.91$~MHz vs NIST
$+29{,}616.951(6)$~MHz at $-0.014\%$ and
$E({}^3\!P_0) - E({}^3\!P_2) = +31{,}844.07$~MHz vs NIST
$+31{,}908.131(6)$~MHz at $-0.20\%$, with the angular content (J-pattern,
spin reduced matrix elements, bipolar Gaunt selection) verified
\emph{sympy-exact} at machine precision via the rank-$k$ $6j$ and
$9j$ identities $(-1)^{L+S+J}\,6j\{1,1,J;1,1,k\}$ at $k = 1$ (SOO)
and $k = 2$ (SS), normalized so $f(J{=}1) = 1$.  Cross-reference:
Paper~14~\S~V and the sympy-verification tests
\texttt{tests/test\_breit\_integrals.py::test\_drake\_f\_SS\_pattern\_from\_6j}
and \texttt{test\_drake\_f\_SOO\_pattern\_from\_6j}.  At the
operator level, the bipolar $(k_1, k_2)$ channels for the He
$(1s)(2p)$ configuration reduce uniquely (Sprint~5 DV) to
$(k_1{=}0, k_2{=}2)$ for the SS direct path and $(k_1{=}1, k_2{=}1)$
for the SS and SOO exchange paths; the SOO direct path is
\emph{Gaunt-forbidden} at the bipolar level (parity selection forces
$k_1 = 0$, then the $(0, k_2, 1)$ triangle inequality requires
$k_2$ odd, contradicting the $(1)(1)$ Gaunt parity constraint).\\
\textbf{Honest scope.}  The small interval $E({}^3\!P_1) - E({}^3\!P_2)$
framework-native at $+2{,}231.16$~MHz vs NIST $+2{,}291.176(15)$~MHz
shows a $-2.62\%$ relative residual.  This is \emph{partial-cancellation
amplification} of the same $\sim$$60$~MHz absolute residual that sits
inside the LS-8a one-loop QED budget on the dominant intervals: SO
($-58{,}396$~MHz) $+$ SS ($-9{,}499$~MHz) $+$ SOO ($+70{,}126$~MHz)
combine to a net $+2{,}231$~MHz, a structural cancellation factor
$(|\text{SO}|+|\text{SS}|+|\text{SOO}|) / |\text{total}| \approx 61.9\times$.
The framework's absolute precision is uniform across the three
intervals; the small interval's fractional residual is amplified by
the cancellation.  This is a feature of how the SO + SS + SOO
operators combine on the $(1s)(2p)\,{}^3\!P$ configuration, not a
framework limitation.  Higher-order corrections (multi-loop QED at
$\alpha^4 \cdot \alpha/\pi$, finite-mass recoil, hyperfine averaging)
are external Layer-2 inputs (sibling status: Drake~1971~\S~IV).
The direct/exchange mixing coefficients
$(\tfrac{3}{50}, -\tfrac{2}{5}, \tfrac{3}{2}, -1)$ in
$A_{SS}$ and $A_{SOO}$ are pure rationals at the angular level but
\emph{not closed from $9j$ alone} on the He triplet: Sprint~5 DV
showed they are convention-dependent identities specific to Drake's
piecewise $M^K_{\rm dir,I} + M^K_{\rm dir,II}$ basis rather than a
direct Racah-algebra output (a fully first-principles derivation
remains open; see \texttt{debug/dd\_drake\_derivation.md} and
\texttt{debug/dv\_drake\_bipolar\_memo.md}).  The bipolar harmonic
projection \emph{itself} is closed; the open question is the basis
change to Drake's specific $M^K$ partitioning.\\
\textbf{Structural reading.}  Sibling of \S~\ref{sec:proj_3j}
(Wigner $3j$ angular coupling): both are angular-content-only
projections living in the pure-rational ring
$\mathbb{Q}[\sqrt{2k+1}]_{k}$ with no dimensional or transcendental
injection.  The structural distinction is multiplicity: Wigner $3j$
couples a \emph{single} pair of angular nodes (one electron's
orbital and one external rank-$k$ multipole), whereas the bipolar
harmonic / Drake combining projection couples \emph{two
single-electron angular tensors} into a joint rank-$K$ object on
the two-electron Hilbert space via the rank-$(k_1, k_2, K)$
$9j$ recoupling.  In dictionary terms, \S~\ref{sec:proj_3j} is
the single-particle special case of this projection at
$(k_1, k_2) = (0, k)$ on the two-electron space (one electron
inert).  The two projections compose: every bipolar reduction
decomposes into a product of two single-particle Wigner-$3j$ Gaunt
factors at the Slater integral level, together with a $9j$ for the
two-electron recoupling and a $6j$ for the J-pattern in
$|LSJ\rangle$.  \emph{Internal multi-focal angular-only finding}: on
configurations like He $(1s)(2p)$ where the two electrons sit at
structurally distinct focal lengths
($Z_{\rm eff}(1s)/Z_{\rm eff}(2p) = 2$), the bipolar Gaunt selection
$(k_1, k_2)$ channels are determined by angular content alone ---
the integer parity and triangle constraints --- and are
\emph{independent of the focal-length ratio}.  Roothaan multipole
termination at $L_{\max} = 2 l_{\max}$ (here $K \le 2$ for
$l_{\max} = 1$) holds verbatim irrespective of the relative
magnitudes of the two electrons' Slater exponents; this is the
\emph{internal} analog of the cross-register Roothaan termination
identified at the multi-focal-composition wall
(\S~\ref{sec:proj_breit_retardation} structural reading).  The
operator-level realization of internal multi-focal architecture is
therefore not a new projection: it is bipolar harmonic / Drake
combining acting at $\lambda_a \neq \lambda_b$ on two electrons in
the same atomic block.  Different focal-length ratio, same angular
projection, same rational ring.

\subsection{Symmetry projection / Young tableau}
\label{sec:proj_symmetry_tableau}
\textbf{Source:} the unrestricted $N$-particle tensor product
$\mathcal{H}^{\otimes N}$ on the single-particle Fock-projected $S^3$
graph (or on any single-particle natural-geometry sector;
\S~\ref{sec:proj_fock}, \S~\ref{sec:proj_bargmann},
\S~\ref{sec:proj_molframe}).  At this stage indistinguishability of
identical fermions has not yet been enforced; the Hilbert space
contains symmetric, antisymmetric, and mixed-symmetry components
without distinction.\\
\textbf{Target:} the $\lambda$-irrep subspace
$\mathcal{H}_\lambda \subset \mathcal{H}^{\otimes N}$, obtained via
the character-based orthogonal projector
\[
  P_\lambda \;=\; \frac{d_\lambda}{|S_N|}\,\sum_{g \in S_N}
                  \chi_\lambda(g)\,\rho(g),
\]
where $\lambda$ is a Young tableau (partition of $N$), $d_\lambda$ is
the dimension of the corresponding irreducible representation of
$S_N$, $\chi_\lambda(g)$ is its character at the permutation $g$, and
$\rho(g)$ is the natural action of $g$ on $\mathcal{H}^{\otimes N}$
by index permutation.  For fully indistinguishable spin-$\tfrac{1}{2}$
fermions in the unrestricted formulation, $\lambda = [1^N]$
(antisymmetric); in the Slater-determinant / spatial-spin coupled
formulation, $\lambda$ is the conjugate of the spin Young diagram
(Paper~16 convention), so that $S_N$
antisymmetry is implemented \emph{jointly} on spatial and spin
indices.\\
\textbf{Variables introduced:} the irrep tableau $\lambda$ itself
(a discrete combinatorial label; equivalently a partition of $N$).
For $N=2$: $\lambda \in \{[2],\,[1,1]\}$ (triplet vs singlet under
spatial-spin coupling).  For $N=4$ closed shell: $\lambda = [2,2]$
is the spin-singlet sector (Paper~17 §VIII, Track AJ).  For general
$N$ with total spin $S$: $\lambda_1 = N/2 + S$ and
$\lambda_2 = N/2 - S$ (Paper~16 convention).\\
\textbf{Dimension introduced:} \emph{none}.  $P_\lambda$ is an
orthogonal projection on the $N$-particle Hilbert space; it does not
introduce a new physical scale, focal length, or coordinate.  The
projector reduces the linear dimension of the working space from
$(\dim\mathcal{H})^N$ to a $\lambda$-irrep-counted subspace, but no
new $[L]$, $[E]$, or $[T]$ enters the framework.\\
\textbf{Transcendental signature:} \textit{none}.  The character
table of $S_N$ is integer-valued (Frobenius character formula); the
projector $P_\lambda$ has rational matrix entries (denominators
dividing $|S_N| = N!$, numerators integer); the dimension $d_\lambda$
is integer (hook-length formula).  No $\pi$, no $\zeta$, no Hurwitz
content enters at this step.  The projection is the unique
\S~\ref{sec:layer2} entry whose entire algebraic content lives in
$\mathbb{Z}/N!$, making it the cleanest example of a
\emph{combinatorial} Layer~1 $\to$ Layer~1 projection.\\
\textbf{Used in:} all $N \geq 2$ computations in the framework, at
multiple levels of explicitness:
\begin{itemize}
\item \textbf{Level 3 (He, $N=2$).}  Singlet $[1,1]$ vs triplet $[2]$
spatial-spin coupling is the projection that splits the He $1^1 S$
ground state from the $2^3 S$ first excited state.  Spatial-spin
coupling for two electrons with total spin $S \in \{0, 1\}$ yields
$\lambda_\text{spatial} = [1,1]$ for $S=0$ (antisymmetric spatial,
symmetric spin: singlet) and $\lambda_\text{spatial} = [2]$ for
$S=1$ (symmetric spatial, antisymmetric spin: triplet).
\item \textbf{Level 4N (LiH, $N=4$).}  The $S_4$ $[2,2]$ Young
diagram projection is the character-based projector onto the
closed-shell singlet sector of the 4-electron mol-frame
hyperspherical Hilbert space on $S^{11}$ (Paper~17 §VIII; Track AJ,
\texttt{geovac/n\_electron\_solver.py}).  The projection reduces the
raw SO(12) angular channel count by a tableau-dependent factor and
makes the 4-electron PES computationally tractable; it is empty at
$l_{\max} = 0$ (no antisymmetric 4-electron state can be built from
$s$-waves), gives 2 channels at $l_{\max} = 1$, 12 channels at
$l_{\max} = 2$, and 40 channels at $l_{\max} = 3$.
\item \textbf{Graph-native CI (all $N$).}  The Slater-determinant
basis used in \texttt{geovac/casimir\_ci.py} is itself the
\emph{Slater-determinant realization} of the $[1^N]$ antisymmetric
projector: each determinant
$\det[\phi_{i_1}(\vec{r}_1)\cdots\phi_{i_N}(\vec{r}_N)]$ is the
image under $P_{[1^N]}$ of the single Hartree product
$\phi_{i_1}(\vec{r}_1)\otimes\cdots\otimes\phi_{i_N}(\vec{r}_N)$,
up to a normalization factor.
\item \textbf{Composed atomic classifier (Paper~16).}  The atomic
structure-type assignment (A/B/C/D/E) for $Z = 1$--$30$ is itself
read off the $S_N$ irrep shape: structure-type letters correspond to
specific tableau patterns under the periodic-table indexing of
Paper~16.  The classifier is the boundary at
which the $S_N$ projection becomes computationally explicit in the
production composed-Hamiltonian pipeline.
\item \textbf{Hyperfine and cross-register couplings (Track NI,
Paper~23 §VI).}  In the composed nuclear--electronic deuterium
Hamiltonian, the Clebsch--Gordan multiplicity $(2I+1)/2$ relating
$F = I+J$ to $F = I-J$ hyperfine levels (D HFS autopsy,
\S\ref{sec:autopsy_d_hfs}) is itself a small-$N$ tableau projection
applied to a coupled spin Hilbert space.  See \S~\ref{sec:proj_3j}
for the angular-momentum coupling content; the present projection
is its $S_N$ analog on the indistinguishable-particle side.
\end{itemize}
\textbf{Empirical anchor:} He $2^1 S$--$2^3 S$ exchange splitting at
$6421.46$~cm$^{-1}$ (NIST experimental) is the cleanest finite-energy
witness of the singlet-vs-triplet ($[1,1]$ vs $[2]$ spatial)
projection in the catalogue.  The splitting is not a small Layer-2
correction; it is the energy difference between two distinct
$\lambda$-irrep sectors of the same single-particle basis, generated
entirely by the exchange-integral content of the two-electron $V_{ee}$
matrix elements under spatial antisymmetrization vs symmetrization.
Sprint Calc-He-2S-2S-splitting and Sprint Hylleraas-$r_{12}$
(2026-05-09) returned this entry at $+11.86\%$ (graph-native CI,
$n_{\max} \leq 11$) and $+209\%$ (single-$\alpha$ Hylleraas-3p,
$\omega = 5$) respectively, with the diagnosis that the residuals
trace to single-exponent basis limitations for the excited
$\lambda$-irrep state rather than to the projection itself
(\S~\ref{sec:offprecision} He singlet-triplet row).  At the
projection level, the splitting is structurally clean: the projector
$P_{[1,1]}$ vs $P_{[2]}$ acts identically on either side; only the
$V_{ee}$ matrix elements respond differently to the two sectors.\\
\textbf{Honest scope.}  At modest $N$ ($N \leq 4$) the framework
implements the character-based projector $P_\lambda$ directly:
explicit construction of the projector matrix in the Slater /
configuration-state-function basis, or via the equivalent Young
symmetrizer in the unrestricted tensor product.  This is exact and
combinatorial.  At large $N$ (heavy atoms, $N \gg 10$), the framework
falls back on \emph{compositional approximations} to the tableau
projection: (i) single-Slater-determinant Hartree--Fock reductions
project onto a single representative of the $[1^N]$ irrep rather
than spanning it, (ii) frozen-core constructions
(\texttt{geovac/neon\_core.py} for [Ne]-core, [Kr]-core, [Xe]-core
species) project the antisymmetric structure onto a product of a
fixed core determinant with a smaller-$N$ valence tableau projection,
(iii) the composed-block architecture for multi-center systems
(\S~\ref{sec:proj_molframe}) treats each block's antisymmetric
subspace independently, neglecting inter-block antisymmetry except
through cross-block ERI / cross-center $V_\text{ne}$ corrections.
Each of these is itself a structured projection onto a coarser
tableau pattern.  The boundary between exact character-projection
and compositional approximation is empirical, set by the largest $N$
at which $|S_N| = N!$ direct construction remains tractable.  As of
v2.40.0 this boundary sits at $N \approx 6$ for the atomic
classifier's explicit construction (extended via lookup tables to
higher $N$ where the periodic-table structure-type pattern repeats),
and at $N = 4$ for the Level~4N solver
(\texttt{geovac/n\_electron\_solver.py}).\\
\textbf{Structural reading.}  This is the projection at which
\textbf{Pauli antisymmetry enters the framework explicitly} rather
than implicitly.  In many quantum-chemistry frameworks
antisymmetrization is folded into the coordinate or operator choice
(e.g.\ second-quantized anticommutators built into the Hamiltonian's
algebraic structure); in GeoVac the dimensionless graph is
symmetry-agnostic by construction, and
Pauli antisymmetry is recovered as a separate, named, character-based
projection at the multi-particle level.  This is the structural
sibling of \S~\ref{sec:proj_spinor} (which introduces the
spin-$\tfrac{1}{2}$ Hilbert space at the single-particle level): the
spinor projection commits to particle statistics (fermionic
half-integer spin), while the present projection commits to the
many-body consequence of those statistics ($N$-particle
antisymmetry).  Composition $P_\lambda \circ
\text{spinor\_lift}$ realizes the standard spatial-spin coupling
scheme; composition $P_\lambda \circ \text{Fock}$ realizes the
$N$-fermion graph Hilbert space on which graph-native CI operates.
Categorically distinct from \S~\ref{sec:proj_3j} (angular-momentum
coupling, $SU(2)$ Clebsch--Gordan, continuous-symmetry character
content) and \S~\ref{sec:proj_wignerD} (rotation between rotated
coordinate frames, $SO(3)$ representation theory).  The shared
abstraction across these three is that all three are character-based
projections onto irreducible representations of finite or compact
groups acting on the Hilbert space; the structural distinction is
the group: $S_N$ for the present projection, $SU(2)$ for
\S~\ref{sec:proj_3j}, $SO(3)$ for \S~\ref{sec:proj_wignerD}.  The
$S_N$ character table has integer entries while the $SU(2)$ and
$SO(3)$ projector matrix elements introduce algebraic-extension and
trigonometric content respectively
(Table~\ref{tab:projection_axes}); the absence of transcendental
content in the symmetry projection is therefore not coincidental
but reflects the discrete finite-group structure underlying Pauli
exclusion.

\subsection{Adiabatic / Born--Oppenheimer projection}
\label{sec:proj_adiabatic_BO}
\textbf{Source:} a single combined Hamiltonian $H$ acting on two
coupled degrees of freedom whose dynamical time-scales differ by a
parametrically large ratio (the small parameter $m_e / M_n \ll 1$ for
electron--nucleus systems, $m_e / M_\text{cluster}$ for inner electrons
versus an outer cluster, etc.).  In the foundational setting this is
the molecular Hamiltonian $H = T_n(\vec{R}) + T_e(\vec{r}) +
V_{nn}(\vec{R}) + V_{ne}(\vec{r}, \vec{R}) + V_{ee}(\vec{r})$ on a tensor
product of nuclear and electronic Hilbert spaces; in the Level~3
hyperspherical setting it is the combined kinetic and potential
operator on the joint $(R, \Omega)$ space of hyperradius and
hyperangle.\\
\textbf{Target:} a \emph{two-stage} decomposition in which the fast
degree of freedom (electronic coordinates $\vec{r}$, hyperangle
$\Omega$, valence electrons) is solved at every fixed value of the
slow degree of freedom (nuclear coordinate $\vec{R}$, hyperradius $R$,
core configuration), producing a parameter family of fast-sector
eigenproblems
\begin{equation*}
  H_\text{fast}(\vec{R}) \, \Phi_\nu(\vec{r}; \vec{R})
  = E_\nu(\vec{R}) \, \Phi_\nu(\vec{r}; \vec{R}),
\end{equation*}
parameterized by the slow variable.  The slow degree of freedom then
sees the fast-sector eigenvalues $E_\nu(\vec{R})$ as an effective
potential, producing the slow-sector Schr\"odinger equation
$[T_n + E_\nu(\vec{R})] \chi(\vec{R}) = E\, \chi(\vec{R})$ on the
slow-only Hilbert space.  The full wavefunction is reconstructed as
$\Psi(\vec{r}, \vec{R}) \approx \Phi_\nu(\vec{r}; \vec{R}) \chi(\vec{R})$
(single-channel adiabatic) or as a sum
$\sum_\nu \Phi_\nu(\vec{r}; \vec{R}) \chi_\nu(\vec{R})$
(coupled-channel; see \S~\ref{sec:proj_coupled_channel}, the
sibling projection that solves the parameterized fast-sector problem
and supplies the coupling between channels).
\\
\textbf{Variables introduced:} the slow continuous parameter $\vec{R}$
itself, which enters the previously fast-sector-only graph
computation as a \emph{continuous} index over a family of
Hamiltonians.  The small parameter controlling the projection's
validity is the mass ratio $m_e / M_n$ (and, more generally, the
ratio of fast and slow time-scales).  Both quantities are content
the bare graph does not carry: the bare graph computation is
formulated at fixed nuclear configuration with no notion of a
parametric ``slow'' variable, and the mass ratio enters only when
the slow-sector kinetic operator $T_n = -\nabla_R^2 / (2 M_n)$ is
restored.\footnote{In mass-independent (Bohr-natural) atomic units
the projection still operates: it is the structural operation of
treating the nuclear configuration as a parameter rather than as a
dynamical variable, even when the mass ratio is not numerically
manifest in the units chosen.  See Paper~4 (archived) reduced-mass
convention discussion --- the adiabatic / Born--Oppenheimer
projection is what makes the electronic problem mass-independent
(at leading order) by separating the slow nuclear sector from the
fast electronic sector; the residual mass-dependence enters at the
diagonal Born--Oppenheimer correction (DBOC) and at non-adiabatic
couplings $P_{\nu\mu}$, both subleading in $m_e / M_n$.}\\
\textbf{Dimension introduced:} preserved at the projection step.  The
fast-sector effective potential $V_\text{eff}(\vec{R}) \equiv
E_\nu(\vec{R})$ has the same energy dimension $[E]$ as the original
Hamiltonian, and the new slow-sector variable $\vec{R}$ carries length
$[L]$ (already present in the projection chain via the stereographic
or molecule-frame hyperspherical projections at \S~\ref{sec:proj_stereo}
or \S~\ref{sec:proj_molframe}).  This projection does
not add a new dimension to the Layer-2 dictionary; it
\emph{reorganizes} the existing dimensional content into a slow~$\vec{R}$
$+$ fast~$\vec{r}$ split.\\
\textbf{Transcendental signature:} \textit{none at the projection step
itself}.  The Born--Oppenheimer scale-separation is parameter-only:
it indexes a family of Hamiltonians by a continuous variable without
injecting $\pi$ or any other transcendental.  Transcendentals appear
\emph{downstream} of this projection, in the solution of the
parameterized fast-sector eigenproblem
(\S~\ref{sec:proj_coupled_channel}): Level~3 hyperspherical $\mu(R)$
is
algebraic over $\mathbb{Q}(\pi, \sqrt{2})$ (CLAUDE.md~\S~12, Track~P1,
v2.0.12); Level~4 molecule-frame hyperspherical $\mu(\rho)$ is
piecewise-smooth in $\rho$, not algebraic over any closed ring
(Track~S, v2.0.13); Level~5 composed-geometry effective potentials
inherit the full transcendental content of every prior projection in
the chain.  The clean separation is: the BO projection introduces the
continuous parameter $\vec{R}$ with no transcendental cost; solving
the parameterized fast-sector eigenproblem then introduces
transcendentals according to the natural geometry of the fast sector
and the operator-order content of $H_\text{fast}(\vec{R})$.\\
\textbf{Used in:} every Level~4 molecular calculation (Paper~15;
H$_2$, HeH$^+$, every two-electron diatomic); every Level~3
hyperspherical calculation that uses the adiabatic effective potential
(Paper~13; He, Li, Be, all atomic species with one slow
hyperradius); every composed-geometry molecular calculation
(Paper~17; LiH, BeH$_2$, H$_2$O, all polyatomics, second-row
diatomics); the \emph{double}-adiabatic separation of Paper~15~\S~VII.C
where the core--valence fiber bundle introduces a third time-scale
(slowest nuclear / slow valence-hyperradius / fast valence-angular,
plus a separate slow core-hyperradius / fast core-angular on the
core fiber); all the ab initio molecular spectroscopy results of
Paper~13~\S~IX (rovibrational H$_2$); all $D_e$ benchmarks reported
across Papers~13, 15, 17; every cross-register / cross-block
calculation in the multi-focal-composition arc (May~2026) that
treats nuclear coordinates as classical parameters in
\texttt{cross\_register\_vne.py}.  This is the most heavily-used
projection in the framework's molecular pipeline; very few molecular
results do not pass through it.\\
\textbf{Empirical anchor:} H$_2$ at 96.0\% $D_e$ (Level~4
mol-frame hyperspherical + Schwartz cusp correction, Paper~15);
LiH $R_\text{eq}$ at 5.3\% error (Level~5 composed, Paper~17);
HeH$^+$ at 93.1\% $D_e$ (Paper~15~\S~VI.E); diagonal Born--Oppenheimer
correction $\Delta E_\text{DBOC} \approx +0.0015$~Ha for He (Paper~13
~\S~VI; explains the apparent variational violation $E = -2.9052$~Ha
vs.\ exact $-2.9037$~Ha at the single-channel adiabatic level by
exactly the magnitude expected from the kinetic energy cost of
angular wavefunction variation with $R$); BeH$^+$ as a transferability
test of the ab initio Phillips--Kleinman + adiabatic composition
(Paper~17).  Each of these is a quantitative test of the BO projection
in combination with the natural-geometry sector that follows it.\\
\textbf{Honest scope.}  The strict-adiabatic single-channel
approximation is the leading order in the small parameter $m_e / M_n$;
it breaks down where two electronic surfaces become near-degenerate
(conical intersections, autoionization channels, non-adiabatic
transitions) and where the kinetic energy of the slow variable is
comparable to electronic-level spacings.  The framework's response to
this breakdown is structural, not perturbative: the 2D variational
solver (Paper~13~\S~XIII, Paper~15~\S~VI.D) treats the slow and fast
variables on equal footing, eliminating the adiabatic approximation
entirely.  Track~AR (CLAUDE.md~\S~2, v2.0.23) showed that the 4-electron
LiH adiabatic solver's apparent $D_e$ violation was an artifact of
the single-channel approximation, not of the natural geometry: the 2D
variational solver gives $E_\text{min} = -7.79$~Ha (variational bound
respected, $D_e$ unbound) versus the adiabatic solver's $D_e = +0.49$~Ha
($5.3\times$ exact).  Track~BQ (CLAUDE.md~\S~2, v2.0.32) further
established that the $R_\text{eq}$ drift $+0.15$--$0.22$~bohr/$l_\text{max}$
in LiH composed geometry is \emph{not} from the adiabatic
approximation: both the adiabatic and 2D variational solvers produce
identical drift, isolating the residual to the composed-geometry PK
architecture rather than the BO projection itself.  The named extension
when adiabatic breaks down is therefore the coupled-channel /
2D-variational sibling at \S~\ref{sec:proj_coupled_channel}.\\
\textbf{Structural reading.}  Categorically distinct from
\S~\ref{sec:proj_restmass}.  The rest-mass projection is a
ring-preserving Bohr-like rescaling of a \emph{single} particle's
spectrum under $m_e \to m_\text{red}$; it modifies the foundational
energy scale (Ry) without introducing a new continuous parameter,
and the spectrum on the projected manifold is the same up to a
rational rescaling factor.  The Born--Oppenheimer projection is a
\emph{scale-separation between two coupled degrees of freedom}: it
takes a single Hilbert space $\mathcal{H}_n \otimes \mathcal{H}_e$ and
produces a parameter family $\{ \mathcal{H}_e \mid \vec{R} \in
\mathcal{M}_\text{nuc} \}$ together with a slow-sector Schr\"odinger
equation on the parameter manifold $\mathcal{M}_\text{nuc}$.  The
small parameter $m_e / M_n$ controls the projection's accuracy but is
not itself the new variable; the new variable is the slow coordinate
$\vec{R}$, which becomes a continuous index on a previously discrete
graph computation.  The two projections answer different structural
questions: rest-mass asks ``what happens to the spectrum when the
single-particle mass changes,'' while Born--Oppenheimer asks ``how
do two coupled degrees of freedom decouple at parametrically separated
time-scales.''
The two are independent and commute (rest-mass rescaling of the
foundational $m_e$ and BO-separation of slow nuclear from fast
electronic degrees of freedom can be applied in either order); they
co-appear in every reduced-mass molecular calculation, where
rest-mass rescales the electronic-sector $m_e$ to $m_\text{red}$ on
the fast side, and Born--Oppenheimer separates the slow nuclear from
the fast electronic sector across the full molecular Hamiltonian.

Structurally, this projection is the \emph{parent} of the
coupled-channel / adiabatic-curve projection at
\S~\ref{sec:proj_coupled_channel} (sibling, immediately downstream):
BO produces the parameterized electronic problem $H_\text{el}(\vec{R})$
with its parameter family of eigenproblems, and the coupled-channel
projection is the next step that solves the parameterized problem
and assembles the adiabatic potential curves
$U_\nu(\vec{R}) = E_\nu(\vec{R})$ as effective potentials for the
slow-sector dynamics, optionally including the non-adiabatic
couplings $P_{\nu\mu}(\vec{R}) = \langle \Phi_\nu | \partial / \partial
\vec{R} | \Phi_\mu \rangle$ that the strict-adiabatic single-channel
approximation drops.  The two projections are conceptually
inseparable in the molecular pipeline --- one does not invoke
Born--Oppenheimer in practice without then solving the parameterized
problem --- but they are taxonomically distinct: BO is the
\emph{scale-separation operation} (parameter-only, no transcendental
content, Layer-1 $\to$ Layer-2 transition), while coupled-channel /
adiabatic-curve is the \emph{solution operation} (introduces the
transcendental content specific to the fast-sector geometry, and
introduces the Berry-connection-like non-adiabatic coupling content
if the single-channel approximation is relaxed).

The double-adiabatic structure of Paper~15~\S~VII.C is a structural
\emph{composition} of this projection with itself at two distinct
time-scale ratios.  In LiH composed geometry the outer BO separates
nuclear (slow) from electronic (fast); within the electronic sector,
an inner BO separates core (slow on the electronic time-scale,
\emph{very} slow on the molecular time-scale) from valence (fast),
producing the core--valence fiber bundle of Paper~17~\S~II.  Iterated
composition of the BO projection is the structural mechanism by
which the framework reaches polyatomic molecules with multiple
distinct dynamical time-scales (Paper~17 LiH/BeH$_2$/H$_2$O composed
geometries; second- and third-row frozen-core diatomics, Sprint~7
v2.19.4).  The projection's own structural content does not change
under iteration; only the count of independent parameter manifolds
$\{\mathcal{M}_\text{nuc}, \mathcal{M}_\text{core}, \mathcal{M}_\text{val},
\ldots\}$ grows.

\subsection{Coupled-channel / adiabatic curve}
\label{sec:proj_coupled_channel}
\textbf{Source:} parameterized electronic Hamiltonian $H_\text{el}(R)$
inherited from the Born-Oppenheimer / adiabatic projection
(\S~\ref{sec:proj_adiabatic_BO}): a family of angular operators on
the Fock-projected $S^3$ graph, indexed by a continuous radial
coordinate ($\rho$ at Level~3 for the He hyperspherical lattice,
internuclear distance $R$ at Level~4 for the molecule-frame
hyperspherical lattice; same continuous coordinate at different
geometric levels per \S~\ref{sec:proj_molframe}).\\
\textbf{Target:} the discrete set of adiabatic potential curves
$\{\mu_\nu(R)\}_{\nu = 1, 2, \ldots}$ obtained by solving the angular
eigenvalue problem at each $R$, together with the non-adiabatic
coupling matrix $P_{\mu\nu}(R) = \langle \Phi_\mu(R) | \partial_R |
\Phi_\nu(R) \rangle$ that couples adjacent channels when the
parameterized angular wavefunctions $\Phi_\nu(R; \hat{\Omega})$ vary
with $R$.  At the matrix-pencil level
\[
  H_\text{ang}(R) \;=\; H_0 + R \cdot V^\text{coupling},
\]
with $H_0$ carrying the free-SO(6) (or SO(3N) at higher~$N$) Casimir
eigenvalues and $V^\text{coupling}$ encoding the nuclear and
electron--electron angular couplings (Paper~13 Eqs.~31--32, §XII).
The eigenvalues $\mu_\nu(R)$ are the adiabatic potential curves; the
coupled-channel radial solver then propagates each radial amplitude
$F_\nu(R)$ in the matrix-valued effective potential built from
$\{\mu_\nu(R), P_{\mu\nu}(R), Q_{\mu\nu}(R)\}$.\\
\textbf{Variables introduced:} the radial coordinate $R$ (or $\rho$)
itself, as the continuous-parameter axis of the channel family.  The
adiabatic curve index $\nu$ is a discrete combinatorial label inherited
from the SO(6) (or SO(3N)) representation theory of the parent angular
problem and is not a new physical variable.\\
\textbf{Dimension introduced:} length (one new $[L]$ scale per
geometric level; Level~3 introduces the hyperradius $\rho$, Level~4
the internuclear distance $R$).  Energy dimension is preserved:
$\mu_\nu(R)$ carries the same $[E]$ scaling as the parent angular
eigenvalues from which it is built, and the curve family is therefore
a one-parameter $[E]$-valued function family on the new $[L]$ axis.
This is the same $[L]$ injection logged at \S~\ref{sec:proj_molframe}
for the Level~4 case; at Level~3 the radial coordinate is the
hyperradius $\rho$ of the SO(6) hyperspherical lattice (Paper~13
§II), structurally analogous but distinct from internuclear $R$.\\
\textbf{Transcendental signature: algebraic-implicit at Level~3,
piecewise-smooth at Level~4 --- and the distinction is structural,
not numerical.}  At Level~3 (He hyperspherical, single-center,
two-electron) the angular Hamiltonian $H_\text{ang}(R)$ is a
\emph{linear matrix pencil} $H_0 + R \cdot V^\text{coupling}$, so its
eigenvalues satisfy the global characteristic polynomial
\[
  P(R, \mu) \;=\; \det\!\bigl(H_0 + R \cdot V^\text{coupling} - \mu I\bigr)
  \;=\; 0
\]
with coefficients in $\mathbb{Q}(\pi, \sqrt{2})[R, \mu]$ (Paper~13
§XII, Track~P1).  The $\pi$ content is inherited from the upstream
calibration of \S~\ref{sec:proj_fock} (normalisation on $[0, \pi/2]$);
the $\sqrt{2}$ content is inherited from the $V_{ee}$ Gaunt integrals
(\S~\ref{sec:proj_3j}); the radial parameter $R$ contributes no new
transcendental content of its own.  The function $\mu_\nu(R)$ is
therefore \emph{algebraic over $\mathbb{Q}(\pi, \sqrt{2})$} (an
algebraic-implicit function in the Paper~18 §IV.4 sense), of degree
$l_\text{max} + 1$ in both variables at angular truncation
$l_\text{max}$.  Branch points of $P(R, \mu) = 0$ occur at zeros of
the discriminant $\Delta(R) = \mathrm{disc}_\mu(P)$, where eigenvalue
sheets meet with square-root singularities --- pointwise numerical
diagonalisation at each $R$ is computational convenience, not
mathematical necessity (Paper~13 §XII.B, Track~P1, v2.0.12).
\smallskip

At Level~4 (mol-frame hyperspherical, two-center, two-electron) the
matrix-pencil form survives but the entries of $V^\text{coupling}$
themselves depend on $\rho$ piecewise via split-region Legendre
factors $\min(s, \rho) / \max(s, \rho)^{k+1}$ from the nuclear charge
function expansion (Paper~15 §V).  These factors are $C^0$ but not
$C^1$ at the region boundaries $s = \rho$, equivalently $\alpha =
\arccos\rho$ and $\alpha = \arcsin\rho$.  The pencil is therefore
\emph{piecewise-linear in $\rho$} rather than linear, and Track~S
(Paper~15 §VI.7, v2.0.13) established the irreducibility result: no
global $P(\rho, \mu) = 0$ exists across the region boundary, and no
algebraic-blow-up, Lie-algebra elevation, or $S^3 \times S^3$ hybrid
restores one.  The function $\mu_\nu(\rho)$ is \emph{piecewise-algebraic
within each region} (sub-eigenvalues on each side of $\rho = 1/2$
satisfy local $P_\text{left}(\rho, \mu) = 0$ and $P_\text{right}(\rho,
\mu) = 0$) but is not algebraic globally.  This is the
\emph{piecewise} tier of Paper~18 Table~II (Paper~18 \S~IV.4),
structurally distinct from the algebraic-implicit tier of Level~3.
\smallskip

The two transcendental regimes share the same operator-level shape
(linear matrix pencil) but differ in how the $R$-dependence enters
the pencil coefficients.  Level~3 admits a single global polynomial
because $V^\text{coupling}$ is $R$-independent; Level~4 does not
because $V^\text{coupling}(\rho)$ carries piecewise $\rho$-dependence
inherited from the two-center nuclear charge function.\\
\textbf{Used in:} Paper~13 (He hyperspherical, $l_\text{max}=2$--$7$
coupled-channel convergence to $0.022\%$ raw at $l_\text{max}=7$,
$0.004\%$ with Schwartz cusp correction at $l_\text{max}=4$ via the 2D
variational solver, Track~DI v2.6.0); Paper~15 (H$_2$ Level~4,
$96.0\%$ $D_e$ at $l_\text{max}=6$ $\sigma+\pi$ with 61 coupled
channels); Paper~17 (composed-geometry diatomics and polyatomics,
where the Level~3 and Level~4 coupled-channel solvers run on each
electron group's natural geometry block); Paper~18 \S~IV.4 (this
projection is where the algebraic-implicit and piecewise tiers live
at operator level); production solver modules
\texttt{geovac/hyperspherical\_adiabatic.py} (Level~3) and
\texttt{geovac/level4\_multichannel.py} (Level~4).\\
\textbf{Empirical anchor:} He hyperspherical at Level~3 reaches
$0.019\%$ error (Track~DI v2.6.0, Sprint~1; 2D variational solver on
$S^5$ with $l_\text{max}=7$, $n_R=35$, self-consistent cusp
correction; the underlying coupled-channel pipeline supplies the
adiabatic curves $\mu_\nu(R)$ that the variational ansatz interpolates
across).  H$_2$ at Level~4 reaches $96.0\%$ $D_e$ at $l_\text{max}=6$
$\sigma+\pi$ with $61$ coupled channels and Schwartz cusp correction
(Paper~15, Track~X v2.0.14).  Graph-native CI on the He pair-state
lattice at $n_\text{max}=9$ ($3{,}927$ configurations, exact algebraic
float integrals) reaches $0.20\%$ (Paper~13, Track~DI Sprint~3C) ---
this is the SAME observable computed by a different projection chain
(direct full-CI on the discrete graph, bypassing
\S~\ref{sec:proj_coupled_channel} entirely) and serves as an
independent transcendental-content cross-check: the residual matches
the small-$Z$ graph-validity-boundary artifact, not the Schwartz
cusp tail, demonstrating that the algebraic-implicit transcendental
content of $\mu_\nu(R)$ is captured cleanly by both the coupled-channel
and the graph-CI projection chains.\\
\textbf{Honest scope.}  This is non-adiabatic correction territory.
The bare adiabatic curve set $\{\mu_\nu(R)\}$ ignores the
$\partial_R$ acting on the parameterized angular wavefunctions
$\Phi_\nu(R; \hat{\Omega})$; capturing that derivative requires the
Hellmann--Feynman P-matrix $P_{\mu\nu}(R) = \langle \Phi_\mu(R) |
V^\text{coupling} | \Phi_\nu(R) \rangle / (\mu_\mu(R) - \mu_\nu(R))$
and the closure-approximated $Q_{\mu\nu}(R) \approx \sum_\kappa
P_{\mu\kappa} P_{\kappa\nu}$ in the coupled radial equations (Paper~13
§XIII.A, Eqs.~70--71).  Both quantities are exact at every $R$ in the
algebraic spectral basis (no finite-difference noise) but the closure
sum runs over an in-principle-infinite set of intermediate channels
$\kappa$; in practice it is truncated together with the included
channel set.

The framework's 2D variational solver
(\texttt{geovac/level3\_variational.py} for He, the
\texttt{level4\_method=\textquotesingle variational\_2d\textquotesingle}
path for Level~4; Paper~15 §VI.4) bypasses adiabatic factorisation
entirely by treating $R$ and the angular coordinate simultaneously in
a single tensor-product basis.  This is structurally a different
projection altogether: it removes the Born-Oppenheimer
\S~\ref{sec:proj_adiabatic_BO} step, so the adiabatic curve
family of the present projection never appears.  It comes at the cost
of giving up the channel-by-channel decomposition that makes the
adiabatic picture chemically interpretable.  Both routes are
legitimate; the choice between them is a Layer-2 control of which
projections the chain commits to.

Algebraic registry status (CLAUDE.md §12): hyperradial overlap and
kinetic matrices $S, K$ are fully algebraic at Levels~3 and~4 via the
three-term Laguerre recurrence (Track~H, v2.0.10; pentadiagonal $M_2$
moment matrix for $S$, tridiagonal derivative kernel for $K$,
machine-precision agreement with quadrature).  The effective radial
potential matrix elements $V^\text{eff}(R)$ are algebraic at
$l_\text{max} = 0$ (single Stieltjes moment seed $e^a \cdot E_1(a)$;
Track~P2, v2.0.13) and numerical-required at $l_\text{max} \geq 1$
because of a radical obstruction --- the discriminant $\sqrt{\Delta(R)}$
enters at $l_\text{max}=1$ and Cardano-type irrationals at
$l_\text{max}=2$.  Pointwise diagonalisation at each $R$ to extract
$\{\mu_\nu(R)\}$ is the residual numerical step that survives all the
upstream algebraic optimisation, and it is irreducible in the sense
that the algebraic curve $P(R, \mu) = 0$ describes the eigenvalues
implicitly without admitting a closed-form root extraction.  Spectral
Laguerre at Level~2 is fully algebraic (one transcendental seed
$e^a \cdot E_1(a)$, Track~J v2.0.10) and provides the cleanest
reference case --- the moment the radial coordinate becomes the
\emph{family index} of an eigenvalue problem rather than the argument
of a single function, the framework crosses from purely algebraic
into algebraic-implicit or piecewise-smooth content.\\
\textbf{Structural reading.}  Categorically distinct from its parent
\S~\ref{sec:proj_adiabatic_BO}: the Born-Oppenheimer projection
produces the parameterised Hamiltonian $H_\text{el}(R)$ and commits to
the slow/fast separation as a structural assumption, but it does not
itself solve the eigenvalue problem at each $R$.  The present
projection makes that commitment: it converts a family of operators
into a family of channel functions plus a non-adiabatic coupling
matrix.

This is the projection that names the \emph{transcendental boundary}
of the framework's algebraic structure.  Levels~1 and~2 (the latter
in $\sigma$ states, where the algebraic Laguerre matrix elements give
a single transcendental seed and no others) are fully algebraic in
the strong sense.  Higher levels are not, and the failure mode is
this projection: the moment one indexes an eigenvalue problem by a
continuous radial coordinate, the eigenvalue itself becomes an
algebraic-implicit function (Level~3) or a piecewise-smooth function
(Level~4) of that coordinate.  The Level~3 case is mild --- the
implicit function admits a global polynomial description and shares
its transcendental content ring $\mathbb{Q}(\pi, \sqrt{2})$ with the
upstream calibration and embedding constants, introducing no new
transcendentals at the $R$-dependence step.  The Level~4 case is
stronger: the piecewise structure is irreducible (Track~S) and the
function does not admit a global polynomial description, although it
remains piecewise-algebraic within each region.

The distinction matters for the broader projection taxonomy because
it splits what Paper~34 \S~\ref{sec:proj_molframe} calls
``piecewise-smooth in $R$ at the adiabatic potential level'' into two
different transcendental classes.  Level~3 is algebraic-implicit and
satisfies $P(R, \mu) = 0$ globally; Level~4 is piecewise-smooth and
does not.  Paper~18 §IV.4 already separates these tiers in its
exchange-constant taxonomy (the ``alg.\ (implicit)'' row for Level~3
hyperspherical, the ``piecewise'' row for Level~4 mol-frame); the
present projection is where that taxonomic split lives at the
projection-chain level.  The Table~\ref{tab:projection_axes} entry for
this projection (algebraic-implicit at Level~3, piecewise-smooth at
Level~4) makes the dual notation explicit; the operator-level
distinction in this subsection is unambiguous regardless of how the
table is summarised.

Sibling to \S~\ref{sec:proj_spectral_action} in the structural sense
that both rely on a continuous parameter ($\Lambda$ for the spectral
action, $R$ for the adiabatic curve) to label a one-parameter family
of finite-rank spectral problems; categorically distinct from
\S~\ref{sec:proj_restmass} (which introduces $m^2 \in \mathbb{Q}$ as a
ring-preserving scalar shift on $D^2$, not as a family index of an
eigenvalue problem) and from \S~\ref{sec:proj_observation} (which
introduces $\beta = 1/T$ as a compactification radius rather than as
a Hamiltonian parameter).  The genuine peer projection is
\S~\ref{sec:proj_molframe}, which provides the radial coordinate that
the present projection then converts into the family index of the
adiabatic curve set.

\subsection{Gauge choice (Coulomb / Lorenz / Feynman)}
\label{sec:proj_gauge_choice}
\textbf{Source:} the framework's $U(1)$ connection on the Fock-projected
$S^3$ Hopf graph (with $\mathrm{SU}(2)$ and $\mathrm{SU}(3)$ siblings
on the same graph and on $S^5$ respectively, per Paper~30 and
Sprint~ST-SU3); equivalently, the photon $1$-cochain in the Hodge
decomposition of edge data, prior to any specific gauge representative
being selected for propagator computations.\\
\textbf{Target:} the same connection with a specific gauge
representative fixed.  Three canonical choices are in play:
\begin{itemize}\setlength{\itemsep}{2pt}
\item \emph{Coulomb gauge} ($\vec{\nabla} \cdot \vec{A} = 0$,
equivalently $d^\dagger A = 0$ on the graph): the photon decomposes
into an instantaneous longitudinal Coulomb part and a transverse
radiative part; the framework's vector-photon construction
(Paper~28~\S vector\_photon\_qed, Paper~33~\S~VI) operates in this
gauge by treating the Hopf-base $1$-form sector as the propagating
photon and computing $G_\gamma = L_1^+$ as the transverse propagator
(see also \S~\ref{sec:proj_wilson}, where the weak-coupling expansion
$L_1 = B^\top B$ is a Coulomb-class gauge fixing of the SU(2) Wilson
construction).
\item \emph{Lorenz gauge} ($\partial_\mu A^\mu = 0$, equivalently the
covariant retarded condition on the spacetime $4$-potential): the
photon propagator carries explicit retardation in time and the
transverse and longitudinal pieces remain coupled.
\item \emph{Feynman gauge} (the propagator $-i g_{\mu\nu}/k^2$): the
simplest covariant choice, used in Paper~28~\S~6 for the one-loop
electron self-energy.
\end{itemize}
The projection step is the act of fixing one of these representatives
\emph{before} intermediate Green's-function, propagator, or matrix-element
computations begin.\\
\textbf{Variables introduced:} \emph{none}.  A gauge choice is a
selection among equivalent classes of representatives, not the
introduction of a continuous parameter.  No physical scale,
focal-length, mass, or coordinate is injected by the gauge fixing
itself.  This places gauge choice alongside \S~\ref{sec:proj_vector_photon}
in the (no-variable) row of Table~\ref{tab:projection_axes}.\\
\textbf{Dimension introduced:} dimension-preserving.  The gauge choice
operates on the existing connection without injecting a new $[L]$,
$[E]$, $[M]$, or $[T]$ scale.  All three gauges produce the same
physical observable with the same physical dimension; what differs is
the intermediate-step organisation of the propagator content.\\
\textbf{Transcendental signature:} \emph{gauge-dependent at
intermediate steps; gauge-invariant on physical observables.}  This is
the distinguishing structural property of the gauge-choice projection.
At the intermediate-propagator level, the three gauges produce
categorically different transcendental content:
\begin{itemize}\setlength{\itemsep}{2pt}
\item In Coulomb gauge --- the framework's implicit choice throughout
Papers~28, 30, and 33 --- the vector-photon construction's per-loop
factor is $1/(4\pi) = \mathrm{Vol}(S^2)/(4 \cdot 4\pi^2)$, identified
in Paper~33~\S~VI as the $S^2$ Weyl exchange constant of the Hopf
base (the same Hopf-base measure $\pi = \mathrm{Vol}(S^2)/4$ that
appears in the M1 master-Mellin sub-mechanism, per Paper~32~\S~VIII
and Paper~18~\S~III.7).  The intermediate transcendental content sits
cleanly in the calibration $\pi$-tier of Paper~18.
\item In Lorenz gauge, the retarded propagator structure would produce
different intermediate angular content, with $\pi$ entering through
the retardation kernel rather than through the static Hopf-base solid
angle.  This route has not been implemented in the framework.
\item In Feynman gauge, the explicit $-i g_{\mu\nu}/k^2$ propagator
gives the textbook covariant intermediate transcendental content, with
$\pi$ appearing through the standard Feynman-parameter denominators
rather than through a Hopf-base normalisation.  Paper~28~\S~6 names
this gauge for the one-loop self-energy formula but the production
machinery operates in Coulomb gauge throughout.
\end{itemize}
On physical observables (cross-sections, transition rates, energy
shifts), all three gauges produce identical results --- this is the
content of the Ward identity (Rule~6 of Paper~33's eight QED selection
rules; see Paper~33~\S~II for the framework's status), and it is
gauge-invariance in its standard QED sense.\\
\textbf{Used in:} the framework's one-loop QED machinery (Paper~28)
implicitly assumes Coulomb gauge throughout the $1$-cochain Hodge
construction; Paper~33's $1{+}6{+}1$ selection-rule partition is
stated in Coulomb gauge; Paper~30's SU(2) Wilson lattice gauge theory
identifies $L_1 = B^\top B$ as the weak-coupling kinetic term via a
Coulomb-class gauge fixing of the link variables.  Feynman gauge is
named at one point in Paper~28~\S~6 (the textbook formula for the
one-loop electron self-energy) but is not used in production
calculations.  No production computation has been performed in Lorenz
gauge.\\
\textbf{Empirical anchor:} the Ward identity itself, i.e.\
gauge-invariance of physical observables.  Standard QED requires
$k_\mu \mathcal{M}^\mu = 0$ for the photon-emission amplitude
$\mathcal{M}^\mu$, and equivalently the equality of physical
predictions across gauge choices.  This is item~R6 in Paper~33's
selection-rule census, where the framework's status is recorded
honestly: ``gauge invariance under the $U(1)$ phase of the photon,
which the $1$-cochain photon has only in temporal gauge and only
approximately'' (Paper~33~\S~II, R6).  The Ward identity is therefore
the standard external anchor for the gauge-choice projection;
gauge-invariance of physical observables is the precise condition
under which the projection itself is structurally well-posed.\\
\textbf{Honest scope.}  This is the single most under-tested entry in
the projection dictionary, and the entry must be tagged as such.

The framework has \emph{not} been tested across gauges.  Every
production calculation in Papers~28, 30, 33, and the Sprint~HF /
Sprint~LS-1..LS-7 / Lamb-shift sub-percent closure of Paper~36 runs
implicitly in Coulomb gauge: the photon is represented as a Hodge
$1$-cochain on the GeoVac edge complex, the propagator is
$G_\gamma = L_1^+$ (the transverse pseudoinverse of the edge
Laplacian), and the per-loop calibration constant $1/(4\pi)$ is
identified through the Hopf-base $S^2$ solid-angle normalisation.
This is a clean, structurally well-motivated choice ---  Coulomb gauge
respects the framework's natural angular-momentum decomposition on
the Fock-projected $S^3$ graph, and the transverse photon sector
sits naturally in the co-exact subspace of the Hodge decomposition
of \S~\ref{sec:proj_wilson} (compare Paper~28~\S~transverse\_photon).
But it \emph{is} a choice.

A genuine validation of the gauge-choice projection at framework level
would require computing a sample one-loop observable in Lorenz gauge
(or Feynman gauge as implemented at full operator level rather than
named for one formula), verifying that the physical result agrees
with the Coulomb-gauge calculation to the precision of the framework,
and confirming that the per-loop calibration constant changes
structurally between the two gauges while the physical observable is
gauge-invariant.  This has not been done.  The Ward identity (R6 in
Paper~33's census) holds only approximately on the $1$-cochain photon
even in temporal gauge, and the framework has not closed the
gauge-invariance test against a competing gauge.

Sub-leading observation: the choice of Coulomb gauge is not arbitrary
on the GeoVac graph.  The Hodge decomposition of the edge complex
distinguishes exact ($d_0$-image), co-exact ($d_1^\top$-image), and
harmonic sectors, and Coulomb gauge is precisely the projection onto
the co-exact (transverse) sector.  The framework's natural data
structure --- $L_0$ on nodes, $L_1 = B^\top B$ on edges,
plaquette Hodge structure from Paper~25 / Paper~30 --- comes equipped
with this decomposition.  Lorenz gauge, by contrast, would require
introducing the time direction explicitly (compactified per
\S~\ref{sec:proj_observation} or extended per the Klein-Gordon
construction of Paper~35), and the framework has not been wired for
covariant time-dependent gauge fixing.  This is a real structural
constraint, not a choice that could be flipped without modifying the
upstream construction.\\
\textbf{Structural reading.}  Layer~$2 \to$ Layer~$2$.  The gauge-choice
projection operates on the projected vector-photon Hilbert space
(\S~\ref{sec:proj_vector_photon}); it does not act on Layer~1 graph
data and does not move the chain across the
\S~\ref{sec:layer2_boundary} boundary.

This is the \emph{second independence witness} in the
(no-variable, no-dimension, transcendental-only) corner of the
variable/dimension/transcendental axis grid.  Until Sprint~2 the only
witness in that corner was \S~\ref{sec:proj_vector_photon} (vector-photon
promotion, which introduces $1/(4\pi)$ without any new variable or
dimension); the gauge-choice projection is its structural sibling.
Both witnesses operate on the same object (the framework's $U(1)$
connection) at successive structural steps: vector-photon promotion
transforms the scalar $1$-cochain into a vector harmonic with
$(q, m_q)$ labels, and the present projection fixes a specific gauge
representative of the resulting connection before propagators are
computed.  The fact that the two witnesses are \emph{both
gauge-related} sharpens the V/D-correlation observation of
\S~\ref{sec:dictionary} and \S~\ref{sec:axiomatization}: the
transcendental-only corner of the axis grid is the gauge sector of the
framework.  It is not an accident that the only two projections
preserving both V and D while moving the transcendental axis are
gauge-construction steps; the structural content of that corner is
``add transcendental content to the gauge sector without adding any
physical variable or dimension.''

Categorically distinct from \S~\ref{sec:proj_vector_photon} despite
the shared axis pattern: vector-photon promotion is the structural
\emph{construction} of the gauge sector (it builds the vector-harmonic
photon from a scalar $1$-cochain), while the gauge-choice projection
is the \emph{selection} of a representative within the already-built
sector.  The two compose: the chain $\textsf{Layer~1 graph} \to
\S~\ref{sec:proj_vector_photon} \to \S~\ref{sec:proj_gauge_choice}$ is
how the framework reaches a computable QED propagator from bare graph
data.  Categorically distinct also from
\S~\ref{sec:proj_multipole_gaunt}, which similarly preserves V and D
(the multipole index $L$ is an internal integer summation label, not
a physical variable; the kernel reorganisation is dimension-preserving)
but lives in the kernel-side angular sector rather than the gauge
sector and is ring-preserving over $\mathbb{Q}[\sqrt{2k{+}1}]_k$ rather
than gauge-dependent in transcendental content.  The multipole / Gaunt
projection is an angular-content-only structural projection on the
Coulomb kernel; the present projection is a gauge-content-only
structural projection on the connection that the Coulomb kernel
contracts against.  Both are necessary to reach a computable
amplitude; both preserve the V/D axes; neither introduces a Layer-2
input.

Sibling in the calibration tier of Paper~18: the per-loop $1/(4\pi)$
that appears in Coulomb gauge is the same calibration $\pi$-content
identified by Paper~33~\S~VI as the $S^2$ Weyl exchange constant.
Under the master-Mellin reading of Paper~32~\S~VIII and
Paper~18~\S~III.7, this calibration $\pi$ is the M1 sub-mechanism
signature ($\mathrm{Vol}(S^2)/4$, the Hopf-base measure).  A
Lorenz-gauge calculation, if performed, would produce different
intermediate $\pi$-content (retardation kernels rather than static
Hopf-base normalisations) but would necessarily land on the same
physical observable.  The gauge-choice projection therefore sits at
the interface of the framework's gauge construction
(\S~\ref{sec:proj_vector_photon} +
\S~\ref{sec:proj_wilson}) and the framework's transcendental taxonomy
(Paper~18), and it is the structural step at which a Layer-2 chain
commits to a specific way of accumulating $\pi$ at each loop.

Open question (flagged for a future sprint): a serious test of the
gauge-choice projection would compute a sample observable (e.g.\ the
one-loop Lamb-shift Bethe-log component of Paper~36, or the
hydrogen $21$\,cm contact term of \S~\ref{sec:autopsy_21cm}) in
Feynman gauge as well as the existing Coulomb-gauge implementation,
and verify gauge invariance to the precision of the framework
($\sim10^{-3}$ on Lamb-shift-class observables, $\sim10^{-4}$ on
hyperfine contact terms).  If gauge invariance fails at framework
precision, the failure localises a specific Layer-2 calibration
content that the Coulomb-gauge construction is silently encoding;
if it succeeds, the gauge-choice projection is verified as a
genuine independence witness in the sense of
Observation~\ref{obs:axes_correlated} rather than a tautological
relabelling.

\subsection{Wick rotation / signature change}
\label{sec:proj_wick_rotation}
\textbf{Source:} the framework's bare spectral-triple machinery on
the Fock-projected $S^3$ graph (or any natural-geometry Euclidean
spectral triple, including the $T_{S^3} \otimes T_{S^1_\beta}$
tensor-product construction of Sprint TD Track~1).  By construction
every operator in the GeoVac inventory ---
Camporesi--Higuchi Dirac on round $S^3$, Bargmann--Segal heat
kernel on $S^5$, Wilson plaquette functionals on Hopf graphs, all
$\zeta_R(2k)$ values entering through $\zeta_{D^2}$ on closed
Riemannian manifolds --- is computed in Euclidean (Riemannian)
signature.\\
\textbf{Target:} the Lorentzian (Minkowski-signature) physical
observable obtained by analytic continuation
$\tau \to i t$, $\beta \to i \beta$ on the imaginary-time circle
(or proper-time circle in the Rindler case).  Concrete instances
under the Wick-rotation chain
[Hartle--Hawking 1976]~$\to$~[Sewell 1982]~$\to$~[Bisognano--Wichmann
1976]~$\to$~[Unruh 1976]:
the Hartle--Hawking KMS state at
$\beta_{\mathrm{cigar}} = 8\pi M$ on the Euclidean Schwarzschild
cigar continues to the Hawking thermal state on the Lorentzian
Schwarzschild exterior at $T_H = 1/(8\pi M)$; the periodic
Euclidean Rindler coordinate $\eta_E$ with period $2\pi$ continues
to the Lorentzian Rindler boost at acceleration $a$, with the
accelerated-observer Unruh temperature $T_U = a/(2\pi)$; the modular
automorphism of the algebra of observables localized in a Lorentzian
Rindler wedge, with respect to the Minkowski vacuum, is identified
by Bisognano--Wichmann with the unitary representation of the Lorentz
boost subgroup preserving the wedge~\cite{tener2025_bw_nonunitary},
and the analytic period of the
modular flow is the same $2\pi = \mathrm{Vol}(S^1)$ that appears on
the Euclidean side.\\
\textbf{Variables introduced:} the signature index --- a discrete
binary choice between Riemannian Euclidean and Lorentzian
Minkowskian.  Structurally analogous to the gauge-choice projection
(\S~\ref{sec:proj_gauge_choice}): a discrete bookkeeping selection
that fixes how the framework's spectral content is read as a physical
observable.  Continuous rapidity parameters relating two Lorentzian
frames at relative velocity $v/c$ are \emph{not} introduced by this
projection; a full Lorentz-boost projection at the operator-system
level (acting on all particles simultaneously, relating two truncated
spectral triples at different observer rapidities) is a separate
structural extension named in \S~\ref{sec:open}.\\
\textbf{Dimension introduced:} dimension-preserving.  Analytic
continuation does not introduce a new $[L]$, $[E]$, $[M]$, or $[T]$
scale; it changes the interpretation of an already-dimensioned
quantity from a Euclidean (imaginary-time-periodic) to a Lorentzian
(real-time-thermal or wedge-modular) reading.  In the V/D/P tagging
of \S~\ref{sec:dictionary}, the projection sits in the same
dimension-preserving class as Sturmian reparameterization
(\S~\ref{sec:proj_sturmian}) and the gauge-choice projection
(\S~\ref{sec:proj_gauge_choice}).\\
\textbf{Transcendental signature:} $2\pi \cdot \mathbb{Q}$ via
$\mathrm{Vol}(S^1)$ on the Euclidean time circle (Hopf-base measure /
M1 sub-mechanism of the master Mellin engine; see Paper~32~\S~VIII
Remark on the Bisognano--Wichmann reading and Paper~35~\S~VIII for
the framework-level structural correspondence).  Critically:\ this
is the \emph{same generator} as the observation / temporal-window
projection (\S~\ref{sec:proj_observation}), which injects
$2\pi \cdot \mathbb{Q}$ per Matsubara mode via compactification of
Euclidean time at finite $\beta = 1/T$.  The two projections share
the M1 transcendental signature but are \emph{structurally distinct}
as mechanisms: observation / temporal-window compactifies a single
Euclidean time direction at finite period $\beta$ and reads out
finite-temperature observables (Casimir at $T \neq 0$,
Stefan--Boltzmann free energy, Matsubara loops), whereas Wick
rotation analytically continues between two signatures and reads out
Lorentzian observables (Hawking radiation, Unruh thermal bath,
Bisognano--Wichmann modular flow) from the same Euclidean data.
Both surface the M1 $2\pi$ on the same $S^1$ factor;
the structural reading of what that $2\pi$ \emph{means} differs:\
a Hopf-bundle measure factor on the Riemannian side; a
modular-flow / Lorentz-boost orbit period on the Lorentzian side
via the published Wick-rotation prescription.\\
\textbf{Used in:} the Hawking-temperature off-precision catalogue
row in \S~\ref{sec:offprecision} (error class C, $M$ external
Einstein-equations input, machinery-witness only); Sprint TD Track~4's
reproduction of $T_H = 1/(8\pi M)$ on the Euclidean Schwarzschild
cigar from M1 alone; Sprint Unruh-pendant's reproduction of
$T_U = a/(2\pi)$ on the Euclidean Rindler wedge from the same M1
mechanism; the structural-correspondence reading of
$\beta = 2\pi / (\text{surface gravity})$ that unifies Hawking,
Sewell, Bisognano--Wichmann, and Unruh as four faces of a single
Wick-rotation theorem.  Any future framework-side observable whose
physical interpretation requires Lorentzian signature (Wightman-class
expectation value of a real-time-ordered operator product, KMS state
on a Lorentzian wedge algebra, real-time Schwinger--Keldysh closed
time contour) consumes this projection.\\
\textbf{Empirical anchor:} the four-witness Wick-rotation theorem,
each face reproducing a shared $2\pi = \mathrm{Vol}(S^1)$ from the
M1 mechanism:
(i)~Hawking temperature $T_H = 1/(8\pi M) = \kappa_g / (2\pi)$ on the
Euclidean Schwarzschild cigar, with surface gravity
$\kappa_g = 1/(4M)$~\cite{hartle_hawking1976};
Sprint~TD Track~4 reproduces this from M1 with sympy residual zero
on all three identities $\beta = 8\pi M$, $T_H = \kappa_g/(2\pi)$,
$\beta = 4M \cdot 2\pi$;
(ii)~Sewell's identification of the Hartle--Hawking thermal state
restricted to the Schwarzschild exterior as a KMS state at
$\beta = 8\pi M$ on a globally hyperbolic Lorentzian manifold with
a bifurcate Killing horizon~\cite{sewell1982};
(iii)~Bisognano--Wichmann identification of the modular automorphism
of the algebra of observables localized in a Rindler wedge, with
respect to the vacuum state, with the unitary representation of the
Lorentz boost subgroup preserving the
wedge~\cite{bisognano_wichmann1976}, with analytic period
$\sigma_{i \cdot 2\pi} = $ identity forced by Tomita--Takesaki KMS;
(iv)~Unruh temperature $T_U = a / (2\pi)$ for a uniformly accelerated
observer at proper acceleration $a$, derived in the same Wick-rotation
construction with $\beta_{\mathrm{Rindler}} = 2\pi / a$ on the
Euclidean Rindler $\eta_E$-circle~\cite{unruh1976}.  Sprint
Unruh-pendant reproduces this verbatim with the same M1 mechanism,
sympy residual zero, lowest bosonic Matsubara mode $= a$, lowest
fermionic mode $= a/2 = \pi T_U$.
All four faces reproduce the shared $2\pi$ from the framework's
M1 Hopf-base measure on the same $S^1$ factor; the published
Wick-rotation prescription supplies the bridge to the Lorentzian-side
modular-flow / boost-orbit vocabulary in each case.\\
\textbf{Honest scope.}  The framework demonstrably reproduces
$\beta = 2\pi / \kappa_g$ on all four Wick-rotation faces from M1
alone, and the published Hartle--Hawking / Sewell / Bisognano--Wichmann
/ Unruh chain identifies that $2\pi$ with the Lorentzian-side
modular-flow / Lorentz-boost orbit period.

\textbf{Sprint L1 closure (2026-05-16):\ literal identification at
operator-system level (Riemannian).}\ The L1-named falsifier
(``compute the modular Hamiltonian $K$ on $\mathcal{T}_{n_{\max}}$
for a wedge-restricted Camporesi--Higuchi state, verify
$\sigma_{2\pi} = $ identity at finite $n_{\max}$'') is now closed
positively at the strongest level.  Sprint L1
(\texttt{debug/l1\_modular\_hamiltonian\_results\_memo.md},
Paper~32~\S~VIII.F) constructs $K$ on the truncated CH spectral triple
via the hemispheric-wedge BW-$\alpha$ geometric realization
($K = J_{\mathrm{polar}}$ with integer spectrum $\mathrm{two\_m}_j$
on the full-Dirac basis) and verifies the period closure
$\sigma_{2\pi}(O) = O$ bit-exactly at $n_{\max} \in \{2, 3, 4, 5\}$
for all four physical witnesses (BW, HH, Sew, Unruh).  Maximum
periodicity residual: $1.8 \times 10^{-16}$ at $n_{\max}=2$ through
$1.6 \times 10^{-15}$ at $n_{\max}=5$ (machine precision scaling
$O(\sqrt{\dim_H} \cdot \varepsilon_{\mathrm{machine}})$).
Cross-witness collapse is bit-exact (max consistency residual
$= 0.0$ at every $n_{\max}$):\ the period closure is independent
of $\kappa_g$ because $\mathrm{e}^{\mathrm{i}\cdot 2\pi \cdot n} = 1$
for integer spectrum.  This lifts the four-witness Wick-rotation
theorem from ``structural correspondence'' to ``literal identification
at the operator-system level (Riemannian)'' at every finite cutoff.
The Lorentzian extension to signature $(3, 1)$ via the
Bizi--Brouder--Besnard~2018 Krein-space spectral triple remains the
named Sprint L2 follow-up (a substantial open program per the
Sprint L0 audit).

\textbf{Sprint L1-tighten (2026-05-16 evening):\ unified BW-$\alpha$/BW-$\gamma$
closure.}\ The L1 closure above is at the geometric BW-$\alpha$ level
(K = $J_{\mathrm{polar}}$ with integer spectrum). Sprint L1-tighten
(\texttt{debug/l1\_tighten\_tomita\_results\_memo.md}, Paper~32~\S~VIII.F
extension) extends this to the full Tomita-Takesaki BW-$\gamma$
construction:\ build $K_{\mathrm{TT}} = -\log \Delta$ from the polar
decomposition $S = J_{\mathrm{TT}} \Delta^{1/2}$ on the GNS
Hilbert-Schmidt space $H_{\mathrm{GNS}} = M_{\dim_W}(\mathbb{C})$,
verify $\sigma_{2\pi}^{\mathrm{TT}}(O) = O$ bit-exactly at the
operator-action level for all four witnesses at $n_{\max} \in \{2, 3,
4, 5\}$. Verdict:\ \emph{UNIFIED\_STRONG closure}. Maximum Tomita-level
residual: $1.09\times 10^{-16}$ ($n_{\max}=2$) through $4.79\times
10^{-15}$ ($n_{\max}=5$).  The BW-$\alpha$ and BW-$\gamma$ flows are
operator-action-level conjugates of each other ($\sigma_t^{\mathrm{TT}}
= \sigma_{-t}^{\alpha}$ bit-exact at general $t$;\ identical at the
period $t = 2\pi$).  The state-dependent $J_{\mathrm{TT}}(X) =
\rho^{1/2} X^* \rho^{-1/2}$ verifies $J_{\mathrm{TT}}^2 = +I$ at the
full Tomita level (not just the canonical $U = I$ sanity case).  The
honest scope above is therefore tightened:\ literal identification at
the operator-system level (Riemannian) holds via \emph{both} the
geometric BW-$\alpha$ \emph{and} the Tomita-Takesaki BW-$\gamma$
constructions, not just one.  The L1 ``BW-$\alpha$-only'' caveat is
removed.

Two further scope distinctions are required.
First, \emph{Wick rotation is not the full Lorentzian extension.}
This projection lifts the framework's metric-functional Riemannian
output (a Mellin trace, a heat-kernel coefficient, a Wilson
expectation value) to its Lorentzian analytic continuation in
cases where the bridge is a single Wick rotation of a temporal
circle.  It does \emph{not} construct a Lorentzian (Minkowski-signature)
spectral triple with Krein-space Hilbert space and fundamental
symmetry $\mathcal{J}^2 = I$; that extension would carry the
framework's signature index from $(m,n) = (3, 0)$ in the
Bizi--Brouder--Besnard~2018~\cite{bizi_brouder_besnard2018}
classification to $(3, 1)$.  A clean Lorentz-boost projection at
the metric / curvature level, acting on \emph{all} particles in the
system simultaneously and relating two truncated metric spectral
triples at different observer rapidities, requires that full
Lorentzian-signature spectral triple plus a Lorentzian analog of
the Latr\'emoli\`ere propinquity used in Paper~38~\cite{paper38}
for the WH1 PROVEN result.  The propinquity literature is purely
Riemannian as of May~2026 (Latr\'emoli\`ere, Toyota,
Hekkelman--McDonald, Leimbach--van~Suijlekom); the parallel synthetic
Lorentzian-geometry program (optimal transport in Lorentzian synthetic
spaces, timelike Ricci curvature lower
bounds~\cite{cavalletti_mondino2024}) provides a metric-geometric
analog on the commutative side but does not yet supply an operator-
algebraic propinquity.  Constructing a Lorentzian propinquity is
original NCG-mathematics work, a substantial open program beyond
sprint scale.  That open extension is flagged in
\S~\ref{sec:open} and is the structural deepening of this
\S~\ref{sec:proj_wick_rotation} entry, not a separate candidate.

Second, \emph{distinct from existing $\gamma$-extensions in the
catalogue.}  The Dirac-sector $\gamma = \sqrt{1 - (Z\alpha)^2}$ in
the Camporesi--Higuchi spinor lift (\S~\ref{sec:proj_spinor}) is a
bound-state radial Lorentz factor for a single-particle hydrogenic
wavefunction in a central potential; it is not a frame
transformation between two observers, and the Wick rotation
discussed here does not introduce or modify it.  The rest-mass
projection (\S~\ref{sec:proj_restmass}) rescales a single particle's
mass under $m_e \to m_\text{red}$ within Euclidean signature; the
Breit retardation projection
(\S~\ref{sec:proj_breit_retardation}) introduces the rest-mass
ratio $m_l/m_n$ as a two-body kinematic control parameter.  None of
these is a signature change.  Wick rotation sits structurally above
all three:\ those projections operate within a fixed Euclidean
spectral triple; this projection toggles between the Euclidean and
Lorentzian readings of the \emph{same} computed quantity, with the
toggle bridged by the published Wick-rotation chain.\\
\textbf{Structural reading.}  Wick rotation sits in the
Layer-2~$\to$~Layer-2 class (cf.\
\S~\ref{sec:layer2_boundary}):\ it operates on already-projected
Euclidean spectral content, not on the bare graph or its initial
projection chain.  Structurally it is a
\emph{metric-functional-level} projection:\ the framework computes a
Mellin trace, heat-kernel coefficient, or other Euclidean
spectral-action functional (Paper~32~\S~VIII case-exhaustion theorem
ensures the transcendental content is M1, M2, or M3); Wick rotation
then reads that functional as a Lorentzian physical observable via
the published Hartle--Hawking / Sewell / Bisognano--Wichmann / Unruh
chain.

The structural-correspondence-not-literal-identification verdict
crystallizes the framework's position on Lorentzian physics in
general:\ at the metric functional level (specifically the M1
$2\pi = \mathrm{Vol}(S^1)$ signature on a Wick-rotated time circle),
the framework's Euclidean output IS the Lorentzian modular-flow /
boost-orbit period, with no extension required.  At the
operator-system level (Tomita--Takesaki modular automorphism acting
on the truncated algebra at finite $n_{\max}$, Krein-space
representation, frame-level Lorentz-boost between truncated metric
spectral triples at different rapidities), the framework currently
sees the same $2\pi$ but cannot internally close the period
identity; that closure is the named open extension.

Three observations distinguish this projection from
\S~\ref{sec:proj_observation} despite shared M1 generator and
shared $2\pi$ transcendental content.
First, mechanism:\ observation / temporal-window compactifies a
single Euclidean time direction at finite period $\beta = 1/T$ and
sums Matsubara modes $\omega_k^t = 2\pi k / \beta$; Wick rotation
analytically continues a Euclidean computation to Lorentzian
signature without changing the periodicity of the underlying time
circle.
Second, target physics:\ observation / temporal-window targets
finite-temperature observables (Casimir at $T \neq 0$, partition
functions, Stefan--Boltzmann); Wick rotation targets
Lorentzian-signature observables read from Euclidean computations
(Hawking radiation, Unruh thermal bath, BW modular flow).
Third, falsifier sharing:\ the operator-level falsifier for
\S~\ref{sec:proj_observation}'s Bisognano--Wichmann footnote
(closing $\sigma_{i\cdot 2\pi}$ on $\mathcal{T}_{n_{\max}}$) is
identical to the falsifier for the four faces of this projection,
because the bridge in every case runs through the same
KMS--Tomita--Takesaki algebra; one operator-system-level proof
lifts all four faces simultaneously.  Same structural correspondence,
same M1 generator, same falsifier, but distinct mechanisms and
distinct target classes of physical observable.  These are not the
same projection.

Cross-references:\ Paper~32~\S~VIII Remark on the Bisognano--Wichmann
reading documents the spectral-triple-side statement and the
case-exhaustion-theorem context; Paper~35~\S~VIII subsection on the
Bisognano--Wichmann reading of the temporal-window projection
develops the cigar identification at Hartle--Hawking
$\beta = 8\pi M$; Paper~35~\S~VIII subsection on the Unruh pendant
develops the Rindler-wedge identification at $\beta = 2\pi / a$;
\S~\ref{sec:proj_observation} carries the original
Bisognano--Wichmann reading footnote on the temporal-window
projection.  The structural ground for the present \S~\ref{sec:proj_wick_rotation}
entry is the Sprint TD Track~4 + Sprint Unruh-pendant + Track~D
documentation chain (May~2026), which promotes the metric-functional-level
content of the Wick-rotation reading from a footnote (\S~\ref{sec:proj_observation})
and an open-question candidate (\S~\ref{sec:open} entry
on the Lorentz boost; see \S~\ref{sec:open}) to a named projection
slot in \S~\ref{sec:layer2}.  The full operator-system-level
Lorentzian extension remains open as a major NCG-mathematics
structural extension and is documented as a deepening of this entry,
not as a separate candidate.

\textbf{Sprint L2-E (2026-05-17):\ Krein-space period closure,
signature-blind at finite cutoff.}\ A Lorentzian-side analog of
the Sprint L1 period closure was constructed (the strong
``literal identification'' reading recorded below was subsequently
\emph{withdrawn} --- see the 2026-06 status update at the end of this
entry).
Sprint L2 (sub-sprints L2-A scoping, L2-B Krein space, L2-C Lorentzian
Dirac, L2-D Connes axiom audit, L2-E Krein-level modular Hamiltonian;
Paper~32~\S~VIII.E.E) builds the Lorentzian Krein space
$\mathcal{K}_{n_{\max}, N_t} = \mathcal{H}_{\rm GV}^{n_{\max}} \otimes
\mathbb{C}^{N_t}$ via van~den~Dungen 2016 Proposition~4.1
(arXiv:1505.01939) and the Bizi--Brouder--Besnard 2018
$(m, n) = (4, 6)$ classification of Lorentzian signature
$(s, t) = (3, 1)$~\cite{bizi_brouder_besnard2018}, then constructs the
hemispheric wedge $W_L = P_W^{\mathrm{spatial}} \otimes P_{t \geq 0}$
and the wedge KMS state $\rho_W^L = e^{-K_L^{\alpha, W}}/Z$ under the
BW choice $H_{\mathrm{local}} = K_L^{\alpha, W}/\beta$.  The
Krein-space period-closure identities close bit-exactly at
every tested $(n_{\max}, N_t) \in \{1, 2, 3\} \times \{1, 11, 21\}$:\
BW-$\alpha$ period closure $\sigma_{2\pi}^{L, \alpha}(O) = O$, BW-$\gamma$
Tomita period closure $\sigma_{2\pi}^{L, \mathrm{TT}}(a) = a$, flow
conjugacy $\sigma_t^{L, \mathrm{TT}} = \sigma_{-t}^{L, \alpha}$ at
general $t$, and six-witness collapse all at residual $\leq 4 \times
10^{-16}$ (machine precision).  At $N_t = 1$ the Lorentzian-side
construction reduces bit-identically to Paper~42's Riemannian
construction (Frobenius residual exactly $0.0$); the $N_t > 1$ content
is temporal-derivative structure layered on the same compact KMS
$\beta = 2\pi$ circle (see the 2026-06 status update below).
Substantive structural finding:\ Paper~42~\S~7.2 open
question O3 (``$H_{\mathrm{local}} \neq D_W$'') is
\textbf{signature-INDEPENDENT at the Riemannian limit} and refined
upward at $N_t > 1$ by temporal-derivative content (the gap between
the spectral-action object and the modular-Hamiltonian generator
widens with temporal cutoff).  \textbf{Net status (as computed, May~2026):} the period-closure
identities $\sigma_{2\pi}^{L}(O) = O$ hold bit-exactly at every tested
$(n_{\max}, N_t)$, and the $N_t = 1$ slice reduces to Paper~42's
Riemannian construction.

\textbf{Status update (2026-06 --- the ``literal identification''
reading is WITHDRAWN).}\ The subsequent Paper~45 $K^+$-annihilation
theorem (2026-06-09) and the compact-boost closure (2026-06-19) show
that these period-closure identities do \emph{not} constitute a
Lorentzian metric identification at finite cutoff.  The truncated
Bisognano--Wichmann boost generator $K = \mathrm{diag}(2m_j)$ has
\emph{integer} spectrum, so $\mathrm{e}^{2\pi\mathrm{i}K} = I$
bit-exactly at every finite cutoff:\ the modular flow is the compact
KMS $\beta = 2\pi$ \emph{circle}, not a non-compact boost, and a
compact circle carries only a Euclidean (forward-triangle) metric.
The Krein $K^+$ compression is signature-blind (the restricted
Lipschitz seminorm is unchanged under a temporal Wick involution), so
the Lorentzian \emph{signature} is metrically invisible to any
finite-cutoff seminorm;\ the Lorentzian reverse-triangle requires the
non-compact continuum boost (continuous spectrum), reached only as
$n_{\max}\to\infty$.  The finite-cutoff statement is therefore
Euclidean/convention, and the earlier ``literal identification at the
Krein operator-system level'' and ``genuine Lorentzian extension''
framing is withdrawn (Paper~45 descope; the group1 synthesis
honest-scope note).  What survives is the norm-resolvent /
period-closure computation, read as the compact-side KMS statement.
No published finite-cutoff Lorentzian propinquity construction exists.  Cross-manifold W2b
($\mathcal{T}_{S^3} \otimes \mathcal{T}_{\mathrm{Hardy}(S^5)}$) is
structurally distinct from the $(3, 1)$ extension closed here and
remains open per Paper~32~\S~VIII.D ``frontier-of-field framing.''
Files:\ \verb|geovac/modular_hamiltonian_lorentzian.py|,
\verb|tests/test_modular_hamiltonian_lorentzian.py| (74 fast + 3 slow,
all pass), \verb|debug/l2_e_modular_hamiltonian_lorentzian_memo.md|.
Paper~32~\S~VIII.E.E and Paper~42~\S~11 carry the synthesis
writeup;\ Paper~43 (Lorentzian-extension math.OA standalone, outlined
in \verb|papers/group1_operator_algebras/paper_43_lorentzian_extension_outline.md|)
is the natural Zenodo deposit for the standalone math.OA form.

\subsection{Apparatus identity / state-side reduction}
\label{sec:proj_apparatus_identity}
\textbf{Source:} the framework's spectral-side output --- a Hamiltonian
$H$ on the Hilbert space $\mathcal{H}$ of a Fock-projected $S^3$ graph
(or its Bargmann--Segal / molecule-frame analogue), with eigenvalues
$\{\lambda_n\}$ and eigenstates $\{|\psi_n\rangle\}$ supplied by the
upstream projection chain.  All twenty-seven preceding entries of
\S~\ref{sec:layer2} (\S~\ref{sec:proj_fock} through the Wick-rotation
entry of Sprint~2) operate on this spectral side: they introduce
variables, dimensions, and transcendentals at the operator level of
$H$, $D$, or $\mathcal{H}$ itself.\\
\textbf{Target:} the state-side observables associated with a density
matrix $\rho$ on $\mathcal{H}$ --- specifically, the von Neumann entropy
$S(\rho) = -\mathrm{Tr}\,\rho \log\rho$, the reduced density matrices
$\rho_A = \mathrm{Tr}_B \rho_{AB}$ for any bipartite split
$\mathcal{H} = \mathcal{H}_A \otimes \mathcal{H}_B$, the mutual
information $I(A{:}B) = S(\rho_A) + S(\rho_B) - S(\rho_{AB})$, and the
classical--quantum apparatus identity
\begin{equation*}
   S_\mathrm{thermo}(T) \;=\; k_B \cdot S_\mathrm{microstate}(\rho_\beta),
   \qquad \rho_\beta = e^{-\beta H} / Z, \; \beta = 1/T,
\end{equation*}
which says that the macroscopic thermodynamic entropy of a system in
contact with a heat bath at temperature $T$ equals (in units of
$k_B$) the von Neumann entropy of the corresponding Gibbs density
matrix.  This projection is the rule that converts spectral-side
data into state-side observables.\\
\textbf{Variables introduced:} the density matrix $\rho$ on
$\mathcal{H}$, or, equivalently for thermal applications, the
inverse-temperature $\beta = 1/T$ that indexes the Gibbs family
$\rho_\beta = e^{-\beta H}/Z$.  The framework's spectral-side
machinery never references $\rho$; the apparatus identity is what
introduces it.  (When this projection follows
\S~\ref{sec:proj_observation}, $\beta$ is already present as the
Matsubara compactification radius; the apparatus identity then names
the state-side use of that $\beta$.  When it stands alone on a
non-thermal density matrix --- e.g.\ the ground-state reduced density
matrix $\rho_A = \mathrm{Tr}_B |\psi_0\rangle\langle\psi_0|$ of
Paper~27 --- $\rho$ is the new variable.)\\
\textbf{Dimension introduced:} preserves.  Von Neumann entropy
$S(\rho) = -\mathrm{Tr}\,\rho\log\rho$ is dimensionless by
construction (the eigenvalues of $\rho$ are pure probabilities and the
logarithm has no unit).  Temperature carries energy dimension via
$k_B T$, but the dimensional content of $T$ is consumed externally
through Boltzmann's constant rather than generated by the projection
itself; the projection's structural role is the dimensionless
operator-to-number map $\rho \mapsto S(\rho)$.\\
\textbf{Transcendental signature: von~Neumann class.}  The functional
$S(\rho) = -\mathrm{Tr}\,\rho\log\rho$ involves the logarithm of an
operator and is generically transcendental at the value level for any
non-trivial $\rho$.  Structurally, however, this transcendental class
is \emph{categorically distinct from the master Mellin engine}
$M_1 \cup M_2 \cup M_3$ that organises every spectral-side
transcendental documented in Paper~34 \S\S~\ref{sec:proj_fock}--%
\ref{sec:proj_wick_rotation} and Paper~32 \S~VIII.  This is not a
framing claim --- it is a verified PSLQ negative.  Sprint~TD Track~5
(CLAUDE.md \S2, 2026-05-08) computed the framework's atomic
correlation entropies $S_\mathrm{full}(\mathrm{He}) =
0.040811051\ldots$, $S_\mathrm{full}(\mathrm{Li}^+) =
0.011211717\ldots$, and $S_\mathrm{full}(\mathrm{He},\, n_{\max}{=}4)
= 0.041879430\ldots$ to 150 dps, then ran PSLQ at 100 dps against a
$12{,}312$-form mechanical basis spanning the M1 (Hopf-base measure),
M2 (Seeley--DeWitt $\sqrt{\pi}$), and M3 (vertex-parity Hurwitz)
sub-rings plus algebraic and rational control classes.  Result: zero
hits across all three systems.  Stress tests at coefficient ceiling
$10^6$ against the ten simplest ring elements ($1, \pi, \pi^2, 1/\pi,
\zeta(2), \zeta(3), G, \log 2, \sqrt{2}, \varphi$) returned zero hits
at $10^{-50}$ tolerance.  Cross-system ratios
$S(\mathrm{He})/S(\mathrm{Li}^+) = 3.6400\ldots$ and the
$n_{\max}$-ratio $0.9744\ldots$ also null.  Six orders of magnitude in
coefficient ceiling do not change the verdict.  GeoVac correlation
entropy is genuinely a state-side von~Neumann quantity, structurally
disjoint from the spectral-side Mellin ring.\\
\textbf{Used in:} every place the framework computes an entropy or a
reduced density matrix.  Paper~27 (entropy as projection artifact;
$S_\mathrm{kin}/S_\mathrm{full} \sim 10^{-14}$ for He at
$n_{\max}=2,3$; $V_{ee}$-only generation of correlation entropy; the
universal $S_B = A(\tilde{w}_B/\delta_B)^\gamma$ scaling with
$\gamma_\infty \approx 1.96$ in the $Z, n_{\max} \to \infty$ limit)
exercises this projection throughout; the apparatus identity is the
structural mechanism that makes ``entropy'' a well-defined Layer-2
observable in the first place.  Sprint~TD Track~2 (apparatus identity
operational verification: $S_\mathrm{thermo} = k_B \cdot
S_\mathrm{microstate}$ to machine precision, $\max$ residual
$4.8\times 10^{-14}$, across H/He/Li$^+$ at $30$ log-spaced
temperatures each).  Future bipartite-entropy / mutual-information
analyses of atomic correlation (Paper~26 \S~VI Z-scaling
extrapolation).  Any future quantum-information-theoretic comparison
between the framework's qubit Hamiltonians (Paper~14) and competing
encodings, where fidelity and trace distance between candidate ground
states sit on the same state-side as $S(\rho)$.\\
\textbf{Empirical anchors.}  Two independent verifications, both at
machine precision.
\smallskip

\textit{Operational identity (Sprint~TD Track~2, 2026-05-08).}  The
classical-quantum apparatus identity $S_\mathrm{thermo}(T) = k_B \cdot
S_\mathrm{microstate}(\rho_\beta)$ is verified to floating-point
machine precision across the framework's hydrogenic-class systems.
Maximum residual across H/He/Li$^+$ at $30$ log-spaced temperatures
each: $4.8 \times 10^{-14}$.  The $T \to 0$ limit reproduces
Paper~27's ground-state correlation entropy $S_\mathrm{full}(\mathrm{GS})$
bit-identically: He $0.040811$ (residual $3.1\times 10^{-10}$),
Li$^+$ $0.011212$ (residual $3.2\times 10^{-10}$), H residual
identically zero by single-determinant ground-state rigidity.  The
$T \to \infty$ limit approaches $\log(\dim \mathcal{H})$ (e.g.\ $\log
14 \approx 2.639$ for hydrogen at $n_{\max}=3$), the Boltzmann maximum
mixing.  The Lindblad concavity bound holds at every $\beta$.  These
verifications confirm that the apparatus identity is operative as a
state-side reduction rule on the framework's hydrogenic systems
without anomaly.
\smallskip

\textit{Master-Mellin-engine disjointness (Sprint~TD Track~5,
2026-05-08).}  PSLQ at 100~dps against a $12{,}312$-form mechanical
basis (frozen before testing, W3-protocol style) returned zero hits
for $S_\mathrm{full}(\mathrm{He})$, $S_\mathrm{full}(\mathrm{Li}^+)$,
and $S_\mathrm{full}(\mathrm{He},\, n_{\max}{=}4)$.  Twenty random
comparable targets at coefficient ceiling $100$ also returned zero
hits (control null distribution genuinely zero).  Stress tests at
coefficient ceiling up to $10^6$ against the simplest ring elements,
and against cross-system ratios, also null.  The verdict is clean and
sharp: the framework's correlation entropies are not algebraic
combinations of the spectral-side master-Mellin transcendentals.\\
\textbf{Honest scope.}  State-side enumeration is currently
\emph{incomplete} in Paper~34's dictionary.  This entry is the first
state-side projection; the natural continuations
(mutual information, conditional entropy, fidelity, trace distance,
relative entropy $S(\rho \| \sigma)$, Wasserstein--Kantorovich
distance on the state space $P(\mathcal{H})$ of Paper~38's propinquity
machinery) are not enumerated here.  The structural status of those
candidate projections is named in the open-question section
(\S~\ref{sec:open}, candidate state-side enumeration); operationally,
they have appeared in the framework only as scaffolding inside
specific calculations (Paper~26 bipartite entanglement, Paper~38
Wasserstein--Kantorovich identification of the propinquity limit) and
have not yet been promoted to dictionary entries with their own
focal-length variable / dimension / transcendental tagging.  The
state-side dictionary should be read as deliberately opened in this
sprint, not as fully populated.

A second honest scope note: the apparatus identity is operationally
verified on the framework's hydrogen-class systems
(H, He, Li$^+$), where the spectral data is dense and the Gibbs
ensemble is small enough to enumerate.  Extension to molecular
spectral triples (Papers 15, 17, 19) and to the nuclear--electronic
composed triple (Paper~23 \S~VI) is mechanical but has not been
exercised; in those settings the relevant $\rho$ is potentially
high-dimensional and the von Neumann entropy may require sampling
or low-rank approximation methods rather than direct
diagonalisation.  This is an engineering scope rather than a
structural one: the projection is well-defined on every spectral
triple the framework produces; the operational cost of computing
$S(\rho)$ scales with the dimension of $\mathcal{H}$.

A third scope note concerns the relationship between this projection
and the foundational fact that pure-state von Neumann entropy of an
ungrouped system is zero.  Apparatus identity / state-side reduction
is non-trivial only when $\rho$ is mixed (Gibbs ensemble) or when a
bipartite trace $\rho_A = \mathrm{Tr}_B \rho_{AB}$ is taken on a pure
state.  This is consistent with Paper~27's framing: entropy is a
projection artifact in the precise sense that it requires choosing
either a thermal-bath partition (apparatus identity in the Gibbs
form) or a Hilbert-space partition (reduced density matrix).  The
state-side reduction is the structural operation that names that
choice and supplies the resulting dimensionless number.\\
\textbf{Structural reading.}  Categorically distinct from
\emph{every} prior entry of \S~\ref{sec:layer2}, in the precise
sense that those entries operate on the spectral data $(\mathcal{A},
\mathcal{H}, D)$ of the spectral triple, while this entry operates
on density matrices $\rho$ on $\mathcal{H}$.  The spectral / state
divide is the deepest taxonomic split in the dictionary so far: it
is sharper than the variable / dimension / transcendental axes
(Theorem~T3 in \S~\ref{sec:axiomatization}), sharper than
the master-Mellin sub-mechanism partition (M1 / M2 / M3), and sharper
than the Layer~1 / Layer~2 partition itself (which lives entirely on
the spectral side: Layer~1 is the bare graph $D$-data, Layer~2 is the
projected $D$-data).  The apparatus identity is structurally
\emph{the} sibling of \S~\ref{sec:proj_spectral_action} (Connes--%
Chamseddine spectral action): both project the framework's
finite-rank spectral data into a single dimensionless number, but
one does so via $\mathrm{Tr}\,f(D/\Lambda)$ on the spectral side
while the other does so via $\mathrm{Tr}\,\rho\log\rho$ on the
state side.  Sprint~TD Track~5's PSLQ negative is the quantitative
statement that these two procedures generate disjoint transcendental
classes: spectral-side traces of operator functions of $D$ live in
the master Mellin engine $M_1 \cup M_2 \cup M_3$ (Paper~32
\S~VIII case-exhaustion theorem; Sprint~TS-E3, May~2026), while
state-side traces of $\rho \log \rho$ do not.

The structural complement to \S~\ref{sec:proj_adiabatic_BO} is also
worth naming explicitly.  The Born--Oppenheimer projection creates a
parameterised family of Hamiltonians $H(\vec{R})$, and a thermal
ensemble of nuclear configurations weighted by Boltzmann factors
$e^{-\beta E(\vec{R})}$ is precisely the density matrix
$\rho_\beta(\vec{R}) = e^{-\beta H(\vec{R})}/Z$ on which the present
projection acts.  In the language of the multi-focal-composition arc
(May~2026, CLAUDE.md \S~2), adiabatic / BO supplies the Layer-1.5
parameter family; the apparatus identity is the Layer-2 reduction
that converts that family into a single thermodynamic number.  The
two projections compose naturally and the composition is exercised
operationally in Sprint~TD Track~2.

Paper~38's Wasserstein--Kantorovich state space $P(S^3)$
(Proposition~``Limit identification'') is the natural metric
container for the state-side complement of the dictionary.  In the
$n_{\max} \to \infty$ propinquity limit, the truncated spectral
triple $\mathcal{T}_{n_{\max}}$ converges to $\mathcal{T}_{S^3}$ with
the round-$S^3$ Wasserstein--Kantorovich distance as the metric on
its state space.  This identifies a clean future-direction target:
state-side projections beyond apparatus identity should sit naturally
in the Wasserstein--Kantorovich complex --- mutual information,
fidelity, and trace distance all have known metric structures on
state spaces of finite-rank operator algebras, and Paper~38's
machinery suggests these are reachable in the GH-convergent limit
without requiring additional NCG framework extension.  The state-side
dictionary is therefore named here as deliberately incomplete but
\emph{not} blocked: extension is structurally available and is the
natural next direction for the dictionary's axis-completeness.

\paragraph{Born rule paragraph (added 2026-05-26 after sprint
\texttt{ahha\_born\_rule\_attempt}).} The apparatus identity
construction produces Born-rule probabilities directly via the
standard GNS construction:\ for a Layer~1 pure state $\psi \in
\mathcal{H}_{\mathrm{GV}}$ and apparatus subalgebra $\mathcal{A}_W$
with apparatus observable $\hat{A} = \sum_a a \, P_a$, the
state-side reduction $\omega_\psi$ assigns probability
$\omega_\psi(P_a) = \langle\psi|P_a|\psi\rangle = \|P_a \psi\|^2 =
|\langle a|\psi\rangle|^2$ to outcome $a$ when $P_a = |a\rangle\langle
a|$.  The squared-modulus structure is inherited from the complex
inner product on $\mathcal{H}_{\mathrm{GV}}$, which is itself
inherited from the standard Peter--Weyl construction on
$L^2(S^3, \Sigma)$ (Paper~32~\S~III, where $\mathcal{H}_{\mathrm{GV}}$
is taken as input data of the spectral triple).  No GeoVac-specific
operation reduces the Born rule to more primitive structures:\ in
this framework, as in standard algebraic quantum mechanics, the Born
rule follows from the Hilbert-space + GNS combination via Gleason's
theorem (1957) and the standard projection-measure formalism.  The
attempt of 2026-05-26 to derive a non-Gleason / non-standard Born
expression from §III.28 plus Pythagorean orthogonality
(Paper~43~\S~10.2) plus the master Mellin engine
(Paper~32~\S~VIII) returned the same Born expression by all three
routes, indicating that the framework is structurally consistent
with Gleason but does not reach further than Gleason at the
substrate level.  The honest framework-internal classification of
the Born rule is therefore that it sits at the
\emph{calibration-data tier} in the sense of Paper~18~\S~IV.6
(inner-factor input data) and CLAUDE.md~\S~1.7 WH8 (the Born measure as
the observation projection's exchange constant, external calibration
input):\ alongside the Yukawa couplings, the fine-structure
constant $\alpha$, the gauge-group lower-bound calibration, and the
multi-loop QED counterterms (LS-8a renormalization gap), the
probability-weighting structure of measurement is external to the
substrate, not derivable from the packing axiom + Fock projection +
spectral-triple data.  The Pythagorean orthogonality theorem of
Paper~43~\S~10.2 contributes a decoherence-flavored constraint on
\emph{which observables} the apparatus's preferred basis selects (the
centralizer of $H_{\mathrm{local}}$ in $\mathcal{A}_W$) but does not
modify the Born probability assignment.  This is a closed structural
position, not an open question; the framework is in the Wigner /
calibration-tier camp on Born rule rather than the Deutsch--Wallace /
Zurek / Bub--Pitowsky derivation camps.  See \S~\ref{sec:open}
``Born rule as calibration data'' for the corresponding §VIII
closure note.

A final structural reading: the apparatus identity / state-side
reduction is the projection that names ``how an observer extracts a
number from the system,'' as a structural complement to the
spectral-side projections that name ``how the framework constructs
the system in the first place.''  In the master-Mellin-engine
partition of Paper~32 \S~VIII and Paper~18 \S~III.7, the spectral
side has three sub-mechanisms (M1, M2, M3); the state-side --- as
documented by Sprint~TD Track~5's PSLQ negative --- is a fourth
mechanism class structurally outside the engine.  Whether the
state-side admits its own internal sub-mechanism partition
(analogous to the M1 / M2 / M3 split on the spectral side) is the
open structural question this entry leaves for future work.  Current
evidence (one verified entropy class, machine-precision apparatus
identity, PSLQ-disjoint from spectral side) is consistent with
either: a single state-side sub-mechanism that produces all
state-side observables uniformly, or multiple sub-mechanisms whose
distinctions are not yet operationally visible.  Either outcome would
sharpen the structural-skeleton-scope statement of CLAUDE.md \S~2
beyond its current spectral-side-only form.

\subsection{Boundary of \S\ref{sec:layer2}: internal graph-derived variables}
\label{sec:layer2_boundary}

The twenty-eight projections enumerated above share a structural
property worth naming explicitly: each introduces a variable that
is \emph{internal to the framework's graph and projection structure}
--- a graph-derived integer or rational invariant (such as $Z$, $n$,
the foundational $\alpha$, or a Wigner $3j$ coefficient), a
continuous parameter that controls a coordinate or focal-length
choice (such as $\rho$, $R$, $\Lambda$, $\beta$, $m$, $r_E$, $r_Z$,
$Q_N$), or a function of prior projection variables (compositional
projections per Observation~\ref{obs:composition}).

This is structurally distinct from \emph{Layer-2 calibration inputs}
that the framework consumes from external sources but does not
autonomously generate.  Examples: nuclear g-factors $g_p, g_d$
(extracted from precision spectroscopy or QCD lattice), multi-loop
QED coefficients (Eides Tab.~7.3, Karshenboim 2005, Penin--Pivovarov
1998 $\alpha^4$ Breit input), Yukawa couplings (per CLAUDE.md
\S~1.7 WH5; Sprint H1 verdict), and external classical-field
strengths (Zeeman $B$, Stark $\vec{\mathcal{E}}$ --- verified by the
2026-05-09 external-field hunting sprint to be Layer-2 inputs, not
\S~\ref{sec:layer2} projections; the bare Landau spectrum is HO-class
and $\pi$-free, structurally a Bargmann--Segal sibling on the
cyclotron plane, with $\pi$ entering only through the 2D phase-space
measure as a calibration~$\pi$ in the M1 Mellin sub-mechanism, not
through Matsubara compactification).

The structural test that distinguishes \S~\ref{sec:layer2}
projection from Layer-2 input is whether the variable's content is
internal graph-derived or external classical-field /
QCD-lattice-derived.  Zemach and charge radii are intrinsic
focal-lengths once the framework's spectral content (not external
calibration) determines the convolved-potential structure given the
radius; nuclear g-factors and Yukawas are external because no
spectral structure of the framework determines their value.  This
boundary is not absolute: promotion of a Layer-2 input to a
projection slot is the structural commitment that the variable
acquires a focal-length role beyond bare scalar consumption, as
Zemach did when \texttt{magnetization\_density.py} was implemented
at the operator level (\S~\ref{sec:proj_magnetization_density}).
Future Layer-2 inputs may follow if the framework develops
operator-level constructions that give them focal-length status.

% ===================================================================
\section{The Three-Axis Dictionary}
\label{sec:dictionary}
% ===================================================================

Table~\ref{tab:projection_axes} consolidates Section~\ref{sec:layer2}
into the three-axis dictionary that this paper proposes as a
companion to Paper~18's transcendental classification.

\begin{table}[h]
\centering\small
\begin{tabular}{p{3.5cm} p{3cm} p{2cm} p{4cm}}
\toprule
Projection & Variables added & Dimension & Transcendental signature \\
\midrule
Fock conformal & $Z$, $E$ & energy & $\kappa = -1/16$;
$\pi$ via $\mathrm{Vol}(S^3)$ \\
Hopf bundle & $\alpha$ & dimensionless & $\pi = \mathrm{Vol}(S^2)/4$ \\
Bargmann--Segal & $\hbar\omega$ & energy & \textit{none} (graph
$\pi$-free) \\
Stereographic & $r$ (length) & length & conformal factor (Coulomb
$1/r$) \\
Sturmian reparam & $Z, n$ (re-use) & preserves & rational
preserved \\
Spectral action & $\Lambda$, $\alpha$ & energy via $\Lambda$ &
$\sqrt{\pi}\times \mathbb{Q}$ SD coefficients;
$\pi^{2k}\times \mathbb{Q}$ observables \\
Camporesi--Higuchi spinor & $\alpha$ via $(Z\alpha)^2$ &
dimensionless & half-integer Hurwitz $\to$ odd~$\zeta$;
Catalan~$G$, $\beta(4)$ at vertex \\
Wigner $3j$ & none & dimensionless &
$\mathbb{Q}[\sqrt{2k{+}1}]_k$ \\
Wigner $D$ rotation & $\vec{R}_{AB}$ & length &
$\mathbb{Q}[\sqrt{2}, \sqrt{3}, \sqrt{6}]$ \\
Wilson plaquette & $\beta = 1/g^2$ & dimensionless & SU(2) Haar \\
Vector-photon promotion & none & dimensionless & $1/(4\pi)$ per loop \\
Molecule-frame hypersph. & $R$ & length & piecewise-smooth in $R$ \\
Drake--Swainson asymptotic subtraction & $K$ (cancels) & dimensionless &
flow tier; $D_{\text{drake}} = 2(2\ell+1)Z^4/n^3$ rational \\
Rest-mass projection & $m$ & mass &
\textit{trivial} (ring-preserving for $m^2 \in \mathbb{Q}$) \\
Observation / temporal-window & $\beta$ (or $T = 1/\beta$) & time &
$2\pi \cdot \mathbb{Q}$ per Matsubara mode; $\pi^{2k}\cdot\mathbb{Q}$ in
integrated quantities; Stefan-Boltzmann $\pi^2/90$ in high-$T$ limit \\
Two-body Dirac / Breit retardation & $m_l/m_n$ & energy &
$\alpha^4 \cdot \mathbb{Q}$ (ring-preserving over $\mathbb{Q}(\alpha)$) \\
Charge-density (Foldy/Friar) & $r_E$ per nucleus & length &
$\mathbb{Q}(\alpha)$ ring-preserving \\
Magnetization-density (Zemach) & $r_Z$ per nucleus & length &
$\mathbb{Q}(\alpha)$ ring-preserving (leading); profile-shape
sub-features at sub-leading \\
Tensor multipole (quadrupole, etc.) & $Q_N$ per nucleus &
length-squared & $\mathbb{Q}$ via Wigner $3j$ at angular level \\
Phillips--Kleinman / core--valence orthogonality &
$\{\phi_c, E_c\}$ core orbitals (Layer-2 input) & dimensionless &
ring-preserving over $\mathbb{Q}(\alpha)$; no transcendental injected \\
Multipole expansion / Gaunt termination &
$L$ (integer; runs $0,\ldots,L_\text{max} = 2 l_\text{max}$) &
preserves & ring-preserving over $\mathbb{Q}[\sqrt{2k{+}1}]_k$;
exact termination, not asymptotic \\
Bipolar harmonic / Drake combining &
$(k_1, k_2, K)$ bipolar triple & preserves &
pure rational over $\mathbb{Q}[\sqrt{2k{+}1}]_k$ (sibling of Wigner~$3j$) \\
Symmetry / Young tableau &
irrep tableau $\lambda$ & dimensionless &
\textit{none} (integer characters; projector entries in $\mathbb{Z}/N!$) \\
Adiabatic / Born--Oppenheimer &
$\vec{R}$ (slow); $m_e / M_n$ small parameter & length &
\textit{none} at projection step (transcendentals appear downstream
in \S~\ref{sec:proj_coupled_channel}) \\
Coupled-channel / adiabatic curve &
$R$ (radial), $\nu$ (channel index) & length &
algebraic-implicit at Level~3 (over $\mathbb{Q}(\pi,\sqrt{2})$);
piecewise-smooth at Level~4 \\
Gauge choice (Coulomb / Lorenz / Feynman) &
none (selection of representative) & preserves &
gauge-dependent at intermediate-propagator level
($1/(4\pi)$ per loop in Coulomb gauge = M1 Hopf-base measure);
gauge-invariant on physical observables (Ward identity) \\
Wick rotation / signature change &
signature (Euclidean $\leftrightarrow$ Lorentzian; binary discrete
choice) & preserves &
$2\pi\cdot\mathbb{Q}$ via $\mathrm{Vol}(S^1)$ (M1 generator);
same generator as observation/temporal-window, structurally distinct
mechanism \\
Apparatus identity / state-side reduction &
$\rho$ density matrix; or $\beta = 1/T$ Gibbs index & preserves &
von~Neumann class ($\mathrm{Tr}\,\rho\log\rho$); PSLQ-disjoint from
master Mellin engine $M_1\cup M_2\cup M_3$ (Sprint TD Track~5) \\
\bottomrule
\end{tabular}
\caption{The projection dictionary.  Each row is a structural
construction that introduces specific physical content into Layer~1
graph data.  The fourth column is the transcendental signature in
the sense of Paper~18; the second and third columns are the new
axes proposed by this paper.  The thirteenth row (Drake--Swainson
asymptotic subtraction) was added in sprint LS-4 to close the 2P
Bethe-log structural floor identified in LS-3; the fourteenth and
fifteenth rows (rest-mass and observation/temporal-window) were
added in sprint KG / Paper~35 (May~2026) to resolve the
mass-vs-observation ambiguity that had been bundled with other
parameter projections; the sixteenth row (two-body Dirac / Breit
retardation) was added in sprint precision-catalogue 2026-05-08
motivated by the Ps 1S-2S sixth-multi-focal-wall instance, where
the $1/M_\text{nucleus}$ expansion fails at equal mass and the
$\alpha^4$ Breit content surfaces one order before the LS-8a
renormalization wall.  The seventeenth, eighteenth, and nineteenth
rows (Foldy/Friar charge density, Zemach magnetization-density,
tensor multipole) were added in the multi-track Roothaan autopsy
sprint of 2026-05-09 as operator-level Layer-2 input projections for
nuclear-structure scalars ($\langle r^2\rangle_E$, $r_Z$, $Q_N$).
Sprint~1 (2026-05-15) added six more rows
(Phillips--Kleinman / core--valence orthogonality, multipole expansion
/ Gaunt termination, bipolar harmonic / Drake combining, symmetry /
Young tableau, adiabatic / Born--Oppenheimer, coupled-channel /
adiabatic curve) as transcription of pre-existing framework
structures, closing the textbook-completeness gap.  Sprint~2
(2026-05-15) added the final three rows (gauge choice
Coulomb/Lorenz/Feynman; Wick rotation / signature change; apparatus
identity / state-side reduction) as axis-completeness additions:
gauge choice is the second independence witness in the
transcendental-only corner alongside vector-photon promotion, Wick
rotation promotes the Bisognano--Wichmann reading from a \S~VIII
open-question candidate to a named §III slot, and apparatus identity
opens the state-side complement of the dictionary (PSLQ-disjoint
from the master Mellin engine per Sprint TD Track~5).}
\label{tab:projection_axes}
\end{table}

\begin{observation}[Two-axis duality]
\label{obs:duality}
The transcendental classification of Paper~18 and the projection
classification of Table~\ref{tab:projection_axes} are not independent.
Each Paper~18 transcendental tier corresponds to specific projections,
and each projection's transcendental signature is a specific Paper~18
tier:
\begin{itemize}\setlength\itemsep{0pt}
\item Calibration $\pi^{2k}$ tier $\leftrightarrow$ spectral action /
Hopf bundle / vector-photon promotion projections.
\item Spinor-intrinsic $(Z\alpha)$ tier $\leftrightarrow$
Camporesi--Higuchi spinor lift.
\item Odd-$\zeta$ tier $\leftrightarrow$ Camporesi--Higuchi spinor
lift in combination with Mellin/heat-kernel evaluation at integer~$s$.
\item Dirichlet~$L$ tier $\leftrightarrow$ vertex topology composed
with Camporesi--Higuchi at quarter-integer Hurwitz shift.
\item Embedding tier $\leftrightarrow$ molecule-frame hyperspherical
or non-trivial bound-state coupling.
\item Combinatorial-rational tier $\leftrightarrow$
Wigner $3j$ / $D$ alone (no projection beyond Layer~1); also the
structural denominator $D_{\text{drake}}(n,\ell) = 2(2\ell+1)Z^4/n^3$
in the Drake--Swainson projection.
\item Flow tier (Paper~18 \S~IV) $\leftrightarrow$ Drake--Swainson
asymptotic subtraction.  This is the projection that takes a
divergent finite-basis ratio (where closure makes the denominator
vanish) and produces a $K$-independent regularized value through
asymptotic subtraction.  The flow content sits in the
$\alpha$-dependent logarithms of the regulated spectral sum, which
cancel against the $K$-cutoff terms.
\item \emph{Trivial / ring-preserving tier} $\leftrightarrow$
rest-mass projection.  The mass projection introduces a parameter
without introducing a transcendental: the bare-graph algebraic
ring is closed under the addition of $m^2$ when $m^2 \in \mathbb{Q}$.
This is the only projection in the dictionary that introduces a
physical variable while preserving the transcendental signature of
the source spectrum.  Paper~35 verifies this explicitly for the
Klein-Gordon spectrum on $S^3 \times \mathbb{R}$.
\item \emph{Calibration $2\pi \cdot \mathbb{Q}$ tier (temporal subtype)}
$\leftrightarrow$ observation / temporal-window projection.  This is
the categorical home of $\pi$ in the framework: every $\pi$ that
appears in a GeoVac observable can be traced to either (a)~a Hopf-base
calibration~$\pi$ (Hopf bundle / vector-photon promotion), (b)~a
spectral-action $\sqrt{\pi}\cdot\mathbb{Q}$ (heat-kernel coefficient),
or (c)~a temporal Matsubara $2\pi/\beta$ (this projection).  The
Casimir constant on $S^3 \times S^1_\beta$ is the cleanest
illustration: spatial-only Casimir is $1/240$ (rational, no $\pi$);
adding $S^1_\beta$ injects $\pi$ exactly through the Matsubara
mechanism.
\item \emph{Layer-2 input tier (two-body kinematic subtype)}
$\leftrightarrow$ two-body Dirac / Breit retardation projection
(\S~\ref{sec:proj_breit_retardation}).  Same ring-preserving
character as the rest-mass projection (no transcendental injected
beyond the $\alpha$ already present in the spectral-action chain),
but distinct mechanism: rest-mass projection rescales single-particle
observables by $m_e \to m_\text{red}$, whereas Breit retardation
introduces the rest-mass ratio $m_l/m_n$ as a two-body kinematic
control parameter.  In the equal-mass leptonic limit
$m_l/m_n \to 1$ the standard $1/M_\text{nucleus}$ recoil expansion
is invalid, and the projection becomes load-bearing for level
energies at $\alpha^4$.  Empirical anchor: Ps 1S--2S framework-native
at $+64.75$ ppm in \S~\ref{sec:offprecision}; closure to sub-MHz
requires the Penin--Pivovarov 1998 $\alpha^4$ Breit input as
Layer-2~content.
\end{itemize}
\end{observation}

\begin{remark}[variable--dimension correlation]
\label{rem:vd_correlation}
The variable axis and dimension axis are not independent: every
projection in Section~\ref{sec:layer2} that introduces a physical
dimension also introduces a physical variable that carries that
dimension.  Of the twenty-eight projections, twelve add a fresh dimension and
all twelve also add the variable that carries it (perfect correlation
on the V/D side); zero projections add a dimension without a
variable.  The six Sprint~1 additions
(\S\S~\ref{sec:proj_phillips_kleinman}--\ref{sec:proj_coupled_channel})
preserve this correlation: four (Phillips--Kleinman, multipole
expansion, bipolar harmonic, symmetry projection) preserve dimension,
and two (adiabatic / Born--Oppenheimer, coupled-channel / adiabatic
curve) re-use the existing length variable $\vec{R}$ already injected
by stereographic and molecule-frame hyperspherical projections.
The three Sprint~2 additions
(\S\S~\ref{sec:proj_gauge_choice}--\ref{sec:proj_apparatus_identity})
further strengthen the V/D-correlation observation: gauge choice and
Wick rotation both preserve V and D while moving the transcendental
axis (joining vector-photon promotion in the (no-variable,
no-dimension, transcendental-only) corner --- giving \emph{three}
independent witnesses there, all gauge-related or signature-related;
the structural content of that corner is the gauge-and-signature
sector of the framework), and apparatus identity introduces $\rho$ as
the first state-side variable while preserving dimension.  The
transcendental axis remains the only genuinely independent axis; V/D
are still perfectly correlated.  The three intra-nuclear projections of
\S~\ref{sec:proj_charge_density}, \S~\ref{sec:proj_magnetization_density},
and \S~\ref{sec:proj_tensor_multipole} preserve this V/D
correlation: each adds an intra-nuclear focal length together with
its associated dimension ($[L]$, $[L]$, $[L]^2$ respectively).  The transcendental axis is the only one that admits an
independence witness --- vector-photon promotion (\S~\ref{sec:layer2},
projection~11) introduces $1/(4\pi)$ without introducing any new
variable or dimension.  This is the structural content of
Section~\ref{sec:axiomatization} (Theorem~3 composition table):
of the three axes, only the transcendental axis is genuinely
independent of the other two.  The sixteenth projection
(two-body Dirac / Breit retardation,
\S~\ref{sec:proj_breit_retardation}) preserves the V/D correlation:
it adds the variable $m_l/m_n$ and the dimension energy together,
and lives in the ring-preserving $\alpha^4 \cdot \mathbb{Q}$ tier.
\end{remark}

% ===================================================================
\subsection{Base-unit decomposition of the projection family}
\label{sec:base_units}
% ===================================================================

Remark~\ref{rem:vd_correlation} establishes that every dimensioned
projection is paired with the variable carrying that dimension, at
the level of compound dimensions (energy, length, time, mass).  This
subsection decomposes the compound dimensions into base-unit
constituents and asks: how often is each base unit injected across
the twenty-eight projections, and which projections inject which units?

We work in atomic units throughout: $\hbar = m_e = e = 1$, so the
mechanical base units are $[L]$, $[M]$, $[T]$, with charge $[Q]$
carried by the integer label $Z$.  Energy is the compound
$[E] = [M][L]^2[T]^{-2}$, but in atomic units acts as an effective
base for the framework's primary observables; we track $[E]$
alongside the strictly-base units below.  We follow the
\S~\ref{sec:layer2} order of projections for cross-reference with
Table~\ref{tab:projection_axes}.

\begin{table}[h]
\centering\small
\begin{tabular}{p{3cm} p{2.2cm} p{0.9cm} p{4.5cm} p{1.6cm}}
\toprule
Projection & Variable & Dim & Base-unit injection & Role \\
\midrule
Fock conformal & $Z$, $E$ & $[E]$ &
foundational anchor for atomic units; sets $a_0$ ($[L]$) and Ry
($[E]$) jointly via $p_0 = Z/n$; $[Q]$ as integer label $Z$ &
foundational \\
Hopf bundle & $\alpha$ (Paper~2) & dim'less & none & pure \\
Bargmann--Segal & $\hbar\omega$ & $[E]$ &
foundational anchor (HO sibling of Fock); sets oscillator length
($[L]$) and $\hbar\omega$ ($[E]$) jointly &
foundational (sibling) \\
Stereographic & $r$ & $[L]$ & $[L]$ as continuous coordinate; no new
scale (coordinatizes the foundational $[L]$) &
coordinate-explicit \\
Sturmian reparam & $Z, n$ (re-use) & --- & none & reparam \\
Spectral action & $\Lambda$, $\alpha$ & $[E]$ & new $[E]$ scale (UV
cutoff $\Lambda$); $\alpha$ dimensionless & scale-extension \\
Camporesi--Higuchi spinor & $\alpha$ via $(Z\alpha)^2$ & dim'less &
none & pure \\
Wigner $3j$ & none & dim'less & none & pure \\
Wigner $D$ rotation & $\vec{R}_{AB}$ & $[L]$ & new $[L]$ scale
(internuclear position vector) & scale-extension \\
Wilson plaquette & $\beta = 1/g^2$ & dim'less & none & pure \\
Vector-photon promotion & none & dim'less & none & pure \\
Mol-frame hyperspherical & $R$ & $[L]$ & new $[L]$ scale
(internuclear distance; same physical length as $|\vec{R}_{AB}|$) &
scale-extension \\
Drake--Swainson & $K$ (cancels) & $[E]$ & transient $[E]$ scale
(cancels in $\beta_\text{low}(K) + \beta_\text{high}(K)$) &
scale-extension (transient) \\
Rest-mass & $m$ & $[M]$ & first $[M]$ scale beyond foundational
$m_e = 1$; one projection per particle species &
scale-extension (new base) \\
Observation / temporal-window & $\beta$ (or $T$) & $[T]$ & first
$[T]$ scale beyond foundational $\tau_0 = \hbar/\mathrm{Ry}$; one
projection per measurement window &
scale-extension (new base) \\
Two-body Dirac / Breit & $m_l / m_n$ & dim'less &
ratio of two prior $[M]$ scales (lepton, nucleon); no new base
content & composition \\
Charge-density (Foldy/Friar) & $r_E$ per nucleus & $[L]$ &
new $[L]$ scale (intra-nuclear, charge distribution); modifies
$V_\text{ne}$ via convolution & scale-extension \\
Magnetization-density (Zemach) & $r_Z$ per nucleus & $[L]$ &
new $[L]$ scale (intra-nuclear, magnetization distribution);
modifies $A_\text{hf}^\text{contact}$ & scale-extension \\
Tensor multipole (rank-$\geq 2$) & $Q_N$ per nucleus & $[L]^2$ &
rank-$L$ multipole moment, dim $[L]^L$; couples to electronic
gradients ($l \geq 1$ on electronic side) &
scale-extension (rank-2) \\
\bottomrule
\end{tabular}
\caption{Base-unit decomposition of the twenty-eight projections (the
six Sprint~1 additions of 2026-05-15 add no new base units; they
preserve or re-use existing dimensions; the three Sprint~2 additions
of 2026-05-15 also add no new base units --- gauge choice, Wick
rotation, and apparatus identity all preserve dimension).
``Variable'' is the new physical content; ``Dim'' repeats the
compound dimension from Table~\ref{tab:projection_axes};
``Base-unit injection'' shows what enters in $[L]$, $[M]$, $[T]$,
$[Q]$, with $[E]$ tracked as the framework's effective base
observable; ``Role'' classifies the projection by its action on the
dimensional content of a projection chain.  Six roles emerge:
\emph{foundational} (sets the atomic-units anchor),
\emph{coordinate-explicit} (re-expresses an anchored unit as a
continuous coordinate), \emph{reparam} (relabels without adding
content), \emph{scale-extension} (adds a new variable in a base
unit, possibly a new base), \emph{pure} (adds transcendental and/or
combinatorial content with no dimensional injection), and
\emph{composition} (variable is a function of prior projection
variables).}
\label{tab:base_units}
\end{table}

The inverse view --- which projections inject each base unit ---
makes the structure visible:

\begin{itemize}\setlength\itemsep{0pt}
\item $[L]$ (scalar): \emph{six} projections.  Fock or
Bargmann--Segal anchors the foundational length ($a_0$ or oscillator
length); stereographic re-expresses it as a continuous coordinate
$r$; Wigner~$D$ rotation introduces nuclear position
$\vec{R}_{AB}$; molecule-frame hyperspherical separation introduces
internuclear distance $R$ (Wigner~$D$ and mol-frame inject the same
physical length scale through different projections);
charge-density (\S~\ref{sec:proj_charge_density}) introduces the
intra-nuclear charge radius $r_E$; magnetization-density
(\S~\ref{sec:proj_magnetization_density}) introduces the
intra-nuclear Zemach radius $r_Z$.  These last three intra-nuclear
length scales are each \emph{categorically distinct} from the
atomic / internuclear scales: they modify different Hamiltonian
operators (Coulomb $V_\text{ne}$, Fermi-contact
$A_\text{hf}^\text{contact}$) with different selection rules, with
empirical separation across hydrogen Lamb shift / 21~cm hyperfine /
deuteron molecular HFS.

\item $[L]^2$ (rank-2 tensor): \emph{one} projection.  Tensor
multipole (\S~\ref{sec:proj_tensor_multipole}) introduces the
nuclear quadrupole moment $Q_N$ at rank-2 leading; higher-rank
multipoles ($L \geq 3$) follow at higher $[L]^L$.  Categorically
distinct from scalar $[L]$ injections: the operator is tensorial
and requires $l \geq 1$ on the electronic side.
\item $[E]$: \emph{four} projections.  Fock or Bargmann--Segal
anchors the foundational energy (Ry or $\hbar\omega$); spectral
action introduces UV cutoff $\Lambda$; Drake--Swainson introduces a
transient subtraction scale $K$ (which cancels in the regularized
result); two-body Breit retardation introduces $\alpha^4 \cdot
\mathrm{Ry}$ as a correction-energy scale, where $\alpha^4$ is
dimensionless and Ry is foundationally anchored.
\item $[M]$: \emph{one} projection (rest-mass).  Each particle
species in the system requires a separate rest-mass projection;
two-particle systems compose two such projections.
\item $[T]$: \emph{one} projection (observation /
temporal-window).  Each measurement window or temperature requires
a separate observation projection.
\item $[Q]$: integer-counter label $Z$ at the Fock projection only.
$[Q]$ does not appear as a continuous scale anywhere in the
projection family --- charge is always counted, never
measured-against-a-scale.
\end{itemize}

\begin{observation}[$L$--$E$ asymmetry]
\label{obs:LE_asymmetry}
Of the mechanical base units (and the charge unit), only $[L]$ and
$[E]$ are injected by multiple projections; $[M]$, $[T]$, and $[Q]$
are each instantiated by exactly one projection beyond the
foundational atomic-units anchor.  After the addition of the three
intra-nuclear projections (charge-density, magnetization-density,
tensor multipole), the asymmetry is no longer symmetric: $[L]$
admits \emph{six} scalar instances plus one rank-2 tensor instance,
while $[E]$ admits four.  The framework records more length scales
(atomic, internuclear, polyatomic, and three intra-nuclear
distribution scales) than energy scales (Rydberg, UV cutoff,
multi-loop subtraction, two-body Breit correction).  The structural
reason: every distinct spatial structure introduces an independent
focal length because different intra-nuclear and intra-molecular
geometries enter different Hamiltonian operators with different
selection rules; spectral content has fewer distinct roles
(foundational, regulator, subtraction, correction).  The conjugate
pair $[L] \leftrightarrow [E]$ remains the only multi-instance
axis, but $[L]$ now dominates: physics has more places to put
length scales than to put energy scales.
\end{observation}

\begin{observation}[Compositional projections]
\label{obs:composition}
The two-body Dirac / Breit retardation projection
(\S~\ref{sec:proj_breit_retardation}) is the first identified
instance of a structural family of \emph{compositional projections}:
projections whose variable is a function (typically a ratio or
product) of variables introduced by prior projections.  All other
projections in Table~\ref{tab:base_units} introduce variables
independent of prior content; compositional projections are
structurally distinct because they are \emph{gated} by the prior
instantiation of two or more distinct values of a foundational
variable in the same projection chain.

We catalogue two confirmed instances of compositional projections:

\begin{enumerate}\setlength\itemsep{0pt}
\item \emph{Kinematic composition.}  Variable: rest-mass ratio
$m_l / m_n$.  Gating: prior instantiation of two distinct $[M]$
scales by separate rest-mass projections (\S~\ref{sec:proj_restmass}).
Manifests at $\alpha^4$ in two-body bound-state energies, one full
order before the $\alpha^5$ multi-loop renormalization wall.
Empirical anchor: positronium 1S--2S framework-native at
$+64.75$~ppm (\S~\ref{sec:catalog}; sixth and structurally cleanest
instance of the multi-focal-composition wall named in
CLAUDE.md~\S~2).

\item \emph{Angular composition.}  Variable: ratio of effective
charges $Z_a / Z_b$ within a single atom (e.g., He 1s
$Z_\text{eff} = 2$ vs.\ 2p $Z_\text{eff} = 1$).  Gating: prior
instantiation of two distinct foundational Fock anchors at different
effective charges within the same electronic Hilbert space (the
``internal multi-focal'' pattern of \S~\ref{sec:catalog}).
Manifests at $\alpha^2$ in fine-structure splittings, where the
bipolar harmonic decomposition $(k_1{=}0, k_2{=}2)$ direct $+$
$(k_1{=}1, k_2{=}1)$ exchange and the Drake combining coefficients
$(\tfrac{3}{50}, -\tfrac{2}{5}, \tfrac{3}{2}, -1)$ are the
structural carrier.  Empirical anchor: He $2{}^3P_0$--$2{}^3P_1$
at $-0.014\%$ and $2{}^3P_0$--$2{}^3P_2$ span at $-0.20\%$
(\S~\ref{sec:catalog}).
\end{enumerate}

The two instances differ in their gating mechanism (kinematic at
the two-body level vs.\ angular at the two-orbital level) but share
the compositional character: each requires the prior projection
chain to have instantiated two distinct values of a foundational
variable, so that the compositional projection's variable (ratio,
product, or function) is well-defined.

The structural prediction the generalization makes: any future
GeoVac observable whose computation requires a function of two or
more foundationally-anchored variables in the same projection chain
belongs to this family, and the order at which it manifests in
$\alpha$-perturbation theory follows the projection at which the
function is gated.  Candidate compositional projections not yet
observed:

\begin{itemize}\setlength\itemsep{0pt}
\item Ratio of two length scales within a single nucleus (e.g.,
$r_\text{charge} / r_Z$, the Friar/Foldy-versus-Zemach distinction
inside one proton; would gate sub-Zemach corrections to hyperfine).
\item Product of three rest-mass scales (three-body kinematic
content beyond Breit retardation; would gate three-body bound-state
corrections at $\alpha^4$ or higher).
\item Ratio of two temperature scales $\beta_a / \beta_b$ in
multi-window observation (e.g., simultaneous interrogation at two
distinct $T$).
\end{itemize}

The compositional-projection family is open-ended: confirmed
instances locate the gating-order pattern (Breit at $\alpha^4$,
internal multi-focal at $\alpha^2$), and additional empirically
anchored instances would refine the structural prediction.
\end{observation}

This decomposition sharpens
Remark~\ref{rem:vd_correlation}: V--D correlation holds at the
compound dimension level, but at the base-unit level $[L]$ and
$[E]$ each admit four instances while $[M]$, $[T]$ admit one each
and $[Q]$ enters only as an integer label.  The base-unit
decomposition is a finer falsifiable test than the compound
dimension alone: an observable that depends on multiple length
scales must engage the foundational projection plus the relevant
length-injecting projections; an observable that depends on
temperature must engage the observation projection; an observable
depending on a species mass ratio must engage two rest-mass
projections plus Breit retardation.  The empirical matches
catalogue (\S~\ref{sec:catalog}) provides a forward test:
each row's projection chain should be traceable to a base-unit
signature consistent with the dimensions of the matched
observable.

% ===================================================================
\section{The Empirical Matches Catalogue}
\label{sec:catalog}
% ===================================================================

We list GeoVac quantities that match standard known physics values
to machine precision or stated uncertainty.  Each row is tagged by
projection(s), variables involved, dimension, transcendental class,
and match precision.  This is a structural index, not an exhaustive
benchmark report; the underlying numerical evidence lives in
\texttt{tests/} and the cited papers.

\begin{longtable}{p{3.5cm} p{2.8cm} p{1.3cm} p{2.2cm} p{1.8cm} p{1.4cm}}
\toprule
Match & Projection(s) & Vars & Dim & Trans.\ class & Match \\
\midrule
\endhead
\textbf{Zero-projection (Layer~1 only)} & & & & & \\
Pauli count $N = 11.10{\cdot}Q$ & --- & none & dimensionless &
integer/rational & exact \\
Magic numbers $2,8,20,28,\ldots$ & HO + LS counting & none &
dimensionless & integer & exact \\
Gaunt selection rules & 3$j$ algebra & none & dimensionless &
rational & exact \\
$\pi$-free certificate (Coulomb) & --- & none & dimensionless &
rational & exact \\
$\pi$-free certificate (Bargmann) & --- & none & dimensionless &
rational & exact \\
Angular sparsity 1.44\% at $\ell_\text{max}=3$ & 3$j$ algebra &
none & dimensionless & rational & exact \\
Discrete $c_1 = 0$ on Hopf-$S^3$ Fock graph at $n_\text{max} = 2, 3$
(Condon--Shortley $+$ $Z_2$ alternating) & --- &
none & dimensionless & integer & exact (Fraction) \\
\midrule
\textbf{Single-projection} & & & & & \\
Hydrogen $E_n = -Z^2/(2n^2)$ Ha & Fock + $\kappa$ & $Z, E$ &
energy & rational & exact symbolic \\
2$p$ doublet $\alpha^2/32$ Ha (Z=1) & spinor lift $\circ$ Fock &
$Z, \alpha$ & energy & $(Z\alpha)^2{\cdot}\mathbb{Q}$ & exact symbolic \\
Dirac fine-structure formula ($n\le 4$) & spinor lift $\circ$ Fock &
$Z, \alpha$ & energy & $(Z\alpha)^2{\cdot}\mathbb{Q}$ & exact \\
$\Pi = 1/(48\pi^2)$ VP coefficient & spectral action &
$\Lambda, \alpha$ & dimensionless & $\pi^2{\cdot}\mathbb{Q}$ &
exact symbolic \\
$\beta(\alpha) = 2\alpha^2/(3\pi)$ one-loop & spectral action &
$\alpha$ & dimensionless & $\pi{\cdot}\mathbb{Q}$ & exact symbolic \\
$D^2$ heat kernel SD coefficients & spectral action & $\Lambda$ &
energy & $\sqrt{\pi}{\cdot}\mathbb{Q}$ & exact symbolic \\
$F = \zeta_R(2) = D_{n^2}(d_\text{max})$ & Fock + Dirichlet at
$s=4$ & graph data & dimensionless & $\pi^2{\cdot}\mathbb{Q}$ &
exact symbolic \\
$\Delta^{-1} = g_3^\text{Dirac} = 40$ & spinor lift & graph data &
dimensionless & integer & exact \\
$B = 42$ Casimir trace at $m=3$ & --- & graph data &
dimensionless & integer & exact \\
$D_\text{even}(4)$ closed form & spinor lift (vertex selection) &
$\alpha$ & dimensionless & $\pi^2, \pi^4, G, \beta(4)$ & PSLQ 80~dps \\
$E_\text{Cas}^{S^3} = 1/240$ (conformal scalar) & Fock + rest-mass
projection (conformal value $m^2 = 1$) & none & dimensionless &
rational & exact symbolic \\
$E_\text{Cas}^{S^3, \text{Dirac}} = 17/480$ (full Dirac) & Fock +
Camporesi--Higuchi spinor lift & none & dimensionless & rational &
exact symbolic \\
$E_\text{Cas}^{\text{HO}, S^5} = -17/3840$ (3D HO, Bargmann--Segal) &
Bargmann--Segal $+$ spectral action ($\zeta$-regularization) & none &
dimensionless & rational & exact symbolic \\
Hydrogen Bohr levels $E_n = -Z^2/(2n^2)$ Ha & Fock $+$ energy-shell
rescaling & $Z, n$ & energy & rational & exact symbolic \\
Static Coulomb-like Green's fn $(L+I)^{-1}$ on Hopf $S^3$ at
$n_\text{max} = 2, 3$ & Fock (Layer-1 matrix inverse, no continuous
integration) & none & dimensionless & rational & exact symbolic at
finite $n_\text{max}$ \\
\midrule
\textbf{Two-projection chain} & & & & & \\
Stefan-Boltzmann $\pi^2/90$ on $S^3 \times S^1_\beta$ (high-$T$) &
Fock $+$ observation/temporal-window $+$ spectral action
($\zeta$-regularization) & $T$ or $\beta$ & energy density &
$\pi^2 \cdot \mathbb{Q}$ & exact symbolic \\
Uehling shift $-27.13$ MHz at 2$S$ & spectral action $\circ$ Fock &
$Z, \alpha$ & energy & $\pi^2{\cdot}\mathbb{Q}$ & $<1\%$ \\
21~cm hyperfine $A_\text{HF} = 1418.84$ MHz Bohr-Fermi (HF-1, $+$ external
$g_p$) & spinor lift (Fermi-contact vertex) $\circ$ Fock $\circ$ rest-mass
projection ($m_p$, proton) & $\alpha, g_p, m_e/m_p$ & energy &
$\alpha^2{\cdot}\mathbb{Q}$ & $-1102$ ppm (= $-a_e$) \\
21~cm hyperfine $A_\text{HF} = 1420.488$ MHz $+$ Parker--Toms $a_e$
(HF-2, $+$ external $g_p$) & spinor lift (Fermi-contact vertex) $\circ$
Fock $\circ$ rest-mass projection ($m_p$, proton) $\circ$ spectral
action (curvature expansion) & $\alpha, g_p, m_e/m_p$ & energy &
$\alpha^2{\cdot}\mathbb{Q}, c_1 \cdot \mathbb{Q}$ & $+58$ ppm \\
21~cm hyperfine $A_\text{HF} = 1420.432$ MHz $+$ external Zemach $r_Z$
(HF-4) & HF-2 chain $\circ$ rest-mass projection $\circ$
magnetization-density (\S~\ref{sec:proj_magnetization_density}) &
$r_Z$ & energy & calibration & $+18$ ppm (Eides Tab.~7.3 consistent) \\
$\mu$H 1S Bohr--Fermi $\nu_F(\mu p) = 182.4433$ meV (Sprint MH Track B,
$+$ external $g_p$) & spinor lift (Fermi-contact vertex) $\circ$ Fock
$\circ$ rest-mass projection ($m_\mu$, muon, lepton swap) $\circ$
rest-mass projection ($m_p$, proton) & $\alpha, g_p, m_\mu/m_e, m_e/m_p$
& energy & $\alpha^2 \cdot \mathbb{Q}$ &
$+2$ ppm vs Eides pure-QED $\nu_F = 182.443$ meV \\
$\mu$H 1S Zemach mass-enhancement factor $185.94 \times$ ep
(Sprint MH Track B; operator-level Track C, May 2026) & rest-mass
projection $\circ$ magnetization-density
(\S~\ref{sec:proj_magnetization_density}) $\circ$ Fock; \texttt{lepton\_mass}
parameter on \texttt{MagnetizationDensitySpec} eliminates the manual
mass-scaling step Track B used in test scripts (Pauli-string sum now
carries correct lepton-mass scaling natively for downstream muonic
VQE) & $r_Z, m_\text{red}$ & dimensionless & calibration &
$0.55\%$ vs Eides muonic target $-7300$ ppm; bit-identical to Track B
manual scaling (intrinsic gap is sub-leading Eides corrections, not
operator construction) \\
$\mu$H 2$S$--2$P$ Uehling $+205.0074$ meV full kernel
(Sprint MH Track A) & spectral action $\circ$ Fock $\circ$ rest-mass
projection ($m_e \to m_\mu$) & $Z, \alpha, m_\mu/m_e$ & energy &
$\pi^2{\cdot}\mathbb{Q}$ & $<1$ ppm vs Antognini~2013 \\
$\mu$d 2$S$--2$P$ Lamb $+202.63$ meV framework $+$ Layer-2
(multi-track Roothaan sprint Track~1, 2026-05-18; operator-level
Roothaan autopsy in \S\ref{sec:autopsy_mud_lamb}) & spectral action
$\circ$ Fock $\circ$ rest-mass projection ($m_e \to m_\mu$, nucleus
$m_p \to m_d$) $+$ Krauth Layer-2 (multi-loop QED $+$ recoil $+$
Friar HO $+$ deuteron polarizability) & $Z, \alpha, m_\mu/m_e,
m_d/m_e, r_d$ & energy & $\pi^2{\cdot}\mathbb{Q}$, calibration &
$-0.12\%$ vs CREMA~2016 $+202.8785(34)$ meV (Pohl) \\
$\mu$d 2$S$--2$P$ Friar leading $-27.85$ meV bare Eides at $r_d
= 2.1413$ fm (Roothaan Track 1) & Fock $\circ$ charge-density
(\S\ref{sec:proj_charge_density}) $\circ$ rest-mass projection &
$Z, \alpha, m_\mu/m_e, m_d/m_e, r_d$ & energy & calibration $+
\pi^2{\cdot}\mathbb{Q}$ via $(Z\alpha)^4$ & $+0.024\%$ vs Krauth~2016
$-27.84$ meV at $r_d$ (bit-tight; scaling
$(m_\text{red}/m_e)^3 (r_d/r_p)^2 = 7.580$ reproduced to 9 digits) \\
$\mu^4$He$^+$ 2$S$--2$P$ Uehling $+1665.77$ meV full kernel
(multi-track sprint Track B, 2026-05-19; operator-level Roothaan
autopsy in \S\ref{sec:autopsy_muhe4_lamb}) & spectral action $\circ$
Fock $\circ$ rest-mass projection ($m_e \to m_\mu$, $Z = 1 \to 2$,
nucleus $m_p \to m_{^4\text{He}}$) & $Z, \alpha, m_\mu/m_e,
m_{^4\text{He}}/m_e$ & energy & $\pi^2{\cdot}\mathbb{Q}$ &
$-0.032\%$ vs Diepold 2018 $+1666.305$ meV at $\beta = 0.682$ (cleanest
single-component muonic VP match at heavy Z; transports kernel
verbatim from $\mu$H/$\mu$d) \\
$\mu^4$He$^+$ 2$S$--2$P$ FNS Foldy $-295.67$ meV bare Eides at $r_\alpha
= 1.6755$ fm (Roothaan Track B) & Fock $\circ$ charge-density
(\S\ref{sec:proj_charge_density}) $\circ$ rest-mass projection &
$Z, \alpha, m_\mu/m_e, m_{^4\text{He}}/m_e, r_\alpha$ & energy &
calibration $+ \pi^2{\cdot}\mathbb{Q}$ via $(Z\alpha)^4$ &
$+2.9\%$ vs Diepold leading, $+4.0\%$ vs Diepold full FNS
($Z = 1 \to 2$ amplifies higher-Friar absorption to $\sim 11$ meV;
scaling $(m_\text{red}/m_e)^3 (r_\alpha/r_p)^2 \cdot Z^4 = 80.45$
reproduced to 9 digits) \\
Mu 1S-2S transition $2{,}455{,}528{,}721$~MHz (rest-mass rescaling
from H experiment, sprint precision-catalogue 2026-05-08) & Fock $\circ$
rest-mass projection (lepton-unchanged, nucleus $m_p \to m_\mu$) &
$\alpha, m_e, m_\mu/m_e$ & frequency &
$\alpha^2{\cdot}\mathbb{Q}$ & $-0.11$ ppm vs Mu-MASS Crivelli~2018 \\
Mu 1S-2S transition $2{,}455{,}528{,}394$~MHz (framework Bohr$+$SE $+$
cross-register recoil) & Fock $\circ$ spectral action $\circ$ rest-mass
projection $\circ$ cross-register $V_{eN}$ (Roothaan) & $\alpha, m_e,
m_\mu/m_e$ & frequency & $\alpha^2{\cdot}\mathbb{Q}$ &
$-0.25$ ppm vs Mu-MASS \\
Mu 2$S_{1/2}$-2$P_{1/2}$ Lamb shift $1047.63$ MHz (framework one-loop
QED with full Uehling, sprint precision-catalogue 2026-05-08) &
Fock $\circ$ spectral action $\circ$ Sturmian $\circ$ rest-mass projection
& $\alpha, m_e, m_\mu/m_e$ & frequency & $\alpha^2{\cdot}\mathbb{Q}$ &
$+0.013\%$ vs Karshenboim 2005 theory $1047.49(5)$ MHz \\
Bethe log $\ln k_0(2S) = 2.786$ (LS-3) & Sturmian (acc.) $\circ$ Fock & $Z$ &
dimensionless & calibration & $-0.92\%$ \\
Bethe log $\ln k_0(1S) = 3.002$ (LS-3) & Sturmian (acc.) $\circ$ Fock & $Z$ &
dimensionless & calibration & $+0.60\%$ \\
D 1S hyperfine $\nu_\text{HFS}(D) = 327.3975$ MHz Bohr--Fermi strict
(sprint precision-catalogue 2026-05-08, $+$ external $g_d$) & spinor
lift (Fermi-contact vertex) $\circ$ Fock $\circ$ rest-mass projection
($m_d$, deuteron); $I=1$ multiplicity $3/2$ on
$A_\text{hf}\,\hat{I}\cdot\hat{S}$ & $\alpha, g_d, m_e/m_d$ & frequency &
$\alpha^2{\cdot}\mathbb{Q}$ & $+40$ ppm vs Wineland--Ramsey 1972
$327.384352522(2)$ MHz \\
T 1S hyperfine $\nu_\text{HFS}(T) = 1515.865$ MHz Bohr--Fermi strict
(multi-track sprint Track 2, 2026-05-18, $+$ external $g_t$; see
\S\ref{sec:autopsy_t_hfs}) & spinor lift (Fermi-contact vertex)
$\circ$ Fock $\circ$ rest-mass projection (universal $m_e/m_p$, $g_t$);
$I=1/2$ multiplicity $1$ shared with H 21cm on
$A_\text{hf}\,\hat{I}\cdot\hat{S}$ & $\alpha, g_t, m_e/m_p$ & frequency &
$\alpha^2{\cdot}\mathbb{Q}$ & $-551$ ppm vs Greene 2017
$1516.701470773(8)$ MHz; cleanest mass-vs-spin separation in catalogue
(I=1/2 like H, $m_t \approx 3 m_p$ like D-extended) \\
Mu 1S hyperfine $\nu_\text{HFS}(\text{Mu}) = 4464.1905$ MHz framework
Bohr--Fermi $+$ Schwinger one-loop (sprint precision-catalogue 2026-05-08,
muonium triple closure) & spinor lift (Fermi-contact vertex) $\circ$ Fock
$\circ$ rest-mass projection (lepton-unchanged, nucleus $m_p \to m_\mu$);
leptonic nucleus, no QCD content & $\alpha, m_e, m_\mu/m_e$ & frequency &
$\alpha^2{\cdot}\mathbb{Q}$ & $+199$ ppm vs Liu 1999 $4463.302765(53)$ MHz;
residual cleanly LS-8a multi-loop QED $+$ recoil (no QCD competing budget) \\
He $2{}^3P_0$--$2{}^3P_1$ fine structure $+29{,}612.91$ MHz (Sprint~3
BF-D + Sprint~4 DD; Drake combining $(\tfrac{3}{50},-\tfrac{2}{5},
\tfrac{3}{2},-1)$, J-pattern from rank-$k$ 6j; sprint precision-catalogue
2026-05-08) & spinor lift $\circ$ Wigner $3j$ ($6j$-coupled) $\circ$
Wigner $3j$ (Gaunt) $\circ$ Fock; internal multi-focal (1s, 2p at
distinct $Z_\text{eff}$), bipolar $(k_1{=}0,k_2{=}2)$ direct +
$(k_1{=}1,k_2{=}1)$ exchange & $\alpha, Z$ & frequency &
$\alpha^2{\cdot}\mathbb{Q}$ & $-0.014\%$ vs NIST $+29{,}616.951(6)$ MHz;
below LS-8a $\alpha^3/\pi$ budget \\
He $2{}^3P_0$--$2{}^3P_2$ span $+31{,}844.07$ MHz (Sprint~3 BF-D
+ Sprint~4 DD; sprint precision-catalogue 2026-05-08) & same chain &
$\alpha, Z$ & frequency & $\alpha^2{\cdot}\mathbb{Q}$ &
$-0.20\%$ vs NIST $+31{,}908.131(6)$ MHz; within LS-8a one-loop QED budget \\
Ps 1S-2S transition $1{,}233{,}687{,}096.1$~MHz (framework-native: Bohr at
$m_\text{red}(ee) = 0.5$ $+$ Eides~\S\,3.2 SE Lamb-shift differential;
sprint precision-catalogue 2026-05-08, equal-mass leptonic 1S-2S; $+$
external Penin--Pivovarov $\alpha^4$ Breit) & Fock $\circ$ rest-mass
projection (equal-mass limit $\lambda_a = \lambda_b = 0.5$) $\circ$
spectral action (Eides) & $\alpha, m_e$ & frequency &
$\alpha^2{\cdot}\mathbb{Q}$ & $+64.75$ ppm vs Fee~1993
$1{,}233{,}607{,}216.4(3.2)$~MHz \\
\midrule
\textbf{Three-projection chain} & & & & & \\
Bethe log $\ln k_0(2P) = -0.0307$ (LS-4) & Drake--Swainson $\circ$ Sturmian
$\circ$ Fock & $Z$ & dimensionless & flow & $+2.4\%$ at $N=40$ \\
Bethe log $\ln k_0(3P) = -0.0406$ (LS-4) & Drake--Swainson $\circ$ Sturmian
$\circ$ Fock & $Z$ & dimensionless & flow & $+6.3\%$ at $N=40$ \\
Bethe log $\ln k_0(3D) = -0.005236$ (LS-4) & Drake--Swainson $\circ$ Sturmian
$\circ$ Fock & $Z$ & dimensionless & flow & $-0.24\%$ at $N=40$ \\
Lamb shift $1054$ MHz (LS-4 N=40 sweet spot) & Fock $\circ$ spectral action
$\circ$ Sturmian $\circ$ Drake--Swainson & $Z, \alpha$ & energy & calibration
$+$ flow & $-0.39\%$ (T$+$A cancellation) \\
Lamb shift $1029$ MHz (LS-3) & Fock $\circ$ spectral action $\circ$
Sturmian (acc.) & $Z, \alpha$ & energy & calibration & $-2.77\%$ (T$+$A)\\
$\mu$H 2$S$--2$P$ Lamb shift $202.17$ meV (framework $+$ literature,
Sprint MH Track A) & Fock $\circ$ spectral action $\circ$ Sturmian
$\circ$ Drake--Swainson $\circ$ rest-mass projection ($m_e \to m_\mu$)
& $Z, \alpha, m_\mu/m_e, r_p$ & energy & calibration $+$ flow &
$-0.10\%$ vs CREMA $202.3706(23)$ meV \\
Heisenberg--Euler 1-loop $\alpha^2/(45\pi m_e^4)$ &
spectral action (Schwinger Mellin) $+$ vector-photon (one loop)
& $\alpha, m_e$ & length${}^4 \cdot$ energy${}^{-1}$ &
$\alpha^2 \cdot \mathbb{Q} / \pi$ & structural (S) \\
Lyman $\alpha$ emission $A_{21}(2P{\to}1S) = 6.2683 \times 10^8$~s$^{-1}$
(Sprint Calc-L, 2026-05-09) & Fock $\circ$ Wigner $3j$ $\circ$
vector-photon promotion & $Z, \alpha, \omega$ & 1/time &
$\alpha^3 \cdot \mathbb{Q}[\sqrt{6}]$ & $+0.055\%$ vs NIST
$6.2649 \times 10^8$~s$^{-1}$; residual is non-relativistic
(Schr\"odinger vs Dirac at $(Z\alpha)^2 \sim 5\times 10^{-5}$),
closes by adding Camporesi--Higuchi spinor lift to the chain
(depth-prediction consistent) \\
H 1S static dipole polarizability
$\alpha_p(\text{H}, 1S) = \tfrac{9}{2}\, a_0^3$
(Sprint Calc-P, 2026-05-09) & Fock $\circ$ Wigner $3j$ $\circ$
Sturmian reparam.\ at $\lambda = Z$ & $Z$ & length${}^3$ &
rational & exact (Dalgarno--Lewis solution
$F|1s\rangle = -\tfrac{1}{2} z(r{+}2)|1s\rangle$ has $p$-wave radial
part lying exactly in span$\{S_2, S_3\}$ at $N_\text{basis} = 2$);
$Z^{-4}$ scaling exact at $Z = 1, 2, 3, 4$.  Bound-states-only sum
on bare Fock graph plateaus at $\sim 81.25\%$ of $\tfrac{9}{2}$
(Bethe--Salpeter 1957 partition); the missing 18.75\% continuum
content is recovered \emph{exactly} by the Sturmian projection.
\textbf{Cleanest L2-count $= 0$ exact example for
Prediction~\ref{pred:l2_presence}}: depth-3 chain at exact match
because the analytic answer is rational with finite-Sturmian-basis
representation; no Layer-2 input is consumed in the chain. \\
He $2{}^1P \to 1{}^1\!S$ oscillator strength $f = 0.2761$
(Drake handbook; Hylleraas--Eckart Track 5 closure, 2026-05-18) &
Fock $\circ$ Wigner $3j$ $\circ$ vector-photon promotion $\circ$
bipolar harmonic / Schwartz P-state $\circ$ Hylleraas $r_{12}$ &
$Z, \alpha, \omega$ & dimensionless &
$\alpha^2 \cdot \mathbb{Q}[\sqrt{6}, \mathrm{poly}(\beta)]$ &
$f = 0.2705$ vs Drake $0.2761$, residual $-2.02\%$ at full Schwartz
two-channel $\omega_s = 3$, $\omega_p = 2$; sub-mHa energy match
on $E(1^1S)$ and $E(2^1P)$.  Closes the 2-electron contact-density
cliff (\S\ref{sec:autopsy_he_oscillator}); validates internal
multi-focal architecture (\S\ref{sec:proj_bipolar_drake}) at the
$2^1P$ transition level. \\
\midrule
\textbf{Observation status} & & & & & \\
$K = \pi(B + F - \Delta) \approx 1/\alpha$ & Fock $\circ$ Hopf
$\circ$ spinor & $\alpha$ & dimensionless & three-sector mix &
$8.8\!\times\!10^{-8}$ \\
\bottomrule
\caption{The empirical matches catalogue.  Rows are grouped by
projection depth.  ``Trans.\ class'' uses Paper~18's tier names.
``Match'' is the precision against standard known values:
``exact'' or ``exact symbolic'' means equality at machine
precision; numerical percentages are residuals.  The bottom row
documents Paper~2's $\alpha$ Observation explicitly as a
multi-projection coincidence; the combination rule itself is a
numerical observation without first-principles derivation
(Paper~18 \S~III.5; Paper~2 status v2.26.1).}
\label{tab:catalog}
\end{longtable}

\begin{remark}[signature $\ne$ value]
\label{rem:signature_value}
Theorem~\ref{thm:dimensional_consistency}'s signature consistency is
a type-level statement: chains that produce the same signature can
produce wildly different numerical values.  For example, the
hydrogen Bohr level $E_n = -Z^2/(2n^2)$ (Fock alone, signature
$\{$energy, no~$\pi\}$) and the hydrogen Lamb shift (Fock $+$
spectral action $+$ Sturmian $+$ Drake--Swainson, signature
$\{$energy, $\pi\}$) both have dimension $\{$energy$\}$, yet differ
by 13 orders of magnitude.  The depth-prediction of
\S~\ref{sec:prediction} addresses the value-level question;
signature arithmetic addresses only the type-level question.
\end{remark}

% -------------------------------------------------------------------
\subsection{Off-precision matches and error-source classification}
\label{sec:offprecision}
% -------------------------------------------------------------------

The catalogue above records matches at machine precision or stated
sub-percent agreement.  Equally informative is the catalogue of
matches that are \emph{off} precision, classified by the source of
residual error.  We propose five distinct error sources, each with
its own scaling behaviour and remediation pathway:

\begin{itemize}
\setlength\itemsep{0pt}
\item \textbf{Truncation (T):} finite $n_\text{max}$, $\ell_\text{max}$,
or basis size $N$.  Residual scales as a power law in the cutoff;
remediated by extending the cutoff.
\item \textbf{Basis quality (B):} structural limit of the chosen
angular or radial basis at fixed cutoff.  E.g.\ PK $\ell_\text{max}=2$
optimal operating point for composed LiH; higher $\ell$ degrades.
Remediated only by changing basis architecture.
\item \textbf{Approximation order (A):} leading-order treatment
omitting higher-order perturbative terms.  Residual is a known
fraction of the standard physics (e.g.\ one-loop QED carries
$\sim 5\%$ residual from $\alpha^5$ multi-loop).
\item \textbf{Calibration mismatch (C):} chained projections where
the intermediate calibration constant for the wrong diagram
topology is applied.  Confirmed in the Schwinger $\alpha/(2\pi)$
graph-native sprint (VP-1, VP-2): the vacuum-polarization projection
constant $C_{\text{VP}}$ does not transfer to vertex correction.
\item \textbf{Structural floor (S):} the framework cannot represent
the relevant physics at the chosen layer.  E.g.\ scalar Fock graph
cannot represent vector-photon QED selection rules without
projection promotion.  Remediated by adding the missing projection
to Layer~2.
\end{itemize}

\begin{table}[h]
\centering\small
\begin{tabular}{p{3.8cm} p{2.5cm} p{1.5cm} p{1.4cm} p{4.5cm}}
\toprule
Match & Projection(s) & Residual & Source & Notes \\
\midrule
He ground state (graph-native CI, $Z=2$) & Fock + Sturmian &
$0.20\%$ at $n_\text{max}=9$ & S & Small-$Z$ graph-validity-boundary
artifact (Sec.~2 of CLAUDE.md, v2.9.2 CUSP-2). \\
He ground state (2D variational + cusp) & Fock + spectral cusp &
$0.019\%$ at $\ell_\text{max}=7$ & T & Track DI Sprint 1; cusp-corrected
$0.004\%$ at $\ell_\text{max}=4$. \\
He ground state (Hylleraas r$_{12}$, $\omega_{\max}=4$) & Fock $\circ$
explicit r$_{12}$ correlation
& $-0.016$~mHa = $0.0006\%$ vs Drake NR at 22 basis functions & T &
Track 1 sprint 2026-05-09 (\texttt{geovac/hylleraas\_r12.py}); cusp ratio
$0.49$ vs Kato $0.5$ at $u=0$ (variationally captured); 3p reproduces
Hylleraas 1929 $-2.90324$~Ha to $+0.79$~mHa.  Ground-state architecture
\emph{categorically distinct} from graph-native CI (no $\kappa = -1/16$
inter-shell coupling artifact, no small-$Z$ graph-validity-boundary
limitation).  Excited states (2$^1$S, 2$^3$S, 2$^1$P) deferred --- single-$\alpha$
trial inadequate, named structural extension is Hylleraas-Eckart
double-zeta. \\
LiH $R_{\text{eq}}$ (composed, PK) & Fock $\circ$ molecule-frame
& $5.3\%$ at $\ell_\text{max}=2$ & B & PK $\ell_\text{max}=2$ optimal
operating point; higher $\ell$ degrades structurally
(Track BQ, Paper~17). \\
BeH$_2$ $R_{\text{eq}}$ & composed + exchange & $11.7\%$ & B &
Same PK structural ceiling as LiH; full 1-RDM exchange.
\\
H$_2$O $R_{\text{eq}}$ & composed + lone pairs & $26\%$ & B &
6{:}1 charge asymmetry overwhelms angular basis at $\ell_\text{max}=2$.
\\
H$_2$ $D_e$ (Level 4) & molecule-frame hyperspherical & $4.0\%$ & T
& $96.0\%$ of experimental $D_e$ at $\ell_\text{max}=6$; CBS
$\sim 97\%$. \\
LiH 4-electron FCI energy & balanced coupled $\circ$ Fock &
$0.20\%$ at $n_\text{max}=3$ & T & Improvement $1.8\% \to 0.20\%$
from $n_\text{max}=2 \to 3$; structural drift in $R_{\text{eq}}$
remains. \\
Schwinger $\alpha/(2\pi)$ from
graph-native vertex & vector-photon $\circ$ graph-native & $33\%$
level overshoot & C & Wrong projection constant; $C_{\text{VP}} \times
F_2$ diverges (sprint VP-1; CLAUDE.md graph-native QED bullet). \\
Anomalous moment $C \times F_2$ asymptotic & graph-native &
$12.7\times$ overshoot at $n_\text{max}=3$ & C & Vector promotion
gets $3.6\times$ closer to scalar baseline but does not close
projection mismatch (sprint VP-1). \\
Bethe logarithm $\ln k_0(2S) = 2.812$ & Sturmian (acc.) $\circ$ Fock &
$-0.92\%$ at $N=40$ & T & Acceleration form $3.3\times$ tighter than
LS-2 velocity at $N=50$; sprint LS-3. \\
Bethe logarithm $\ln k_0(1S) = 2.984$ & Sturmian (acc.) $\circ$ Fock &
$+0.60\%$ at $N=20$ & T & Acceleration form crosses Drake at smaller $N$
than velocity; sprint LS-3. \\
Bethe logarithm $\ln k_0(2P) = -0.0307$ & Drake--Swainson $\circ$
Sturmian $\circ$ Fock & $+2.4\%$ at $N=40$ & T & \textbf{LS-4 closes
the LS-3 structural floor.} New Drake--Swainson projection
(\S~\ref{sec:proj_drake_swainson}) produces $K$-independent value via
$\ln k_0 = J_v / D_{\text{drake}}$ with structural denominator
$D_{\text{drake}}(n,\ell) = 2(2\ell+1)Z^4/n^3$.  Convergence
monotonic: $+2.4\%$ at $N=40 \to +1.24\%$ at $N=55$. \\
Bethe logarithm $\ln k_0(3P) = -0.0406$ & Drake--Swainson $\circ$
Sturmian $\circ$ Fock & $+6.3\%$ at $N=40$ & T & Slowest-converging
$\ell>0$ case; sprint LS-4. \\
Lamb shift $1057.85$ MHz & Fock $\circ$ spectral action $\circ$
Sturmian (acc.) $\circ$ (Drake--Swainson, l=1) & $-0.39\%$ at LS-4
sweet spot $N=40$ & T cancels A & At $N=40$, $T_{2S} = +28.8$~MHz
truncation cancels $A = -32.78$~MHz one-loop ceiling; converges to
LS-1 baseline $-3.10\%$ at $N\to\infty$.  Sprint LS-4. \\
Lamb shift $1057.85$ MHz (LS-3 baseline) & Fock $\circ$ spectral action
$\circ$ Sturmian (acc.) & $-2.77\%$ (LS-3 prediction $1029$ MHz) &
\textbf{C} (resolved by LS-6a) $+$ A & Sprint LS-6a (May 2026)
re-derived the same one-loop physics in the Eides \S~3.2 canonical
convention and found the LS-1 implementation had inadvertently
subtracted the Uehling kernel constant from the canonical $10/9$ SE
coefficient (Uehling already enters separately as a VP contribution).
The convention fix supersedes this row: see the LS-6a row below for
the clean result. The $\sim +24.7$~MHz residual that had been
classified A is structurally a $\sim +27.13$~MHz one-loop convention
correction (\textbf{C} resolved), leaving the genuine A-tier
multi-loop residual of $\sim +5.65$~MHz. \\
Lamb shift $1057.85$ MHz (LS-6a, Eides \S~3.2 convention) & Fock
$\circ$ spectral action $\circ$ Sturmian (acc.) $\circ$ (Drake--Swainson,
$\ell=1$) & $-0.534\%$ (LS-6a prediction $1052.19$~MHz) & A (small) +
non-loop sectors (reframed by LS-7) & Eides \S~3.2 canonical one-loop
SE convention.  Convention shift from LS-1: $+27.13$~MHz, exactly
$(4/15) \cdot \alpha^3 Z^4 / (\pi n^3) \cdot \text{Ha-to-MHz}$ at
$n=2$, $Z=1$ (the Uehling kernel constant $10/9 - 38/45 = 4/15$).
Components: SE 2$S_{1/2}$ $+1066.44$~MHz (Bethe $+953.40$~MHz +
canonical $10/9$ piece $+113.04$~MHz); VP 2$S_{1/2}$ $-27.13$~MHz;
SE 2$P_{1/2}$ $-12.88$~MHz; VP 2$P_{1/2}$ $0$.  $5.8\times$
reduction in absolute error vs LS-3.  \textbf{Sprint LS-7 reframed
the residual decomposition}: the LS-5 / LS-6a attribution to
``multi-loop QED in Eides Tab.~7.4 band $[+3.5, +10.7]$~MHz'' was a
misreading of which Eides table is the proper multi-loop reference.
Eides Tab.~7.3 ($\alpha^5$ multi-loop QED proper) totals only
$\sim +1.20$~MHz; the remaining $\sim +4.4$~MHz of the $-5.65$~MHz
residual is non-loop physics (recoil $-2.40$, finite nuclear size
$+1.18$, hyperfine averaging $\sim +5.0$, sum $\sim +4.98$).  The
framework's one-loop QED machinery on $S^3$ is therefore even more
complete than originally framed: the multi-loop test target for
Paper~35 Prediction~1 is sharpened to $\sim +1.20$~MHz, with the
non-loop sectors covered by the Paper~34~\S\,III.14 rest-mass
projection and Paper~23's nuclear-electronic composed Hamiltonian
(not a test of Paper~35).  Sprint provenance:
\texttt{debug/ls6a\_eides\_convention\_memo.md},
\texttt{debug/ls7\_two\_loop\_se\_memo.md}. \\
Two-loop SE prefactor $(\alpha/\pi)^2 (Z\alpha)^4 / n^3$ on
Dirac-$S^3$ (LS-7, structural) & Fock $\circ$ spectral action
$\circ$ Sturmian (acc.) $\circ$ Camporesi--Higuchi $\circ$ iterated
proper time & structural prefactor confirmed to symbolic exactness;
dimensionless bracket $C_{2S} = +3.63$ used from Eides literature
(LS-8a target for native derivation) & A (sharpened) & Iterated CC
spectral action on Dirac-$S^3$ produces the standard two-loop SE
prefactor with the $1/\pi^2$ tracing to two iterated Schwinger
proper-time integrations, confirming Paper~35 Prediction~1 (the
$\pi$-as-temporal-projection prediction) \emph{at the structural
level} for two-loop content.  Weak test (the same prefactor appears in
flat-space two-loop QED); the strong test is LS-8a's native
derivation of $C_{2S}$.  No structural obstruction encountered;
\texttt{geovac/qed\_two\_loop.py} infrastructure is adequate.
Sprint provenance: \texttt{debug/ls7\_two\_loop\_se\_memo.md}. \\
He $2{}^3P_1$--$2{}^3P_2$ small interval $+2{,}231.16$~MHz (Drake-pattern
Breit--Pauli + spinor; Sprint~3 BF-D + Sprint~4 DD + Sprint~5 CP) &
spinor lift $\circ$ Wigner $3j$ ($6j$-coupled, Drake) $\circ$ Wigner $3j$
(Gaunt) $\circ$ Fock &
$-2.62\%$ vs NIST $+2{,}291.176(15)$~MHz & A & Partial-cancellation
amplification: absolute residual $-60$~MHz on a $2.3$~GHz interval.
The dominant intervals (P$_0$-P$_1$ at $-0.014\%$, P$_0$-P$_2$ span at
$-0.20\%$) are in \S~\ref{sec:matches}; this small interval has the
same absolute residual but $\sim$10$\times$ larger fractional residual
because of the structural cancellation SO ($-58$~GHz) $+$ SS ($-9.5$~GHz)
$+$ SOO ($+70$~GHz) $\approx +2.2$~GHz.  Residual is below the LS-8a
$\alpha^3/\pi$ one-loop QED budget on the absolute scale; multi-electron
SS/SOO higher-order corrections (Drake 1971 \S~IV) are the standard
remediation.  Sprint precision-catalogue 2026-05-08.  Convention exposure
on the relativistic Bethe-log sub-component (Drake 1990s vs Pachucki--Yerokhin
2009) catalogued in V.D.6 (\S\ref{sec:conv_he_bethe_log}; V.D-prediction
sprint, 2026-05-10): the Drake-vs-PY 3$\sigma$ disagreement on $\nu_{01}$
sits inside the same $\alpha^3(Z\alpha)^2$ multi-loop budget that this
row attributes the framework's residual to. \\
Li 2$^2$P doublet (Sprint 5 CP) & Breit--Pauli + Slater $Z_{\text{eff}}$
& $+8.9\%$ & A & Standard convention with $Z_\text{val}=1$;
core polarization worsens accuracy (parity-forbidden
$2s2p \leftrightarrow 2p^2$ mixing confirmed). \\
21~cm $A_\text{HF} = 1418.84$ MHz (HF-1 reduced-mass, $+$ external $g_p$)
& spinor lift (Fermi-contact vertex) $\circ$ Fock $\circ$ rest-mass
projection ($m_p$, proton) & $-1102$ ppm (= $-a_e$) & A & Bohr-Fermi level;
$g_e = 2$ (no anomalous moment).  Residual is exactly $a_e \approx
\alpha/(2\pi)$, addressed by HF-2.  Sprint HF, 2026-05-07. \\
21~cm $A_\text{HF} = 1420.488$ MHz (HF-2, $+a_e$ Schwinger, $+$ external
$g_p$) & HF-1 chain $+$ graph-native $F_2$ on Dirac-$S^3$ & $+58$ ppm & A
& Schwinger
asymptote $a_e = \alpha/(2\pi)$ derived structurally on Dirac graph;
\textbf{Parker--Toms first-order curvature correction $c_1 = R/12 = 1/2$
verified at $0.5\%$ precision against the spectral sum}, no fits.  See
Paper~28 \S\ref{sec:curvature_coefficients}.  Residual $+58$ ppm is
recoil + Zemach + multi-loop. \\
21~cm $A_\text{HF} = 1420.432$ MHz (HF-4, $+r_Z$ external) & HF-2 chain
$+$ Zemach Layer-2 input ($r_Z = 1.045$ fm) & $+18$ ppm & A + S(Zemach)
& Convention catch: standard form is $\Delta\nu/\nu = -2 Z\alpha m_e r_Z$
(natural units), $-2 Z r_Z[\text{bohr}]$ (atomic units), giving
$-39.5$ ppm at $r_Z = 1.045$ fm.  Framework has no
magnetization-distribution operator on the proton register
(structural floor S); $r_Z$ is external focal length.  Final $+18$ ppm
matches Eides Tab.~7.3 attribution to Bodwin--Yennie recoil-beyond-
reduced-mass + multi-loop QED + nuclear polarizability. \\
Recoil from cross-register cross-spatial coupling (HF-3) & Track NI
nuclear-electronic & not derivable & S & \texttt{finite\_size\_correction}
returns scalar at classical \texttt{R\_PROTON\_BOHR}; 0/71 Pauli strings
couple spatial qubits across registers; the 12 cross-register terms are
all spin-spin hyperfine.  Framework couples discrete spin labels but
not spatial coordinates.  Sprint HF, 2026-05-07. \\
Multi-loop $a_2$ on Dirac-$S^3$ (HF-5) & iterated CC spectral action
$\circ$ Camporesi--Higuchi & $a_2$ extraction blocked & S (renormalization)
& Bare two-loop topologies (VP-on-photon, SE-on-electron, crossed) UV
diverge as $\sim N^{3.79}$.  Framework reproduces $(\alpha/\pi)^2$
prefactor, three-topology decomposition, SO(4) selection at all four
vertices, UV degree count, but does not autonomously generate $Z_2/Z_3/
\delta m$ counterterms.  Same structural wall as LS-8a (Lamb shift),
vertex sector. \\
$\mu$H 1S HFS $\nu_\text{HFS}(\mu p) = 181.32$ meV (BF + Schwinger +
leading Zemach, Sprint MH Track B) & rest-mass projection $\circ$
spinor $\circ$ Fermi contact $\circ$ Fock; magnetization-density
Layer-2 input $r_Z = 1.045$~fm & $-7710$ ppm vs Krauth full theory
182.725 meV & A (one-loop closure only) + S (LS-8a wall, muonic regime)
& Dominant residual is electron-VP (Uehling) in the muonic potential,
the leading muonic QED correction (Eides Tab.~7.4 / Karshenboim 2005:
$\Delta_\text{el-VP} \approx +1.5$ meV).  Same LS-8a wall as HF-5; in
muonic atoms the electron loop is enhanced because the muonic Bohr
radius probes the electron Compton wavelength.  Inner-factor input
data tier per Paper~18 \S~IV.6.  Sprint MH Track B, 2026-05-08. \\
Ps 1S HFS $116.79$~GHz (Fermi-contact only, framework-native) & spinor
lift (Fermi-contact vertex) $\circ$ Fock $\circ$ rest-mass projection
($m_p \to m_e$, equal-mass limit) & $-42.58\%$ vs Ishida 2014
$203.389$~GHz & A & Framework-native part captures only the magnetic
dipole-dipole contribution ($4/12$ of the LO HFS).  The
remaining $3/12$ is the annihilation channel $e^+e^- \to \gamma \to e^+e^-$,
a virtual-loop second-quantized effect not generated by the framework's
bare action -- structurally analogous to the K\"all\'en--Sabry two-loop
VP input for muonic H. Multi-focal architecture at $\lambda_a = \lambda_b$
verified: Roothaan multipole termination at $L_\text{max} = 2 \ell_\text{max}$
preserved (Gaunt selection is angular, not small-parameter).
Sprint precision-catalogue, 2026-05-08. \\
Ps 1S HFS $204.39$~GHz (LO total $= 7/12 \cdot m_e c^2 \alpha^4 / h$,
framework $+$ Layer-2 annihilation) & spinor lift (Fermi-contact vertex)
$\circ$ Fock $\circ$ rest-mass projection $+$ Layer-2 annihilation &
$+0.49\%$ vs Ishida 2014 & A & Residual is multi-loop QED
(Czarnecki--Melnikov--Yelkhovsky 2000, $\alpha^5 m_e c^2/h$ corrections).
LS-8a wall confirmed in the equal-mass leptonic regime: same
renormalization gap as $\mu$H Track A and HF-5.  Sprint precision-catalogue,
2026-05-08. \\
D 1S hyperfine $\nu_\text{HFS}(D) = 327.4779$ MHz (BF $+$ recoil
$+$ Schwinger $a_e$ $+$ leading Zemach with $r_Z(D) = 2.593$~fm
external; sprint precision-catalogue 2026-05-08) & spinor $\circ$ Fermi
contact $\circ$ Fock $+$ reduced-mass $+$ vertex correction
$+$ magnetization-density Layer-2 input
& $+286$ ppm vs Wineland--Ramsey 1972 $327.384352522(2)$ MHz & A
(framework-native) $+$ S (LS-8a wall, multi-loop QED) $+$ S (deuteron
polarizability, QCD-internal NN dynamics)
& Strict Bohr-Fermi at $+40$ ppm verifies $I=1$ cross-register angular
coupling structure: $A_{hf}\,\hat{I}\cdot\hat{S}$ Hamiltonian unchanged,
$3/2$ multiplicity is Clebsch--Gordan, no new mechanism vs $I=1/2$.
Cumulative residual $+286$ ppm budgeted to deuteron polarizability
($+44$ ppm Pachucki--Yerokhin 2010, QCD-internal), multi-loop QED
($\sim$few ppm, LS-8a wall), recoil NLO, finite-size charge.
Operator-level Zemach $-98.0$ ppm with $r_Z(D) = 2.593$~fm matches
analytic $-2 Z m_e r_Z$ to $10^{-14}$ ppm; profile (Gaussian vs
exponential) independence preserved at leading order, same as Sprint MH
Track C / hydrogen 21cm (Sprint HF Track 4).  Multi-focal architecture
verified at I=1: deuteron quadrupole moment couples only to gradients,
sub-ppm for s-states (negligible).  Sprint precision-catalogue,
2026-05-08.  Operator-level four-component Roothaan autopsy in
\S\ref{sec:autopsy_d_hfs} (May~9 2026): $\hat{I}\cdot\hat{S}$
multiplicity ratio D/H $= 3/2$ at machine precision (12 digits);
\S\ref{sec:proj_magnetization_density} reproduces Eides analytic LO
to $\sim 1.45 \times 10^{-14}\%$ of LO shift; cumulative $+285.6$~ppm
chain residual matches PY~2010 Layer-2 budget magnitude.  Polarizability
sub-component convention exposure (Friar--Payne~2005 vs Pachucki--Yerokhin
2010+) catalogued in V.D.5 (\S\ref{sec:conv_d_polarizability}; V.D-prediction
sprint, 2026-05-10), refining \S\ref{sec:conv_rZG_DHFS}. \\
T 1S hyperfine $\nu_\text{HFS}(T) = 1516.6969$ MHz (BF $+$ Schwinger
$a_e$ $+$ recoil $+$ leading Zemach with $r_Z(t) = 1.762$~fm external;
multi-track sprint Track 2, 2026-05-18) & Fock $\circ$ spinor (Fermi-contact
vertex) $\circ$ spectral action ($a_e$) $\circ$ rest-mass projection
($m_p \to m_t$) $\circ$ magnetization-density Layer-2 input
& $-3.0$ ppm vs Greene 2017 $1516.701470773(8)$ MHz & B
(basis-quality: Zemach scalar from Carlson~2008 $r_Z(t) = 1.762$~fm,
convention drift $\pm 5$~ppm Carlson vs Sick~2014; sub-ppm absolute
LS-8a $+$ W1a $+$ W3 budget) & \textbf{Cleanest LS-8a multi-loop QED
isolation in the hadronic-nucleus catalogue.}  Tritium polarizability
$\sim +0.7$~ppm is $\sim 300\times$ smaller than deuteron's
$\sim +200$~ppm because the 2n+1p binding ($B(t) = 8.482$~MeV) is
much tighter than the deuteron's n+p binding ($B(d) = 2.225$~MeV),
making the tritium nucleus essentially compact and rigid.  T fills
the I=1/2 hadronic non-proton cell of the mass $\times$ spin coverage
matrix that was previously empty.  Operator-level autopsy in
\S\ref{sec:autopsy_t_hfs}: $\hat{I}\cdot\hat{S}$ structurally
identical to H~21cm (3 Pauli terms, $X{\otimes}X + Y{\otimes}Y +
Z{\otimes}Z$); \S\ref{sec:proj_magnetization_density} reproduces
Eides analytic LO to bit-identical machine precision under both
Gaussian and exponential profiles; T/H Zemach ratio matches T/H
radii ratio at machine precision.  Cumulative chain residual sits
inside the projected Karshenboim 2005 Layer-2 budget $+12 \pm 5$~ppm
(LS-8a $\sim$$+6$ ppm $+$ W1a recoil NLO $\sim$$+3$ ppm $+$ W3
polarizability $\sim$$+0.7$ ppm $+$ finite-size $\sim$$+1.5$ ppm,
sum $\sim$$+12$ ppm, framework gives $-3$~ppm within convention
band).  No new convention mismatch surfaced beyond the well-known
$r_Z(t)$ value drift (Carlson 2008 vs Sick 2014, $\pm 5$~ppm). \\
Mu 1S-2S transition (test of rest-mass projection at unchanged lepton)
& Fock $\circ$ rest-mass projection (m$_p \to m_\mu$) & $-0.11$ ppm vs
Mu-MASS & A & Pure $m_\text{red}$ rescaling from H experimental;
sub-ppm verification of the rest-mass projection theorem (Paper~34
\S~III.14) at varied nucleus mass.  Differential-recoil contribution
(framework + cross-register $V_{eN}$ Roothaan) closes to $-0.25$ ppm.
Sprint precision-catalogue, 2026-05-08.  \\
$\mu$H 2$S$--2$P$ Lamb framework-native $200.50$ meV (no literature
inputs, Sprint MH Track A) & Fock $\circ$ spectral action $\circ$
Sturmian $\circ$ Drake--Swainson $\circ$ rest-mass projection
($m_e \to m_\mu$) & $-0.92\%$ vs CREMA $202.3706$ meV & A + S
(LS-8a wall, muonic regime)
& Framework reproduces dominant Uehling at $<1$ ppm vs Antognini and
Friar moment leading-order at $4.3\%$.  $-0.92\%$ residual decomposes
exactly into missing K\"all\'en--Sabry two-loop VP ($+1.51$ meV),
multi-loop ($+0.19$ meV), recoil next-to-leading ($-0.045$ meV), and
nuclear polarizability ($+0.013$ meV).  Same LS-8a renormalization
wall as HF-5 and Sprint MH Track B, vertex sector with electron loop
in muonic potential.  Plus literature-input total: $-0.10\%$ vs CREMA.
\textbf{Headline architectural result}: full Uehling kernel integration
(rather than contact-density approximation) handles the muonic
$\beta = 2 m_e a_\mu = 1.475$ regime where Bohr radius $\sim$ electron
Compton wavelength.  Sprint MH Track A, 2026-05-08. \\
Mu 2$S_{1/2}$-2$P_{1/2}$ Lamb $1047.63$ MHz vs Mariam~1982 experiment
$1054(22)$ MHz & Fock $\circ$ spectral action $\circ$ Sturmian $\circ$
rest-mass projection & $-0.6\%$ (within $\pm 22$ MHz exp uncertainty)
& A & Framework-native at $+0.013\%$ vs Karshenboim~2005 theory
$1047.49(5)$~MHz; $-0.6\%$ residual against Mariam dominated by 1982-era
experimental uncertainty, not framework limitation.  Karshenboim's
``Recoil NLO $-11.30$ MHz'' is absorbed into the framework's
$m_\text{red}$ Eides bracket (the $+1074.13$ MHz framework SE differs
from Karshenboim's $+1085.84$ MHz SE itemization by exactly the recoil
NLO).  Naively adding Karshenboim's recoil to the framework total
double-counts.  See \S~\ref{sec:convention_exposures} for catalogue
context (entry \S~\ref{sec:conv_karshenboim_recoil}).
Sprint precision-catalogue, 2026-05-08. \\
He$^+$ 2$S_{1/2}$-2$P_{1/2}$ Lamb $14041.13(17)$ MHz (van
Wijngaarden--Holuj--Drake 2000) & Fock $\circ$ spectral action $\circ$
Sturmian $\circ$ (Drake--Swainson, $\ell=1$) $\circ$
\S\ref{sec:proj_charge_density} & $-1.29\%$ (framework $+$ Drake-Yan
Layer-2 lit total $13\,859.58$ MHz; framework-native $13\,841.11$ MHz at
$-1.43\%$) & A & \textbf{Z-scaling test of Paper~36 architecture
(Z=1 $\to$ Z=2)}.  Framework's projection-chain machinery $Z$-portable:
three Z-scaling structural tests pass (Lamb$_\text{SE}$ ratio $13.22$
matches Eides bracket prediction $13.18$; Lamb$_\text{VP}$ ratio $15.88$
matches $Z^4 = 16$ contact-density; Lamb$_\text{FNS}$ ratio $63.60$
bit-exact match to $Z^4 \times (r_E({}^4\text{He})/r_p)^2 = 63.52$).  The
$-1.43\%$ framework-native residual is structurally the Z-amplified analog
of Paper~36's $-0.55\%$ at Z=1: the $\alpha^3$ Eides-bracket truncation
amplifies by $Z^4 \times \ln$-factor to $\sim -200$ MHz at Z=2.  Dominant
Layer-2 input is Drake--Yan 1992 $\alpha(Z\alpha)^6 \ln(Z\alpha)^{-2}$
relativistic-Bethe-log correction ($+14.7$~MHz), below precision visibility
at Z=1 but structurally visible at Z=2.  Cleanest W3 inner-factor in
catalogue: ${}^4$He polarizability $+0.06$~MHz / $4 \times 10^{-9}$
relative (tightly bound nucleus suppresses virtual photonuclear
excitations).  $I({}^4\text{He}) = 0$ identically zeroes the
\S\ref{sec:proj_magnetization_density} channel (no HFS, no Zemach).
Autopsy: \S\ref{sec:autopsy_he_plus_lamb}.  Source memo:
\texttt{debug/he\_plus\_lamb\_autopsy\_track\_a\_memo.md}.  Sprint:
2026-05-19. \\
Mu 1S hyperfine framework Bohr--Fermi $+$ Schwinger one-loop
$4464.1905$ MHz vs Liu~1999 $4463.302765(53)$ MHz & spinor lift
(Fermi-contact vertex) $\circ$ Fock $\circ$ rest-mass projection
(lepton-unchanged, nucleus $m_p \to m_\mu$) & $+199$ ppm & A &
\textbf{Cleanest LS-8a isolation in the catalogue}: leptonic nucleus
(point antimuon, no QCD) means residual is purely multi-loop QED $+$
cross-register recoil NLO, no Zemach / polarizability / hadronic VP
budget interfering.  Karshenboim 2005 itemized total of corrections
($+4.703$ MHz, $\sim +1054$ ppm relative to BF strict) decomposes as
$\alpha^2(Z\alpha) \sim +0.17$ MHz (two-loop QED), $\alpha(Z\alpha)^2(m/M)^2
\sim +0.7$ MHz (recoil $+$ QED), hadronic VP $\sim +0.2$ MHz, plus the
dominant Schwinger one-loop already captured by the framework.  The
$+199$ ppm framework residual is the structural depth of the LS-8a wall
at $\alpha^2(Z\alpha)$ for Bohr--Fermi-class observables in clean
isolation: no QCD competing budget, both leptons receiving full
$\alpha/(2\pi)$ at one loop.  Adding Karshenboim's full Layer-2
itemization closes the gap to $<1$~ppm.  Sprint precision-catalogue,
2026-05-08. \\
Ps 1S-2S framework-native $+64.75$ ppm vs Fee~1993
$1{,}233{,}607{,}216.4(3.2)$~MHz (sprint precision-catalogue 2026-05-08,
equal-mass leptonic 1S-2S) & Fock $\circ$ rest-mass projection $\circ$
spectral action (Eides) & $+64.75$ ppm & A & \textbf{Sixth instance of
the multi-focal-composition wall, and the cleanest known.}  At equal
mass $\lambda_a = \lambda_b$ the standard $1/M_\text{nucleus}$ recoil
expansion is invalid; the dominant residual is the $\alpha^4$ Breit /
two-body-Dirac correction at the $\sim 80$~GHz scale, NOT multi-loop
QED ($\sim$10 MHz, three orders of magnitude smaller).  Same structural
wall as Sprint H1 Yukawa, LS-8a $Z_2/\delta m$, HF-3 recoil, HF-4
Zemach, HF-5 multi-loop $a_e$: framework couples discrete labels via
Wigner symbols cleanly, does not natively compose two Fock-style
projections at equal mass.  Cleanest known instance because the failure
surfaces at $\alpha^4$ -- one full order before the LS-8a renormalization
wall -- isolating two-body-projection failure from renormalization
failure.  This residual is the empirical anchor for the projection family
in \S~\ref{sec:proj_breit_retardation} (two-body Dirac / Breit
retardation, the sixteenth projection, added 2026-05-08 motivated by
this match).  Layer-2 $\alpha^4$ Breit input from Penin--Pivovarov 1998
(PRL~80, 2101) closes the chain to sub-MHz; cumulative match is partly
explained-by-construction (Breit value inferred as residual).  Honest
verdict: framework-native is at $+65$~ppm; all closure beyond Bohr$+$SE
Lamb is Layer-2.  Sprint precision-catalogue, 2026-05-08; operator-level
five-component autopsy (\S\ref{sec:autopsy_ps_1s2s}, 2026-05-09)
locates the wall at the kernel-vs-bound-state boundary
of \S~\ref{sec:proj_breit_retardation}, not at the kernel level itself
(closed-form $\mathbb{Q}[\log p]$ matrix elements verified;
state-specific Bethe--Salpeter evaluation is the named structural
extension). \\
He $2{}^1S_0$--$2{}^3S_1$ exchange splitting $7183.3$~cm$^{-1}$ vs NIST
$6421.47$~cm$^{-1}$ (sprint precision-catalogue 2026-05-08, ss-only
sub-block at $n_\text{max}=11$; refreshed 2026-05-08 with fixed
\texttt{compute\_rk\_float} evaluator, trajectory bit-identical to
displayed 0.1~cm$^{-1}$ precision, largest shift 0.05~cm$^{-1}$) &
graph-native CI $\circ$ Fock &
$+11.86\%$ & B & \textbf{First multi-electron correlation-driven catalogue
entry.}  No relativistic, no QED, no cross-register content; pure $V_{ee}$
exchange in the angular momentum eigenbasis.  Splitting is \emph{less}
accurate than absolute energies (34$\times$ ratio: $\sim 12\%$ on splitting
vs $\sim 2.7$--$2.8\%$ on individual states) -- common-mode errors on
$2^3S$ and $2^1S$ have \emph{opposite} phase impact, no cancellation.
\textbf{Variational bound violated for both excited states} (not for
$1^1S$): graph $\kappa = -1/16$ adjacency adds constant inter-shell
coupling regardless of orbital extent, overbinding loosely-bound
Rydberg-like states.  Same mechanism as the small-$Z$
graph-validity-boundary artifact ($Z_c \approx 1.84$, He at $Z=2$).
Single-exponent hydrogenic basis at $k = Z = 2$ cannot represent
$Z_\text{eff} \approx 1$ screening of 2s; multi-focal composition
(Phase~C-W1c, May~2026) is the structural extension.  Sprint
precision-catalogue, 2026-05-08.
\textbf{Hylleraas r$_{12}$ Track 1 attempt (2026-05-09):} single-$\alpha$
Hylleraas explicit-correlation trial in
\texttt{geovac/hylleraas\_r12.py} does NOT close this residual --- the
trial gives splitting $\sim 19{,}900$~cm$^{-1}$ at $\omega_{\max}=5$
($+209\%$ vs NIST), worse than the graph-native CI baseline.  The
2$^3$S excited state goes BELOW Drake by $\sim 50$~mHa due to incomplete
representation of the asymmetric Z$_\text{eff}$ structure
($1s$ at $Z_a\!\approx\!2$, $2s$ at $Z_b\!\approx\!1$); a single-$\alpha$
trial $e^{-\alpha s}$ cannot represent two distinct exponential decay
scales without a $\beta t$-dependent factor.  The structural fix is
the standard \textbf{Hylleraas--Eckart double-zeta} extension (1933);
catalogue residual stays at $+11.86\%$ pending that extension.
Source: \texttt{debug/hylleraas\_r12\_track1\_memo.md} \S3. \\
Self-consistent $r_Z$ global extraction (Sprint Calc-rZG, 2026-05-09;
Layer-2 corrected 2026-05-09; extended-fit diagnostic 2026-05-09)
& §III.18 magnetization-density global fit across H 21cm $+$
$\mu$H 1S HFS $+$ D 1S HFS & \emph{Original 3-obs fit} (Layer-2 corrected,
$\sigma_{\mu H}$=1500 ppm, ``LS-8a wall scale''): $r_Z(p) = 1.180(166)$ fm
[$+0.80\sigma$ above Eides $1.045(20)$, $+1.00\sigma$ above lattice
$1.013(10)(12)$], $r_Z(D) = 2.583(1588)$ fm [$-0.01\sigma$ from
Friar--Payne $2.593(16)$], $\chi^2/\text{dof} = 0.43/1$.
\emph{Extended-fit diagnostic} (sprint Calc-rZG-extended, $\sigma_{\mu H}$
hardened to literature kernel scale ${\sim}367$ ppm): $r_Z(p) =
1.257(51)$~fm, $\sigma$ closes to $51$~mfm but extracted central
value sits $+3.85\sigma$ above Eides and $+4.55\sigma$ above lattice.
& C & \textbf{Framework $r_Z(p)$ extraction is currently kernel-limited,
not Layer-2-itemization-limited}.  Original 3-obs fit at $\sigma=166$~mfm
(driven by $\sigma_{\mu H}=1500$~ppm conservative) lands within $1\sigma$
of all three literature determinations after Layer-2 budget correction
(2026-05-09 bug fix; original rZG memo specified Layer-2$_D$ = $-150$ ppm
based on a mis-itemization, corrected to $-286$ ppm per Pachucki--Yerokhin
2010 + Friar--Payne 2005 closure analysis; Layer-2$_H$ tightened from
$0$ to $-18.5$ ppm per Eides Tab.~7.3); $\chi^2/\text{dof} = 0.43/1$
confirms internal fit consistency.  \textbf{The original $r_Z(D) = 6.18$
fm artifact was traced to a Layer-2 budget misspecification, NOT to a
recoil double-counting in the production cross-register $V_{eN}$
kernel.}  Diagnostic verification
(\texttt{debug/rzg\_bug\_diagnosis\_compute.py}): the production
\texttt{geovac.cross\_register\_vne.cross\_register\_recoil\_correction}
reproduces leading-order Bethe--Salpeter recoil at $2.86\%$ for
hydrogen, $2.03\%$ for deuterium, and $8.18\%$ for muonium without
any species-specific code path; it is species-agnostic and
correct.  Six new D and muonium tests added to
\texttt{tests/test\_cross\_register\_vne.py} document and protect
this property going forward.  \textbf{Extended-fit diagnostic
finding (2026-05-09 follow-on)}: tightening $\sigma_{\mu H}$ from
1500~ppm to either 150~ppm (Krauth 2017 multi-loop QED itemization) or
to 367~ppm (5\% kernel approximation in muonic systems) drives
$\sigma(r_Z(p))$ below 50 mfm but extracts a central value
$r_Z(p)\sim 1.18$--$1.26$~fm, INCONSISTENT with both Eides
$1.045(20)$ and lattice $1.013(16)$ at $>3.85\sigma$.  The +0.22~fm
offset between the muH-alone extraction ($r_Z=1.265$~fm) and Eides
($1.045$~fm) is the structural fingerprint of recoil-mixing in the
literature Zemach formula that the framework's leading-order Eides
kernel $-2Zm_\text{red}r_Z$ does not yet capture
(Sprint MH Track B \S2.4--2.5 W1b operator-level extension point).
\textbf{Verdict: framework currently \emph{cannot} resolve the
Eides-vs-lattice 30~mfm tension at the operator-level precision of
the leading-order Eides Zemach kernel.}  Path to sub-30~mfm: extend
\texttt{magnetization\_density.py} to include recoil-mixing
(Friar--Payne 2005 \S III treatment), expected to close the +0.22~fm
muH-alone offset and bring $\sigma(r_Z(p))$ to ${\sim}30$~mfm.
\textbf{Genuine prospective contribution}: §III.18-based global fit
operator-level-localizes the ``Layer-2 itemization is the load-bearing
systematic for atomic $r_Z$ extraction'' finding --- a unified
analytical lens not currently produced by standard atomic-physics
compilations.  Note that the diagnostic exposes a SECOND load-bearing
systematic: the kernel approximation order of the framework's W1b
operator dominates over Layer-2 itemization in muonic regime
extraction.  §III.18 itself reproduces Eides analytical
$-2 Z\alpha m_e r_Z$ to floating-point precision.  Source memos:
\texttt{debug/calc\_track\_rZG\_global\_zemach\_memo.md},
\texttt{debug/rzg\_bug\_diagnosis\_memo.md} (bug fix),
\texttt{debug/calc\_track\_rZG\_extended\_memo.md} (extended-fit
diagnostic).  \textbf{Sprint W1b recoil-mixing extension
(2026-05-09 follow-on)}: \S III.18 operator extended with the
NLO recoil-mixing prefactor $m_l/(m_l + m_n)$ (per arXiv:2604.06930
eq.~95) and the Friar moment correction $(1/2)(Z\alpha m_\text{red})^2
\langle r^2\rangle_{(2)}$ (per Friar 1979 / Eides \S6.2), gated behind
opt-in \texttt{include\_recoil\_mixing} flag with backward-compatibility
preserved bit-identical (37 baseline tests, 20 new tests, 57/57 pass).
With NLO kernel enabled, the muH 1S HFS Zemach line at $r_Z=1.045$
fm reproduces Krauth~2017's full-theory line to $\sim 1\%$ (vs $\sim
3\%$ overshoot at LO).  \textbf{rZG-extended-v2 fit result}: with
NLO kernel and Layer-2 re-anchored at the Eides reference,
$\sigma(r_Z(p))$ tightens from 51~mfm to \textbf{16~mfm} (below the
30~mfm Eides-vs-lattice gap), but the central value \emph{remains}
at $r_Z(p) \approx 1.286$~fm --- only $\sim 10\%$ of the v1
$+0.22$~fm offset closes.  \textbf{Diagnostic verdict revised}: the
$+0.22$~fm offset is NOT kernel-approximation; it is a Layer-2
itemization-convention mismatch between Krauth~2017 (used for muH
Layer-2) and Eides~2024 Tab.~7.3 (used for H~21cm Layer-2), with
the convention difference at $\sim 0.01\%$ of $\nu_F$ propagating
to $\sim 25$~mfm in extracted $r_Z(p)$.  Next-step recommendation:
Layer-2 convention reconciliation sprint, NOT further kernel
extensions.  See \S~\ref{sec:convention_exposures} for catalogue
context (entry \S~\ref{sec:conv_rZG_DHFS}).  Source memo:
\texttt{debug/w1b\_recoil\_mixing\_extension\_memo.md}. \\
Hawking temperature $T_H = 1/(8\pi M)$ on Schwarzschild cigar &
Fock $\circ$ observation/temporal-window
(\S~\ref{sec:proj_observation}) at $\beta = 8\pi M$ &
machinery-witness only & C &
Sprint TD Track~4 (\texttt{debug/sprint\_td\_track4\_memo.md},
May~2026) reproduces $\beta = 8\pi M$ from the framework's M1
Hopf-base measure / $\mathrm{Vol}(S^1_\tau)$ mechanism alone (sympy
residual zero on $\beta = 4M\cdot 2\pi$ cigar regularity, $T_H =
\kappa_g/(2\pi)$, lowest bosonic Matsubara mode $1/(4M) = \kappa_g$
surface gravity).  $M$ is external Einstein-equations input, not a
framework calibration mismatch in the usual sense; the row is a
machinery-witness, not a precision match.  Track~D
(\texttt{debug/bisognano\_wichmann\_track\_d\_memo.md}, May~2026)
documents the structural correspondence between the framework's M1
$2\pi$ and the Bisognano--Wichmann modular-flow / Lorentz-boost
orbit period via the standard Wick-rotation chain (Hartle--Hawking
1976~\cite{hartle_hawking1976}; Sewell 1982~\cite{sewell1982};
Bisognano--Wichmann 1976~\cite{bisognano_wichmann1976}).  See
Paper~32~§VIII Remark on the Bisognano--Wichmann reading; Paper~35
\S~VIII Bisognano--Wichmann subsection for the framework-level
identification.  Operator-level
literal identification ($\sigma_{i\cdot 2\pi} = $ identity on the
truncated metric spectral triple) is the principal Track~D
falsifier, a follow-up beyond sprint scale. \\
Unruh temperature $T_U = a/(2\pi)$ on Rindler wedge &
Fock $\circ$ observation/temporal-window
(\S~\ref{sec:proj_observation}) at $\beta_U = 2\pi/a$ &
machinery-witness only & C &
Unruh pendant (\texttt{debug/unruh\_pendant\_memo.md}, May~2026,
a sprint-scale pendant follow-up to Track~D) reproduces $\beta_U = 2\pi/a$
from the framework's M1 Hopf-base measure /
$\mathrm{Vol}(S^1_{\eta_E})$ mechanism alone (sympy residual zero on
$T_U$, lowest non-zero bosonic Matsubara mode equals $a$ --- the Unruh
analog of Hawking's surface gravity --- and lowest fermionic
Matsubara mode equals $a/2 = \pi T_U$).  Wick-rotation construction:\
Rindler $(\eta, \rho)$ wedge metric, $\eta \to i\eta_E$ gives
Euclidean polar $\rho^2\, \mathrm{d}\eta_E^2 + \mathrm{d}\rho^2$ with
conical regularity at $\rho = 0$ forcing $\eta_E$-period $2\pi$;
proper-time period $\beta_U = 2\pi/a$.  $a$ is external
proper-acceleration input, not a framework calibration mismatch in
the usual sense; row is a machinery-witness, not a precision match.
Same structural-correspondence scope as the Hawking row above:\ under
the unified Wick-rotation chain (Hartle--Hawking
1976~\cite{hartle_hawking1976}; Sewell 1982~\cite{sewell1982};
Bisognano--Wichmann 1976~\cite{bisognano_wichmann1976};
Unruh 1976~\cite{unruh1976}), the framework's M1 $2\pi$ on
$S^1_{\eta_E}$ IS the modular-flow / boost-orbit period of the
accelerated observer.  See Paper~32~§VIII
\texttt{rem:bisognano\_wichmann\_reading} (Unruh paragraph) and
Paper~35 \S~VIII \texttt{subsec:unruh\_pendant} for the framework-level
identification.  The principal Track~D operator-level falsifier
($\sigma_{i\cdot 2\pi} = $ identity on the truncated metric spectral
triple for a wedge-restricted state) covers all three
Wick-rotation faces (Hawking, BW, Unruh) simultaneously, since the
chain runs through the same KMS-Tomita-Takesaki algebra in every case;
\textbf{Sprint L1 closure (2026-05-16) lifts the structural correspondence to literal
identification at the operator-system level (Riemannian)} via the
machine-precision $\sigma_{2\pi}(O) = O$ closure at $n_{\max} \in
\{2, 3, 4, 5\}$ (see row below and Paper~32~\S~VIII.F). \\
Modular periodicity $\sigma_{2\pi}(O) = O$ on $\mathcal{T}_{n_{\max}}$ &
\S~\ref{sec:proj_wick_rotation} (Wick rotation) at
$\beta = 2\pi / \kappa_g$;\ four-witness collapse to canonical BW &
$1.8\times 10^{-16}$ at $n_{\max}=2$ rising to
$1.6\times 10^{-15}$ at $n_{\max}=5$ (machine precision) & C &
Sprint L1 (\texttt{debug/l1\_modular\_hamiltonian\_results\_memo.md},
2026-05-16) constructs the modular Hamiltonian $K$ on the truncated
Camporesi--Higuchi spectral triple $\mathcal{T}_{n_{\max}}$ via the
hemispheric-wedge BW-$\alpha$ geometric realization $K = J_{\mathrm{polar}}$
and verifies the period closure $\sigma_{2\pi}(O) = O$ bit-exactly at
every tested $n_{\max} \in \{2, 3, 4, 5\}$ for all four physical
witnesses (BW, HH, Sew, Unruh).  Residual scales as machine precision
$O(\sqrt{\dim \mathcal{H}_{n_{\max}}} \cdot
\varepsilon_{\mathrm{machine}})$;\ cross-witness collapse is bit-exact
(max consistency residual $= 0.0$).  Row is a machinery-witness for
the L1 falsifier closure, not a precision match.  Lifts the
four-witness Wick-rotation theorem from structural correspondence
to literal identification at the operator-system level (Riemannian).
Lorentzian extension to signature $(3, 1)$ via BBB Krein lift is the
named Sprint L2 follow-up (a substantial open program per the
Sprint L0 audit).
See Paper~32~\S~VIII.F and \S~\ref{sec:proj_wick_rotation} honest-scope
paragraph. \\
Bekenstein--Hawking probe:\ scaling of
$S(\rho_W) = -\mathrm{Tr}(\rho_W \log \rho_W)$ at finite $n_{\max}$ on the
hemispheric wedge KMS state $\rho_W = e^{-K_\alpha^W}/Z$ &
\S~\ref{sec:proj_wick_rotation} (Wick rotation) at BW canonical
$\beta = 2\pi$;\ tests whether modular-Hamiltonian closure extends to
BH area-law entanglement &
clean negative on naive $S \sim n_{\max}^2$:\
$S \approx 1.944\, \log(n_{\max}) + 0.565$, $R^2 = 0.99991$ across
$n_{\max} \in \{2, \dots, 7\}$ & C &
Sprint BH-Phase0 (\texttt{debug/bh\_phase0\_diagnostic\_memo.md},
2026-05-22) tests whether the modular-Hamiltonian construction of
Papers~42--46 carries a Bekenstein--Hawking area-law signature at
finite cutoff.  Verdict:\ NO at the wedge KMS state level.  $S(\rho_W)$
scales logarithmically, NOT polynomially.  Structural reading:\
$S(\rho_W) \approx \log(n_{\mathrm{equator}})$ where
$n_{\mathrm{equator}} = n_{\max}(n_{\max}+1)$ is the count of states
with smallest $|m_j| = 1/2$;\ the BW canonical state is effectively
maximally mixed on the lowest-$K_\alpha$ shell with exponentially
suppressed higher-shell contributions.  The slope $\approx 2 =
\dim(\partial W) = \dim S^2_{\mathrm{equator}}$ matches the boundary
dimension, giving a Cardy--Calabrese-flavor log entropy along the
$d = 2$ boundary, NOT a BH area-law $S \sim A/(4G)$.  Asymptotic
limits saturate bit-exactly:\ $S(s \to 0^+) = \log(\dim W)$ and
$S(s \to \infty) = \log(n_{\mathrm{equator}})$ at $n_{\max} = 7$.
The $s$-parameter scan across $s \in \{0.01, \dots, 2\pi\}$ shows
power-law exponents $\alpha \in [0.64, 0.72]$, never close to $2$ or
$3$ --- no temperature regime produces polynomial scaling.  Structural
implication:\ the operator-system truncation preserves the four-witness
Wick-rotation theorem (period closure, KMS, $J^2$ signatures) but does
NOT preserve continuum QFT vacuum entanglement.  The Bekenstein--Hawking
area law and the modular-Hamiltonian period closure are structurally
distinct features;\ the framework captures the latter and not the
former.  Consistent with the structural-skeleton-scope pattern
(Sprint H1, LS-8a, HF-3/4/5, multi-focal wall, W3).  Named follow-on
candidate:\ \texttt{BH-Phase0-route-A} (sprint-scale)
testing the reduced density matrix of the GeoVac Dirac ground state
on the same wedge.  The original May-16 BH scoping (Connes--Chamseddine
spectral action on Euclidean Schwarzschild cigar with one-parameter
$\Lambda^2 \leftrightarrow 1/G$ calibration) remains the longer-route
alternative (beyond sprint scale).  See Paper~32~\S~VIII
\texttt{rem:bh\_phase0\_negative} for the spectral-triple-side reading. \\
HD molecule $J=1$ deuteron quadrupole coupling
$\chi_d \approx 224.54$~kHz (Code-Ramsey 1971;
Komasa--Pachucki 2020 ab~initio $224.51(5)$~kHz) &
\S~\ref{sec:proj_multipole_gaunt} (multipole at $R_\text{eq}$)
$\circ$ \S~\ref{sec:proj_3j} ($J=1, m_J=0$ rigid-rotor)
$\circ$ \S~\ref{sec:proj_tensor_multipole} ($Q_d$ scalar Layer-2 input)
$\circ$ \S~\ref{sec:proj_adiabatic_BO} (BO) &
$-39.9\%$ (OoM, 135 vs 224~kHz) & B$+$S &
First operator-level test of
\S~\ref{sec:proj_tensor_multipole} (rank-2 tensor multipole),
Track~4 of the multi-track Roothaan autopsy sprint (2026-05-18).
Order-of-magnitude estimate using only quantities the framework
currently exposes analytically:\ bare proton multipole at
$R_\text{eq} = 1.4011$~bohr gives $\chi_d^\text{nuc} = 195.25$~kHz
($-13\%$ of experiment, right sign and order); adding free-atom
1s electronic screening overshoots in the wrong direction giving
$\chi_d^\text{tot} = 135.01$~kHz ($-40\%$).  Bare-nuclear-only
contribution reaches 87\% of experiment alone, confirming that
the framework's existing multipole machinery captures the dominant
LO contribution.  The $40\%$ residual is specific and localized:\
molecular bond-charge redistribution requires the position-space
electronic density at the deuteron coordinate, an operator the
framework does not currently expose (Level~4
\texttt{geovac/level4\_multichannel.py} returns hyperspherical
$(R, \alpha, l, m)$ channel amplitudes, not
$\psi(\vec{r}_1, \vec{r}_2)$).  Error class B (basis architecture)
$+$ S (missing operator).  First molecular precision-catalogue entry.
\textbf{Named follow-on}:\ Sprint \texttt{HD-EFG-position-space}
(beyond sprint scale; analog of the W1a-D Roothaan recoil and
\S~\ref{sec:proj_magnetization_density} operator-level extensions).
Reference target Komasa--Pachucki 2020 ab~initio.  Sources:\
\texttt{debug/HD\_rotational\_HFS\_autopsy\_track4.py},
\texttt{debug/HD\_rotational\_HFS\_scoping\_track4\_memo.md}.
See \S~\ref{sec:autopsy_HD_rotational_HFS} for the full autopsy
table and the named follow-on architecture. \\
\midrule
He-3 $2{}^3S_1$ HFS $\nu_\text{HFS}(\text{He-3 } 2^3S_1) = 6653.4$ MHz
(BF + Schwinger $a_e$ + recoil + leading Zemach with $r_Z(\text{He-3})
= 1.965$~fm at Slater-screened $Z_\text{eff,2s} = 1.15$; Track-3,
2026-05-18) &
\S\ref{sec:proj_fock} $\circ$ \S\ref{sec:proj_spinor} $\circ$
\S\ref{sec:proj_3j} ($I=1/2$, mult 3/2 at $J=1$) $\circ$
\S\ref{sec:proj_bipolar_drake} (multi-electron triplet projection,
factor 1/2) $\circ$ \S\ref{sec:proj_phillips_kleinman} (Slater
screening of 2s by 1s) $\circ$ \S\ref{sec:proj_spectral_action}
$\circ$ \S\ref{sec:proj_restmass} ($m_p \to m_\text{He3}$) $\circ$
\S\ref{sec:proj_magnetization_density} ($Z=2$-scaled) &
$-1.28\%$ vs Schuessler--Fortson--Dehmelt 1969 / Prior--Wang 1977
$6739.701177(16)$~MHz & C (multi-electron screening prescription,
PK convention-dependent) + S (multi-electron QED Pachucki--Drachman
beyond one-loop) &
\textbf{First multi-electron HFS autopsy in catalogue.} The dominant
Layer-2 wall is \S\ref{sec:proj_phillips_kleinman} multi-electron
screening: spread across hydrogenic-vs-Slater-vs-full-screen
prescriptions is \textbf{10.5\% = 105\,468~ppm}, well above both
experimental precision (2.4 ppb) and LS-8a multi-loop QED budget
($\sim$tens of ppm). The framework-native architecture is identical
to H/D HFS; what changes is the dominant Layer-2 wall, which shifts
from LS-8a multi-loop QED (single-electron HFS) to
\S\ref{sec:proj_phillips_kleinman} multi-electron screening
(multi-electron HFS). New architectural ingredient: multi-electron
triplet contact density carries structural factor 1/2 from
$|J=S=1, M_S=1\rangle$ projection (a
\S\ref{sec:proj_bipolar_drake} bipolar-harmonic content for HFS).
He-3$^+$ 1s sanity check (single-electron hydrogen-like ion) matches
experiment at $+0.05\%$, anchoring convention.
\S\ref{sec:proj_magnetization_density} operator at $I=1/2$ He-3
reproduces Eides LO at machine precision (same operator as H/D HFS;
$I$- and multi-electron-independent).
New \S V.D candidate class-(i) convention exposure: Track 5 (D HFS)
BF formula mass-slot convention is system-dependent (works at I=1
only); standard convention ($g_N = \mu/I$, $m_p$ in mass slot) is
required at I=1/2. \textbf{Named follow-on}: ab initio
\S\ref{sec:proj_phillips_kleinman} via
\texttt{geovac.neon\_core.screened\_psi\_origin\_squared} on
no-frozen-core Z=2 system would give the framework-preferred unique
value within the 10.5\% spread. Autopsy:
\S\ref{sec:autopsy_he3_hfs}. Source memo:
\texttt{debug/He3\_HFS\_autopsy\_track3\_memo.md}. \\
\midrule
Li-7 $2{}^2S_{1/2}$ HFS $\nu_\text{HFS}({}^7\text{Li}) = 83.04$ MHz
(BF + Schwinger $a_e$ + recoil + leading Zemach with
$r_Z({}^7\text{Li}) = 3.71$~fm; Track-5, 2026-05-18) &
\S\ref{sec:proj_fock} $\circ$ \S\ref{sec:proj_spinor} $\circ$
\S\ref{sec:proj_3j} ($I=3/2$, mult 2) $\circ$
\S\ref{sec:proj_spectral_action} $\circ$ \S\ref{sec:proj_restmass}
$\circ$ \S\ref{sec:proj_magnetization_density} ($Z_\text{eff}$-scaled) &
$-89.7\%$ vs Beckmann--Boklen--Elke 1974 $803.504086(34)$~MHz &
B (single-zeta hydrogenic basis) + S (multi-electron contact-density
cliff, $\sim 10\times$) &
First multi-electron HFS in catalogue; first $I=3/2$ entry.  Cliff
factor $\sim 9.7\times$ on $A_\text{hf}$ from hydrogenic-$Z_\text{eff}$
underestimating $|\psi_{2s}(0)|^2$ vs Clementi--Roetti 1974 SCF
($\sim 8\times$ enhancement) plus \S\ref{sec:proj_phillips_kleinman} PK
orthogonality and \S\ref{sec:proj_bipolar_drake} bipolar core
polarization (1--2\%).  Effective $Z_\text{eff}$ for contact density
$\approx 2.73$ vs Slater 1.3 (orbital extent screened, contact density
not).  Two cliffs, one mechanism: structural sibling of the
\S\ref{sec:autopsy_cs_hfs} Z\textgreater 20 cliff (single-zeta basis
no radial nodes).  Operator-level $\hat{I}\cdot\hat{S}$ at $I=3/2$
gives eigenvalues $\{+3/4,-5/4\}$, mult 2 at machine precision.
Pauli encoding $Q_\text{tot}=3$, 8 terms.
\S\ref{sec:proj_magnetization_density} operator at $I=3/2$ reproduces
Eides LO at machine precision (same operator as H 21cm and D HFS,
4 Pauli terms, $I$-independent).  New \S~V.D class~(iii) exposure:
framework-architecture limitation (\S\ref{sec:conv_alkali_hfs}).
\textbf{Named follow-on}: Hylleraas--Eckart double-zeta for $2{}^2S$
HFS (closes Li atomic accuracy floor $1.10\%$ $\to$ sub-$0.1\%$
\emph{and} the contact-density cliff simultaneously;
\texttt{geovac/hylleraas\_r12.py} extension).  Autopsy:
\S\ref{sec:autopsy_li7_hfs}.  Source memo:
\texttt{debug/Li7\_HFS\_autopsy\_track5\_memo.md}. \\
\midrule
$\mu^4$He$^+$ 2$S$--2$P$ Lamb $+1358.66$ meV framework $+$ Layer-2
(multi-track sprint Track B, 2026-05-19; full Roothaan autopsy in
\S\ref{sec:autopsy_muhe4_lamb}) &
\S\ref{sec:proj_fock} $\circ$ \S\ref{sec:proj_spectral_action}
$\circ$ \S\ref{sec:proj_sturmian} $\circ$ \S\ref{sec:proj_restmass}
$\circ$ \S\ref{sec:proj_charge_density}; \emph{no}
\S\ref{sec:proj_magnetization_density} ($^4$He has $I=0$) &
$-1.44\%$ vs CREMA 2021 $1378.521(48)$ meV (Krauth 2021,
arXiv:2102.05728) &
C (Diepold conventions: Friar HO absorbed into leading
$\sim +11.4$ meV, $r_\alpha$ extraction $\sim +7$ meV) +
A (Eides SE recoil-mixing $\sim +1$ meV, same gap as $\mu$H/$\mu$d) &
\textbf{Cleanest pure-Lamb test in muonic catalogue} (no Zemach,
no HFS contamination, no $\hat{I}\cdot\hat{S}$).  Most extreme
mass-hierarchy $\times$ Z combination tested: $m_\text{red} \sim
201\,m_e$, $Z = 2$, $\beta = 0.682$ (vs $1.475$ at $\mu$H, $1.400$
at $\mu$d).  Framework-native VP at $-0.032\%$ vs Diepold 2018
$+1666.305$ meV --- bit-identical kernel transports verbatim
across $Z = 1 \to 2$.  Contact-form Uehling fails by
$\mathbf{8.76\times}$ at $\beta = 0.682$ (vs $3.55\times$ at $\mu$H);
full kernel integration mandatory at heavy-Z muonic.  Three
attributable convention surfaces fully explain the $-19.86$ meV gap:
Friar HO absorption ($\sim +11.4$ meV; flagged \S~V.D candidate),
$r_\alpha$ extraction signal ($\sim +7$ meV; Diepold reference 1.681
fm vs CREMA-extracted $1.67824(13)(82)$ fm), SE bracket recoil-mixing
($\sim +0.9$ meV; same Phase C-W1a target as $\mu$H $+24\%$, $\mu$d
$+41\%$).  Source memo:
\texttt{debug/mu\_he4\_lamb\_autopsy\_track\_b\_memo.md}. \\
\bottomrule
\end{tabular}
\caption{Off-precision matches with error-source classification.
Residual error is the absolute relative deviation from the standard
physics value.  Source codes: T (truncation), B (basis quality),
A (approximation order), C (calibration mismatch), S (structural
floor).  Combinations indicate multiple independent contributions
to residual error.  This table is a complement to
Sec.~\ref{sec:catalog}; together they document the full empirical
profile of the framework.}
\label{tab:offprecision}
\end{table}

\begin{remark}[Error sources are projection-specific]
\label{rem:err_sources}
The error-source classification is not orthogonal to the projection
dictionary.  Truncation (T) errors are properties of the spectral
sums that the Sturmian or spectral-action projections invoke.
Calibration mismatch (C) is a property of how multiple projections
chain.  Structural floors (S) appear where a needed projection is
absent from Layer~2 entirely.  Approximation order (A) is the
``flow'' tier of Paper~18: it is the residual after Layer~1 and
Layer~2 do their work, and it indicates which higher-order projection
or perturbation order should be added next.
\end{remark}

% -------------------------------------------------------------------
\subsection{Roothaan autopsies: multi-component decompositions}
\label{sec:autopsies}
% -------------------------------------------------------------------

The matches catalogue rows in \S\ref{sec:catalog} and
\S\ref{sec:offprecision} each capture a single observable as a
single row.  For observables that admit a multi-component literature
decomposition --- e.g.\ Eides Tab.~7.3 for the hydrogen Lamb shift,
where the standard $1057.85$ MHz value decomposes into one-loop SE,
one-loop VP, multi-loop QED, finite nuclear size, recoil NLO, and
hyperfine averaging --- the single-row treatment hides the structure
that the projection-chain dictionary makes visible.  This subsection
hosts \emph{Roothaan autopsies}: decomposition tables that show each
measured observable as a sum of focal-length-tagged components, with
each component's projection chain made explicit.

The Roothaan autopsy is structurally the same kind of object Roothaan
(1951) made for two-electron molecular integrals: a multi-component
decomposition where each piece has its own focal length and its own
closed-form structure.  Applying the same discipline to multi-component
\emph{observables} makes the projection-chain dictionary's reach
visible at the observable level, not just at the catalogue-row level.
The autopsy partitions each component into framework-native (FN ---
computable from \S~\ref{sec:layer2} projections without external
numerical constants beyond $\alpha$) versus Layer-2 (L2 --- parameter
at operator level, projection chain framework-native), giving an
explicit operator-level reading of the structural-skeleton scope
statement (CLAUDE.md \S~1.7).  Several autopsies surface
literature convention mismatches (class~(a) findings of the
multi-observable focal-length decomposition program); these are
collected separately in \S~\ref{sec:convention_exposures}.

\subsubsection{Hydrogen 1S Lamb shift autopsy}
\label{sec:autopsy_lamb}

\textbf{Reference.} $\nu(2S_{1/2} - 2P_{1/2}) = 1057.845(9)$ MHz
(Lundeen--Pipkin 1981; NIST ASD).

\begin{table}[h]
\centering\small
\begin{tabular}{p{4cm} r p{6cm} c}
\toprule
Component & MHz & Projection chain & Status \\
\midrule
SE 2$S_{1/2}$ (Bethe + Eides 10/9) & $+1066.44$ &
\S\ref{sec:proj_fock} $\circ$ \S\ref{sec:proj_sturmian} $\circ$
\S\ref{sec:proj_spectral_action} $\circ$ \S\ref{sec:proj_drake_swainson}
& FN \\
VP 2$S_{1/2}$ (Uehling) & $-27.13$ &
\S\ref{sec:proj_fock} $\circ$ \S\ref{sec:proj_spectral_action} & FN \\
SE 2$P_{1/2}$ & $+12.88$ &
\S\ref{sec:proj_fock} $\circ$ \S\ref{sec:proj_sturmian} $\circ$
\S\ref{sec:proj_spectral_action} $\circ$ \S\ref{sec:proj_spinor}
$\circ$ \S\ref{sec:proj_drake_swainson} & FN \\
VP 2$P_{1/2}$ & $0$ &
\S\ref{sec:proj_fock} (parity zero) & FN \\
$\alpha^5$ multi-loop QED & $+1.20$ &
\S\ref{sec:proj_fock} $\circ$ \S\ref{sec:proj_sturmian} $\circ$
$(\S\ref{sec:proj_spectral_action})^2$ $\circ$ \S\ref{sec:proj_spinor}
$+$ Layer-2 $C_{2S}$ & L2 \\
FNS $r_p$ contribution & $+1.18$ &
\S\ref{sec:proj_fock} $\circ$ \S\ref{sec:proj_spinor} $\circ$
\textbf{\S\ref{sec:proj_charge_density}} $+$ Layer-2 $r_p$ & L2 \\
Recoil NLO (beyond reduced-mass) & $-2.40$ &
\S\ref{sec:proj_fock} $\circ$ \S\ref{sec:proj_restmass} $\circ$
\S\ref{sec:proj_breit_retardation} $+$ Layer-2 $R_{nl}$ & L2 \\
Hyperfine averaging (F-state mixing) & $\sim +5.00$ &
\S\ref{sec:proj_fock} $\circ$ \S\ref{sec:proj_spinor} $\circ$
\S\ref{sec:proj_3j} $\circ$ \S\ref{sec:proj_restmass} & FN \\
\midrule
\textbf{Sum} & $+1057.17$ & & \\
Experimental & $+1057.845(9)$ & & \\
\textbf{Residual} & $+0.68$ & within Eides Tab.~7.3 quoting precision & \\
\bottomrule
\end{tabular}
\caption{Hydrogen 1S Lamb shift Roothaan autopsy.  Each component is
tagged to its \S\ref{sec:layer2} projection chain.  All entries are
contributions to the $2S_{1/2}-2P_{1/2}$ splitting: the $2P_{1/2}$
self-energy is a $-12.88$~MHz level shift, entering the splitting as
$+12.88$~MHz.  Status: FN
$=$ framework-native; L2 $=$ Layer-2 input.  Source memo:
\texttt{debug/calc\_track\_LAR\_lamb\_autopsy\_memo.md}.}
\label{tab:autopsy_lamb}
\end{table}

\textbf{Framework-native subtotal:} $+1057.19$ MHz ($99.9\%$ of
measurement).  \textbf{Layer-2 net:} $+1.20 + 1.18 - 2.40 \approx
-0.02$ MHz --- empirical near-cancellation specific to the H 1S
Lamb shift, \emph{not} a generic property.  The framework-native
components alone reproduce the measurement at $\approx 6 \times
10^{-4}$ precision; the Layer-2 inputs cancel in this observable.

\textbf{Multi-focal-composition wall visible.}  The two non-cancelling
Layer-2 components correspond to:
\begin{itemize}\setlength\itemsep{0pt}
\item Component 5 (multi-loop QED) is the LS-8a renormalization-gap
wall (CLAUDE.md \S~1.7 LS-8a entry).  The framework reproduces the
two-loop SE structural prefactor $(\alpha/\pi)^2 (Z\alpha)^4/n^3$
but cannot autonomously generate the $Z_2/\delta m$ counterterms.
\item Component 7 (recoil NLO) is the
\S\ref{sec:proj_breit_retardation} two-body kinematic wall.  The
framework's Breit retardation projection produces the $\alpha^4$
kinematic structure but the $R_{nl}$ coefficients enter as
Layer-2 input.
\item Component 6 (FNS) is now \emph{framework-native at the
operator level} (today's \S\ref{sec:proj_charge_density} promotion);
the remaining Layer-2 input is the scalar $r_p$ value, not the
operator structure consuming it.  This is the cleanest beneficiary
of the 2026-05-09 promotion.
\end{itemize}

\textbf{Structural reading.}  The autopsy demonstrates Paper~34
\S\ref{sec:layer2}'s reach at finer grain than single-row catalogue
entries: eight components, eight different projection chains, every
Eides Tab.~7.3 row reproduced.  The structural-skeleton scope
statement is made concrete --- framework-native pieces reproduce the
observable to $\approx 6 \times 10^{-4}$, with the Layer-2 inputs
(renormalization, Component 5; two-body kinematics, Component 7)
netting to $-0.02$~MHz, and FNS (Component 6) now in projection-chain
language rather than as a black-box external constant.

\subsubsection{Hydrogen 21~cm hyperfine autopsy}
\label{sec:autopsy_21cm}

\textbf{Reference.} $\nu_\text{HF} = 1\,420\,405\,751.768$~Hz (Hellwig
1970; CODATA / NIST), the most precisely measured atomic transition.

\begin{table}[h]
\centering\small
\begin{tabular}{p{4cm} r p{6cm} c}
\toprule
Component & MHz & Projection chain & Status \\
\midrule
Bohr--Fermi Dirac (point nucleus, $g_e=2$, no recoil) & $+1421.1595$ &
\S\ref{sec:proj_fock} $\circ$ \S\ref{sec:proj_spinor} $\circ$
\S\ref{sec:proj_3j} & FN \\
$+$ Schwinger $a_e$ (Parker--Toms-verified at $+0.5\%$) & $\times (1+\alpha/2\pi)$ &
\S\ref{sec:proj_fock} $\circ$ \S\ref{sec:proj_spinor} $\circ$
\S\ref{sec:proj_spectral_action} & FN (with calibration) \\
$+$ Reduced-mass $(1+m_e/m_p)^{-3}$ & $\times 0.998366$ &
\S\ref{sec:proj_fock} $\circ$ \S\ref{sec:proj_restmass} & FN \\
$+$ Zemach $r_Z=1.045$~fm via \S\ref{sec:proj_magnetization_density}
operator-level & $\times (1 - 39.495 \text{~ppm})$ &
\S\ref{sec:proj_fock} $\circ$ \S\ref{sec:proj_spinor} $\circ$
\textbf{\S\ref{sec:proj_magnetization_density}} & FN at op-level $+$ L2 ($r_Z$) \\
\midrule
\textbf{Final $A_\text{HF}$} & $+1420.4318$ & & \\
Experimental & $+1420.4058$ & (12-digit precision) & \\
\textbf{Residual} & $+0.026$ & $= +18.4$~ppm; within Eides Tab.~7.3 itemized $+12$--$+18$~ppm band & \\
\bottomrule
\end{tabular}
\caption{Hydrogen 21~cm hyperfine Roothaan autopsy.  Each component is
tagged to its \S\ref{sec:layer2} projection chain.  Status: FN $=$
framework-native; L2 $=$ Layer-2 input.  Source memo:
\texttt{debug/h21\_autopsy\_v1\_memo.md} (Sprint Calc-H21-Autopsy v1, May~2026).}
\label{tab:autopsy_21cm}
\end{table}

\textbf{Operator-level Zemach.}  Component 4 exercises
\S\ref{sec:proj_magnetization_density} at the operator level via the
production module \texttt{geovac.magnetization\_density}.  The bilinear
matrix element $\langle \hat{r}_Z\rangle$ on a Sturmian register at
the L=0 multipole reduces to $M_1[\rho_M] = r_Z$ at leading order in
$r_Z/a_0 \sim 2\times 10^{-5}$; the operator output $-39.495$~ppm at
$r_Z = 1.045$~fm vs Eides Tab.~7.3 reference $-39.500$~ppm gives
residual $+4.7\times 10^{-3}$~ppm $= 0.012\%$ of the LO shift --- the
operator collapses to the Eides scalar at sub-percent precision
automatically through the multipole reduction, without any external
substitution.  Profile independence verified: Gaussian and exponential
$\rho_M$ give bit-identical $A_\text{HF}$ at LO since both are
calibrated to first moment $M_1 = r_Z$.

\textbf{Framework-native subtotal:} $A_\text{HF}^\text{native}\!=\!1420.4318$~MHz
($99.998\%$ of measurement).  \textbf{Layer-2 net:} $-39.495 + 18.37
\approx -21.1$~ppm of $\nu_F$.  The Layer-2 input is the scalar $r_Z$
value; the operator structure consuming it is framework-native via
\S\ref{sec:proj_magnetization_density}.  This is the cleanest
beneficiary of the W1b operator-level extension (Sprint
Calc-rZG-extended, 2026-05-09), which promoted Zemach from
analytic-substitution to operator-level.

\textbf{Multi-focal-composition wall visible.}  The $+18.37$~ppm
residual decomposes per Eides Tab.~7.3 as:
\begin{itemize}\setlength\itemsep{0pt}
\item Multi-loop QED ($\alpha^2(Z\alpha)$, $\sim+6$~ppm) is the
LS-8a renormalization gap (CLAUDE.md \S~1.7 LS-8a entry).
\item Recoil NLO beyond reduced-mass (Bodwin--Yennie, $\sim +6$~ppm)
is the W1a-D Roothaan kernel-level recoil-mixing wall.
\item Nuclear polarizability $\Delta_\text{pol}$ ($+1.4$~ppm) and
hadronic VP ($+0.1$~ppm) are W3 inner-factor calibration data
(QCD-internal).
\item Zemach NLO recoil-mixing $m_e/(m_e+m_p) \approx 5 \times 10^{-4}$
is structural noise in the electronic regime ($+0.02$~ppm) but the
\emph{dominant} systematic in the muonic regime
($f_\text{recoil} \approx 0.092$, structural $\sim 10\%$).  The W1b
NLO opt-in flag covers this cleanly; cf. Sprint MH for the muonic
analog.
\item Convention drift on $r_Z$ ($\pm 5$~ppm; Eides 1.045 vs
Karshenboim 1.054) is the largest class~(a) literature-itemization
sensitivity in this observable.
\end{itemize}

\textbf{Structural reading.}  The autopsy demonstrates Paper~34's
multi-observable focal-length decomposition program (CLAUDE.md
\S~1.8) at the operator level on the most precisely measured atomic
transition: four components, four different projection chains, $+18$
ppm residual fully attributable to LS-8a (multi-loop QED) $+$ W1a-D
(Roothaan recoil) $+$ W3 (nuclear polarizability), exactly as
classified in CLAUDE.md \S~1.7.  The structural-skeleton scope
statement is verified: the framework reproduces the observable to
$\sim 2 \times 10^{-5}$ precision via four named projections, and
the residual budget decomposes cleanly into already-named walls
without surfacing any new class~(a) literature convention mismatch
beyond the known $r_Z$-value drift.

\subsubsection{Muonic hydrogen 2$S_{1/2}$--2$P_{1/2}$ Lamb shift autopsy}
\label{sec:autopsy_muh_lamb}

\textbf{Reference.} $\Delta E = 202.3706(23)$ meV (CREMA 2010,
Pohl~et~al.; muonic convention $E(2P_{1/2}) - E(2S_{1/2})$).

The CREMA 2010 muonic hydrogen Lamb shift decomposes under five
named projection-chain components plus three companion Layer-2
inputs.

\begin{table}[h!]
\centering\small
\caption{Five-component Roothaan decomposition of the muonic hydrogen
Lamb shift (muonic convention $E(2P_{1/2}) - E(2S_{1/2})$).  The
framework-native subtotal $+200.50$ meV comes from full Uehling
integration via \S\ref{sec:proj_spectral_action}, the self-energy
bracket under rest-mass projection
(\S\ref{sec:proj_restmass}~+~\S\ref{sec:proj_sturmian}), and the
operator-level Friar moment via the magnetization-density module
(\S\ref{sec:proj_magnetization_density}; cf.\
\S\ref{sec:proj_charge_density} for the parent charge-density
projection).  The $+1.67$ meV gap to the $+202.17$ meV
framework-plus-literature total is closed by Antognini~2013
literature inputs covering multi-loop QED (K\"all\'en--Sabry,
mixed VP-VP/VP-SE, $\alpha^7$) and QCD-internal proton
polarizability (W3 inner-factor).  Source memo:
\texttt{debug/muH\_Lamb\_autopsy\_v1\_memo.md}.}
\label{tab:autopsy_muh_lamb}
\begin{tabular}{p{4.5cm} r p{4.8cm} l}
\toprule
Component & Value (meV) & Projection chain & Status \\
\midrule
1.~Full Uehling VP & $+205.0074$ &
\S\ref{sec:proj_fock} $\circ$ \S\ref{sec:proj_spectral_action}
& FN ($-0.10$ ppm vs Antognini) \\
2.~Self-energy (Eides bracket) & $-0.8295$ &
\S\ref{sec:proj_fock} $\circ$ \S\ref{sec:proj_sturmian}
$\circ$ \S\ref{sec:proj_restmass}
& FN ($+24\%$ vs Antognini, recoil-mixing gap) \\
3.~Friar moment (closed-form) & $-3.6752$ &
\S\ref{sec:proj_fock} $\circ$ \S\ref{sec:proj_charge_density}
$/$ \S\ref{sec:proj_magnetization_density}
& FN ($-4.3\%$ vs Antognini) \\
4.~K\"all\'en--Sabry 2-loop VP & $+1.5081$ & Layer-2 input &
LS-8a renormalization wall \\
5.~Proton polarizability & $+0.0129$ & Layer-2 input &
W3 inner-factor (QCD) \\
Companions (recoil $+$ multi-loop) & $+0.144$ & Layer-2 input &
Antognini 2013 catalogue \\
\midrule
\textbf{Total} & $+202.17$ & \multicolumn{2}{l}{Residual $-0.10\%$
vs CREMA $+202.3706(23)$ meV} \\
\bottomrule
\end{tabular}
\end{table}

\textbf{Headline framework computation.}  The full Uehling integration
matches the Antognini~2013 / Pachucki~1996 reference at $-0.10$ ppm
--- the cleanest single-component match in the muonic catalogue.
The $\beta = 2/(m_{\text{red},\mu}\,\alpha) = 1.475$ muonic regime,
where the muonic Bohr scale and the $e^+ e^-$ Compton wavelength are
comparable, is handled correctly by the full kernel where the
contact-form Uehling overshoots by $\sim 3.6\times$.  This is the
muonic-regime analog of the electronic Uehling row in
\S\ref{sec:autopsy_lamb} (Component 2 of \S\ref{tab:autopsy_lamb}),
where contact form is sufficient because $\beta_\text{ep} \approx 274$.

\textbf{Convention-mismatch surface.}  Antognini~2013 and Krauth~2017
itemizations agree on the bottom-line total but differ at $\sim 100$
ppm in per-component split (VP one-loop and SE-recoil aggregation).
Specifically, Krauth absorbs $\sim 0.021$~meV of higher-order VP into the
one-loop line ($+205.0282$~meV), while Antognini itemizes it separately
as ``VP-VP mixed'' ($+0.150$~meV), with one-loop reported at
$+205.0074$~meV.  The framework's $-0.10$~ppm match against Antognini
is $1000\times$ tighter than the $102$~ppm Antognini-vs-Krauth gap.
At this precision level the choice of Layer-2 reference compilation
matters; this autopsy anchors throughout to Antognini~2013.  See
\S~\ref{sec:convention_exposures} for catalogue context (entry
\S~\ref{sec:conv_muH_VP}).

\textbf{Operator-level Friar moment.}  Computed via the
\S\ref{sec:proj_magnetization_density} module
\texttt{geovac.magnetization\_density} with
\texttt{lepton\_mass}~=~$m_{\text{red},\mu}$.  The operator-level
result $-4.33$ meV vs closed-form $-3.68$ meV reflects the Gaussian
profile's $\langle r^2\rangle / \langle r\rangle^2 = 4/\pi \approx
1.27$ factor --- a profile-convention issue between calibrating to
the first moment $\langle r\rangle$ versus the RMS radius
$\sqrt{\langle r^2\rangle}$.  Eides Eq.~2.35 closed-form derivation
takes $\langle r^2\rangle^{1/2} = r_p$ (RMS radius); the operator-level
realization via \S\ref{sec:proj_magnetization_density} takes
$\langle r\rangle = r_p$ (first moment).  The discrepancy is exactly
the Gaussian profile factor.  The flagged extension is a sibling
\texttt{MagnetizationDensitySpec} calibrated to RMS radius for the
Lamb-shift FNS channel, distinct from the existing first-moment
calibration for the HFS Zemach channel.  See
\S~\ref{sec:convention_exposures} for catalogue context (entry
\S~\ref{sec:conv_friar_profile}).

\textbf{Structural reading.}  The autopsy realizes the
structural-skeleton-scope statement at the operator level for a
cross-register multi-focal observable in the muonic regime.  Three
framework-native components reproduce $+200.50$ meV (the
five-projection-chain skeleton); three Layer-2 inputs close the
remaining $+1.67$ meV gap, attributable cleanly to LS-8a multi-loop
renormalization (Component 4) and W3 inner-factor (Component 5).
The same architecture that closes the H 1S Lamb shift at
\S\ref{sec:autopsy_lamb} closes $\mu$H 2$S$--2$P$ here, with
Uehling-kernel kinematics (contact-form $\to$ full kernel) and
Friar-profile calibration (first moment $\to$ RMS) being the only
two ingredients that change between regimes.

\subsubsection{Muonic deuterium 2$S_{1/2}$--2$P_{1/2}$ Lamb shift autopsy}
\label{sec:autopsy_mud_lamb}

\textbf{Reference.}  $\Delta E = 202.8785(34)$ meV (CREMA 2016, Pohl
et al., \emph{Science} \textbf{353}, 669; muonic convention
$E(2P_{1/2})-E(2S_{1/2})$).

The CREMA 2016 muonic deuterium Lamb shift decomposes under the
same five-component skeleton as the muonic hydrogen
autopsy~(\S\ref{sec:autopsy_muh_lamb}), with three architectural
substitutions: (a) nuclear charge radius $r_p \to r_d = 2.1413(25)$
fm (factor $2.547$ in linear size, $6.487$ in $r^2$); (b) reduced
mass $m_{\text{red},\mu p} \to m_{\text{red},\mu d} = 195.74\,m_e$
(factor $1.053$, mildly larger via the heavier nucleus); (c)
dominant Layer-2 input shifts from K\"all\'en--Sabry two-loop VP
(at $\mu$H) to deuteron polarizability ($\sim 130\times$ larger
than proton polarizability, dominant at $\mu$d).

\begin{table}[h!]
\centering\small
\caption{Five-component Roothaan decomposition of the muonic
deuterium Lamb shift (muonic convention $E(2P_{1/2})-E(2S_{1/2})$).
The framework-native subtotal $+198.91$ meV comes from full Uehling
integration via \S\ref{sec:proj_spectral_action}, the self-energy
bracket under rest-mass projection
(\S\ref{sec:proj_restmass}~+~\S\ref{sec:proj_sturmian}), and the
closed-form Friar moment at the deuteron's measured RMS charge
radius via \S\ref{sec:proj_charge_density}.  The $+3.71$ meV gap
to the $+202.63$ meV framework-plus-literature total is closed by
Krauth~2016 literature inputs covering multi-loop QED
(K\"all\'en--Sabry, iterated Uehling, $\alpha^7$), higher-Friar
moments and recoil-mixing, NLO recoil corrections, and---dominant
contribution---deuteron polarizability (Carlson--Gorchtein--
Vanderhaeghen 2014).  Source memo:
\texttt{debug/muD\_Lamb\_autopsy\_track1\_memo.md}.}
\label{tab:autopsy_mud_lamb}
\begin{tabular}{p{4.7cm} r p{4.8cm} l}
\toprule
Component & Value (meV) & Projection chain & Status \\
\midrule
1.~Full Uehling VP & $+227.6347$ &
\S\ref{sec:proj_fock} $\circ$ \S\ref{sec:proj_spectral_action}
& FN ($-0.50\%$ vs Krauth) \\
2.~Self-energy (Eides bracket) & $-0.8737$ &
\S\ref{sec:proj_fock} $\circ$ \S\ref{sec:proj_sturmian}
$\circ$ \S\ref{sec:proj_restmass}
& FN ($+41\%$ vs Krauth, same recoil-mixing gap as $\mu$H widened
at heavier nucleus) \\
3.~Friar moment leading (closed-form Eides) & $-27.8466$ &
\S\ref{sec:proj_fock} $\circ$ \S\ref{sec:proj_charge_density}
& FN ($+0.024\%$ vs Krauth at $r_d$, bit-tight) \\
4.~Multi-loop QED (KS + iterated + mixed + 3-loop) & $+1.8578$ &
Layer-2 input & LS-8a renormalization wall \\
5.~Friar HO + recoil-mixing & $+0.1030$ &
Layer-2 input & higher Friar moments \\
6.~Recoil corrections (NLO + HO) & $+0.0637$ &
Layer-2 input & W1a sub-leading recoil \\
7.~\textbf{Deuteron polarizability} & $\mathbf{+1.6900}$ &
Layer-2 input & \textbf{W3 inner-factor (QCD, dominant)} \\
\midrule
\textbf{Total} & $+202.63$ & \multicolumn{2}{l}{Residual
$-0.12\%$ vs CREMA $+202.8785(34)$ meV} \\
\bottomrule
\end{tabular}
\end{table}

\textbf{Headline framework computation.}  The Friar moment leading
order via the closed-form Eides Eq.~2.35 matches the Krauth 2016
tabulated leading line at the deuteron's measured radius at
$+0.024\%$ --- the cleanest framework-vs-literature match for the
nuclear-structure leading line across the muonic catalogue.
The scaling
\begin{equation}
\frac{\Delta E_\text{FNS}(\mu d)}{\Delta E_\text{FNS}(\mu p)}
= \left(\frac{m_{\text{red},\mu d}}{m_{\text{red},\mu p}}\right)^3
  \left(\frac{r_d}{r_p}\right)^2
= 1.168 \times 6.487 = 7.580
\end{equation}
is reproduced by the framework's bare Eides calculation to 9 digits.
This confirms the rest-mass projection (\S\ref{sec:proj_restmass})
transports the Friar moment leading-order calculation cleanly
across nuclear-isotope variation.

\textbf{Same sub-leading SE recoil-mixing gap as $\mu$H.}  The
framework's self-energy leading-order rest-mass scaling overshoots
Krauth's catalogued SE Lamb-shift contribution by $+0.25$ meV at
$\mu$d ($+41\%$ fractional) vs $+0.16$ meV at $\mu$H ($+24\%$
fractional).  The \emph{absolute} gap grows by $1.5\times$ between
isotopes, while the \emph{fractional} gap grows from $+24\%$ to
$+41\%$ because Krauth's tabulated SE shrinks slightly more.  Same
physical mechanism (sub-leading $\alpha(Z\alpha)^4 \times
(m_\text{red}/m_n)$ recoil-mixing pieces, Eides~\S~4.2, Pachucki~1996,
omitted at the leading-order rest-mass projection); the gap is
diagnostically robust under nuclear-isotope variation.
Phase~C-W1a-physics extension at the SE-bracket level would close
both observables simultaneously.

\textbf{Dominant Layer-2 shifts from LS-8a to W3 at $\mu$d.}  In
$\mu$H, the dominant Layer-2 input is the K\"all\'en--Sabry two-loop
VP ($+1.51$ meV, LS-8a renormalization wall).  In $\mu$d, the
dominant Layer-2 input is the deuteron polarizability ($+1.69$ meV,
W3 inner-factor wall) --- $\sim 130\times$ larger than the proton
polarizability ($+0.013$ meV in $\mu$H) because the deuteron is a
weakly bound $n+p$ system whose virtual photonuclear excitations sit
at MeV scale.  This structural feature (different walls dominate at
different nuclei) is direct empirical evidence for the
multi-focal-composition wall taxonomy: different observables probe
different walls at different relative weights, even within the same
projection chain.

\textbf{Convention-mismatch surface.}  Carlson--Gorchtein--Vanderhaeghen
2014 catalogues the deuteron polarizability at $+1.690 \pm 0.020$
meV; the older Carlson--Vanderhaeghen 2011 evaluation reported
$+1.94$ meV --- a $\sim 13\%$ spread on the \emph{dominant} Layer-2
line at $\mu$d.  At the framework's $-0.12\%$ Lamb-shift residual
scale this is the controlling uncertainty.  \emph{Flagged as a new
V.D-prediction candidate}, class~(i) inter-compilation itemization,
sibling to \S\ref{sec:conv_d_polarizability} (which addresses
sub-component itemization of deuteron HFS polarizability across
Friar--Payne~2005 and Pachucki--Yerokhin~2010+).

\textbf{Structural reading.}  The autopsy validates the multi-focal
architecture under nuclear-isotope substitution at the sub-percent
precision level on a precision-frontier observable.  Three
framework-native components reproduce $+198.91$ meV; four Layer-2
inputs close the remaining $+3.71$ meV gap, with the dominant
contribution now coming from the W3 inner-factor (deuteron
polarizability, QCD-internal NN dynamics) rather than the LS-8a
multi-loop QED that dominated at $\mu$H.  The rest-mass projection
(\S\ref{sec:proj_restmass}) transports the entire skeleton (Uehling
integrator, SE bracket, Friar leading) verbatim; only the nuclear
charge-radius input ($r_p \to r_d$) and the polarizability
Layer-2 input change between isotopes.  Same architecture, same
sub-percent residual ($-0.10\%$ at $\mu$H, $-0.12\%$ at $\mu$d),
different dominant Layer-2 wall.

\subsubsection{He$^+$ 2$S_{1/2}$--2$P_{1/2}$ Lamb shift autopsy (Z-scaling test of Paper 36)}
\label{sec:autopsy_he_plus_lamb}

\textbf{Reference.}  $\nu(2S_{1/2} - 2P_{1/2}) = 14\,041.13(17)$ MHz
(van Wijngaarden--Holuj--Drake, \emph{Phys. Rev. A} \textbf{63}, 012505
(2000); theoretical reference $14\,041.18(13)$ MHz from
Pachucki~\emph{Phys. Rev. A} \textbf{63}, 042503 (2001) and Drake--Yan
\emph{Phys. Rev. A} \textbf{46}, 2378 (1992) consistent to $\sim 50$~kHz).
Convention $E(2S_{1/2}) - E(2P_{1/2}) > 0$ matches Paper~36 normal-H.

The He$^+$ Lamb shift is the cleanest $Z$-scaling test of the Paper~36
one-loop QED architecture (\S\ref{sec:autopsy_lamb}): identical lepton
(electron), identical projection chains, only the nuclear charge $Z=2$
and the nuclear input parameters $m_\alpha = 7294.30\,m_e$, $r_E({}^4\text{He})
= 1.6755$ fm change.  ${}^4$He has $I=0$, so the entire
\S\ref{sec:proj_magnetization_density} channel is identically zero ---
no HFS, no Zemach radius, no hyperfine averaging.  The Uehling parameter
$\beta = 2/(Z\,m_{\text{red}}\,\alpha) = 137.05$ sits in the same large-$\beta$
regime as H ($\beta = 274$), so contact-form vacuum polarization is valid
to $\sim$ppm and the full Uehling kernel matches contact form at 0.5\%.

\begin{table}[h!]
\centering\small
\caption{Five-component Roothaan decomposition of the He$^+$
2$S_{1/2}$--2$P_{1/2}$ Lamb shift (normal-H convention
$E(2S_{1/2}) - E(2P_{1/2})$).  The framework-native subtotal
$+13\,841.11$ MHz comes from the same projection-chain machinery that
closed the H Lamb shift (\S\ref{sec:autopsy_lamb}) with $Z = 2$
substituted throughout: Bohr-Fock Dirac (zero by 2$S_{1/2}=2P_{1/2}$
degeneracy), one-loop self-energy via Eides Sec.~3.2 bracket, one-loop
vacuum polarization via full Uehling kernel, and Foldy-Friar finite
nuclear size via \S\ref{sec:proj_charge_density}.  The $-200$~MHz
deficit against experiment is the Z-amplified bracket-completion gap
($\alpha^3$ truncation amplified by $Z^4 \times \ln$ from the $-0.55\%$
residual at Z=1), closed at Layer-2 by Drake-Yan 1992 multi-loop QED
($\alpha^2(Z\alpha)^4 + \alpha(Z\alpha)^5 + \alpha(Z\alpha)^6
\ln(Z\alpha)^{-2}$).  Source memo:
\texttt{debug/he\_plus\_lamb\_autopsy\_track\_a\_memo.md}.}
\label{tab:autopsy_he_plus_lamb}
\begin{tabular}{p{5.5cm} r p{4.7cm} l}
\toprule
Component & Value (MHz) & Projection chain & Status \\
\midrule
1.~Bohr-Fock Dirac baseline & $+0.0000$ &
\S\ref{sec:proj_fock} & FN (zero by Dirac 2$S_{1/2}=2P_{1/2}$ degeneracy) \\
2.~One-loop SE (Eides Sec.~3.2 bracket) & $+14\,258.5252$ &
\S\ref{sec:proj_fock} $\circ$ \S\ref{sec:proj_sturmian} $\circ$
\S\ref{sec:proj_spectral_action} $\circ$ \S\ref{sec:proj_drake_swainson}
& FN ($\alpha^3 Z^4$ truncation) \\
3.~One-loop VP (full Uehling kernel) & $-426.2096$ &
\S\ref{sec:proj_fock} $\circ$ \S\ref{sec:proj_spectral_action}
& FN ($\beta = 137$, contact form $0.5\%$ check) \\
4.~Foldy/Friar FNS ($r_E({}^4$He$)^2$) & $+8.7913$ &
\S\ref{sec:proj_fock} $\circ$ \S\ref{sec:proj_charge_density} $+$
Layer-2 $r_E$ & FN \\
5a.~$\alpha^2(Z\alpha)^4$ two-loop QED & $+6.71$ &
Layer-2 input & LS-8a renormalization wall \\
5b.~$\alpha(Z\alpha)^5$ higher-order & $+0.18$ &
Layer-2 input & LS-8a \\
5c.~$\alpha(Z\alpha)^6 \ln$ relativistic Bethe-log (Drake-Yan)
& $+14.7$ & Layer-2 input & bracket-completion gap at $Z = 2$ \\
6.~Recoil NLO $m_e^2/m_\alpha$ & $-3.18$ &
Layer-2 input & W1a sub-leading \\
7.~${}^4$He polarizability & $+0.06$ &
Layer-2 input & W3 inner-factor (smallest in catalogue) \\
\midrule
\textbf{Framework-native subtotal} & \textbf{+13\,841.11} & (1+2+3+4) & \\
Framework $+$ literature total & $+13\,859.58$ & (1+2+3+4+5a+5b+5c+6+7) & \\
\textbf{Experimental} & \textbf{+14\,041.13(17)} & van Wijngaarden 2000 & \\
\textbf{Theory} & \textbf{+14\,041.18(13)} & Pachucki 2001 & \\
\textbf{Residual (framework$+$lit)} & $\mathbf{-181.55}$ MHz / $-1.29\%$
& Z-amplified bracket-completion gap & \\
\bottomrule
\end{tabular}
\end{table}

\textbf{Z-scaling structural verification (the load-bearing test).}
Three component ratios He$^+$/H test whether the projection-chain
machinery transfers cleanly Z=1 $\to$ Z=2:

\begin{itemize}\setlength\itemsep{0pt}
\item Lamb$_\text{SE}$ ratio = $13.218$ (vs naive $Z^4 = 16$).  The
deviation from $16$ is structurally the $(4/3) \ln(1/(Z\alpha)^2)$
piece of the Eides bracket, which shrinks by $(4/3) \ln(4) \approx
+1.85$ at Z=2.  Framework predicts $16 \times \text{bracket}(Z=2)/
\text{bracket}(Z=1) = 16 \times 0.8237 \approx 13.18$, matching observed
$13.22$ to $0.3\%$.

\item Lamb$_\text{VP}$ ratio = $15.878$ (vs naive $Z^4 = 16$).  The
$0.8\%$ deviation is the $\beta$-dependent kernel integral structure
modulating the leading Z$^4$ contact-density scaling.

\item Lamb$_\text{FNS}$ ratio = $63.599$ (vs predicted $Z^4 \times
(r_E({}^4\text{He})/r_p)^2 = 16 \times (1.9925)^2 = 63.52$).
\textbf{Bit-exact match} to 0.1\%, confirming the
\S\ref{sec:proj_charge_density} projection $Z$-scales without
architectural modification.
\end{itemize}

All three structural tests pass.  The projection chains are
\textbf{$Z$-portable}: substituting $Z=2$ into the existing pipeline
gives He$^+$ with no architectural changes.  This is what Paper~34's
projection taxonomy was designed to guarantee --- projections describe
physics, not specific $Z$ values.

\textbf{Residual attribution.}  The framework-native subtotal is
$-200.02$ MHz / $-1.43\%$ below experiment, which is structurally the
same as Paper~36's $-5.65$ MHz / $-0.55\%$ at Z=1 amplified by
$Z^4 \times \ln$-amplification.  The dominant Layer-2 contribution is
the $\alpha(Z\alpha)^6 \ln(Z\alpha)^{-2}$ Drake-Yan 1992 relativistic-
Bethe-log correction ($+14.7$ MHz), which is below precision visibility
at Z=1 ($\sim +0.05$ MHz, absorbed into Paper~36's residual) but
structurally visible at Z=2.  The remaining $\sim 180$ MHz of the deficit
comes from higher-order $\alpha(Z\alpha)^n_{\ln}$ mixed terms that are
extensively itemized in Drake-Yan 1992 and Pachucki 2001 but live
below the $\alpha^3$-truncation visibility at H precision.

\textbf{Cleanest W3 inner-factor contribution in the catalogue.}
${}^4$He polarizability contributes $+0.06$ MHz / $4 \times 10^{-9}$
relative, the smallest W3 contribution in the catalogue (cf.\
$\mu$d Lamb where deuteron polarizability is $+1.69$ meV, the
dominant Layer-2 line).  Mechanism: ${}^4$He is tightly bound (binding
energy $\sim 28$ MeV per nucleon vs deuteron $1.1$ MeV), so virtual
photonuclear excitations are far above the Lamb-shift scale.  He$^+$
Lamb is therefore the cleanest QED-at-Z=2 test with nuclear-structure
contamination at the $10^{-9}$ level --- the regime where the
\S\ref{sec:proj_spectral_action} LS-8a multi-loop QED becomes the
dominant Layer-2 contribution rather than nuclear-structure
inputs.

\textbf{Convention-mismatch surface.}  Different compilations
(Drake-Yan 1992, Pachucki 2001, Eides Tab.~7.6) split the leading
SE bracket and the higher-order $\alpha(Z\alpha)^n \ln$ corrections
differently at $Z = 2$.  The framework's $\alpha^3$ Eides bracket is
unambiguous; what depends on compilation is the partition of the
Layer-2 residual between ``SE $\alpha^3$ deficit'' (in the
Pachucki convention) and ``$\alpha(Z\alpha)^6 \ln$ Drake-Yan
relativistic-Bethe-log correction'' (in the Eides convention).  Same
total, different per-component itemization.  See
\S~\ref{sec:convention_exposures} for catalogue context.

\textbf{Structural reading.}  The autopsy demonstrates three things.
First, the framework's projection-chain machinery is \emph{Z-portable}
at the structural level: all three Z-scaling tests pass cleanly, with
the SE bracket's logarithmic Z-dependence reproduced exactly.
Second, the $-1.43\%$ native residual is the Z-amplified analog of
Paper~36's $-0.55\%$ Z=1 residual --- same physical content ($\alpha^3$
leading-order truncation), different size by $Z^4 \times \ln$.
Third, He$^+$ joins the multi-focal precision catalogue as the cleanest
QED-at-Z=2 test, with nuclear-structure Layer-2 inputs at $10^{-9}$
relative.  This is the canonical Z-scaling validation of the Paper~36
one-loop QED closure.

\subsubsection{Helium $2{}^3P$ fine-structure autopsy (first internal multi-focal entry)}
\label{sec:autopsy_he_2_3p}

\textbf{Reference.}  $\nu(2{}^3P_0 - 2{}^3P_1) = +29\,616.951(6)$~MHz,
$\nu(2{}^3P_1 - 2{}^3P_2) = +2\,291.176(15)$~MHz,
$\nu(2{}^3P_0 - 2{}^3P_2) = +31\,908.131(6)$~MHz (NIST compilation;
Pachucki--Yerokhin 2010 theoretical at sub-kHz vs experiment).

This is the first \emph{internal multi-focal} Roothaan autopsy: two electrons
on the same nucleus at structurally distinct effective nuclear charges
($Z_\text{eff}(1s) = 2$ full-nuclear, $Z_\text{eff}(2p) = 1$ full-shield).
Distinct from the cross-register multi-focal autopsies of \S\ref{sec:autopsy_lamb}
(electron $+$ I$=$1/2 nucleus) and \S\ref{sec:autopsy_muh_lamb} (electron $+$
muonic nucleus), which couple particles on different registers.

The operator-level Drake (1971) decomposition gives three components:
\begin{equation}
E({}^3P_J) = \tfrac{\zeta_{2p}}{2}\, X(J) + A_{SS}\, f_{SS}(J)
              + A_{SOO}\, f_{SOO}(J)
\end{equation}
with
$X(J) = J(J+1) - L(L+1) - S(S+1)$,
$f_{SS}(J) = (-2, +1, -\tfrac{1}{5})$ from the rank-2 6j symbol
$(-1)^{L+S+J}\cdot 6j\{1,1,J;1,1,2\}$ at $L = S = 1$ (Sprint 4 DD,
sympy-exact),
$f_{SOO}(J) = (+2, +1, -1)$ from the rank-1 6j,
and Drake combining coefficients
$A_{SS} = \alpha^2 (\tfrac{3}{50}\, M^2_\text{dir} - \tfrac{2}{5}\, M^2_\text{exch})$,
$A_{SOO} = \alpha^2 (\tfrac{3}{2}\, M^1_\text{dir} - M^1_\text{exch})$
(Sprint 3 BF-D, intrinsic-tier rationals).

\begin{table}[h]
\centering\small
\begin{tabular}{p{2.6cm} r r r r p{4.5cm} c}
\toprule
Interval & $E_{SO}$ (MHz) & $E_{SS}$ (MHz) & $E_{SOO}$ (MHz) & Total (MHz) & Projection chain & Status \\
\midrule
$P_0 - P_1$ & $-29\,198.090$ & $+23\,747.770$ & $+35\,063.226$ & $+29\,612.906$ &
\S\ref{sec:proj_fock} $\circ$ \S\ref{sec:proj_spinor} $\circ$
\S\ref{sec:proj_sturmian} (SO); $+$ (SS, SOO chains via
\S\ref{sec:proj_tensor_multipole} $\circ$ \S\ref{sec:proj_3j}) & FN \\
$P_1 - P_2$ & $-58\,396.180$ & $-9\,499.108$ & $+70\,126.452$ & $+2\,231.165$ &
(same as $P_0$-$P_1$) & FN \\
$P_0 - P_2$ & $-87\,594.270$ & $+14\,248.662$ & $+105\,189.679$ & $+31\,844.070$ &
(same as $P_0$-$P_1$) & FN \\
\midrule
\textbf{NIST} & & & & & & \\
$P_0 - P_1$ & & & & $+29\,616.951(6)$ & & \\
$P_1 - P_2$ & & & & $+2\,291.176(15)$ & & \\
$P_0 - P_2$ & & & & $+31\,908.131(6)$ & & \\
\midrule
\textbf{Residual} & & & & & & \\
$P_0 - P_1$ & & & & $-4.05$ MHz / $-0.014\%$ & & \\
$P_1 - P_2$ & & & & $-60.01$ MHz / $-2.62\%$ & partial cancellation $61.9\times$ & \\
$P_0 - P_2$ & & & & $-64.06$ MHz / $-0.20\%$ & & \\
\bottomrule
\end{tabular}
\caption{He $2{}^3P$ fine-structure operator-level Roothaan autopsy.
The angular content (Drake $J$-pattern, spin reduced m.e., bipolar
Gaunt selection) is sympy-exact at machine precision; residuals live
entirely in the radial sector (multi-electron correlation, higher-order
Breit-Pauli, $\alpha^3/\pi$ QED, recoil).  Source memo:
\texttt{debug/he\_2\_3P\_autopsy\_v1\_memo.md}.}
\label{tab:autopsy_he_2_3p}
\end{table}

\textbf{Bipolar (k$_1$, k$_2$) decomposition.}
For the (1s)(2p) configuration with $L_\text{max} = 2 l_\text{max} = 2$:
SS (K$=$2) admits direct (k$_1=0$, k$_2=2$) and exchange (k$_1=1$, k$_2=1$)
channels only; SOO (K$=$1) admits exchange (k$_1=1$, k$_2=1$) only (the
direct path is empty by parity selection).  All higher (k$_1$, k$_2$)
channels are Gaunt-forbidden.  This is the INTERNAL analog of the
cross-register Roothaan termination (B-W1a-diag finding): the
multipole termination at $L_\text{max} = 2 l_\text{max}$ holds
\emph{independent of the focal-length ratio} $Z_\text{eff}(1s) /
Z_\text{eff}(2p) = 2$ --- the channel structure is angular content,
not radial content.

\textbf{Partial-cancellation mechanism ($P_1$-$P_2$ amplification).}
The small interval is the difference SO ($-58.4$ GHz) $+$ SS
($-9.5$ GHz) $+$ SOO ($+70.1$ GHz) $= +2.2$ GHz, a structural
cancellation by factor $61.9\times$.  The framework's $60$ MHz
absolute residual on $P_1 - P_2$ is the same order as the $64$ MHz
residual on $P_0 - P_2$; the $-2.62\%$ fractional residual is the
$60$ MHz absolute divided by the $14\times$-smaller small interval.
Generic feature of multi-component fine-structure splittings (Li
$2{}^2P$ and Be $2s2p\,{}^3P$ show analogous partial-cancellation
amplification, Sprint 5 CP).

\textbf{Three-class diagnosis (per CLAUDE.md \S~1.8 directive).}
Class A (literature convention mismatch): NEGATIVE at LO Breit-Pauli ---
Drake (1971) and Pachucki--Yerokhin (2010) agree on leading-order
Breit-Pauli combining coefficients; \textbf{but POSITIVE at the
$\alpha^3(Z\alpha)^2$ multi-loop level} (V.D-prediction sprint,
2026-05-10) where Pachucki--Yerokhin~2009 reevaluated the relativistic
Bethe-log term to resolve a 3$\sigma$ disagreement on $\nu_{01}$ vs
Drake's earlier itemization (cross-ref \S\ref{sec:conv_he_bethe_log}).
The convention exposure sits inside the framework's LS-8a
$\alpha^3(Z\alpha)^2$ multi-loop wall, NOT at the LO Breit-Pauli sector
that the framework computes sympy-exactly.  Class B (framework gap,
RADIAL): POSITIVE --- residual
attributes to (i) multi-electron correlation in the (1s)(2p)
configuration (largest source, $\sim 45$ MHz on $P_0$-$P_2$ scale,
not in framework's single-configuration hydrogenic basis), (ii)
$\alpha^3(Z\alpha)^2$ QED on Breit-Pauli (LS-8a vertex sector,
$\sim 15$ MHz), (iii) $\alpha^3$ recoil (W1a cross-register
kinematic wall, $\sim 2$ MHz), (iv) higher-order $\alpha^4 + \alpha^5
+ \alpha^6$ NRQED (LS-8a multi-loop, sub-MHz).  Class C (focal-length
cataloguing): POSITIVE --- partial-cancellation amplification is
generic across atomic fine-structure splittings.

\textbf{Structural reading.}  The autopsy realizes the
structural-skeleton-scope statement at the operator level for an
internal multi-focal observable.  Five projection chains, three
operators, sympy-exact angular content, residuals in the radial
sector --- the same architecture that closes hydrogen Lamb shift at
\S\ref{sec:autopsy_lamb} closes He $2{}^3P$ fine structure here.
The internal multi-focal angular-only prediction (Observation 3)
holds at machine precision: angular Roothaan termination is
independent of focal-length ratio.

\subsubsection{Deuterium 1S hyperfine autopsy}
\label{sec:autopsy_d_hfs}

\textbf{Reference.} $\nu_\text{HFS}(D) = 327\,384\,352.522$~Hz
(Wineland--Ramsey 1972), the most precisely measured I=1 atomic HFS
transition.  Distinguished from \S\ref{sec:autopsy_21cm} by the
nuclear-spin axis: I=1 (deuteron) vs I=1/2 (proton).  Distinguished
from \S\ref{sec:autopsy_he_2_3p} by being a cross-register
(electron $+$ I=1 nucleus) rather than internal (two electrons on
the same I=1/2 nucleus) multi-focal observable.

\begin{table}[h]
\centering\small
\begin{tabular}{p{4.5cm} r p{6cm} c}
\toprule
Component & MHz & Projection chain & Status \\
\midrule
Bohr--Fermi Dirac (point nucleus, $g_e=2$, no recoil), I=1
multiplicity 3/2 & $+327.397464$ &
\S\ref{sec:proj_fock} $\circ$ \S\ref{sec:proj_spinor} $\circ$
\S\ref{sec:proj_3j} ($\hat{I}\cdot\hat{S}$, 3/2 from CG) & FN \\
$+$ Schwinger $a_e$ (Parker--Toms-verified at $+0.5\%$) &
$\times (1+\alpha/2\pi)$ &
\S\ref{sec:proj_fock} $\circ$ \S\ref{sec:proj_spinor} $\circ$
\S\ref{sec:proj_spectral_action} & FN (with calibration) \\
$+$ Reduced-mass $(1+m_e/m_d)^{-3}$ & $\times 0.999183$ &
\S\ref{sec:proj_fock} $\circ$ \S\ref{sec:proj_restmass} (variable
$m_n$: $m_p \to m_d$) & FN \\
$+$ Zemach $r_Z(D)=2.593$~fm via \S\ref{sec:proj_magnetization_density}
operator-level & $\times (1 - 98.001 \text{~ppm})$ &
\S\ref{sec:proj_fock} $\circ$ \S\ref{sec:proj_spinor} $\circ$
\textbf{\S\ref{sec:proj_magnetization_density}} (at I=1) & FN at op-level
$+$ L2 ($r_Z$ scalar) \\
\midrule
\textbf{Final $\nu_\text{HFS}(D)$} & $+327.477853$ & & \\
Experimental & $+327.384353$ & (Wineland--Ramsey 12-digit precision) & \\
\textbf{Residual} & $+0.094$ & $= +285.6$~ppm; matches PY 2010 Layer-2
budget magnitude & \\
\bottomrule
\end{tabular}
\caption{Deuterium 1S hyperfine Roothaan autopsy.  Each component
tagged to its \S\ref{sec:layer2} projection chain. Status: FN $=$
framework-native; L2 $=$ Layer-2 input.  Source memo:
\texttt{debug/d\_hfs\_autopsy\_track5\_memo.md} (multi-track sprint
Track 5, May 9 2026).}
\label{tab:autopsy_d_hfs}
\end{table}

\textbf{Operator-level Zemach at I=1.}  Component~4 exercises
\S\ref{sec:proj_magnetization_density} at the operator level via
\texttt{geovac.magnetization\_density.hydrogen\_zemach\_eides\_leading\_order}
with $r_Z(D) = 2.593$~fm (Friar--Payne~2005).  The bilinear matrix
element $\langle\hat{r}_Z\rangle$ on the Sturmian register at L=0
multipole reduces to $M_1[\rho_M] = r_Z$ to machine precision; the
operator output $-98.001197$~ppm vs Eides analytic $-2 Z m_e r_Z =
-98.001197$~ppm gives reproduction residual
$+1.42 \times 10^{-14}$~ppm $= 1.45 \times 10^{-14}\%$ of the LO shift
--- \emph{machine precision}.  No architecture change vs the H 21cm
operator-level Zemach in \S\ref{sec:autopsy_21cm}; only the scalar
parameter $r_Z(D) = 2.593$~fm replaces $r_Z(p) = 1.045$~fm.  Profile
(Gaussian vs exponential) independence at machine precision.  The
Pauli encoding is 4 non-identity terms ($II, Z_e, Z_p, Z_e Z_p$),
the same minimal sparse encoding as for H 21cm --- confirming the
\S\ref{sec:proj_magnetization_density} operator does not depend on
nuclear spin (it operates on the spatial register only).

\textbf{Operator-level $\hat{I}\cdot\hat{S}$ Hamiltonian at I=1.}
The $\hat{I}\cdot\hat{S}$ operator at $I=1$, $J=1/2$ has eigenvalues
$\{+1/2, -1\}$ on a 6-dim joint nuclear-electronic Hilbert space
$(2I+1)(2J+1) = 3 \cdot 2 = 6$.  The F=3/2 to F=1/2 splitting is
$\frac{1}{2} - (-1) = 3/2$, exactly the Clebsch--Gordan output.  At
$I=1/2$ the eigenvalues are $\{+1/4, -3/4\}$ with splitting $1$.
Multiplicity ratio D/H = $3/2 = 1.500000000000$ at machine precision.
The Pauli encoding extends from $Q_\text{tot}=2$ qubits, 3 terms
(I=1/2) to $Q_\text{tot}=3$ qubits, 10 terms (I=1) --- one extra
qubit for the larger nuclear-spin space, no new structural mechanism.
\textbf{Operator-level verdict: leading\_order\_I\_independent}, the
elevation of the May 8 precision-catalogue verdict from energy level
to operator level.

\textbf{Framework-native subtotal:} $\nu_\text{HFS}^\text{native}(D) =
327.4779$~MHz ($99.971\%$ of measurement).  \textbf{Layer-2 net:}
$+285.6$~ppm of $\nu_F$ (the framework chain overshoots; PY 2010
Layer-2 corrections close it).  The Layer-2 inputs are the scalar
$r_Z(D)$ value and the PY 2010 Layer-2 budget itemization
($-286$~ppm convention; alternate Eides Tab 7.3 gives $-150$~ppm,
the source of the well-known $\sim 5$~fm extraction artifact in the
global rZG fit, see \S\ref{sec:conv_rZG_DHFS}).  The operator
structure consuming the inputs is framework-native via
\S\ref{sec:proj_magnetization_density}, identical to the H 21cm
autopsy.

\textbf{Why D Layer-2 is deeper than H Layer-2.}  The cumulative
residual decomposes per Pachucki--Yerokhin~2010:
\begin{itemize}\setlength\itemsep{0pt}
\item Deuteron polarizability ($\sim +200$~ppm) is the dominant
component, $\sim 30\times$ the proton's contribution.  This is
QCD-internal NN dynamics (W3 inner-factor wall, Paper~18~\S~IV.6),
fundamentally because the deuteron is a weakly bound n+p system
with much larger spatial extent ($r_Z(D) = 2.5\,r_Z(p)$).
\item Multi-loop QED ($\alpha^2(Z\alpha)$, $\sim +40$~ppm) is the
LS-8a renormalization gap, $\sim 7\times$ H scaling because of the
larger nuclear extent's $\langle r\rangle$ couplings.
\item Recoil NLO beyond reduced-mass ($\sim +30$~ppm) is the W1a-D
Roothaan kernel-level recoil-mixing wall.
\item Finite-size charge ($\sim +10$~ppm) and higher Friar moments
($\sim +5$~ppm) are Layer-2 inputs via \S\ref{sec:proj_charge_density}
and \S\ref{sec:proj_magnetization_density} NLO opt-in.
\item Zemach NLO recoil-mixing $m_e/(m_e+m_d) \approx 2.7 \times 10^{-4}$
is structural noise in the electronic regime ($+0.028$~ppm).
\item Convention drift ($\pm 25$~ppm; Eides Tab~7.3 vs Pachucki--Yerokhin
2010) is the largest pre-existing literature-itemization
sensitivity, documented in \S\ref{sec:conv_rZG_DHFS}.
\end{itemize}

\textbf{Structural reading.}  The framework-native architecture is
identical for H 21cm and D 1S HFS --- same operators, same projection
chains, same scope-boundary statements.  What changes at I=1 is:
(a) the Clebsch--Gordan multiplicity (3/2 vs 1, structural angular
content), (b) the nuclear-spin Hilbert space dimension (3 vs 2,
adding one qubit and 7 Pauli terms to the I·S encoding), (c) the
Layer-2 budget magnitude ($-286$ ppm vs $-18$ ppm, scaling with the
nuclear extent and complexity).  The operator-level autopsy verifies
that all three changes are structural consequences of nuclear physics,
not framework modifications.  This is the I=1 verification of the
multi-focal architecture: the same framework that reproduces H 21cm
at $+18$~ppm reproduces D 1S HFS at the depth of the deuteron Layer-2
budget, with framework-native components (BF, $a_e$, recoil, Zemach)
at the same operator-level precision.  No new convention mismatch
surfaced beyond the pre-existing \S\ref{sec:conv_rZG_DHFS} exposure.

\subsubsection{Tritium 1S hyperfine autopsy}
\label{sec:autopsy_t_hfs}

\textbf{Reference.}  $\nu_\text{HFS}(T) = 1\,516\,701\,470.773(8)$~Hz
(Greene 2017 hydrogen-maser update of Mathur, Crampton, Kleppner,
Ramsey 1967), the most precisely measured I=1/2 atomic HFS transition
with a non-proton nucleus.  Distinguished from \S\ref{sec:autopsy_21cm}
by the mass-hierarchy axis at fixed nuclear-spin slot
($m_t \approx 3 m_p$ at I=1/2; H 21cm has I=1/2 at $m_p$).
Distinguished from \S\ref{sec:autopsy_d_hfs} by sharing the I=1/2
nuclear-spin axis with H 21cm rather than the I=1 of D 1S HFS.  This
is the cleanest separation in the catalogue of nuclear-mass-effect
($m_t$ vs $m_p$) from nuclear-spin-effect (I=1/2 same as proton).

\begin{table}[h]
\centering\small
\begin{tabular}{p{4.5cm} r p{6cm} c}
\toprule
Component & MHz & Projection chain & Status \\
\midrule
Bohr--Fermi Dirac (point nucleus, $g_e=2$, no recoil), I=1/2
multiplicity 1 & $+1515.865482$ &
\S\ref{sec:proj_fock} $\circ$ \S\ref{sec:proj_spinor} $\circ$
\S\ref{sec:proj_3j} ($\hat{I}\cdot\hat{S}$, multiplicity 1 for I=1/2) & FN \\
$+$ Schwinger $a_e$ (Parker--Toms-verified at $+0.5\%$) &
$\times (1+\alpha/2\pi)$ &
\S\ref{sec:proj_fock} $\circ$ \S\ref{sec:proj_spinor} $\circ$
\S\ref{sec:proj_spectral_action} & FN (with calibration) \\
$+$ Reduced-mass $(1+m_e/m_t)^{-3}$ & $\times 0.999454$ &
\S\ref{sec:proj_fock} $\circ$ \S\ref{sec:proj_restmass} (variable
$m_n$: $m_p \to m_t$) & FN \\
$+$ Zemach $r_Z(t)=1.762$~fm via \S\ref{sec:proj_magnetization_density}
operator-level & $\times (1 - 66.594 \text{~ppm})$ &
\S\ref{sec:proj_fock} $\circ$ \S\ref{sec:proj_spinor} $\circ$
\textbf{\S\ref{sec:proj_magnetization_density}} (at I=1/2, $r_Z(t)$) & FN at op-level
$+$ L2 ($r_Z(t)$ scalar) \\
\midrule
\textbf{Final $\nu_\text{HFS}(T)$} & $+1516.696899$ & & \\
Experimental & $+1516.701471$ & (Greene 2017, 12-digit precision) & \\
\textbf{Residual} & $-0.005$ & $= -3.0$~ppm; inside Karshenboim 2005
Layer-2 budget $+12 \pm 5$ ppm & \\
\bottomrule
\end{tabular}
\caption{Tritium 1S hyperfine Roothaan autopsy.  Each component
tagged to its \S\ref{sec:layer2} projection chain. Status: FN $=$
framework-native; L2 $=$ Layer-2 input.  Source memo:
\texttt{debug/T\_HFS\_autopsy\_track2\_memo.md} (multi-track sprint
Track 2, May 18 2026).}
\label{tab:autopsy_t_hfs}
\end{table}

\textbf{Operator-level Zemach at $r_Z(t) = 1.762$~fm.}  Component~4
exercises \S\ref{sec:proj_magnetization_density} at the operator level
via \texttt{geovac.magnetization\_density.hydrogen\_zemach\_eides\_leading\_order}
with $r_Z(t) = 1.762$~fm (Carlson~2008 review).  The bilinear matrix
element $\langle\hat{r}_Z\rangle$ on the Sturmian register at L=0
multipole reduces to $M_1[\rho_M] = r_Z$ to machine precision; the
operator output $-66.594$~ppm vs Eides analytic $-2 Z m_e r_Z = -66.594$~ppm
gives bit-identical reproduction at IEEE 754 float64 precision --- no
architecture change vs the H 21cm operator-level Zemach in
\S\ref{sec:autopsy_21cm} or the D 1S HFS Zemach in
\S\ref{sec:autopsy_d_hfs}; only the scalar parameter $r_Z(t)$ replaces
$r_Z(p)$.  Profile (Gaussian vs exponential) independence at machine
precision.  The Pauli encoding is 4 non-identity terms
($II, Z_e, Z_p, Z_e Z_p$), the same minimal sparse encoding as for H
21cm and D 1S HFS --- confirming the
\S\ref{sec:proj_magnetization_density} operator does not depend on
nuclear identity (it operates on the spatial register only with $r_Z$
as a scalar parameter; the same operator handles H, D, T at the same
4-Pauli encoding).

\textbf{Operator-level $\hat{I}\cdot\hat{S}$ Hamiltonian at I=1/2
(shared with H 21cm).}  The $\hat{I}\cdot\hat{S}$ operator at $I=1/2$,
$J=1/2$ has eigenvalues $\{+1/4, -3/4\}$ on a 4-dim joint
nuclear-electronic Hilbert space $(2I+1)(2J+1) = 2 \cdot 2 = 4$.  The
F=1 to F=0 splitting is $\frac{1}{4} - (-\frac{3}{4}) = 1$, the same
multiplicity as hydrogen.  Multiplicity ratio T/H = $1$ at machine
precision (target 1).  The Pauli encoding is $Q_\text{tot}=2$ qubits,
3 non-identity terms ($X \otimes X + Y \otimes Y + Z \otimes Z$), the
\emph{standard hydrogen hyperfine encoding} --- structurally identical
between T and H at qubit-encoding level.  Distinguished from the D 1S
HFS encoding which requires $Q_\text{tot}=3$ qubits and 10 Pauli terms
due to the larger I=1 nuclear-spin Hilbert space.
\textbf{Operator-level verdict: T 1S HFS shares the I=1/2 operator
structure with H 21cm exactly}, with only numerical Pauli
coefficients differing (proportional to $g_t$ and $r_Z(t)$ rather than
$g_p$ and $r_Z(p)$).

\textbf{Framework-native subtotal:} $\nu_\text{HFS}^\text{native}(T) =
1516.6969$~MHz ($99.9997\%$ of measurement).  \textbf{Layer-2 net:}
$-66.6 + 64.6 \approx -2$~ppm of $\nu_F$, with framework-native
$-3.0$~ppm residual sitting INSIDE the projected Karshenboim 2005
Layer-2 budget of $+12 \pm 5$ ppm.  The Layer-2 inputs are the scalar
$r_Z(t)$ value (Carlson~2008 central $1.762$~fm) and the Karshenboim
2005 itemization for sub-leading multi-loop QED + recoil NLO.  The
operator structure consuming the inputs is framework-native via
\S\ref{sec:proj_magnetization_density}, identical to the H 21cm and
D 1S HFS autopsies.

\textbf{Why T 1S HFS is the cleanest LS-8a isolation in the hadronic
catalogue.}  The cumulative residual decomposes per Karshenboim 2005:
\begin{itemize}\setlength\itemsep{0pt}
\item Multi-loop QED ($\alpha^2(Z\alpha)$, $\sim +6$~ppm) is the
LS-8a renormalization gap, identical mechanism to H 21cm (both Z=1).
\item Recoil NLO beyond reduced-mass ($\sim +3$~ppm) is the W1a-D
Roothaan kernel-level recoil-mixing wall, slightly smaller than H's
$+5.85$~ppm because the tritium recoil-mixing factor is smaller
($m_e/(m_e+m_t) \approx 1.82 \times 10^{-4}$ vs H's $5.45 \times 10^{-4}$).
\item \textbf{Tritium polarizability ($\sim +0.7$~ppm)} is the
W3 inner-factor calibration tier --- $\sim 300\times$ smaller than
deuteron's $\sim +200$~ppm because the tritium 2n+1p binding
($B(t) = 8.482$~MeV) is much tighter than the deuteron's n+p binding
($B(d) = 2.225$~MeV), making the tritium nucleus essentially compact
and rigid under external fields.
\item Hadronic VP $\sim +0.1$~ppm, finite-size charge $\sim +1.5$~ppm,
higher Friar moments $\sim +0.5$~ppm are Layer-2 inputs via
\S\ref{sec:proj_charge_density} and \S\ref{sec:proj_magnetization_density}
NLO opt-in.
\item Zemach NLO recoil-mixing $m_e/(m_e+m_t) \approx 1.82 \times 10^{-4}$
is structural noise in the electronic regime ($+0.013$~ppm).
\item Convention drift ($\pm 5$~ppm; Carlson 2008 vs Sick 2014 $r_Z(t)$)
is the largest literature-itemization sensitivity in this observable.
\end{itemize}

\textbf{Cleanest mass-vs-spin separation in the catalogue.}  Tritium
fills the I=1/2 hadronic non-proton cell that was previously empty in
the precision catalogue:

\begin{center}
\small
\begin{tabular}{c|c|c}
& I=1/2 & I=1 \\
\hline
$m_N = m_p$ & H 21cm ($+18$~ppm) & --- \\
$m_N \neq m_p$ (hadronic) & \textbf{T 1S HFS ($-3$~ppm)} & D 1S HFS ($+286$~ppm) \\
$m_N \neq m_p$ (leptonic) & Mu 1S HFS ($+199$~ppm) & --- \\
\end{tabular}
\end{center}

The framework's structural-skeleton scope statement now applies in
tritium with the cleanest isolation: T and H share the I=1/2
nuclear-spin axis exactly (identical Pauli encoding, identical I·S
eigenstructure); T and D share the non-proton nucleus axis (both
$m_N \neq m_p$); what separates T from BOTH is the combination
``I=1/2 like H, mass $\sim 3 m_p$ like D-extended''.  This makes T
1S HFS the cleanest LS-8a multi-loop QED isolation point for an
electron$+$hadronic-nucleus system after Mu 1S HFS (which has no QCD
nucleus at all).  No new convention mismatch surfaced beyond the
expected $r_Z(t)$ value drift between Carlson 2008 and Sick 2014.

\textbf{Structural reading.}  The framework-native architecture is
identical for H 21cm, D 1S HFS, and T 1S HFS at the operator level
--- same operators, same projection chains, same scope-boundary
statements.  What changes between cells of the mass $\times$ spin
matrix is:
(a) the Clebsch--Gordan multiplicity (1 vs 3/2 between I=1/2 and I=1),
(b) the nuclear g-factor $g_N^\text{atomic}$ as numerical Pauli
coefficient,
(c) the recoil factor $(1+m_e/m_N)^{-3}$ as numerical multiplicative
factor,
(d) the Zemach scalar $r_Z$ as input to \S\ref{sec:proj_magnetization_density},
(e) the Layer-2 budget magnitude (small for compact nuclei T/H,
large for extended nuclei D/$\mu$H).  All five changes are structural
consequences of nuclear physics, not framework modifications.
\textbf{The operator-level autopsy verifies that the multi-focal
architecture scales across the full mass $\times$ spin matrix of
light-nucleus HFS observables.}

\subsubsection{Helium-3 $2{}^3\!S_1$ hyperfine structure autopsy}
\label{sec:autopsy_he3_hfs}

\textbf{Reference.} $\nu_\text{HFS}(\text{He-3},\,2{}^3\!S_1) =
6\,739\,701\,177(16)$~Hz (Schuessler--Fortson--Dehmelt 1969;
Prior--Wang 1977; Rosner--Pipkin 1970), the F=3/2 to F=1/2 splitting
of the metastable $2{}^3\!S_1$ state. \textbf{First multi-electron HFS
autopsy in the catalogue.} The 2$^3\!S_1$ state has electronic
configuration (1s, 2s) with triplet coupling (S=1, L=0, J=1), and the
He-3 nucleus has I=1/2 with \emph{negative} magnetic moment
($\mu(\text{He-3}) = -2.127625\,\mu_N$, $g_\text{He3} = -4.25525$).

Distinguished from \S\ref{sec:autopsy_21cm} / \S\ref{sec:autopsy_d_hfs}
/ \S\ref{sec:autopsy_t_hfs} by being a multi-electron HFS at J=1 rather
than single-electron HFS at J=1/2.  Distinguished from
\S\ref{sec:autopsy_he_2_3p} by being an HFS (cross-electron-nucleus
spin-spin) rather than fine-structure (intra-electron spin-orbit +
spin-spin); both are multi-electron observables with He nucleus as
Layer-2 input but probe categorically different operator structures.

\begin{table}[h]
\centering\small
\begin{tabular}{p{4.5cm} r p{6cm} c}
\toprule
Component & MHz & Projection chain & Status \\
\midrule
Bohr--Fermi 2$^3\!S_1$ (multi-electron, Slater-screened
$Z_\text{eff,2s}=1.15$, $\hat{I}\!\cdot\!\hat{J}$ multiplicity 3/2
at I=1/2, J=1) & $+6650.31$ &
\S\ref{sec:proj_fock} $\circ$ \S\ref{sec:proj_spinor} $\circ$
\S\ref{sec:proj_3j} $\circ$ \S\ref{sec:proj_bipolar_drake}
(multi-electron triplet) $\circ$
\S\ref{sec:proj_phillips_kleinman} (Slater screen) & FN at op-level
(PK convention-dependent) \\
$+$ Schwinger $a_e$ & $\times(1+\alpha/2\pi)$ &
\S\ref{sec:proj_spectral_action} & FN (with calibration) \\
$+$ Reduced-mass $(1+m_e/m_\text{He3})^{-3}$ & $\times 0.999454$ &
\S\ref{sec:proj_restmass} ($m_p \to m_\text{He3}$) & FN \\
$+$ Zemach $r_Z(\text{He-3})=1.965$~fm via
\S\ref{sec:proj_magnetization_density} (Z=2 scaling) &
$\times (1-148.5\text{~ppm})$ &
\S\ref{sec:proj_fock} $\circ$ \S\ref{sec:proj_spinor} $\circ$
\textbf{\S\ref{sec:proj_magnetization_density}} (at I=1/2) &
FN at op-level $+$ L2 ($r_Z$ scalar) \\
\midrule
\textbf{Final $\nu_\text{HFS}$(He-3 $2{}^3\!S_1$)} & $+6653.41$ & & \\
Experimental & $+6739.70$ & (Schuessler/Prior, 14-digit precision) & \\
\textbf{Residual} & $-86.29$ & $= -1.28\%$ ($-12\,804$~ppm); dominant
wall is \S\ref{sec:proj_phillips_kleinman} multi-electron screening,
NOT LS-8a multi-loop QED & \\
\bottomrule
\end{tabular}
\caption{Helium-3 $2{}^3\!S_1$ hyperfine structure Roothaan autopsy.
\textbf{First multi-electron HFS in the catalogue.} The Slater-screened
contact density is a sprint-scope placeholder; the actual ab initio
\S\ref{sec:proj_phillips_kleinman} calculation would give the
framework-preferred unique value within the 10.5\% spread documented
below. Source memo: \texttt{debug/He3\_HFS\_autopsy\_track3\_memo.md}
(multi-track sprint Track 3, 2026-05-18).}
\label{tab:autopsy_he3_hfs}
\end{table}

\textbf{Multi-electron contact density (new architectural content).}
The 2$^3\!S_1$ spatial wavefunction is the antisymmetric Slater
determinant $\Psi_T = [\varphi_{1s}(r_1)\varphi_{2s}(r_2) -
\varphi_{2s}(r_1)\varphi_{1s}(r_2)]/\sqrt{2}$. The Fermi-contact
operator $\sum_i \delta^3(r_i)\,\mathbf{s}_i\!\cdot\!\mathbf{I}$
evaluated on the $|J=S=1, M_S=1\rangle$ state gives
$\langle \sum_i \delta^3(r_i)\,s_{i,z}\rangle_{M_S=1} =
(|\varphi_{1s}(0)|^2 + |\varphi_{2s}(0)|^2)/2$. The factor 1/2 is a
\emph{structural multi-electron spin projection} that does not appear
in single-electron HFS calculations: mechanism is the antisymmetric
Slater determinant ($\langle\delta^3(r_i)\rangle =
(|\varphi_{1s}|^2+|\varphi_{2s}|^2)/2$ per electron) weighted by
$\langle s_{i,z}\rangle = 1/2$ in the $M_S=1$ state. This is the
\S\ref{sec:proj_bipolar_drake} bipolar-harmonic content for HFS
observables (sibling of the He $2{}^3\!P$ spin-spin + spin-other-orbit
Drake combining coefficients).

\textbf{Operator-level \S\ref{sec:proj_magnetization_density} at He-3.}
Component~4 exercises the magnetization-density operator at I=1/2,
He-3 nucleus, with $r_Z(\text{He-3}) = 1.965$~fm (Sick~2014).
Reproduces Eides analytic leading-order Zemach to machine precision
(reproduction residual $0.0$~ppm, $0.0\%$ of the LO shift); profile
(Gaussian vs exponential) independence preserved at machine precision;
4 non-identity Pauli terms ($II, Z_e, Z_p, Z_e Z_p$) --- same minimal
sparse encoding as for H 21cm and D 1S HFS, confirming
\S\ref{sec:proj_magnetization_density} does not depend on multi-electron
structure (it acts on the spatial register only).

\textbf{Operator-level $\hat{I}\!\cdot\!\hat{J}$ at I=1/2, J=1.} The
$\hat{I}\!\cdot\!\hat{J}$ operator at $I=1/2, J=1$ has eigenvalues
$\{+1/2, -1\}$ on a 6-dim joint nuclear-electronic Hilbert space
$(2I+1)(2J+1) = 2 \cdot 3 = 6$. The F=3/2 to F=1/2 splitting is
$1/2 - (-1) = 3/2$. The D HFS case (I=1, J=1/2) has the
\emph{identical} 6-dim Hilbert space with the same $\{+1/2, -1\}$
eigenvalues (Wigner $3j$ I-J swap symmetry): multiplicity ratio
He-3/D = $1.000000000000$ at machine precision.
\textbf{Operator-level verdict: $I\leftrightarrow J$ swap symmetric}
--- the F-multiplet structure depends only on $(2I+1)(2J+1)$ and the
total angular momentum content, not on which factor is the nucleus
vs the electron-pair.

\textbf{\S\ref{sec:proj_phillips_kleinman} multi-electron screening
detectability (THE HEADLINE FINDING).} The 2$^3\!S_1$ contact density
depends sensitively on the screening prescription for the outer 2s
electron:

\begin{center}\small
\begin{tabular}{l c c r r}
\toprule
Path & $Z_\text{eff,1s}$ & $Z_\text{eff,2s}$ & $\nu_\text{HFS}$ (MHz) &
Residual (ppm) \\
\midrule
Hydrogenic (unscreened) & 2.00 & 2.00 & $7311.34$ & $+84\,817$ \\
Slater (1s $\sigma=0.85$) & 2.00 & 1.15 & $6653.41$ & $-12\,804$ \\
Full screen & 2.00 & 1.00 & $6600.52$ & $-20\,651$ \\
\bottomrule
\end{tabular}
\end{center}

\textbf{Spread across screening prescriptions: 105\,468~ppm = 10.5\%.}
This is the \emph{first precision-catalogue entry where the framework's
multi-electron screening prescription is the dominant Layer-2
systematic}, well above both the ppm-scale experimental precision
(16~Hz / 6.7~GHz = 2.4~ppb) and the LS-8a multi-loop QED budget
($\sim$tens of ppm). \S\ref{sec:proj_phillips_kleinman} multi-electron
screening is \emph{structurally detectable} at He-3 $2{}^3\!S_1$ HFS
precision --- the autopsy converts the existing
\S\ref{sec:proj_phillips_kleinman} dictionary entry from
``architectural'' (handles core-valence orthogonality in molecular
systems) to ``load-bearing for HFS precision''. The framework-native
ab initio path uses
\texttt{geovac.neon\_core.screened\_psi\_origin\_squared} on a
no-frozen-core Z=2 system; that is the named follow-on closure target.

\textbf{Convention finding (new \S\ref{sec:layer2_d}-class candidate).}
The Bohr--Fermi formula's mass slot uses $m_e/m_p$ \emph{always},
because the nuclear magneton $\mu_N = e\hbar/(2 m_p)$ is referenced
to the proton mass regardless of the actual nucleus. Track 5 (D HFS,
\S\ref{sec:autopsy_d_hfs}) used $m_d$ in this slot with a doubled
g-factor convention ($g_d^\text{atomic} = 2\mu_d/I$) that happens to
give the correct splitting by a \emph{factor cancellation} specific to
I=1 (the $m_d/m_p$ factor cancels the 2/2 g-factor doubling). For
I=1/2 nuclei (He-3, H, Mu, T) the cancellation does NOT hold; the
standard convention ($g_N = \mu/I$, $m_p$ in mass slot) is required.
Sanity check: He-3$^+$ 1s (single-electron hydrogen-like ion, no
multi-electron factor) matches experiment at \textbf{$+0.05\%$}
(framework $-8661.25$ MHz vs experimental $-8665.65$ MHz). The Track 5
convention transferred naively to He-3 is off by $3/2$; the standard
convention generalizes correctly across all I-values. This is a
candidate \S\ref{sec:layer2_d} class-(i) inter-track itemization
convention exposure.

\textbf{Framework-native subtotal:} $\nu_\text{HFS}^\text{native}
(\text{He-3 } 2{}^3\!S_1) = 6653.41$~MHz at Slater-screened contact
density. \textbf{Layer-2 net:} $-12\,804$~ppm of $\nu_\text{exp}$ at
this prescription. Decomposition:
\begin{itemize}\setlength\itemsep{0pt}
\item \S\ref{sec:proj_phillips_kleinman} multi-electron screening of
2s by 1s ($\pm 50\,000$~ppm spread across hydrogenic/Slater/full-screen):
dominant; the framework's own multi-electron Layer-2 prescription is
the load-bearing systematic.
\item Multi-electron QED Pachucki--Drachman corrections (LS-8a multi-
particle generalization, few hundred ppm): structural framework scope
ends at one-loop on single-particle Dirac-S$^3$.
\item He-3 nuclear polarizability (W3 inner-factor, few ppm): smaller
than deuteron polarizability because alpha-particle-like binding is
stronger than n+p; QCD-internal nuclear structure.
\item Multi-loop QED $\alpha^2(Z\alpha)$ (LS-8a, $\sim$tens of ppm):
the one-loop closure architecture wall.
\end{itemize}

\textbf{Structural reading.} The He-3 2$^3\!S_1$ HFS autopsy
\emph{restructures the catalogue's understanding of where the structural
skeleton's residuals sit for HFS observables}. Single-electron HFS
(H 21cm, D, T, $\mu$H, Mu) all have LS-8a multi-loop QED as the
dominant Layer-2 wall, with framework residuals ranging from
$+2$~ppm ($\mu$H) to $+286$~ppm (D). Multi-electron HFS (this entry;
candidate future entries: Li-7, Cs-133) has
\S\ref{sec:proj_phillips_kleinman} multi-electron screening as the
dominant wall, with framework residuals 50--500$\times$ deeper
($-1.28\%$ here). The framework's one-loop QED scope is preserved at
the same operator-level precision; what changes is which Layer-2 input
dominates. \textbf{This is the first precision-catalogue entry where
\S\ref{sec:proj_phillips_kleinman} is the load-bearing dictionary
entry, validating its Sprint~1 dictionary-completion inclusion
(May~15 2026).}

\subsubsection{Helium $2{}^1\!P \to 1{}^1\!S$ oscillator strength}
\label{sec:autopsy_he_oscillator}

\textbf{Reference.}  $f_\text{length} = 0.27616$ (Drake handbook;
Theodosiou 1984), oscillator strength for the He $2{}^1P \to 1{}^1S$
electric-dipole transition; $\omega_\text{exact} = 0.7799$ Ha.

The first multi-electron extension of Sprint Calc-L (hydrogen Lyman~$\alpha$,
$+0.055\%$ at depth-3) and Sprint Calc-P (hydrogen polarizability,
machine-exact at $N_{\rm basis}=2$ in the Sturmian basis).  The
structural question: does the Sturmian continuum-closing property
extend to multi-electron transition matrix elements?

The framework-native length-form oscillator strength is computed via
\[
f = 2\omega\, |\langle 1{}^1\!S | z | 2{}^1\!P_{m=0}\rangle|^2
\]
in the graph-native CI sector at fixed $M_L$, with the $1{}^1\!S$ state in
the $(l_1=l_2=0)$ ss singlet sub-block, the $2{}^1\!P$ state in the
$(l_1=0, l_2=1, M_L=0)$ sp singlet sub-block, and the transition dipole
assembled by Slater--Condon expansion of the multi-electron CI vectors.

\begin{table}[h]
\centering\small
\begin{tabular}{c c c c c c}
\toprule
$n_{\max}$ & configs & $\omega$ (Ha) & $|\langle z\rangle|^2$ & $f_\text{length}$ & err vs Drake \\
\midrule
2 & 5 & 0.9044 & 1.107 & 2.002 & $+625\%$ \\
3 & 12 & 0.7045 & 0.420 & 0.592 & $+114\%$ \\
4 & 22 & 0.6930 & 0.397 & 0.550 & $+99\%$ \\
5 & 35 & 0.6891 & 0.323 & 0.446 & $+61\%$ \\
6 & 51 & 0.6893 & 0.334 & 0.460 & $+67\%$ \\
7 & 70 & 0.6893 & 0.322 & 0.444 & $+61\%$ \\
\midrule
exact & --- & 0.7799 & 0.177 & 0.27616 & 0 \\
\bottomrule
\end{tabular}
\caption{Coulomb Sturmian convergence for He $2{}^1P \to 1{}^1S$
oscillator strength at single exponent $k_\text{orb} = Z = 2$.  Both
$\omega$ and $|\langle z\rangle|^2$ plateau by $n_{\max} = 6$ at
\emph{wrong limits} ($\omega$ low by $-12\%$, $|\langle z\rangle|^2$
high by $+82\%$).  Source memo:
\texttt{debug/he\_oscillator\_v1\_memo.md}.}
\label{tab:autopsy_he_oscillator}
\end{table}

Both $\omega$ and $|\langle z\rangle|^2$ \emph{plateau} by $n_{\max} = 6$ at
\emph{wrong limits}: $\omega$ converges $-12\%$ low and $|\langle z\rangle|^2$
converges $+82\%$ high.  The $2{}^1\!P$ state is variationally violated by
$\sim 0.08$ Ha ($E(2{}^1\!P)_\text{GeoVac} = -2.204$ vs Drake $-2.124$ Ha) ---
the same $\kappa$-induced over-binding mechanism documented for
the He $2{}^1\!S - 2{}^3\!S$ splitting (small-$Z$
graph-validity-boundary artifact, $Z_c \approx 1.84$; He at $Z = 2$
is just above; cf.\ \S\ref{sec:offprecision} He singlet-triplet row).
The angular (Wigner $3j$) and vector-photon
(\S\ref{sec:proj_vector_photon}) projections are exact; the failure
is in the CI sector.

\textbf{Structural reading.}  Sturmian's continuum-closing is a
\emph{single-focal-length} property.  For multi-electron transitions where
initial and final states have structurally different effective focal lengths
(here, doubly-screened $1s$ and singly-screened $2p$), the single-exponent
Sturmian basis cannot represent both, and the framework's CI eigenvectors
overlap incorrectly even at large $n_{\max}$.  The natural extension is the
\emph{multi-focal} architecture (\verb|geovac/cross_register_vne.py| for
chemistry; internal split-$Z$ for atomic excited states is a named
follow-on track per Sprint Calc-He recommendation).

The $+61\%$ residual at $n_{\max} = 7$ does not improve at higher $n_{\max}$
within the single-exponent architecture and is therefore tagged error code
\texttt{B} (basis architecture) in
Table~\ref{tab:catalog_off}.  This entry is a \emph{structural} negative
result that complements the He $2{}^1\!S - 2{}^3\!S$ off-precision row:
both expose the small-$Z$ graph-validity-boundary at He as the dominant
limitation for excited-state precision, separable from the framework's
clean closures on hydrogenic and ground-state observables.

\textbf{Internal multi-focal Phase D production result (May 9 follow-on).}
The extended-angular-CI architecture in
\texttt{geovac/internal\_multifocal.py} (~1521 lines, 28 fast tests $+$ 2
slow, 191/191 zero regression) lands at $f = 0.286$ ($+3.4\%$ vs Drake
$0.276$) on physically-motivated multi-exponent panels at well-conditioned
$\mathrm{cond}(S) = 400$.  The $E(2{}^1\!P)$ state is variationally
violated by $\sim 0.08$~Ha (3.8\%) due to $\kappa$-induced inter-shell
over-binding, and the variational floor sits at $+14\%$.  An exponent
sweet spot ($E1$ panel) gives $-0.88\%$ but is NOT robust ($f$ varies
$0.06$--$0.44$ across plausible exponent panels).  Sub-1\% accuracy
requires explicit $r_{12}$ correlation (Hylleraas) on a P-state basis.

\textbf{Hylleraas r$_{12}$ closure status (Track 1 follow-on, May 9).}
The Hylleraas r$_{12}$ explicit-correlation production module
\texttt{geovac/hylleraas\_r12.py} (Track 1 sprint, May 2026) reproduces
the He 1$^1$S ground state to $-0.016$~mHa vs Drake exact NR at
$\omega_{\max} = 4$ (22 basis functions) and the Hylleraas 1929 published
3p value to $+0.79$~mHa.  Cusp condition variationally captured (cusp
ratio at $u=0$ converges $0.29 \to 0.49$ toward Kato $0.5$).  The
single-exponent Hylleraas trial does not close the asymmetric-Z$_{\text{eff}}$
excited-state issue at modest basis size; the named structural closure
was \textbf{Hylleraas--Eckart double-zeta} (Hylleraas--Eckart 1933).

\textbf{Hylleraas--Eckart full Schwartz two-channel closure (May 2026,
the structural endpoint of the autopsy).}  The Hylleraas--Eckart double-$\alpha$
implementation arc (CLAUDE.md \S2 Track 5 backlog) closes with the full
Schwartz 1961 two-channel 2$^1$P trial
\[
\Psi_{2^1P}^{M=0} = \sum_q c_q^{+} (z_1+z_2)\,e^{-\alpha s}\cosh(\beta t)\,
                       s^{l_q} t^{2m_q} u^{n_q}
                 + \sum_q c_q^{-} (z_1-z_2)\,e^{-\alpha s}\sinh(\beta t)\,
                       s^{l_q} t^{2m_q} u^{n_q}
\]
delivering Drake-class oscillator strength.  At $\omega_s = 3$,
$\omega_p = 2$ (14 P-state basis functions total: 7 sym $+$ 7 antisym),
the framework gives $E(1^1\!S) = -2.903659$~Ha (vs Drake $-2.903724$,
$+0.064$~mHa), $E(2^1\!P) = -2.123744$~Ha (vs Drake $-2.123843$,
$+0.099$~mHa), $\Delta E = 0.7799$~Ha (sub-mHa match), and
\textbf{$f = 0.2705$ vs Drake $0.2761$, residual $\boxed{-2.02\%}$}.
The dipole decomposes as $D_\text{sym} = +0.269$ and
$D_\text{antisym} = +0.148$ (35\% antisym contribution; constructive
addition); the antisym channel is genuinely required for Drake-class
accuracy, not basis-size-limited.  Sub-mHa energy convergence at the
modest basis (14 functions) confirms the variational architecture is
saturated; the residual $-2\%$ reflects the $\omega_s = 3$, $\omega_p = 2$
truncation.  Production code: \texttt{geovac/hylleraas\_eckart\_pstate.py}
(universal SO(3)-averaged 3D quadrature kinetic for the 3 non-trivial
channel pairs $(-,-)$, $(+,-)$, $(-,+)$, plus the algebraic Hartree
kinetic for $(+,+)$ from Track 3.5).  Sprint memo:
\texttt{debug/hylleraas\_eckart\_track5\_closure\_memo.md}.  Driver:
\texttt{debug/he\_2p\_oscillator\_full\_channel.py}.  Data:
\texttt{debug/data/he\_2p\_oscillator\_full\_omega3\_2.json}.  Tests:
\texttt{tests/test\_hylleraas\_eckart\_pstate.py}::\texttt{TestUniversalQuadratureKinetic}
(6 new), \texttt{TestFullChannelOscillatorStrength} (2 new).

The closure validates two structural claims:
(i) the \textbf{internal multi-focal architecture} (two electrons at
distinct effective $Z$'s) is verified operationally at the $2^1P$
transition level, not just at the angular-content level (\S\ref{sec:proj_bipolar_drake});
(ii) the \textbf{2-electron contact-density cliff} identified in the
2026-05-09 multi-track sprint (this autopsy at $+61\%$ with graph-native
CI; the He $2^1S$--$2^3S$ splitting at $+11.9\%$) is closed at Drake-class
accuracy by Hylleraas--Eckart full Schwartz on the 2-electron problem.

\textbf{Wigner--Eckart $f$-formula normalization.}  For $L=0 \to L'=1$
absorption, summing over final $M_L$ states gives
$f = (2/3)\Delta E\,|\langle L'\|r\|L\rangle|^2 = 2\Delta E\,|D_z|^2$
where $D_z = \langle L'=1, M=0 | r_z | L=0\rangle$.  The factor of $2$
(rather than $2/3$ on $|D_z|^2$) absorbs the $M_L$ sum.  Verified
against hydrogen 1S$\to$2P at $f = 0.4162$.

\textbf{Honest scope (closure does NOT extend to):}
(i) the Li $2^2S_{1/2}$ HFS cliff (\S\ref{sec:autopsy_li7_hfs}, multi-electron
3-body system requires Hylleraas--CI hybrid); (ii) the Cs $Z > 20$ cliff
(\S\ref{sec:autopsy_cs_hfs}, heavy-atom screening cliff requires BBB93/KTT
$+$ Bohr--Weisskopf, structurally distinct mechanism per \S\ref{sec:autopsy_li7_hfs}
closure-path discussion).  The ``three cliffs, one mechanism'' framing
that surfaced in the multi-track sprint was tighter than the math
supports: Hylleraas--Eckart cleanly closes the 2-electron subset
(He 1$^1$S, He 2$^1$S--2$^3$S splitting, He 2$^1$P$\to$1$^1$S oscillator
strength, and prospective He-3 2$^3$S$_1$ HFS), but the Li and Cs cliffs
are separate downstream sprints.

\subsubsection{Cesium 6S$_{1/2}$ hyperfine (atomic clock)}
\label{sec:autopsy_cs_hfs}

\textbf{Reference.}  $\Delta\nu = 9\,192\,631\,770$ Hz (BIPM 1967,
SI-second-defining transition), equivalently
$A(6S_{1/2}) = \Delta\nu/4 = 2298.158$ MHz for the
Cs-133 $6S_{1/2}$, $F=4 \leftrightarrow F=3$ hyperfine transition.

\begin{table}[h]
\centering\small
\begin{tabular}{p{6.5cm} r r l}
\toprule
Component & $A$ (MHz) & Residual & Status \\
\midrule
C1: Bohr--Fermi strict ($g_e = 2$, point nucleus)
  & $783.0$ & $-65.9\%$ & FN \\
C2: Screened $|\psi_{6s}(0)|^2$ via [Xe] \texttt{FrozenCore}
  $Z_\text{eff}(r)$ (\S\ref{sec:proj_fock} $+$ \S\ref{sec:proj_pk})
  & \multicolumn{2}{c}{$1.328~\text{bohr}^{-3}$ (Richardson extrap.)} & FN \\
C3: Schwinger $a_e$ (\S\ref{sec:proj_vector_photon})
  & $783.9$ & $-65.9\%$ & FN \\
C4: Casimir relativistic enhancement $F_R = 1.555$
  (\S\ref{sec:proj_spinor})
  & $\mathbf{1219.1}$ & $\mathbf{-47.0\%}$ & FN \\
\midrule
C5: Zemach $r_Z = 6.7$~fm (\S\ref{sec:proj_magnetization_density})
  & $-0.1$ & & L2 \\
C5: Bohr--Weisskopf $+1.2\%$ (Roberts--Ginges 2015)
  & $+14.6$ & & L2 \\
L: Multi-loop QED $\sim +100$~ppm (\S\ref{sec:proj_breit_retardation})
  & $+0.1$ & & L2 \\
\midrule
\textbf{Framework-native total} & $\mathbf{1219.1}$ & $\mathbf{-47.0\%}$ & \\
\textbf{With Layer-2 corrections} & $\mathbf{1233.8}$ & $\mathbf{-46.3\%}$ & \\
\bottomrule
\end{tabular}
\caption{Cs-133 $6S_{1/2}$ hyperfine five-component Roothaan autopsy.
The $-47\%$ residual is dominated by two identified gaps:
(a) Clementi--Raimondi single-zeta screening produces
$|\psi(0)|^2 \approx 1.33~\text{bohr}^{-3}$
(outer-shell overshoot $2.68\times$ from missing radial nodes);
(b) Casimir $F_R = 1.555$ is leading-order only
(full Bohr--Weisskopf $\approx 2.6$ at $Z=55$).
Source: \texttt{debug/data/cs\_hfs\_v2.json}.}
\label{tab:autopsy_cs_hfs}
\end{table}

The framework's projection-chain machinery (\texttt{atomic\_classifier},
[Xe] \texttt{FrozenCore}, Bohr--Fermi formula, generic
$A \cdot \vec{I} \cdot \vec{J}$ Pauli wrapper) extends to $Z=55$
without structural failure.  The $-47\%$ residual is cleanly
attributed to two named kernel gaps, not to architectural limitations.

\textbf{Z-scaling assessment.}  The heavy-atom regime does not expose
structural limitations beyond those already identified at $Z=20$--$56$
in Sprint 7b screened SO splittings (Clementi--Raimondi single-zeta
screening qualitative; W1c-residual core-polarization correction
needed for sub-percent precision).

\textbf{Z-cliff diagnostic refinement (Track 2, 2026-05-09 follow-on).}
A diagnostic-only sprint (post-multi-track Roothaan autopsy + screening
kernel upgrade closeout) ran five hypothesis probes to localize the
dominant cause of the $-47\%$ Cs HFS framework-native residual.  All
four candidate alternative causes were RULED OUT cleanly: (b) the FD
solver \texttt{\_solve\_screened\_radial\_log} is faithful to $<1.2\%$
in $|\psi(0)|^2$ across hydrogenic Z=1..80 with monotone n\_grid
convergence (probe~b); (c) no factor-of-X convention bug in the
$Z_\text{eff}(r) \to \psi \to A_\text{HF}$ chain (probe~c); (d) the
multi-zeta machinery in \texttt{geovac/multi\_zeta\_orbitals.py}
itself is bug-free (STO normalization, MultiZetaOrbital reduction,
density integration to $N_\text{core}$ all correct, probe~d).  The
\textbf{dominant cause is (a) sharpened by a structural mechanism (d):
the missing radial nodes in the CR67 single-zeta hydrogenic basis.}
Probe~(a) extended on Cs Z=55 across all 11 shells of the [Xe] core
shows the cliff is shell-resolved: inner shells (1s, 2s, 2p) overshoot
by mean $1.01\times$ (well-fit), outer shells (4d, 5s, 5p) overshoot
by mean $2.68\times$ (4d alone overshoots $4.14\times$).  Probe~(d)
reveals the structural mechanism: the all-positive-coefficient
single-zeta hydrogenic basis (and the heuristic two-zeta extension)
produces s-orbitals that violate radial orthogonality at the level of
$|\langle R_{2s}|R_{3s}\rangle| = 0.86$ --- without negative
coefficients, no radial nodes can form, and the orbital amplitude
piles up near the inner-zeta peak.  Real RHF outer-shell orbitals
have $n - \ell - 1$ radial nodes whose physical role is to push
amplitude OUT to the correct radial extent by orthogonalizing against
inner shells.  The two-zeta heuristic FAILED in the wrong direction
($-47\% \to -90\%$) because it inherited the no-nodes problem and
made the screening transition steeper without fixing the position.
Probe~(e) v2 localizes the cliff onset cleanly at K (Z=19): H Z=1
ratio $|\psi_\text{FW}(0)|^2 / |\psi_\text{HF lit}(0)|^2 = 0.999$
(perfect), Na Z=11 = 1.243 (mild over), K Z=19 = 0.309 (3.2$\times$
under, cliff onset), Rb Z=37 = 0.193 (deepest), Cs Z=55 = 0.625
(partial recovery).  Source memo:
\texttt{debug/z\_cliff\_diagnostic\_track2\_memo.md}; probe scripts in
\texttt{debug/z\_cliff\_probe\_\{a,a\_extended,b,c,d,e,e\_v2\}.py}.

\textbf{BBB93/KTT GO with eyes open}: The Track~2 diagnostic confirms
the BBB93/Koga-Tatewaki-Thakkar full multi-zeta tabulation IS the
correct address for the dominant cause (negative coefficients in
BBB93 produce the radial nodes that all-positive heuristic two-zeta
cannot).  Expected outcome: BBB93 closes $\sim 20\%$ of the residual
($-47\% \to \sim -25$ to $-30\%$ for Cs), the full sub-percent closure
requiring \emph{both} the screening kernel upgrade AND the full
Bohr--Weisskopf relativistic enhancement via spinor lift
(\S\ref{sec:proj_spinor}).  Self-consistent Hartree--Fock iteration
in \texttt{geovac/neon\_core.py} is the principled long-term
alternative that closes the cliff for ALL Z (not just [Xe] core).

\textbf{Diagnostic-instrument angle.}  The framework's projection-chain
decomposition would isolate three known atomic-structure
dependencies of the percent-level Cs PNC uncertainty:
\begin{itemize}\setlength\itemsep{0pt}
\item radial wavefunction at the origin (Bohr--Weisskopf via
\S\ref{sec:proj_magnetization_density}),
\item core polarization (relativistic enhancement via
\S\ref{sec:proj_spinor}),
\item multi-loop QED in the heavy-atom regime (two-body Dirac/Breit
Layer-2 input via \S\ref{sec:proj_breit_retardation}),
\end{itemize}
each tied to a different Layer-2 input.  This is the prospective
catalogue's first heavy-atom test of the framework's projection-chain
machinery in a regime where simple hydrogenic formulas break down,
and connects to atomic-clock-grade precision ($10^{-15}$ relative,
the SI-second definition).  Source memo:
\texttt{debug/cs\_hfs\_v1\_memo.md}.

\subsubsection{Lithium-7 $2{}^2S_{1/2}$ hyperfine autopsy (first multi-electron HFS at $I=3/2$)}
\label{sec:autopsy_li7_hfs}

\textbf{Reference.}  $\nu_\text{HFS}({}^7\text{Li}, 2{}^2S_{1/2},
F=2 \leftrightarrow F=1) = 803{,}504{,}086.0(34)$~Hz
(Beckmann, Boklen, Elke 1974; confirmed Walls 2003).  Nuclear:
$I=3/2$, $g_{{}^7\text{Li}} = 2.170951(4)$ (CODATA $\mu/I/\mu_N$
convention), $m_{{}^7\text{Li}} \approx 12{,}786.4\,m_e$.

This is the \emph{first multi-electron HFS} in the catalogue at
$I=3/2$, exercising \emph{three} projection-dictionary entries
\emph{simultaneously}: \S\ref{sec:proj_magnetization_density}
(Zemach with $r_Z({}^7\text{Li}) = 3.71$ fm),
\S\ref{sec:proj_phillips_kleinman} (PK / core-valence orthogonality
between the unpaired 2s valence and the $1s^2$ closed-shell core),
and \S\ref{sec:proj_bipolar_drake} (multi-electron Fermi contact
via 1s--2s--2p CI mixing).  All previous HFS catalogue entries
are single-electron (H 21cm, $\mu$H, D HFS, Mu HFS, Ps HFS); Li-7
is the first three-electron HFS.

\begin{table}[h]
\centering\small
\begin{tabular}{p{4.8cm} r p{5.5cm} c}
\toprule
Component & MHz & Projection chain & Status \\
\midrule
Bohr--Fermi Dirac at $I=3/2$, $J=1/2$, hydrogenic $2s$
with $Z_\text{eff} = 1.279$ (CR67), multiplicity $\Delta\nu = 2 A_\text{hf}$ &
$+82.977$ &
\S\ref{sec:proj_fock} $\circ$ \S\ref{sec:proj_spinor} $\circ$
\S\ref{sec:proj_3j} ($\hat{I}\cdot\hat{S}$ at $I=3/2$;
eigenvalues $\{+3/4, -5/4\}$, mult 2) & FN (with $Z_\text{eff}$ L2) \\
$+$ Schwinger $a_e$ & $\times(1+\alpha/2\pi)$ &
\S\ref{sec:proj_fock} $\circ$ \S\ref{sec:proj_spectral_action} & FN (with calib) \\
$+$ Reduced-mass $(1+m_e/m_{{}^7\!\text{Li}})^{-3}$ & $\times 0.999765$ &
\S\ref{sec:proj_restmass} at $m_n = m_{{}^7\!\text{Li}}$ & FN \\
$+$ Zemach $r_Z({}^7\text{Li})=3.71$~fm via
\S\ref{sec:proj_magnetization_density} ($Z_\text{eff}$-scaled) &
$\times(1 - 178~\text{ppm})$ &
\textbf{\S\ref{sec:proj_magnetization_density}} at $I=3/2$ &
FN at op-level $+$ L2 ($r_Z$ scalar) \\
\midrule
\textbf{Final framework-native $\nu_\text{HFS}$} & $+83.04$ & & \\
Experimental & $+803.504$ & (Beckmann--Boklen--Elke 12-digit precision) & \\
\textbf{Residual} & $-720.5$ & $= -89.7\%$
($\sim 9.7\times$ cliff factor) & \\
\bottomrule
\end{tabular}
\caption{Lithium-7 $2{}^2S_{1/2}$ hyperfine Roothaan autopsy.  Each
component tagged to its \S\ref{sec:layer2} projection chain.  The
framework-native chain at hydrogenic-$Z_\text{eff}$ sits $\sim 9.7\times$
below experiment --- an order of magnitude cliff, NOT the
$\sim 1.10\%$ framework atomic accuracy floor (Paper FCI-A).  Source
memo: \texttt{debug/Li7\_HFS\_autopsy\_track5\_memo.md}.}
\label{tab:autopsy_li7_hfs}
\end{table}

\textbf{Operator-level $\hat{I}\cdot\hat{S}$ at $I=3/2$.}  Eigenvalues
$\{+3/4, -5/4\}$ on the 8-dim joint nuclear-electronic space
$(2I+1)(2J+1) = 4 \cdot 2 = 8$.  $F=2$ to $F=1$ splitting is
$+3/4 - (-5/4) = 2$, the expected CG output at $I=3/2$, $J=1/2$
(from $E_F = A/2[F(F+1) - I(I+1) - J(J+1)]$: $E_2 = +3A/4$,
$E_1 = -5A/4$, $\Delta E = 2A$).  Pauli encoding extends to
$Q_\text{tot} = 3$ qubits (2 nuclear binary for $I=3/2$ in 4-dim
space $+$ 1 electronic), 8 non-identity terms ---
\emph{operator-level verdict: CG multiplicity correct to machine precision}.

\textbf{Operator-level \S\ref{sec:proj_magnetization_density} at $I=3/2$.}
Component~4 applies the same operator architecture as in the H 21cm
and D HFS autopsies (\S\ref{sec:autopsy_21cm} and \S\ref{sec:autopsy_d_hfs}):
\texttt{geovac.magnetization\_density} reproduces the Eides leading-order
Zemach scalar at hydrogenic $Z=1$ to machine precision, profile (Gaussian
vs exponential) independence preserved, Pauli count $= 4$ (same as H and D,
confirming \S\ref{sec:proj_magnetization_density} is $I$-independent at the
operator level).  The Li 2s contact density scales the operator output by
$Z_\text{eff} = 1.279$ for the leading-order Zemach shift.

\textbf{Multi-electron contact-density cliff (headline finding).}
The framework-native chain sits at $-89.7\%$ residual.  This is
\emph{not} the $\sim 1.10\%$ Li atomic accuracy floor (Paper FCI-A
graph-native CI at $n_\text{max}=4$); the cliff is an order of magnitude
deeper.  $Z_\text{eff}$ sensitivity scan:

\begin{itemize}\setlength\itemsep{0pt}
\item $Z_\text{eff} = 1.279$ (CR67): $\nu = 82.98$~MHz ($-89.7\%$)
\item $Z_\text{eff} = 1.350$: $\nu = 97.58$~MHz ($-87.9\%$)
\item $Z_\text{eff} = 3.000$ (no screening): $\nu = 1070.80$~MHz ($+33.3\%$)
\item Experimental $\nu = 803.5$~MHz $\Leftrightarrow$
$Z_\text{eff}^\text{effective} \approx 2.73$
\end{itemize}

The effective $Z_\text{eff}$ for Li 2s \emph{contact density} sits
near $Z=2.7$, not $Z=1.3$, despite the orbital's spatial extent being
screened to $Z_\text{eff} \sim 1.3$.  This is the well-known
\emph{Fermi contact enhancement} in alkali atoms: the
\S\ref{sec:proj_phillips_kleinman} PK constraint
$\langle 2s | 1s_\text{core}\rangle = 0$ pulls the 2s radial node
inward and dramatically enhances $|\psi_{2s}(0)|^2$ over the hydrogenic
$Z_\text{eff}$ approximation.  Clementi--Roetti 1974 SCF
$|\psi_{2s}(0)|^2$ is $\sim 8\times$ the hydrogenic-$Z_\text{eff}=1.279$
value, and \S\ref{sec:proj_bipolar_drake} bipolar-harmonic core
polarization (Lindgren 1985; Yan \& Drake 2003) contributes a further
$\sim 1$--$2\%$ on top.

\textbf{Why this is not the H/D HFS regime.}  Single-electron HFS
observables (H 21cm at $+18$~ppm, $\mu$H 1S HFS at $+2$~ppm, D 1S HFS at
$+286$~ppm) hit \emph{nuclear-structure} Layer-2 walls (W3 inner-factor
deuteron polarizability, multi-loop QED).  Multi-electron HFS hits a
prior \emph{electronic-structure} wall: the framework's
single-zeta hydrogenic basis cannot represent
\S\ref{sec:proj_phillips_kleinman} core-valence orthogonality at the
amplitude-at-origin level.  Until this is engineered, neither $r_Z$
nor multi-loop QED is independently testable in multi-electron HFS.

\textbf{Two cliffs, one mechanism.}  The Li-7 multi-electron
contact-density cliff (this autopsy) and the Cs HFS Z\textgreater 20 cliff
(\S\ref{sec:autopsy_cs_hfs} Track-2 diagnostic) are \emph{two manifestations
of the same structural limitation}: the all-positive-coefficient
single-zeta hydrogenic basis cannot represent the PK-orthogonalized SCF
radial wavefunction.  The Z\textgreater 20 cliff manifests as a
$\sim 50\%$ undershoot on Cs HFS (single-zeta gives $-47\%$, heuristic
two-zeta worse at $-90\%$, BBB93 multi-zeta would close $\sim 20\%$).
The Z$=3$ cliff manifests as a $\sim 90\%$ undershoot on Li HFS because
the 2s valence is the \emph{first} shell beyond the closed core where
PK orthogonality matters most.  Same mechanism, same direction, much
larger magnitude at Z$=3$ than at Z$=55$.

\textbf{New §V.D-class exposure (class iii).}  This autopsy surfaces
a new exposure class distinct from the four existing §V.D entries
(all literature convention mismatches): the \emph{multi-electron
contact-density cliff} is a framework-ARCHITECTURE limitation, not a
literature convention nor a framework precision ceiling.  Magnitude:
$\sim 10\times$ on $A_\text{hf}$ for Li-7; scales similarly for Na, K,
Rb, Cs valence HFS.  See \S\ref{sec:conv_alkali_hfs}.

\textbf{Named structural follow-on.}  \emph{Hylleraas-Eckart double-zeta
for Li $2{}^2S$ HFS:} extend \texttt{geovac/hylleraas\_r12.py} from the
He singlet ground state (currently $+0.0006\%$ at $\omega=4$; see
\S\ref{sec:autopsy_he_oscillator}) to 3-electron Hylleraas with
$r_{12}, r_{13}, r_{23}$.  Expected closure: framework Li atomic accuracy
$1.10\% \to$ sub-$0.1\%$ \emph{simultaneously} with closure of the
$|\psi_{2s}(0)|^2$ cliff (Hylleraas explicitly satisfies PK orthogonality
at the amplitude level, no $Z_\text{eff}$ fitting required).  Downstream
impact: unlocks separable testing of \S\ref{sec:proj_magnetization_density}
(Zemach via $r_Z$), \S\ref{sec:proj_phillips_kleinman} (PK orthogonality),
\S\ref{sec:proj_bipolar_drake} (bipolar core polarization), and multi-loop
QED in multi-electron HFS.  The alkali HFS frontier (Li, Na, K, Rb, Cs)
then becomes a testbed for the framework's projection-chain dictionary
at $Z \in [3, 55]$.

\subsubsection{Positronium 1S--2S two-photon transition (first equal-mass two-body autopsy)}
\label{sec:autopsy_ps_1s2s}

\textbf{Reference.}  $\nu(1^3 S_1 \to 2^3 S_1) = 1{,}233{,}607{,}216.4(3.2)$
MHz (Fee, Chu, Mills, Mader, Mills, Chichester, \emph{Phys.\ Rev.\ A}
\textbf{48}, 192, 1993; $2.6$~ppb).

The first \S~V.C autopsy of an \emph{equal-mass two-body} system, and the
first operator-level test of the just-added 16th projection
(\S~\ref{sec:proj_breit_retardation} Breit retardation, added 2026-05-08)
since its introduction.  Continues the multi-observable focal-length
decomposition program (\S~1.8 directive) at the cleanest known instance of
the multi-focal-composition wall.

\begin{table}[h!]
\centering\small
\caption{Five-component Roothaan decomposition of the positronium 1S--2S
two-photon transition.  Framework-native subtotal ($1 + 3$) $=$
$1{,}233{,}687{,}096.0775$ MHz reproduces the May-8 Ps 1S-2S sprint
baseline at $+64.75$~ppm vs Fee 1993; the $+80$~GHz overshoot is the
missing $\alpha^4$ Breit retardation content (Component 2 = first
operator-level test of \S~\ref{sec:proj_breit_retardation}).
Cumulative ($1+2+3+4+5$) closes to $+0.0075$~MHz, within Fee 1993
$\pm 3.2$~MHz uncertainty.  Source memo:
\texttt{debug/ps\_1s2s\_autopsy\_track4\_memo.md}.}
\label{tab:autopsy_ps_1s2s}
\begin{tabular}{p{4.6cm} r p{4.5cm} l}
\toprule
Component & Value (MHz) & Projection chain & Status \\
\midrule
1.~Bohr at $m_\text{red}(ee) = 0.5$ & $+1{,}233{,}690{,}735.0941$ &
\S\ref{sec:proj_fock} $\circ$ \S\ref{sec:proj_sturmian} $\circ$
\S\ref{sec:proj_restmass} & FN \\
2.~$\alpha^4$ Breit / two-body Dirac & $-79{,}861.9$ &
\S\ref{sec:proj_fock} $\circ$ \S\ref{sec:proj_sturmian} $\circ$
\S\ref{sec:proj_restmass} $\circ$
\textbf{\S\ref{sec:proj_breit_retardation}} $+$ Layer-2 BS expansion
& L2 (op-level kernel: FN) \\
3.~Eides~\S~3.2 SE Lamb-shift differential & $-3{,}639.0166$ &
\S\ref{sec:proj_fock} $\circ$ \S\ref{sec:proj_sturmian} $\circ$
\S\ref{sec:proj_restmass} $\circ$ \S\ref{sec:proj_spectral_action}
$\circ$ \S\ref{sec:proj_drake_swainson} & FN \\
4.~Annihilation channel ($e^+e^- \to \gamma$) & $-9.3100$ &
Layer-2 input (LS-8a vertex sector) & L2 \\
5.~Multi-loop $\alpha^6 + \alpha^7$ & $-8.4600$ &
Layer-2 input (LS-8a renormalization) & L2 \\
\midrule
\textbf{Framework-native subtotal} (1+3) & $+1{,}233{,}687{,}096.08$ &
\multicolumn{2}{l}{$+64.75$~ppm vs Fee 1993} \\
\textbf{Cumulative} (1+2+3+4+5) & $+1{,}233{,}607{,}216.41$ &
\multicolumn{2}{l}{$+5\times 10^{-6}$ ppm; within Fee~1993 $\pm 3.2$ MHz} \\
\bottomrule
\end{tabular}
\end{table}

\textbf{Operator-level rest-mass projection (Component 1, FN).}  The
ratio $\nu_\text{Ps}/\nu_\text{H} = 0.5002723085$ matches
$m_\text{red}(ee)/m_\text{red}(ep) = 0.5002723085$ to relative residual
$0.00 \times 10^{0}$ (bit-identical).  This is the operator-level
verification of \S~\ref{sec:proj_restmass} ring-preservation in the
equal-mass limit at the same observable type as Mu 1S-2S (May-8 Track 1
sprint, $-0.11$ ppm).  Together the two cover the rest-mass projection
across the lepton-mass-ratio axis (Mu, $\lambda_l/\lambda_n \sim 1/207$;
Ps, $\lambda_l/\lambda_n = 1/1$).

\textbf{Operator-level Breit retardation (Component 2, partial FN).}  The
\S~\ref{sec:proj_breit_retardation} kernel is implemented as production
module \texttt{geovac.breit\_integrals} (26 tests, exact
Fraction/sympy arithmetic).  Closed-form matrix elements at $Z=1$, $k=0$:

\begin{table}[h]
\centering\small
\begin{tabular}{l c}
\toprule
Orbital pair & $R^0_\text{BP}(Z=1)$ \\
\midrule
$(1s,1s; 1s,1s)$ & $-5 + 8\log 2 \approx 0.5452$ \\
$(1s,1s; 2s,2s)$ & $4/81 \approx 0.0494$ \\
$(1s,2s; 1s,2s)$ & $-4\log 2 - 19/9 + 9\log(3)/2 \approx 0.0601$ \\
$(2s,2s; 2s,2s)$ & $-175/256 + \log 2 \approx 0.0096$ \\
\bottomrule
\end{tabular}
\end{table}

\noindent
The kernel lives in $\mathbb{Q}[\log 2, \log 3]$ (Paper~18 embedding-log
content; the Mellin-regularized inner integral produces $\log p$ for
small primes $p$ from the negative-power singularity treatment).

Three operator-level subtests:
\begin{itemize}\setlength\itemsep{0pt}
\item \textbf{Test 1 (kernel implementation): PASS.}  Closed-form
expressions for all four tested orbital pairs; production-grade module.
\item \textbf{Test 2 (Roothaan multipole termination at equal mass):
PASS.}  For $(1s,1s)$ and $(2s,2s)$ products $l_a = l_b = 0$, only $k=0$
contributes by Wigner-3j triangle inequality (\S~\ref{sec:proj_3j}).
The termination is angular content, \emph{not} a small-parameter
expansion --- it holds independent of mass ratio, re-confirming the
Ps HFS sprint finding for the 1S-2S observable.
\item \textbf{Test 3 (state-specific bound-state energy at $\alpha^4$):
PARTIAL.}  The framework implements (a) the radial kernel, (b) the
angular projection via \S~\ref{sec:proj_3j}, (c) the spinor sector
via \S~\ref{sec:proj_spinor}.  The full bound-state matrix element
$\langle \text{Ps}, 1S | H_B | \text{Ps}, 1S\rangle - \langle
\text{Ps}, 2S | H_B | \text{Ps}, 2S\rangle$ in the equal-mass two-body
Dirac normalization requires the Bethe--Salpeter quasi-energy expansion
machinery, which is structurally outside the framework's single-particle
Eides~\S~3.2 SE bracket.  This is the named structural extension of
\S~\ref{sec:proj_breit_retardation}.
\end{itemize}

\noindent
\textbf{\S~\ref{sec:proj_breit_retardation} verdict: OPERATOR-LEVEL
PARTIALLY VERIFIED.}  Kernel-level + Roothaan termination at equal mass
$=$ OK at operator level; full bound-state energy evaluation at $\alpha^4$
requires Bethe--Salpeter machinery as the named extension target.

\textbf{Layer-2 magnitude consistency check.}  The framework-extracted
$\alpha^4$ Breit residual (calibrated against Penin--Pivovarov 1998
complete theoretical value) gives fraction
$-79{,}861.9 / (m_e c^2 \alpha^4 / h) = -0.2279$, vs canonical
Karshenboim 2005 \S~4 closed-form $-(11/48) = -0.2292$; agreement
$0.55\%$.  Match consistent with the calibration-by-construction
caveat documented in \texttt{debug/precision\_catalogue\_positronium\_1s2s\_memo.md}.

\textbf{Crystallization of the multi-focal-composition wall at $\alpha^4$.}
This autopsy identifies the operator-level signature of the wall.  The
framework handles the $\alpha^2$ Bohr level (Component 1) and the
$\alpha^3$ self-energy (Component 3) as native, with full operator-level
realization.  At $\alpha^4$ (Component 2), the wall surfaces: the Breit
retardation kernel is implemented and tested at operator level; the
\emph{bound-state matrix-element evaluation} in two-body Dirac
normalization is the named structural extension.

This is the cleanest known instance of the multi-focal-composition wall
because: (i) the wall surfaces at $\alpha^4$ --- one full order before
the LS-8a renormalization wall at $\alpha^5$ or higher; (ii) the wall
is \emph{isolated} from the renormalization wall (Components 4, 5 at
the $-9$ to $-17$ MHz scale, while Component 2 is at the $-80{,}000$ MHz
scale --- three orders of magnitude separation); (iii) the wall is at
the operator level without being at the renormalization level (kernel
implemented in closed form, angular structure preserved, multipole
termination at equal mass preserved --- only the bound-state evaluation
machinery is the named extension).

The multi-focal-composition wall (CLAUDE.md \S~1.7, Sprint HF May 2026)
is therefore made visible at \emph{the operator level} on the cleanest
possible observable, transitioning from a structural-skeleton-scope
\emph{statement} to an operator-level \emph{structural identification}
with a named extension target.

\textbf{Structural reading.}  Same architecture that closes the
H 1S Lamb shift at \S~\ref{sec:autopsy_lamb} and $\mu$H $2S$--$2P$ at
\S~\ref{sec:autopsy_muh_lamb} closes Ps 1S-2S here, with the equal-mass
two-body Dirac kinematic content (Component 2) being the structural
extension specific to $\lambda_a = \lambda_b$ that has no analog in
the unequal-mass cases.  Together with the existing Ps 1S HFS sprint
(framework-native Fermi-contact at 4/7 of LO, annihilation channel as
3/7 of LO Layer-2), the equal-mass leptonic axis is now spanned at
both observable types (HFS spin-coupling + 1S-2S level energy) and the
dominant Layer-2 mechanism shifts cleanly between observable types
(annihilation for HFS, $\alpha^4$ Breit for 1S-2S).

\subsubsection{HD molecule $J=1$ rotational hyperfine (first molecular precision-catalogue entry)}
\label{sec:autopsy_HD_rotational_HFS}

\textbf{Reference.}  $\chi_d^\text{HD}(J=1) = 224.540(20)$~kHz
(Code-Ramsey~1971), the deuteron electric quadrupole coupling in
HD ground-state vibrational $v=0$ at rotational $J=1$.  Komasa--Pachucki
2020 ab~initio reference at CCSDT-F12 $+$ nonadiabatic ro-vibrational
averaging: $224.51(5)$~kHz.  Precision target $\sim 0.02\%$
($\sim 50$~Hz absolute).

\textbf{First operator-level test of \S\ref{sec:proj_tensor_multipole}
(\S III.19, rank-2 tensor multipole).}  Pre-track, \S III.19 was
the only Layer-2 projection with no operator-level test.  Track~4
of the multi-track Roothaan autopsy sprint (2026-05-18) produced
an order-of-magnitude Roothaan autopsy, the first
molecular precision-catalogue entry, and the first concrete
operator-level test of \S~\ref{sec:proj_tensor_multipole}.

The framework currently exposes the following ingredients of the
$\chi_d$ chain analytically:
\begin{itemize}\setlength\itemsep{0pt}
\item Bare nuclear multipole of the partner H proton at the
deuteron position via \S~\ref{sec:proj_multipole_gaunt}:
$q_\text{nuc} = 2 Z_H / R^3$ at $R = R_\text{eq} = 1.4011$~bohr.
\item $J=1, m_J=0$ rigid-rotor Wigner-3j angular reduction via
\S~\ref{sec:proj_3j}: $\langle V_{zz}\rangle = -(2/5)\,V_{zz}^\text{mol}$
(Townes--Schawlow standard).
\item $Q_d = 0.285699(15)(18)$~fm$^2$ as scalar Layer-2 input
(Komasa--Pachucki 2020), in the same projection class as the
$\langle r^2\rangle_E$ and $r_Z$ scalars of \S~\ref{sec:proj_charge_density}
and \S~\ref{sec:proj_magnetization_density}.
\item Born--Oppenheimer separation via \S~\ref{sec:proj_adiabatic_BO}.
\end{itemize}
What the framework does NOT currently expose: the molecular
electronic density $\rho_e(\vec{r})$ at the deuteron position
in position space.  Level~4 (\texttt{geovac/level4\_multichannel.py})
returns hyperspherical $(R, \alpha, l, m)$ channel amplitudes,
not $\psi(\vec{r}_1, \vec{r}_2)$, and consequently the EFG
operator $V_{zz}^\text{elec} = \partial^2 V_e/\partial z^2$ cannot
yet be integrated against the molecular density.

\begin{table}[h]
\centering\small
\begin{tabular}{p{5.5cm} r r}
\toprule
Component & $\chi_d^\text{mol}$ (kHz) & $\chi_d^\text{J=1}$ (kHz) \\
\midrule
Bare proton at $R_\text{eq}$ (\S\ref{sec:proj_multipole_gaunt}) &
$+488.13$ & $-195.25$ \\
Free H-atom 1s electron (\S\ref{sec:proj_fock}) &
$-150.52$ & $+60.21$ \\
\midrule
\textbf{Free-atom OoM total} (with $|J=1, m_J=0\rangle$ via
\S\ref{sec:proj_3j}) & $+337.61$ & $\textbf{135.01}$ \\
\midrule
Experimental & $-$ & $+224.54$ \\
Komasa--Pachucki 2020 ab~initio & $-$ & $+224.51(5)$ \\
\textbf{Framework OoM residual} & $-$ &
$\mathbf{-89.5}$~kHz, $\mathbf{-39.9\%}$ \\
\bottomrule
\end{tabular}
\caption{HD $J=1$ deuteron quadrupole coupling order-of-magnitude
autopsy.  $Q_d = 0.285699(15)(18)$~fm$^2$ as Layer-2 scalar input.
The bare-nuclear-only contribution alone reaches 87\% of experiment
(195 vs 224~kHz); the 40\% residual is the absence of molecular
bond-charge redistribution from the free-atom electronic picture.
Source: \texttt{debug/HD\_rotational\_HFS\_autopsy\_track4.py}.}
\label{tab:autopsy_HD_rotational_HFS}
\end{table}

\textbf{What the OoM result establishes.}  The framework's existing
multipole machinery (\S\ref{sec:proj_multipole_gaunt}) captures
the dominant LO contribution to within 13\% of experiment, with
the right sign, the right magnitude, and the right angular
structure (J=1 rigid-rotor Wigner-3j reduction is exact and
framework-native).  The 40\% residual after adding the free-atom
electronic screening is \emph{specific and well-localized}: it
attributes to the absence of position-space molecular electron
density at the deuteron.  This is not a framework gap at the
projection-class level; it is one missing operator (electronic
density at a point), structurally analogous to how the
\S~\ref{sec:proj_magnetization_density} operator-level extension
(May~2026) bridged $r_Z$ as a scalar Layer-2 input to a full
operator-level computation.

\textbf{Named follow-on sprint:} \texttt{HD-EFG-position-space}
(beyond sprint scale; analog of the W1a-D Roothaan recoil and
\S\ref{sec:proj_magnetization_density} operator-level extensions).
Five components, each precedent-bounded:
(1) Expose Level~4 wavefunction at a nuclear coordinate
(coordinate transform of existing $(R, \alpha)$ amplitudes);
(2) Add EFG operator $V_{zz}(\vec{R}_D) = \int (3 z_e^2 - r_e^2)/r_e^5
\, \rho_e(\vec{r}_e)\,d^3 r_e$;
(3) $|J=1, m_J=0\rangle$ rotational averaging via existing Wigner-3j
infrastructure;
(4) $v=0$ vibrational averaging over Level-4 hyperradial wavefunction;
(5) $T^{(2)}_{ij}$ rank-2 spin tensor on $|I=1\rangle$ (extension of
\texttt{hyperfine\_a\_pauli\_for\_atomic\_hfs} from rank-1 to rank-2).
Reference target: Komasa--Pachucki 2020 (J. Chem. Phys. 152, 154301).
Realistic accuracy: $\sim 5\%$ at operator-level (sub-percent requires
CCSDT-F12-class correlation that Level~4 does not provide; the
residual would then attribute to non-adiabatic / electron-correlation
Layer-2 input class, in the same way Komasa--Pachucki itemize their
ab~initio chain).

\textbf{Structural reading.}  This is the catalogue's first
molecular precision-catalogue entry (all prior \S~\ref{sec:autopsies}
autopsies are atomic) and the first operator-level test of
\S~\ref{sec:proj_tensor_multipole}.  The OoM result confirms the
\S III.19 entry's architectural claims: $Q_d$ as Layer-2 scalar
$\times$ Wigner-3j angular ($J=1$ rigid-rotor) $\times$ multipole
expansion $\times$ Born--Oppenheimer separation is consistent with
the experimental observable at order-of-magnitude precision.
\S III.19 ceases to be \emph{zero-operator-level-tests}; it now
has one OoM-level autopsy with a named follow-on track for
quantitative closure.  The pattern matches the
\S~\ref{sec:proj_magnetization_density} (Zemach) operator-level
extension trajectory: scoping autopsy first, full operator-level
sprint as named follow-on, sub-percent closure with explicit
attribution of the remaining Layer-2 budget.

\subsubsection{Muonic helium-4 ion 2$S_{1/2}$--2$P_{1/2}$ Lamb shift autopsy}
\label{sec:autopsy_muhe4_lamb}

\textbf{Reference.} $\Delta E = 1378.521(48)$ meV (CREMA 2021, Krauth
et al., \emph{Nature} \textbf{589}, 527; arXiv:2102.05728; muonic
convention $E(2P_{1/2}) - E(2S_{1/2})$).

The CREMA 2021 muonic helium-4 ion Lamb shift completes the muonic-Z=2
cell of the precision catalogue.  It is the most extreme mass-hierarchy
$\times$ Z combination tested so far: muon on a spinless ($I=0$)
high-Z nucleus.  Decomposes under the same five-component skeleton
as the muonic hydrogen autopsy~(\S\ref{sec:autopsy_muh_lamb}) and
muonic deuterium autopsy~(\S\ref{sec:autopsy_mud_lamb}), with three
architectural substitutions: (a) nuclear charge $Z = 1 \to 2$
(Sommerfeld $\beta = 2/(Z m_\text{red} \alpha)$ scales as $1/Z$;
$\beta = 1.475 \to 0.682$); (b) reduced mass $m_{\text{red},\mu p} \to
m_{\text{red},\mu^4\text{He}} = 201.069\,m_e$ (factor 1.082);
(c) nuclear charge radius $r_p \to r_\alpha = 1.6755(28)$ fm
(factor 1.993, FNS scaling $\sim 80\times$ via $(Z\alpha)^4
m_\text{red}^3 r_E^2/12$); and one major structural simplification:
(d) \emph{nuclear spin $I = 0$}, so the magnetization-density operator
(\S\ref{sec:proj_magnetization_density}) is identically zero and no
Zemach term is computed.  This is the \emph{cleanest pure-Lamb test
in the muonic catalogue}: no HFS contamination, no $\hat I \cdot \hat S$
mixing, no nuclear-spin convention questions.

\begin{table}[h!]
\centering\small
\caption{Five-component Roothaan decomposition of the muonic helium-4
ion Lamb shift (muonic convention $E(2P_{1/2})-E(2S_{1/2})$).  The
framework-native subtotal $+1358.24$ meV comes from full Uehling
integration via \S\ref{sec:proj_spectral_action}, the self-energy
bracket under rest-mass projection
(\S\ref{sec:proj_restmass}~+~\S\ref{sec:proj_sturmian}), and the
closed-form Foldy charge-density at the $\alpha$-particle's measured
RMS charge radius via \S\ref{sec:proj_charge_density}.  No
magnetization-density (\S\ref{sec:proj_magnetization_density}) is
activated, since $^4$He has $I=0$.  Layer-2 inputs from
Diepold~2018 / Krauth~2021 (multi-loop QED, recoil NLO,
$\alpha$-particle polarizability) close the remaining gap to the
CREMA 2021 measurement modulo two convention surfaces (SE bracket
recoil-mixing, Friar HO absorption) documented below.  Source memo:
\texttt{debug/mu\_he4\_lamb\_autopsy\_track\_b\_memo.md}.}
\label{tab:autopsy_muhe4_lamb}
\begin{tabular}{p{4.7cm} r p{4.8cm} l}
\toprule
Component & Value (meV) & Projection chain & Status \\
\midrule
1.~Full Uehling VP & $+1665.7731$ &
\S\ref{sec:proj_fock} $\circ$ \S\ref{sec:proj_spectral_action}
& FN ($-0.032\%$ vs Diepold 2018) \\
2.~Self-energy (Eides bracket) & $-11.8583$ &
\S\ref{sec:proj_fock} $\circ$ \S\ref{sec:proj_sturmian}
$\circ$ \S\ref{sec:proj_restmass}
& FN ($+8.4\%$ vs Diepold, recoil-mixing gap shrunken at higher Z) \\
3.~Foldy FNS (closed-form Eides) & $-295.6714$ &
\S\ref{sec:proj_fock} $\circ$ \S\ref{sec:proj_charge_density}
$\circ$ \S\ref{sec:proj_restmass}
& FN ($+2.9\%$ vs Diepold leading, $+4.0\%$ vs Diepold full FNS) \\
4.~Multi-loop QED (KS + 3L VP + VPVP/VPSE mixed + 2L SE + $\alpha^7$) & $+0.787$ &
Layer-2 input & LS-8a renormalization wall \\
5.~Recoil corrections (relativistic + binding) & $-5.941$ &
Layer-2 input & W1a sub-leading recoil \\
6.~\textbf{$\alpha$-particle polarizability} (+ hadronic VP) & $\mathbf{+4.061}$ &
Layer-2 input & \textbf{W3 inner-factor (QCD; small)} \\
7.~FS-NS cross + residual higher-Friar & $+1.510$ &
Layer-2 input & FNS cross-channel sub-leading \\
\midrule
\textbf{Total} & $+1358.66$ & \multicolumn{2}{l}{Residual
$-1.44\%$ vs CREMA $+1378.521(48)$ meV (Krauth 2021)} \\
\bottomrule
\end{tabular}
\end{table}

\textbf{Headline framework computation.}  The full Uehling integration
at $\beta = 0.682$ (vs $\beta = 1.475$ at $\mu$H, $\beta = 1.400$ at
$\mu$d) matches the Diepold~2018 / Pachucki~1996 reference at
$-0.032\%$ --- the cleanest single-component muonic VP match in
the catalogue at heavy-Z, transporting bit-identical kernel
infrastructure from the $\mu$H regime to the $\mu^4$He$^+$ regime.
This is the architectural transport theorem of the muonic Lamb
catalogue: the same Uehling-kernel integral closes verbatim under
$Z = 1 \to 2$ and $m_\text{red} = 185.84 \to 201.07$.

\textbf{Contact-form Uehling failure escalates at heavy Z.}  At
$\beta = 0.682$, the contact-form Uehling overshoots full kernel by a
factor of \textbf{8.76$\times$}, vs 3.55$\times$ at $\mu$H ($\beta = 1.475$).
The contact approximation is fully out-of-regime when the bound-state
Bohr radius becomes comparable to or smaller than the $e^+ e^-$
Compton wavelength.  Full numerical kernel integration is
\emph{mandatory} for the $\mu^4$He$^+$ Lamb shift.  This sharpens
the architectural claim of \S\ref{sec:autopsy_muh_lamb}: contact-form
Uehling is correct only at $\beta \gtrsim 100$ (electronic regime);
transitional at $\beta \sim 1$ ($\mu$H, $\mu$d); fails by an order of
magnitude at $\beta < 1$ (heavy-Z muonic).

\textbf{Same SE bracket recoil-mixing gap, shrunken at higher Z.}  The
framework's one-loop SE Eides bracket gives $-11.86$ meV vs Diepold's
$-10.94$ meV, a $+8.4\%$ overshoot.  This is the same recoil-mixing gap
identified in Sprint MH Track A ($\mu$H $+24\%$) and the multi-track
Roothaan sprint Track~1 ($\mu$d $+41\%$).  The fractional gap shrinks
monotonically with $Z$ because the leading-log $(4/3)\ln(1/(Z\alpha)^2)$
term in the Eides bracket grows with $Z$ while the missing sub-leading
$\alpha(Z\alpha)^4 \times (m_\text{red}/m_n)$ recoil-mixing pieces
(Eides \S 4.2; Pachucki 1996) become a smaller percentage of the bracket.
The absolute SE shift grows as $Z^4 \cdot m_\text{red} \cdot (\text{bracket
ratio})$ (verified bit-exact in the Z-scaling cross-check below).

\textbf{Foldy FNS overshoot reveals higher-Friar absorption.}  The
framework's bare Eides Eq.~2.35 leading-order Foldy correction
overshoots Diepold's tabulated "FNS leading" by $+2.9\%$ at $r_\alpha = 1.6755$
fm, and Diepold's \emph{total} FNS (leading $+$ HO $+$ recoil) by
$+4.0\%$.  Diepold's "FNS leading" line implicitly internalizes
third-Zemach-moment and Darwin relativistic corrections (Friar 1979
Ann.\ Phys.\ 122:151) of order $(Z\alpha)^2$ at $Z=2$, totaling
$\sim 11.4$ meV (the dominant component of the $-19.86$ meV total
residual).  The framework's bare leading is structurally complete to
its operator order ($\propto \langle r^2 \rangle$); the higher-Friar
absorption is a Layer-2 input that the bare-Eides formula doesn't
natively generate.  \emph{Flagged as a new V.D-prediction candidate}:
class-(ii) closed-form-vs-operator-level (Eides bare leading vs
Diepold-tabulated FNS-leading-with-absorbed-Friar-HO), sibling to
\S\ref{sec:conv_friar_profile} (which addresses the magnetization-side
profile-convention issue).  At Z=2 the higher-Friar piece is 11.4 meV
on a 295 meV leading; at $\mu$H Z=1 the same higher-Friar piece is
$\sim 0.13$ meV on a 3.68 meV leading (sub-leading by $(Z\alpha)^2$
scaling).

\textbf{Z-scaling verification.}  The component-by-component Z-scaling
of framework values from $\mu$H to $\mu^4$He$^+$ reproduces the
predicted forms exactly:
\begin{equation}
\frac{\Delta E_\text{SE}(\mu^4\text{He}^+)}{\Delta E_\text{SE}(\mu p)}
= Z^4 \cdot \frac{m_{\text{red},\mu^4\text{He}}}{m_{\text{red},\mu p}}
  \cdot \frac{\text{bracket}(Z=2)}{\text{bracket}(Z=1)}
= 17.31 \cdot 0.826 = 14.30 \quad (\text{observed: } 14.30)
\end{equation}
\begin{equation}
\frac{\Delta E_\text{FNS}(\mu^4\text{He}^+)}{\Delta E_\text{FNS}(\mu p)}
= Z^4 \cdot \left(\frac{m_{\text{red},\mu^4\text{He}}}{m_{\text{red},\mu p}}\right)^3
  \cdot \left(\frac{r_\alpha}{r_p}\right)^2
= 16 \cdot 1.267 \cdot 3.969 = 80.45 \quad (\text{observed: } 80.45)
\end{equation}
Both ratios are reproduced to seven-digit precision under the rest-mass
projection \S\ref{sec:proj_restmass}, confirming the §III.14 ring-
preservation theorem transports the entire Eides-bracket and Foldy-FNS
machinery verbatim across $Z = 1 \to 2$ and isotope substitution.

\textbf{Dominant Layer-2 wall shifts again.}  In $\mu$H the dominant
Layer-2 line was K\"all\'en--Sabry two-loop VP ($+1.51$ meV, LS-8a
renormalization wall).  In $\mu$d it shifted to the deuteron polarizability
($+1.69$ meV, W3 inner-factor; \S\ref{sec:autopsy_mud_lamb}).  In
$\mu^4$He$^+$, the dominant Layer-2 lines are the \emph{recoil corrections}
($-5.94$ meV, W1a sub-leading) and the $\alpha$-particle polarizability
($+3.10$ meV alone, $+4.06$ meV with hadronic VP combined; W3 inner-
factor).  The $\alpha$-particle polarizability per nucleon is much smaller
than the deuteron polarizability per nucleon (the $\alpha$-particle is
tightly bound at $28.3$ MeV, with no easily accessible photonuclear
excitations compared to the deuteron's near-threshold $\sim 2.2$ MeV);
this is the structural reason the W3 wall shrinks from $\mu$d to
$\mu^4$He$^+$ despite the heavier nucleus.

\textbf{r$_\alpha$ extraction sensitivity.}  The framework's FNS
coefficient at $Z=2$ in the muonic convention gives $d \Delta E /
d r_\alpha = -352.94$ meV/fm.  The CREMA 0.048 meV measurement
uncertainty propagates to $\sigma(r_\alpha) = 0.14$ amf --- the
statistical precision of muonic atoms as nuclear-radius probes.  Taking
the framework's $-19.86$ meV residual at face value would shift the
extracted radius by $-0.056$ fm to $r_\alpha = 1.619$ fm vs CREMA's
$1.67824(13)(82)$ fm; this is \emph{not} a framework extraction (which
would use Diepold's Friar-HO-absorbed leading line, eliminating the
$+11.4$ meV convention gap), but the size of the convention-mismatch
correction at Z=2 muonic when applied to the bare Eides Eq.~2.35
leading directly.

\textbf{Structural reading.}  The autopsy validates the multi-focal
architecture under the most extreme mass-hierarchy / Z combination in
the precision catalogue.  Three framework-native components reproduce
$+1358.24$ meV; four Layer-2 inputs (one less than $\mu$H since no
Zemach is needed) close the remaining gap, with component-level
convention attributions (Friar HO absorption $+11.4$ meV; SE recoil-
mixing $+0.92$ meV; r$_\alpha$ extraction signal $\sim 7$ meV) summing
to the $-19.86$ meV residual within 1 meV.  None of these is a framework
architectural failure; all are documented convention surfaces with sibling
entries in $\mu$H and $\mu$d.  The rest-mass projection
(\S\ref{sec:proj_restmass}) transports the entire five-projection-chain
skeleton bit-identically; only the nuclear charge ($Z = 1 \to 2$), nuclear
mass ($m_p \to m_{^4\text{He}}$), and nuclear charge radius ($r_p \to
r_\alpha$) inputs change between the muonic catalogue entries.  The
spinless-nucleus simplification eliminates one projection
(\S\ref{sec:proj_magnetization_density}) entirely; the remaining four
projections give the same architecture at all three muonic systems.
The same architecture that closes the H 1S Lamb shift at
\S\ref{sec:autopsy_lamb} closes the muonic Lamb shifts at
\S\ref{sec:autopsy_muh_lamb}, \S\ref{sec:autopsy_mud_lamb}, and
$\mu^4$He$^+$ here, with full Uehling kernel kinematics (contact-form
$\to$ full kernel) and Foldy r$^2$ scaling ($r_p \to r_d \to r_\alpha$)
being the only two ingredients that change between the muonic regimes.

\subsubsection{Hydrogen Lyman-$\alpha$ $1\,^{2}S_{1/2} \to 2\,^{2}P$ oscillator strength autopsy}
\label{sec:autopsy_lyman_alpha}

\textbf{Reference.} NIST tabulates Lyman-$\alpha$ in two fine-structure
components: $f(1\,S \to 2\,P_{1/2}) = 0.13881$ and
$f(1\,S \to 2\,P_{3/2}) = 0.27764$, summing to $0.41645$.  The
non-relativistic spin-averaged textbook value is $0.4162$.  The
framework computes the non-relativistic spin-averaged value
directly, so the relevant reference is the NIST sum.

The Lyman-$\alpha$ oscillator strength is the cleanest possible
dipole-transition test in atomic physics:\ one electron, no nuclear
structure, no multi-electron correlation, no continuum integration.
The framework-native value comes out as an exact rational in $\mathbb{Q}$
after the squaring step ($|\langle d \rangle|^2$ removes the
intermediate $\sqrt{2}, \sqrt{3}, \sqrt{6}$).

\begin{table}[h!]
\centering\small
\caption{Six-component Roothaan decomposition of the Lyman-$\alpha$
oscillator strength.  Components~1--4 produce the framework-native
$f = 2^{13}/3^9 = 0.41620$ exactly.  Component~5 (reduced mass)
multiplies by $m_e/\mu_{\mathrm{red}} \approx 1.000545$, giving
$f^{\text{framework}+\mu} = 0.41642$.  Component~6 (LS-8a-class)
closes the remaining $+75$~ppm to NIST.}
\label{tab:autopsy_lyman_alpha}
\begin{tabular}{cllll}
\toprule
\# & Component & Value & Layer & Projection \\
\midrule
1 & Bohr energy gap $E_2 - E_1$ & $3/8\, E_h$ exact & 1 & \S\ref{sec:proj_fock_conformal} \\
2 & Radial dipole $\langle R_{2p}|r|R_{1s}\rangle$ & $128\sqrt{6}/243\,a_0$ & 1 & Algebraic Hydrogenic integration \\
3 & Angular factor $\langle Y_{10}|\cos\theta|Y_{00}\rangle$, $\sum_{m_f}$ & $1/\sqrt{3} \times 3$ & 1 & \S\ref{sec:proj_wigner3j} \\
4 & Oscillator-strength prefactor $(2/3) \times (1/g_i)$ & $2/3$ & 2 & \S\ref{sec:proj_vector_photon} \\
5 & Reduced mass $\mu/m_e$ in $m_e \cdot \omega \cdot |\langle r \rangle|^2$ scaling & $+545$ ppm & 2 & \S\ref{sec:proj_restmass} \\
6 & Relativistic Dirac + radiative QED & $\sim +75$ ppm & 2 & LS-8a class \\
\bottomrule
\end{tabular}
\end{table}

\textbf{Walking the algebra.}  The dipole matrix element factorizes
into radial and angular pieces.  Hydrogenic wavefunctions
$R_{1s}(r) = 2 e^{-r}$ and $R_{2p}(r) = (1/(2\sqrt{6})) r e^{-r/2}$
(atomic units, $Z=1$) give the radial integral
$\langle R_{2p} | r | R_{1s}\rangle = 128\sqrt{6}/243$ (sympy
verified by direct integration).  The angular factor
$\langle Y_{10} | \cos\theta | Y_{00}\rangle = 1/\sqrt{3}$ closes
to give $\langle 2p_0 | z | 1s\rangle = 128\sqrt{2}/243$ — exact in
$\mathbb{Q}[\sqrt{2}]$.  Squaring removes the irrational:\
$|\langle 2p_0 | z | 1s\rangle|^2 = 32768/59049 \in \mathbb{Q}$.
The sum over $m_f \in \{-1, 0, +1\}$ multiplies by $3$ (spherical
symmetry), giving $\sum_{m_f} |\langle 2p_m | r | 1s\rangle|^2 =
3 \times 32768/59049 = 32768/19683$.  Energy gap $\Delta E = 3/8\,E_h$
exact.  Combining via the oscillator-strength formula:
\begin{equation*}
f^{\mathrm{framework}}_{1s \to 2p} \;=\; \frac{2}{3} \cdot \frac{3}{8} \cdot \frac{32768}{19683} \;=\; \frac{8192}{19683} \;=\; \frac{2^{13}}{3^9} \;\approx\; 0.41620.
\end{equation*}
This is the framework-native value in exact rational form.  Every
intermediate transcendental ($\sqrt{6}, \sqrt{3}, \sqrt{2}$) cancels
out at the squaring step.

\textbf{Layer-2 residual attribution.}  Reduced-mass scaling
$f \propto m_e \cdot \omega \cdot |\langle r \rangle|^2 \propto m_e \cdot \mu \cdot (m_e/\mu)^2 = m_e \cdot (m_e/\mu)$
multiplies the framework value by $m_e/\mu_{\mathrm{red}} = 1 + m_e/m_p \approx 1.000545$.  Framework $+$ reduced mass:\ $0.41642$.  NIST $0.41645$
gives a residual of $+75$~ppm, attributable to relativistic Dirac
corrections at order $\alpha^2 Z^2 \approx 5 \times 10^{-5}$
(about $50$~ppm) plus sub-Lamb radiative content.

\textbf{Structural reading.}  This is the cleanest substrate-side
autopsy in the catalogue.  Layer~1 produces $f = 2^{13}/3^9$ exact
in $\mathbb{Q}$, with all transcendentals cancelling at the
squaring step.  Layer~2 contributes $\sim 600$~ppm split between
reduced mass (\S\ref{sec:proj_restmass}, dominant) and LS-8a-class
radiative.  The autopsy demonstrates the generic pattern of dipole
oscillator strengths in the framework:\ amplitudes can be irrational
algebraic, but $|\mathrm{amplitude}|^2$ is in $\mathbb{Q}$.

\subsubsection{Hydrogen $2P_{1/2}$--$2P_{3/2}$ fine-structure splitting autopsy}
\label{sec:autopsy_h_2p_fs}

\textbf{Reference.} NIST $E(2P_{3/2}) - E(2P_{1/2}) = 10{,}969.044$ MHz
(includes all QED).  This is arguably the cleanest precision QED
test in atomic physics:\ no nuclear-structure budget, no
multi-electron complication, no continuum integration --- just
spinor Dirac in a Coulomb field.

\begin{table}[h!]
\centering\small
\caption{Six-component Roothaan decomposition of the H $2P$
fine-structure splitting.  Layer~1 produces the closed-form
$\Delta E = \alpha^2/32\,E_h$ exactly.  CODATA $\alpha$ gives the
infinite-mass value $10{,}949.3$ MHz; reduced-mass correction gives
$10{,}943.3$ MHz.  Differential QED radiative closes the remaining
$+25.7$ MHz to NIST.}
\label{tab:autopsy_h_2p_fs}
\begin{tabular}{cllll}
\toprule
\# & Component & Value & Layer & Projection \\
\midrule
1 & Camporesi--Higuchi spinor Dirac on $S^3$, $|\lambda_n^D| = n + 3/2$ & Spectrum exact & 1 & \S\ref{sec:proj_spinor_lift} \\
2 & Fock projection $p_0 = Z/n$, maps to H bound states & Exact & 1 & \S\ref{sec:proj_fock_conformal} \\
3 & Sommerfeld/Dirac expansion (T8 sympy-verified) & Closed form & 1 & Algebraic Tier-3 \\
4 & Coefficient $1/32 \in \mathbb{Q}$ from $j$-dependent bracket & $1/32$ & 1 & $j$-angular structure \\
5 & $\alpha^2$ factor (CODATA) & $5.32514 \times 10^{-5}$ & 2 & External Class-1 calibration \\
6 & Reduced mass $\mu/m_e$ in Dirac equation & $\approx 0.99946$ & 2 & \S\ref{sec:proj_restmass} \\
7 & Differential QED radiative on 2P states & $\approx +25.7$ MHz & 2 & LS-8a class \\
\bottomrule
\end{tabular}
\end{table}

\textbf{Walking the algebra.}  Sprint Dirac-on-S³ Tier~3 Track~T8
(April 2026) verified that Darwin $+$ mass-velocity $+$ spin-orbit
on hydrogen sum exactly to the Sommerfeld/Dirac formula:
\begin{equation*}
E_{n,j} \;=\; -\frac{m_e c^2 (Z\alpha)^2}{2n^2}\left[1 + \frac{(Z\alpha)^2}{n^2}\left(\frac{n}{j+1/2} - \frac{3}{4}\right) + O(\alpha^4)\right].
\end{equation*}
For $n=2$, $Z=1$:\ bracket factor is $5/4$ at $j=1/2$ and $1/4$ at
$j=3/2$.  Subtracting and simplifying (sympy):\
\begin{equation*}
\Delta E(2P_{3/2} - 2P_{1/2}) \;=\; \frac{m_e c^2 \alpha^4}{32} \;=\; \frac{\alpha^2}{32}\, E_h.
\end{equation*}
The coefficient $1/32$ is intrinsic (in $\mathbb{Q}$, from the
$j$-dependent angular structure); the $\alpha^2$ factor is Class-1
calibration.  With CODATA $\alpha$:\ $\Delta E = 10{,}949.3$ MHz
(infinite mass).  Reduced-mass scaling
$\Delta E \propto \mu \cdot \alpha^4$ multiplies by
$\mu_{\mathrm{red}}/m_e = 1 - m_e/m_p \approx 0.99946$, giving
$10{,}943.3$ MHz.

\textbf{Layer-2 residual attribution.}  NIST $10{,}969.0$ MHz minus
framework $10{,}943.3$ MHz $= +25.7$ MHz residual ($+2349$ ppm).
This is the differential QED radiative correction between
$2P_{3/2}$ and $2P_{1/2}$:\ one-loop self-energy and vacuum
polarization act differently on the two $j$ components (the famous
Lamb shift of $2P_{1/2}$ relative to Dirac is part of this).
Order-of-magnitude check:\ $\alpha \cdot \alpha^2 \cdot E_h / 32
\sim 80$ MHz naively; observed $+26$ MHz reflects the partial
cancellation between $2P_{1/2}$ and $2P_{3/2}$ Lamb shifts (similar
magnitude, opposite-sign differential).  The framework structurally
cannot reach this --- multi-loop QED counterterms ($Z_2$, $\delta m$,
vertex renormalization) are required per the LS-8a closure pattern.
Same wall as Lamb shift, Born rule (\S\ref{sec:proj_apparatus_identity}
``Born rule paragraph''), Yukawas.

\textbf{Structural reading.}  Simplest non-trivial QED test in the
framework.  Layer~1 (substrate) produces exact
$\Delta E = (1/32) \alpha^2 E_h$ in $\mathbb{Q} \times$ Class-1
calibration.  Layer~2 contributes reduced mass ($-545$ ppm) plus
differential QED ($+2349$ ppm, dominant).  The $1/32$ coefficient
generalises across any $\Delta E(nj' - nj)$ derived from Sprint T8:\
Sommerfeld is exact in the framework, only the rational changes
per $(n, j, j')$.  Canonical ``framework gets Dirac exactly,
multi-loop is the wall'' autopsy entry.

\subsubsection{Helium $2\,^{1}S - 2\,^{3}S$ exchange splitting Roothaan walk
(first three-architecture comparison)}
\label{sec:autopsy_he_singlet_triplet}

\textbf{Reference.} NIST $E(2\,^{1}S) - E(2\,^{3}S) = 6{,}421.46$
cm$^{-1}$ $= 0.029257\,E_h = 192.51$ GHz.  The singlet sits above
the triplet (standard Hund's-rule pattern).  This observable
isolates the exchange integral cleanly:
$\Delta E(2\,^{1}S - 2\,^{3}S) = 2\,K(1s, 2s)$ at leading order
(direct Coulomb $J(1s, 2s)$ cancels between singlet and triplet).

\textbf{Texture difference from preceding autopsies.}  Lyman-$\alpha$
(\S\ref{sec:autopsy_lyman_alpha}) and H $2P$ FS
(\S\ref{sec:autopsy_h_2p_fs}) have ONE clean framework value at
the substrate level.  This observable has THREE different framework
values from three different basis architectures, spanning a factor
of $\sim 30$ between best and worst.  The autopsy structure is
therefore ``three-architecture reduction'' rather than ``single
substrate value $+$ Layer-2 correction.''

\begin{table}[h!]
\centering\small
\caption{Three-architecture framework reduction on the He
$2\,^{1}S - 2\,^{3}S$ splitting.  Single-$\alpha$ architectures
(A, B) fail by factors $1.1$ and $3.1$ respectively; double-$\alpha$
(C) closes to $1.4\%$ at sprint-scale basis truncation.}
\label{tab:autopsy_he_singlet_triplet}
\begin{tabular}{clllr}
\toprule
\# & Architecture & Basis & $\Delta E$ (cm$^{-1}$) & Error vs NIST \\
\midrule
A & Graph-native CI, single-$\alpha$ ($k = Z = 2$, $\kappa = -1/16$) at $n_{\max} = 8$ & $1{,}218$ configs & $7{,}183$ & $+11.86\%$ \\
B & Hylleraas $r_{12}$ single-$\alpha$ ($\zeta = Z = 2$ for both 1s and 2s) & $\omega = 5$ & $\sim 19{,}900$ & $+209\%$ \\
C & Hylleraas--Eckart double-$\alpha$ ($\zeta_1 = 2$, $\zeta_2 \approx 1$; Eckart 1933) & Variational & $\sim 6{,}331$ & $-1.4\%$ \\
\bottomrule
\end{tabular}
\end{table}

\textbf{Walking the architectures.}  Architecture A (graph-native CI
with $\kappa = -1/16$ inter-shell, exact Slater integrals from
\texttt{hypergeometric\_slater}, single exponent $k = Z = 2$):\
$\kappa$ inter-shell coupling adds constant amplitude between 1s
and 2s regardless of orbital extent.  For excited He
($1s, 2s$), $\kappa$ overbinds the diffuse 2s, breaking the
$K(1s,2s)$ computation.  The $2^{3}S$ state is variationally
violated at this architecture (CI eigenvalue below exact).
Same mechanism as the small-$Z$ graph-validity-boundary artifact
($Z_c \approx 1.84$, He at $Z = 2$ just above), specifically for
excited states where the radial extent matters.  Net effect on the
splitting:\ $+11.86\%$ overestimate.

Architecture B (Hylleraas $r_{12}$ single-$\alpha$, $\zeta = Z = 2$
for both 1s and 2s, no $\kappa$, explicit $r_{12}$):\ the 2s orbital
with $\zeta = 2$ is wildly too compact --- physically 2s should
have $\zeta_{\mathrm{eff}} \approx 1$ because the 1s electron
screens the nucleus.  With $\zeta = 2$ on both, $K(1s, 2s)$ is
roughly $3\times$ too large; the splitting overshoots by $+209\%$.
Framework's Sprint Calc Track~1 first-pass identified this as the
structural failure mode for excited states.

Architecture C (Hylleraas--Eckart double-$\alpha$, Eckart 1933):\
each orbital gets its own variational exponent.  $\zeta_{\mathrm{inner}}
= 2$ on 1s (full nuclear charge); $\zeta_{\mathrm{outer}} \approx 1$
on 2s (screened nuclear charge).  Variational bound preserved.
Result $-1.4\%$ (slight undershoot at $\omega = 5$ basis
truncation).  The residual is dominated by \emph{basis truncation},
not Layer-2 calibration --- extending to $\omega = 8$ would close
further (Drake's exhaustive calculation reaches sub-cm$^{-1}$).

\textbf{Structural lesson.} $K(1s, 2s)$ is fundamentally a
\emph{two-exponent problem} because the two electrons see different
effective nuclear charges.  Single-exponent architectures (A and B)
fail by factors of $3$ or $1.1$ depending on whether $\kappa$
adjacency or explicit $r_{12}$ correlation is included.
Two-exponent architecture (C) succeeds at $1.4\%$.  The framework's
\texttt{hylleraas\_r12.py} Sprint Calc Track~1 was DEFERRED 2026-05-18
in favor of Track~5 (oscillator strength); the scoping memo
confirmed Gate~2 at $-1.4\%$ for double-$\alpha$.

\textbf{Layer-1 vs Layer-2 split (Architecture C).}  Layer~1
produces $\Delta E \approx 6{,}331$ cm$^{-1}$ at $\omega = 5$
basis truncation.  The $1.4\%$ residual is \emph{algorithmic}
(basis truncation), not calibration --- closeable by increasing
$\omega$.  This is qualitatively different from Lyman-$\alpha$ and
H $2P$ FS, where the framework substrate-side value was exact and
Layer~2 residual was unavoidable LS-8a class.  Here, LS-8a-class
Layer~2 content (relativistic Breit, radiative QED) is sub-$100$
ppm --- basis truncation is the dominant residual.

\textbf{Connection to W1e wall pattern.}  Texture matches the W1e
chemistry wall in atomic form.  W1e (NaH 2026-05-23) discovered
that multi-determinant correlation on a two-electron NaH valence
active space cannot close binding even with $\kappa$ adjacency,
screened valence, multi-zeta, cross-block $h_1$, and
Phillips--Kleinman --- because multi-determinant correlation lives
on the N-representable density-matrix manifold that the substrate's
spectral-triple machinery does not natively engage.  The He
$2\,^{1}S - 2\,^{3}S$ autopsy is the same lesson applied to atomic
precision physics:\ the framework's substrate-level CI fails by
$+12\%$ because $\kappa$-adjacency cannot capture the differential
exponent structure; closing to $1.4\%$ requires the architectural
extension (Hylleraas--Eckart double-$\alpha$).  The framework can
close it, but only by extending the basis architecture, not by
improving substrate-level CI.

\textbf{Structural reading.}  Canonical ``framework architecture
choice matters for excited multi-electron states'' autopsy entry.
Three architectures produce three answers spanning factor $\sim 30$.
The structural lesson is that excited-state He requires a basis
that can host two different effective nuclear charges simultaneously
--- single-exponent bases fail by factors of $3$ or $12$ depending
on whether $\kappa$ adjacency or explicit $r_{12}$ correlation is
included.  This is NOT an LS-8a-class wall:\ the production-architecture
($-1.4\%$) residual is basis-truncation and would close to sub-cm$^{-1}$
at $\omega = 8$.  The bridge to the chemistry-side W1e wall is the
deepest cross-domain structural pattern in the multi-focal-composition
arc.

\subsubsection{Muonic hydrogen 1S hyperfine structure autopsy
(first muonic HFS; $r_Z(p)$ cross-link)}
\label{sec:autopsy_muh_hfs}

\textbf{Reference.}  $E_\text{HFS}(\mu\text{H},\,1S) = 182.725$ meV
$= 44\,183$~GHz (Karshenboim--Ivanov 2002, \emph{Phys.\ Lett.\ B}
\textbf{524}, 259; no direct experimental measurement --- muon lifetime
$\tau_\mu = 2.2\,\mu$s prevents conventional precision HFS spectroscopy).
\textbf{First muonic HFS entry in the catalogue.}  Creates the
cross-link for $r_Z(p)$ identified in the cross-observable consistency
matrix (Finding~5, May~31 2026).

The projection chain is \emph{identical} to
\S\ref{sec:autopsy_21cm} (H 21cm): Fock $\circ$ spinor $\circ$ 3$j$
$\circ$ rest-mass $\circ$ magnetization-density, with the single
architectural substitution $m_e \to m_\mu$ via
\S\ref{sec:proj_restmass}.

\begin{table}[h]
\centering\small
\begin{tabular}{p{4.5cm} r p{5.5cm} c}
\toprule
Component & GHz & Projection chain & Status \\
\midrule
BF strict ($g_\mu = 2$, point nucleus) & $+44\,115$ &
\S\ref{sec:proj_fock} $\circ$ \S\ref{sec:proj_spinor} $\circ$
\S\ref{sec:proj_3j} $\circ$ \S\ref{sec:proj_restmass}
& FN \\
$+$ Schwinger $a_\mu$ ($g_\mu = 2.0023318$) & $\times 1.001166$ &
\S\ref{sec:proj_spectral_action} & FN (with calibration) \\
$+$ Reduced mass $(m_p/(m_\mu+m_p))^3 = 0.7261$ &
(folded into BF) &
\S\ref{sec:proj_restmass} & FN \\
$+$ Zemach $r_Z = 1.045$~fm (Eides 2024) & $\times(1 - 7340\text{~ppm})$ &
\S\ref{sec:proj_fock} $\circ$ \S\ref{sec:proj_spinor} $\circ$
\textbf{\S\ref{sec:proj_magnetization_density}} & FN at op-level $+$ L2 ($r_Z$) \\
\midrule
\textbf{Final (Eides $r_Z$)} & $+43\,842$ & & \\
\textbf{Final (Karshenboim $r_Z = 1.054$)} & $+43\,839$ & & \\
Literature (Karshenboim--Ivanov) & $+44\,183$ & & \\
\textbf{Residual} & $-341$ & $= -0.77\%$; LS-8a $+$ recoil-mixing budget & \\
\bottomrule
\end{tabular}
\caption{Muonic hydrogen 1S HFS Roothaan autopsy.  First muonic HFS
entry in the catalogue; creates the $r_Z(p)$ cross-link.  The
$-0.77\%$ residual is the same LS-8a $+$ recoil-mixing budget as
electronic H 21cm (\S\ref{sec:autopsy_21cm}), scaled to muonic mass.
Source: \texttt{debug/muh\_hfs\_and\_na\_cliff.py}.}
\label{tab:autopsy_muh_hfs}
\end{table}

\textbf{Zemach sensitivity enhancement: $186\times$.}  The Zemach
correction at muonic mass is $-7340$~ppm (at $r_Z = 1.045$~fm) versus
$-39.5$~ppm at electronic mass --- an enhancement factor
$m_{\text{red},\mu p}/m_{\text{red},ep} = 185.84/0.999 = 186$.  The
Eides ($r_Z = 1.045$~fm) vs Karshenboim ($r_Z = 1.054$~fm) convention
split produces a $64$~ppm $= 2{,}792$~MHz difference in the muonic HFS,
versus $\sim 0.3$~ppm $\sim 56$~kHz at electronic H 21cm.  \textbf{The
convention split is resolvable at muonic sensitivity}: $64$~ppm is well
above any plausible framework residual budget.

\textbf{Cross-link with H 21cm (\S\ref{sec:autopsy_21cm}).}  The two
autopsies share the same Layer-2 input ($r_Z(p)$) consumed through the
same projection (\S\ref{sec:proj_magnetization_density}), at mass ratios
differing by $186\times$.  A simultaneous fit constraining $r_Z(p)$ from
both channels would:
\begin{itemize}\setlength\itemsep{0pt}
\item resolve the Eides/Karshenboim convention split (currently
$\pm 5$~ppm on H 21cm, below the $+18$~ppm framework residual;
$\pm 64$~ppm on $\mu$H HFS, resolvable);
\item test whether the LS-8a multi-loop QED budget scales correctly
from electronic to muonic mass (same physics, different kinematics);
\item provide the first muonic-electronic HFS consistency check in
the precision catalogue.
\end{itemize}

\textbf{Alkali contact-density cliff: quantitative Z-sequence.}
The cross-observable consistency analysis (Finding~4, May~31 2026)
computed the hydrogenic-$Z_\text{eff}$ BF A-constant across all five
stable alkali ground-state HFS systems, testing the cliff severity
prediction from \S\ref{sec:autopsy_li7_hfs} and
\S\ref{sec:autopsy_cs_hfs}.

\begin{table}[h]
\centering\small
\begin{tabular}{l c c c r r r r}
\toprule
Atom & $Z$ & $n$s & $Z_\text{eff}$(CR67) & $A_\text{fw}$ (MHz) &
$A_\text{exp}$ (MHz) & Cliff & $Z_\text{eff}^\text{needed}/Z_\text{eff}^\text{CR67}$ \\
\midrule
Li-7    &  3 & 2s & 1.279 &  21.0 &  401.8 & $-94.8\%$  & $19\times$ \\
Na-23   & 11 & 3s & 2.507 &   9.6 &  885.8 & $-98.9\%$  & $92\times$ \\
K-39    & 19 & 4s & 2.162 &   0.3 &  230.9 & $-99.9\%$  & $851\times$ \\
Rb-85   & 37 & 5s & 2.771 &   0.3 & 1011.9 & $-100.0\%$ & $3{,}635\times$ \\
Cs-133  & 55 & 6s & 3.475 &   0.3 & 2298.2 & $-100.0\%$ & $8{,}306\times$ \\
\bottomrule
\end{tabular}
\caption{Alkali ground-state HFS cliff sequence using
hydrogenic-$Z_\text{eff}$ contact density (CR67 single-zeta).
The rightmost column shows the ratio of true contact density to the
$Z_\text{eff}^3/(\pi n^3)$ estimate --- the quantitative measure of
core-penetration contact enhancement from PK orthogonality
(\S\ref{sec:proj_phillips_kleinman}).  The enhancement grows as
$\sim Z^{2.5}$ across the alkali series.  Source:
\texttt{debug/muh\_hfs\_and\_na\_cliff.py}.}
\label{tab:alkali_cliff_sequence}
\end{table}

The $Z_\text{eff}^\text{needed}/Z_\text{eff}^\text{CR67}$ ratio column
is the quantitative measure of the PK contact-enhancement cliff:
$19\times$ at Li (the existing \S\ref{sec:autopsy_li7_hfs} finding),
growing to $8{,}306\times$ at Cs.  The growth scaling $\sim Z^{2.5}$
reflects the increasing number of inner shells against which the valence
$ns$ orbital must orthogonalize --- each orthogonalization pulls a
radial node inward, piling amplitude at the origin.  The existing Cs
autopsy (\S\ref{sec:autopsy_cs_hfs}) achieved $-47\%$ using the full
\texttt{FrozenCore} screened wavefunction (which captures \emph{some}
core penetration); the $-100\%$ here uses the bare
hydrogenic-$Z_\text{eff}$ contact density (which captures \emph{none}).
The gap between $-47\%$ and $-100\%$ at Cs quantifies the partial
recovery from the \texttt{FrozenCore} infrastructure.

% -------------------------------------------------------------------
\subsection{Literature convention exposures (running catalogue)}
\label{sec:convention_exposures}
% -------------------------------------------------------------------

The framework's projection-chain decomposition acts as a unique
diagnostic instrument for precision atomic and molecular physics:
multi-observable global fits using a single structural dictionary
(the \S~\ref{sec:layer2} projection family) expose itemization
choices in standard-physics compilations that single-observable
analyses cannot see.  This is the directive's class~(a) finding
type --- ``literature convention mismatches'' between independent
compilations (Eides~2024, Karshenboim~2005, Krauth~2017,
Pachucki--Yerokhin~2010, Antognini~2013, Friar--Payne~2005)
on Layer-2 itemization at the sub-percent precision the framework
operates at.  Frameworks that consume single-observable Layer-2
inputs in their native conventions cannot expose these mismatches
structurally.  This subsection consolidates the six exposures
catalogued to date (V.D.1--V.D.4 from 2026-05-09 multi-track Roothaan
autopsy sprint; V.D.5--V.D.6 from 2026-05-10 V.D-prediction sprint);
each future sprint that surfaces a
class~(a) finding adds a new \S~V.D.$N$ subsection with the
five-element template \emph{Observable~$|$~Source compilations~$|$~
Convention difference~$|$~Magnitude~$|$~Sprint reference $+$ Closure verdict.}

\begin{table}[h]
\centering\small
\caption{Literature convention exposures catalogued through 2026-05-10.
Each row is documented in detail in the indicated subsection
(originating sprint memo) and cross-references the
\S~\ref{sec:layer2} projection chain whose operator-level
decomposition exposed it.  V.D.5 and V.D.6 added 2026-05-10 by the
V.D-prediction sprint testing Pattern C from the 2026-05-09 synthesis review.}
\label{tab:convention_exposures}
\begin{tabular}{p{0.6cm} p{2.6cm} p{3.0cm} p{2.6cm} p{1.6cm} p{1.6cm}}
\toprule
ID & Observable & Source compilations & Convention difference &
Magnitude & Sprint ref. \\
\midrule
\S\ref{sec:conv_rZG_DHFS} & D 1S HFS Layer-2 budget &
Eides~2024 Tab.~7.3 vs Pachucki--Yerokhin~2010 &
itemization of recoil NLO + multi-loop QED + finite-size correction &
$\sim 0.01\%$ of $\nu_F$ $\to$ $\sim 25$ mfm in $r_Z(p)$ &
Calc-rZG (\S\ref{sec:offprecision}) \\
\S\ref{sec:conv_muH_VP} & $\mu$H 2$S$-2$P$ VP one-loop &
Antognini~2013 vs Krauth~2017 &
absorption of higher-order VP into one-loop line vs separate
``VP-VP mixed'' itemization &
$\sim 100$ ppm in per-component split (totals agree) &
Sprint MH Track A / Calc-muH Lamb autopsy v1 \\
\S\ref{sec:conv_karshenboim_recoil} & $\mu$ Lamb shift SE bracket &
Framework Eides~\S 3.2 vs Karshenboim~2005 &
$m_\text{red}$ rescaling absorbs recoil NLO at bracket level vs
separately-itemized ``Recoil NLO'' line &
$\sim 11$ MHz at $\mu$ Lamb (exact match between framework
$+1074.13$ and Karshenboim $+1085.84 - 11.30$) &
Sprint precision-catalogue Mu Lamb 2026-05-08 \\
\S\ref{sec:conv_friar_profile} & $\mu$H Friar moment &
Eides Eq.~2.35 closed-form vs operator-level
\S\ref{sec:proj_magnetization_density} &
$\langle r^2\rangle^{1/2} = r_p$ (RMS) vs $\langle r\rangle = r_p$
(first moment) for the Gaussian profile &
$\langle r^2\rangle/\langle r\rangle^2 = 4/\pi \approx 1.27$
$\to$ $+12.7\%$ on Friar moment &
Sprint MH / Calc-muH Lamb autopsy v1 \\
\S\ref{sec:conv_d_polarizability} & D 1S HFS polarizability sub-component &
Friar--Payne~2005 vs Pachucki--Yerokhin~2010+ &
low-energy aggregate (FP) vs explicit polarizability $+$ Zemach split (PY) &
$\sim 240$~ppm of $E_F$; $\sim 80$~kHz absolute &
V.D-prediction sprint, 2026-05-10 \\
\S\ref{sec:conv_he_bethe_log} & He $2{}^3P$ fine structure (29~GHz; $\alpha$-determination) &
Drake~1990 / Drake--Yan~1995 vs Pachucki--Yerokhin~2009 (PRA \textbf{79}, 062516) &
Drake's original relativistic Bethe-log itemization was numerically incorrect; PY reevaluated and resolved a 3$\sigma$ disagreement on $\nu_{01}$ &
$\sim$kHz on 29~GHz; $\sim 30$~ppb on $\alpha$ &
V.D-prediction sprint, 2026-05-10 \\
\S\ref{sec:conv_d_polarizability_mud_lamb} &
$\mu$d 2$S$--2$P$ Lamb polarizability &
Carlson--Gorchtein--Vanderhaeghen~2014 vs Carlson--Vanderhaeghen~2011 &
2014 evaluation reports nuclear polarizability $+1.690(20)$ meV
(QCD-internal NN dynamics);
2011 evaluation reported $+1.94$ meV.
The two convergent QCD evaluations differ at ~$13\%$ on the
\emph{dominant} Layer-2 contribution to $\mu$d Lamb shift &
$\sim 13\%$ on Layer-2 $\to$ $\sim 0.25$~meV; dominates the
framework's own $-0.12\%$ residual on the Lamb shift &
Multi-track Roothaan sprint Track~1, 2026-05-18 \\
\S\ref{sec:conv_alkali_hfs} &
Alkali $n$s valence HFS (Li-7 2$^2$S; generalizes Na, K, Rb, Cs) &
Single-zeta hydrogenic ($Z_\text{eff}$) vs SCF + multi-electron CI &
Hydrogenic $Z_\text{eff}$ misses $|\psi_{ns}(0)|^2$ by $\sim 8\times$
for Li 2s; \emph{architecture-class (iii) exposure}, not a
literature-convention mismatch.  Same mechanism as the
\S\ref{sec:autopsy_cs_hfs} Z\textgreater 20 cliff, two manifestations
of single-zeta basis missing PK orthogonality &
$\sim 10\times$ on $A_\text{hf}$ at Z=3; scales similarly for Na,
K, Rb, Cs HFS &
Multi-track Roothaan sprint Track~5, 2026-05-18 \\
\S\ref{sec:conv_multifocal_bimodule} &
Multi-focal-composition wall taxonomy (CLAUDE.md \S~1.7) &
Single-observation reading vs.\ Hilbert $C^*$-bimodule reading
(Latr\'emoli\`ere modular propinquity) &
Wall splits into bimodule cross-shifts (HANDLED: HF-3 recoil, HF-4
Zemach, HF-5 Parker--Toms, cross-register $V_{eN}$, FNS) vs.\
module endomorphisms (NOT HANDLED: LS-8a $Z_2-1$, $\delta m$;
Sprint H1 Yukawa; inner-factor calibration data).  \emph{Architecture
class (iv) exposure}, framework-level structural taxonomy refinement &
Qualitative structural finding;
quantitative tests are LS-8a (multi-loop QED) and
W1c Layer-3 (alkali-hydride binding) &
Sub-sprint M-Z, 2026-05-23 \\
\bottomrule
\end{tabular}
\end{table}

\subsubsection{D 1S HFS Layer-2 itemization (Eides vs Pachucki--Yerokhin)}
\label{sec:conv_rZG_DHFS}

\textbf{Observable.}  Deuterium 1S hyperfine splitting
$\nu_\text{HFS}(D) = 327\,384\,352.522$ kHz (Wineland--Ramsey 1972),
used in the Sprint Calc-rZG global fit (\S\ref{sec:offprecision}, row
``Self-consistent $r_Z$ global extraction'') to extract
$r_Z(D)$ via the framework's \S\ref{sec:proj_magnetization_density}
operator.

\textbf{Source compilations.}  Eides~2024 Tab.~7.3 (used as the
canonical Layer-2 budget for H 21cm in this paper) and
Pachucki--Yerokhin~2010 (the standard atomic-physics itemization for
deuterium HFS).

\textbf{Convention difference.}  The two compilations use different
groupings for the recoil-NLO + multi-loop QED + finite-size
contributions to $\nu_F^D$.  Eides Tab.~7.3 itemizes the leading
Zemach line at $-150$~ppm; Pachucki--Yerokhin~2010, when fully
expanded with the Friar--Payne~2005 closure analysis, places the
itemization at $-286$~ppm.  Both compilations close to the same
total $\nu_\text{HFS}(D)$ at the precision the deuterium HFS
measurement reports; the difference is in how sub-leading recoil and
multi-loop contributions are aggregated.

\textbf{Magnitude.}  $\sim 0.01\%$ of $\nu_F$, propagating to
$\sim 25$~mfm in extracted $r_Z(p)$ via the global fit.

\textbf{Sprint reference and closure.}  The original Sprint Calc-rZG
fit (2026-05-09 morning) used the Eides $-150$~ppm itemization for
deuterium; this gave $r_Z(D) = 6.18(132)$~fm, the well-known
$\sim 5$~fm artifact.  The diagnostic memo
\texttt{debug/rzg\_bug\_diagnosis\_memo.md} verified that the
production cross-register $V_{eN}$ kernel
(\texttt{geovac.cross\_register\_vne}) reproduces leading-order
Bethe--Salpeter recoil at $2.86\%$ for hydrogen, $2.03\%$ for
deuterium, and $8.18\%$ for muonium, species-agnostically and
correctly.  The artifact came from the Layer-2 budget specification
for D HFS, not from the framework kernel.  Re-anchoring to the
Pachucki--Yerokhin~2010 budget at $-286$~ppm closed the extraction
to $r_Z(D) = 2.583(1588)$~fm (within $0.01\sigma$ of Friar--Payne
$2.593(16)$~fm).  \textbf{Closure verdict:} resolved by Layer-2
budget re-anchoring.  The original ``recoil double-counting''
diagnosis was a misdiagnosis; the issue was Layer-2 itemization
specification, not a kernel bug.

\textbf{Operator-level autopsy refinement (May 9 2026).}  The Track 5
operator-level autopsy of D 1S HFS (\S\ref{sec:autopsy_d_hfs})
operates exclusively at the Pachucki--Yerokhin~2010 / Friar--Payne~2005
$-286$~ppm Layer-2 budget convention.  Component-by-component the
framework reproduces (a) the $\hat{I}\cdot\hat{S}$ Hamiltonian at I=1
on the 6-dim joint nuclear-electronic Hilbert space at machine
precision (multiplicity ratio D/H $= 3/2$ to 12 digits); (b) the
\S\ref{sec:proj_magnetization_density} operator-level Zemach at
$r_Z(D) = 2.593$~fm to $\sim 1.45 \times 10^{-14}\%$ of the LO shift;
and (c) the cumulative $+285.6$~ppm chain residual matches the
$-286$~ppm PY 2010 Layer-2 budget magnitude to within the $\pm 25$~ppm
itemization drift documented above.  No new convention mismatch
surfaced at the operator level.  This sharpens the $\sim 0.01\%$ /
$\sim 25$~mfm magnitude estimate to: framework-side architecture is
identical for H and D (same operators, same chains, same scope
boundaries); the Layer-2 budget magnitudes differ by $\sim 16\times$
because the deuteron's nuclear-structure complexity (NN binding
dynamics inside a weakly bound n+p system) makes its polarizability,
multi-loop QED, and recoil NLO contributions much larger than the
proton's at the same precision.  The convention exposure is
intrinsically a Layer-2 itemization split, NOT a framework
architecture issue.

\textbf{Cross-references.}  Originating documentation:
\S\ref{sec:offprecision} Sprint Calc-rZG row.  Operator-level
elaboration: \S\ref{sec:autopsy_d_hfs}.  Source memos:
\texttt{debug/calc\_track\_rZG\_global\_zemach\_memo.md},
\texttt{debug/rzg\_bug\_diagnosis\_memo.md},
\texttt{debug/d\_hfs\_autopsy\_track5\_memo.md}.

\subsubsection{$\mu$H 2$S_{1/2}$--2$P_{1/2}$ VP one-loop split (Antognini vs Krauth)}
\label{sec:conv_muH_VP}

\textbf{Observable.}  Muonic hydrogen Lamb shift VP one-loop
component, the dominant ($\sim +205$~meV out of $+202.37$~meV total)
contribution to the CREMA 2010 measurement.  Documented in the
\S\ref{sec:autopsy_muh_lamb} Roothaan autopsy.

\textbf{Source compilations.}  Antognini~2013 (the canonical
muonic-hydrogen reference for CREMA Lamb-shift theory; used as the
Layer-2 anchor in \S\ref{sec:autopsy_muh_lamb}) and Krauth~2017 (the
canonical muonic-hydrogen reference for the parallel HFS measurement
and the most thorough recent itemization).

\textbf{Convention difference.}  Antognini reports VP one-loop at
$+205.0074$~meV with ``VP-VP mixed'' itemized separately at
$+0.150$~meV.  Krauth absorbs $\sim 0.021$~meV of higher-order VP
into the one-loop line and reports $+205.0282$~meV.  Both compilations
agree on the bottom-line $\mu$H Lamb total to printed precision; the
difference is in per-component itemization.

\textbf{Magnitude.}  $\sim 100$~ppm in per-component split (totals
agree at much tighter precision).

\textbf{Sprint reference and closure.}  Sprint MH Track A
(\S\ref{sec:offprecision}, row ``$\mu$H 2$S$--2$P$ Lamb shift'') and
Sprint Calc-muH Lamb autopsy v1 (\S\ref{sec:autopsy_muh_lamb}) both
operationally surfaced the convention difference: the framework's
full-Uehling kernel matches Antognini~2013 at $-0.10$~ppm, which is
$1000\times$ tighter than the Antognini-vs-Krauth $\sim 100$~ppm
itemization gap.  This means the framework's match precision is
sensitive to the choice of reference compilation for Layer-2 inputs.
\textbf{Closure verdict:} anchored to Antognini~2013 throughout the
$\mu$H Lamb autopsy.  No reconciliation against Krauth attempted at
this time; flagged for future sprint if Krauth-anchored cross-checks
become relevant.

\textbf{Cross-references.}  Originating documentation:
\S\ref{sec:autopsy_muh_lamb} Convention-mismatch surface paragraph.
Source memo: \texttt{debug/muH\_Lamb\_autopsy\_v1\_memo.md} (LCM tag).

\subsubsection{$\mu$ Lamb shift SE-recoil aggregation (framework vs Karshenboim)}
\label{sec:conv_karshenboim_recoil}

\textbf{Observable.}  Muonium 2$S_{1/2}$-2$P_{1/2}$ Lamb shift
self-energy bracket.  Documented in the \S\ref{sec:offprecision} row
``Mu 2$S_{1/2}$-2$P_{1/2}$ Lamb $1047.63$~MHz vs Mariam~1982
experiment $1054(22)$~MHz.''

\textbf{Source compilations.}  Framework Eides~\S 3.2 SE bracket
(used in this paper's $\mu$ Lamb-shift framework calculation) and
Karshenboim~2005 (the canonical muonium-Lamb reference).

\textbf{Convention difference.}  The framework's $m_\text{red}$
rescaling at the Eides bracket level absorbs the leading
recoil-correction beyond reduced-mass automatically, giving
SE bracket $= +1074.13$~MHz.  Karshenboim itemizes ``SE one-loop''
at $+1085.84$~MHz with a separately-itemized ``Recoil NLO''
line at $-11.30$~MHz.  The two reach exact agreement at the bracket
level: $+1085.84 - 11.30 = +1074.54$~MHz vs framework $+1074.13$~MHz
(within bracket-evaluation precision).

\textbf{Magnitude.}  $\sim 11$~MHz on the $\mu$ Lamb-shift SE bracket
($\sim 1\%$ of the SE component, $\sim 0.01$ of the total Lamb shift);
the difference is exactly the recoil-NLO itemization line.

\textbf{Sprint reference and closure.}  Sprint precision-catalogue
2026-05-08 (Mu Lamb shift row, \S\ref{sec:offprecision}) operationally
surfaced this: the framework total at $1047.63$~MHz matches
Karshenboim~2005 theory at $1047.49(5)$~MHz to $+0.013\%$.  Naively
adding Karshenboim's separately-itemized ``Recoil NLO $-11.30$~MHz''
to the framework total double-counts the same recoil correction that
the framework's $m_\text{red}$ rescaling already absorbs.
\textbf{Closure verdict:} resolved by recognizing the framework's
$m_\text{red}$-bracket convention is structurally equivalent to
Karshenboim's SE one-loop $+$ Recoil NLO sum, and not adding the
recoil NLO twice.

\textbf{Cross-references.}  Originating documentation:
\S\ref{sec:offprecision} Mu 2$S_{1/2}$-2$P_{1/2}$ Lamb row.

\subsubsection{Friar moment profile convention (Eides closed-form vs operator-level)}
\label{sec:conv_friar_profile}

\textbf{Observable.}  Muonic hydrogen Friar moment, the
$O((Z\alpha)^2)$ finite-nuclear-size correction in the
\S\ref{sec:autopsy_muh_lamb} Lamb-shift autopsy (Component~3).

\textbf{Source compilations.}  Eides Eq.~2.35 closed-form derivation
(used as the Antognini~2013 reference value $-3.84$~meV at $r_p =
0.84087$~fm) and the framework's operator-level realization via the
\S\ref{sec:proj_magnetization_density} module
\texttt{geovac.magnetization\_density}.

\textbf{Convention difference.}  Eides Eq.~2.35 is derived assuming
the input radius is $\langle r^2\rangle^{1/2}$ (RMS charge radius);
the operator-level
\S\ref{sec:proj_magnetization_density} module is calibrated to the
first moment $\langle r\rangle = r_p$ (consistent with its
existing HFS Zemach use case).  For a Gaussian profile, the two
calibrations differ by exactly the factor
$\langle r^2\rangle / \langle r\rangle^2 = 4/\pi \approx 1.27$.

\textbf{Magnitude.}  $+12.7\%$ overshoot on the Friar moment
(operator-level $-4.33$~meV vs Eides closed-form $-3.68$~meV at
fixed $r_p$).  This is exactly $1.27\times$ the closed-form value to
$\sim 0.5\%$ relative precision.

\textbf{Sprint reference and closure.}  Sprint Calc-muH Lamb autopsy v1
(2026-05-08, \S\ref{sec:autopsy_muh_lamb} Operator-level Friar moment
paragraph) surfaced the discrepancy by computing the Friar moment
both ways: closed-form Eides reproduction matches Antognini~2013 at
$-4.3\%$, while operator-level realization overshoots Antognini by
$+12.7\%$, with the ratio between the two paths equal to the Gaussian
profile factor $4/\pi$ to sub-percent precision.
\textbf{Closure verdict:} flagged for a sibling
\texttt{MagnetizationDensitySpec} calibrated to RMS radius for the
Lamb-shift FNS channel, distinct from the existing first-moment
calibration for the HFS Zemach channel.  The HFS Zemach channel is
unaffected (cf. \S\ref{sec:autopsy_21cm}, where the operator-level
Zemach result $-39.495$~ppm matches Eides $-39.500$~ppm at
$0.012\%$).  Resolution is mechanical (sprint-scale).

\textbf{Cross-references.}  Originating documentation:
\S\ref{sec:autopsy_muh_lamb} Operator-level Friar moment paragraph.
Source memo: \texttt{debug/muH\_Lamb\_autopsy\_v1\_memo.md}
(profile-convention finding).

\subsubsection{D 1S HFS polarizability sub-component (Friar--Payne vs Pachucki--Yerokhin)}
\label{sec:conv_d_polarizability}

\textbf{Observable.}  Deuterium 1S hyperfine splitting
$\nu_\text{HFS}(D) = 327\,384\,352.522$ kHz, with focus on the
\emph{deuteron polarizability} sub-component of the Layer-2 budget
that closes the framework's $+285.6$~ppm chain residual
(\S\ref{sec:autopsy_d_hfs} Track~5 autopsy).  This refines the broader
\S\ref{sec:conv_rZG_DHFS} entry to the dominant polarizability piece.

\textbf{Source compilations.}  Friar--Payne~2005 (PRC \textbf{72},
014002; building on Friar~1979 PRC \textbf{20}, 325) versus
Pachucki--Yerokhin~2010 (the modern itemization framework adopted in
Kalinowski--Pachucki--Yerokhin~2018 PRA \textbf{98}, 062513 for
muonic D and applied analogously for electronic D), and the recent
Patkos--Pachucki--Yerokhin~2024/2025 chiral-EFT update
(arXiv:2508.18776).

\textbf{Convention difference.}  Friar--Payne~2005 aggregates the
``low-energy nuclear-structure $+$ polarizability $+$ dispersion $+$
Friar-moment recoil'' contributions into a single low-energy
aggregate term (e.g., $\delta E_\text{Low}(eD) = 87.3$~kHz in
Friar--Payne's reporting style).  Pachucki--Yerokhin~2010+ separates
polarizability ($\sim +240 \times 10^{-6}\,E_F \approx +78.6$~kHz)
and Zemach ($\sim -100 \times 10^{-6}\,E_F \approx -32.7$~kHz)
explicitly as distinct itemized lines, using chiral-EFT-derivable
nuclear-structure operators rather than Friar--Payne effective
Friar moments.  Both compilations close to similar bottom-line
totals at the experimental precision (Wineland--Ramsey~1972 $\sim$few
ppb), but per-component itemization differs structurally.  The
disagreement reaches operational consequences: Patkos--Pachucki--Yerokhin
2024/2025 documents a $5\sigma$ disagreement with previous
calculations on muonic D HFS, while the underlying convention
difference applies to electronic D HFS as well.

\textbf{Magnitude.}  $\sim 240$~ppm of $E_F$ for the polarizability
term itself; the convention split (bundled-vs-separated) is
$\sim$tens-to-hundreds of ppm of the observable; $\sim 80$~kHz
absolute on the 327~MHz observable.  Comparable to V.D.1's
$\sim 25$~mfm in $r_Z(p)$ precedent; both at the $10^{-4}$ to
$10^{-3}$ relative scale.

\textbf{Sprint reference and closure.}  V.D-prediction sprint, 2026-05-10
(\texttt{debug/v\_d\_prediction\_sprint\_memo.md} Candidate~A); refines
\S\ref{sec:conv_rZG_DHFS} to the polarizability-specific sub-component,
addressing the dominant Layer-2 budget item.  \textbf{Closure verdict:}
the framework's projection chain consumes both compilations' polarizability
input as a Layer-2 scalar (single-channel
\S\ref{sec:proj_magnetization_density} entry), so the convention
choice propagates directly into the $+285.6$~ppm residual budgeted in
\S\ref{sec:autopsy_d_hfs}.  Currently anchored to Pachucki--Yerokhin~2010+
itemization (matching the framework's Track~5 autopsy assumption); the
Friar--Payne~2005 convention would shift the budget bucketing without
changing bottom-line totals.

\textbf{Cross-references.}  Originating documentation: this sprint memo.
Refines \S\ref{sec:conv_rZG_DHFS} (broader D HFS itemization) to the
polarizability-specific axis.  Cross-references
\S\ref{sec:autopsy_d_hfs} ``Why D Layer-2 is deeper than H Layer-2''
paragraph (deuteron polarizability is the dominant component there).

\textbf{Honest caveat: overlap with V.D.1.}  V.D.1 documents the
broader Eides Tab.~7.3 vs Pachucki--Yerokhin~2010 D HFS itemization
mismatch.  V.D.5 refines V.D.1 to the polarizability-specific
sub-component on a distinct convention axis (Friar--Payne~2005 vs
Pachucki--Yerokhin~2010+).  The two entries differ in which
compilation pair they compare; kept separate per V.D protocol since
the convention axes differ.  Future audit may merge V.D.5 into V.D.1
as a sub-paragraph if the two axes prove operationally equivalent.

\subsubsection{He $2{}^3P$ relativistic Bethe-log reevaluation (Drake vs Pachucki--Yerokhin)}
\label{sec:conv_he_bethe_log}

\textbf{Observable.}  Helium $2{}^3P$ fine-structure intervals
$\nu_{01} = 29\,616\,946.2(1.6)$~kHz and $\nu_{12} =
2\,291\,177.3(1.6)$~kHz (current theoretical precision), used
historically as a determination of the fine-structure constant
$\alpha$ to $\sim 30$~ppb.  Documented in the
\S\ref{sec:autopsy_he_2_3p} Roothaan autopsy.

\textbf{Source compilations.}  Drake~1990 / Drake--Yan~1995
(Drake--Yan~PRA \textbf{46}, 2378, 1992; Drake--Yan~PRL \textbf{74},
4791, 1995; Drake in \emph{Long-Range Casimir Forces}, ed. F.~S.~Levin
and D.~A.~Micha, Plenum, NY, 1993): original itemization of relativistic
Bethe-log term and combining coefficients $(3/50, -2/5, 3/2, -1)$
that GeoVac uses sympy-exactly (Sprint~4 DD).
Pachucki--Yerokhin~2009 (PRA \textbf{79}, 062516, 2009 ``Reexamination
of helium fine structure''; cf. Pachucki--Yerokhin~2011 Can.~J.~Phys.
\textbf{89}, 95, arXiv:1011.1370): reevaluated the relativistic
Bethe-log term and resolved a 3$\sigma$ disagreement on $\nu_{01}$
that Drake's earlier itemization had with experiment.

\textbf{Convention difference.}  Pachucki and Yerokhin
(PRA \textbf{79}, 062516, 2009) explicitly state they ``eliminated this
[3$\sigma$] disagreement [vs Drake on $\nu_{01}$] by reevaluating the
relativistic Bethe logarithm term.''  Pachucki and Sapirstein
``pointed out several computational mistakes and inconsistencies in
earlier calculations, and eventually the results turned out to be
consistent after Drake revised the sign'' (per literature retrieval).
This is a textbook example of inter-compilation itemization correction
at the $\alpha^3 (Z\alpha)^2$ multi-loop level, sitting inside the
LS-8a renormalization-wall budget that the framework does not capture.

The relativistic Bethe-log term is part of the $\alpha^3(Z\alpha)^2$
multi-loop sector that the GeoVac framework attributes to its LS-8a
structural wall.  The framework's residual on $P_0 - P_1$ is
$-0.014\%$, on $P_0 - P_2$ is $-0.20\%$, on $P_1 - P_2$ is $-2.62\%$
from partial-cancellation amplification (see \S\ref{sec:autopsy_he_2_3p}).
The Drake-vs-Pachucki--Yerokhin convention difference lives inside
this same budget.

\textbf{Magnitude.}  $\sim$kHz on the 29~GHz $\nu_{01}$ interval
(Pachucki--Yerokhin's revised value $29\,616\,946.2$~kHz vs Drake's
older value disagreeing at the 3$\sigma \sim$few~kHz level).  In
ppm: $\sim 0.1$~ppm on $\nu_{01}$.  In $\alpha$-determination units:
the helium $2{}^3P$ fine structure determines $\alpha$ to $\sim 27$-$31$~ppb;
the Bethe-log reevaluation was the load-bearing fix for the discrepancy
that earlier $\alpha$-determination work had against alternative
determinations.

\textbf{Sprint reference and closure.}  V.D-prediction sprint, 2026-05-10
(\texttt{debug/v\_d\_prediction\_sprint\_memo.md} Candidate~C); refines
the \S\ref{sec:autopsy_he_2_3p} Three-class diagnosis paragraph's
``Class~A: NEGATIVE'' statement, which was correct at LO Breit-Pauli
only, not at the $\alpha^3(Z\alpha)^2$ multi-loop level where the
relativistic Bethe-log reevaluation lives.  \textbf{Closure verdict:}
this entry is the strongest of the V.D-prediction sprint candidates
because the literature explicitly documents both the disagreement
(Drake--Pachucki--Yerokhin original 3$\sigma$ split) and its resolution
(Pachucki--Yerokhin~2009 Bethe-log reevaluation, with Drake's subsequent
agreement after sign revision per Pachucki--Sapirstein).  It is not an
inferred convention split but an explicitly noted correction in the
published record.

\textbf{Cross-references.}  Originating documentation: this sprint memo.
Refines \S\ref{sec:autopsy_he_2_3p} ``Three-class diagnosis'' paragraph
(``Class~A: NEGATIVE'' statement now read as ``Class~A: NEGATIVE at LO
Breit-Pauli, POSITIVE at $\alpha^3(Z\alpha)^2$ multi-loop level'').
Cross-references \S\ref{sec:offprecision} He $2{}^3P$ catalogue rows.

\subsubsection{$\mu$d 2$S$--2$P$ Lamb deuteron polarizability
(Carlson--Gorchtein--Vanderhaeghen vs Carlson--Vanderhaeghen)}
\label{sec:conv_d_polarizability_mud_lamb}

\textbf{Observable.}  Muonic deuterium 2$S_{1/2}$--2$P_{1/2}$ Lamb
shift $\Delta E = 202.8785(34)$ meV (Pohl~et~al., CREMA 2016).
The dominant Layer-2 contribution to the theoretical prediction
is the deuteron nuclear polarizability, a QCD-internal evaluation
of virtual photonuclear excitations between the bound $n+p$ system.

\textbf{Source compilations.}  Carlson--Gorchtein--Vanderhaeghen 2014
(\emph{Phys. Rev. A} \textbf{89}, 022504): $+1.690(20)$ meV.
Carlson--Vanderhaeghen 2011 (earlier evaluation, distinct
sub-channel decomposition): $+1.94$ meV.  Two distinct QCD-internal
evaluations, both anchored to deuteron photoabsorption data, with
the 2014 update including refined sub-leading contributions
(elastic-inelastic interference, subtraction-function refinement,
finite-nuclear-size corrections to the polarizability kernel).

\textbf{Convention difference.}  The 2014 evaluation explicitly
splits the nuclear-structure correction into three sub-contributions
(elastic intermediate states, inelastic photoabsorption, subtraction
function); the 2011 evaluation used a broader-aggregate treatment
that produced a $\sim 0.25$~meV larger value.  Both evaluations
converge to the same physical mechanism (virtual photoabsorption
between $n+p$ via the muonic Lamb shift kinetic energy); they
differ in how the QCD content is itemized.

\textbf{Magnitude.}  $\sim 13\%$ spread on the \emph{dominant}
Layer-2 contribution at $\mu$d Lamb shift ($+1.69$ meV vs
$+1.94$ meV; $\sim 0.25$ meV absolute).  At the framework's
$-0.12\%$ Lamb-shift residual, the convention choice between the
two QCD evaluations sets whether the framework reads as
``sub-percent closure'' (with CGV~2014 anchoring, $-0.12\%$) or
``over-shoots experiment by $+0.12\%$'' (with CV~2011 anchoring,
$+0.12\%$).  The framework's own residual is at the same precision
as the spread between the two QCD evaluations.

\textbf{Three-class tag.}  Class~(i) inter-compilation itemization,
distinct from \S\ref{sec:conv_d_polarizability} (deuteron HFS
polarizability sub-component, Friar--Payne vs Pachucki--Yerokhin):
this entry is between two QCD nuclear-structure evaluations of
the deuteron polarizability rather than between atomic-physics
itemization frameworks.  Cross-references the
\S\ref{sec:autopsy_mud_lamb} ``Convention-mismatch surface''
paragraph and W3 inner-factor wall (CLAUDE.md~\S1.7).

\textbf{Verdict (multi-track Roothaan sprint Track~1, 2026-05-18):
CONFIRMED.}  The two QCD evaluations explicitly disagree at $\sim
13\%$ on the dominant Layer-2 line of a precision-frontier
observable.  This is the first explicit Class~(i) convention
exposure originating in a QCD-internal sub-component evaluation
rather than an atomic-physics itemization compilation.  Refines
the broader W3 inner-factor calibration question to a specific
operational instance: until the deuteron polarizability convention
is reconciled at the framework-residual scale, multi-observable
extractions involving the deuteron ($\mu$d Lamb, D 1S HFS,
isotope-shift extractions of $r_d$) cannot reach beyond the
$\sim 13\%$ Layer-2-determined precision.

\textbf{Cross-references.}  Originating documentation:
\texttt{debug/muD\_Lamb\_autopsy\_track1\_memo.md}.  Refines
\S\ref{sec:autopsy_mud_lamb} ``Convention-mismatch surface''
paragraph.  Sibling to \S\ref{sec:conv_d_polarizability}
(D HFS polarizability sub-component; different observable, same
W3 inner-factor wall).  Cross-references CLAUDE.md~\S1.7
multi-focal-composition wall taxonomy W3.

\subsubsection{Alkali $n$s valence HFS multi-electron contact-density
cliff (class~(iii): framework architecture)}
\label{sec:conv_alkali_hfs}

\textbf{Observable.}  Alkali-atom ground-state hyperfine splittings on
single-zeta hydrogenic basis; concretely Li-7 $2{}^2S_{1/2}$
HFS = $803.504086(34)$~MHz (Beckmann--Boklen--Elke 1974), with
generalization to Na, K, Rb, Cs (all alkali $n$s valence HFS share
the same mechanism).  Originating sprint: multi-track Roothaan
sprint Track~5, 2026-05-18 (\S\ref{sec:autopsy_li7_hfs}).

\textbf{Source of mismatch.}  Hydrogenic-$Z_\text{eff}$ approximation
(Clementi--Raimondi~1963 for alkalis: $Z_\text{eff} = 1.279$ for
Li 2s) vs SCF + multi-electron CI (Clementi--Roetti~1974;
Yan \& Drake~2003) at the level of the contact density
$|\psi_{ns}(0)|^2$ \emph{specifically}.  Spatial extent
($\langle r \rangle$, polarizability) of the orbital is well-approximated
by $Z_\text{eff}$; contact density is not.

\textbf{Convention difference.}  The hydrogenic-$Z_\text{eff}$ contact
density $|\psi_{2s}(0)|^2 = Z_\text{eff}^3/(8\pi)$ at
$Z_\text{eff} = 1.279$ gives $|\psi_{2s}(0)|^2 \approx 0.083$
in atomic units.  Clementi--Roetti~1974 SCF Li 2s gives
$|\psi_{2s}(0)|^2 \approx 0.665$, a factor of $\sim 8\times$
enhancement.  The mechanism is the
\S\ref{sec:proj_phillips_kleinman} core-valence orthogonality
constraint $\langle 2s | 1s_\text{core}\rangle = 0$ \emph{exactly},
which pulls the 2s radial node inward and dramatically enhances
amplitude at the origin.  This is the well-known Fermi contact
enhancement in alkali atoms.

\textbf{Magnitude.}  $\sim 10\times$ on $A_\text{hf}$ for Li-7
($A_\text{hf}^\text{framework, hydro-} = 41.5$ MHz vs
$A_\text{hf}^\text{exp} = 401.75$ MHz).  Scales similarly for
Na, K, Rb, Cs $n$s valence HFS:
the same factor-$\sim 10\times$ undershoot is the structural prediction
of the framework's current single-zeta hydrogenic architecture for
\emph{all} alkali valence HFS observables.

\textbf{Class.}  This is a class~(iii) exposure --- \emph{framework
architecture limitation}, distinct from class~(i) inter-compilation
itemization (entries V.D.1, V.D.2, V.D.3, V.D.5, V.D.6, V.D.7)
and class~(ii) closed-form vs operator-level convention (entry V.D.4
Friar profile RMS-vs-first-moment).  Class~(iii) exposes a structural
property of the GeoVac graph-native basis choice (single-zeta
hydrogenic with empirical $Z_\text{eff}$) that the projection-chain
dictionary alone cannot fix; closing it requires extending the
graph-native basis to multi-zeta with explicit PK orthogonality at
the orbital level.

\textbf{Two cliffs, one mechanism.}  This exposure shares a structural
mechanism with the \S\ref{sec:autopsy_cs_hfs} Track-2 Z\textgreater 20
cliff (single-zeta basis with no radial nodes cannot represent the
PK-orthogonalized SCF radial wavefunction).  The Cs cliff manifests
as a $-47\%$ undershoot on $A_\text{hf}$ at $Z=55$; the Li cliff
manifests as a $-89.7\%$ undershoot at $Z=3$.  The two cliffs are
not independent walls but two manifestations of the same structural
limitation, scaling with $Z$ inversely to expectation: the Z=3 case
is \emph{more} affected than the Z=55 case because the 2s valence is
the \emph{first} shell beyond the closed $1s^2$ core where PK
orthogonality matters most acutely.

\textbf{Named structural closure.}  Hylleraas--Eckart double-zeta
extension for $n=2$ multi-electron systems
(\texttt{geovac/hylleraas\_r12.py}, currently He singlet at $\omega=4$
giving $+0.0006\%$ on $E_\text{tot}$ --- \S\ref{sec:autopsy_he_oscillator}).
The Hylleraas basis explicitly satisfies PK orthogonality at the
amplitude level (no $Z_\text{eff}$ fitting required), and closure of
the multi-electron contact-density cliff and the Li atomic accuracy
floor ($1.10\%$, Paper FCI-A) is expected \emph{simultaneously}.
Sibling closure to the BBB93/KTT multi-zeta tabulation closure for
the \S\ref{sec:autopsy_cs_hfs} Z\textgreater 20 cliff (also named in
\S\ref{sec:autopsy_cs_hfs}).

\textbf{Cross-references.}  Originating documentation:
\texttt{debug/Li7\_HFS\_autopsy\_track5\_memo.md}.  Cross-references
\S\ref{sec:autopsy_li7_hfs} (Li-7 HFS autopsy) and
\S\ref{sec:autopsy_cs_hfs} (Cs HFS scoping + Z\textgreater 20 cliff
diagnostic).  Generalizes to Na ($Z=11$), K ($Z=19$), Rb ($Z=37$),
Cs ($Z=55$) valence HFS as a precision-frontier challenge for the
framework's atomic-spectroscopy reach.

\subsubsection{Multi-focal-composition wall: bimodule cross-shifts
vs.\ module endomorphisms (Sub-sprint M-Z, 2026-05-23)}
\label{sec:conv_multifocal_bimodule}

\textbf{Observable.}  The multi-focal-composition wall named in
CLAUDE.md \S~1.7: five independent observables (Sprint H1 Yukawa
selection, LS-8a multi-loop QED renormalization, HF-3 recoil, HF-4
Zemach, HF-5 multi-loop $a_e$) hitting one structural wall on the
framework's ability to compose multiple Fock-style projections.

\textbf{Source compilations / framing perspectives.}  The
single-observation reading (each observable a separate failure mode)
versus the Hilbert $C^*$-bimodule reading of the modular propinquity
literature (Latr\'emoli\`ere~2019, 2021;
arXiv:1608.04881, arXiv:1811.04534).

\textbf{Convention difference.}  Single-observation reading treats
each instance as an independent kernel-approximation gap; the bimodule
reading splits the wall cleanly into two structural pieces:
\emph{bimodule cross-shifts} (the framework's production code does
compose them: HF-3 recoil via \texttt{geovac/cross\_register\_vne.py},
HF-4 Zemach via \texttt{geovac/magnetization\_density.py}, HF-5
Parker--Toms first-order curvature on Dirac-$S^3$, FNS via Foldy--Friar,
cross-register $V_{eN}$) versus \emph{module endomorphisms} (the
framework does NOT compose them: LS-8a renormalization counterterms
$Z_2-1$ and $\delta m$, Sprint H1 Yukawa $Y$, inner-factor calibration
data).  The cross-shift piece is empirically handled (e.g., HF-3 at
sub-ppm match with Bethe--Salpeter recoil); the endomorphism piece is
the LS-8a wall precisely.  This is an architecture-class~(iv) exposure
(framework-level structural taxonomy refinement); not a
literature-convention mismatch in the strict sense.

\textbf{Magnitude.}  Qualitative structural finding at this level.
Quantitative tests live in the constituent sprints: LS-8a yields the
two-loop QED renormalization gap (Paper~35 \S~VII.3 Refined Prediction~1);
the W1c Layer~3 finding (Paper~19 \S~\ref{sec:w1c_residual}, 2026-05-23
$\alpha$-PES test) is the most recent concrete demonstration that the
module-endomorphism piece can be invisible at finite basis size when
the FCI variational state does not occupy the affected subspace.

\textbf{Sprint reference and closure.}
Sub-sprint M-Z (\texttt{debug/sprint\_modular\_propinquity\_mZ\_bethe\_log\_memo.md},
2026-05-23) surfaced the bimodule reading; the modular propinquity
synthesis memo
(\texttt{debug/sprint\_modular\_propinquity\_synthesis\_memo.md})
consolidates it as the load-bearing structural finding of the
modular sprint's chemistry-side thread.  The Track $\beta$ synthesis
(\texttt{debug/sprint\_modular\_alpha\_arc\_synthesis\_memo.md})
extends the consolidation with the W1c Layer~3 finding from the
follow-on $\alpha$ arc.

\subsubsection{Three-bucket M-Z partition hypothesis (basis-closable
cross-shift as intermediate class) --- FALSIFIED via NaH max\_n=3 test}
\label{sec:conv_three_bucket_falsified}

\textbf{Observable.} The W1c $\times$ M-Z partition bridge sprint
(\texttt{debug/sprint\_w1c\_mz\_partition\_analysis\_memo.md}, 2026-05-23)
proposed a three-bucket refinement of the M-Z partition
(\S\ref{sec:conv_multifocal_bimodule}): cross-shift (handled) /
\emph{basis-closable cross-shift} (handled at sufficient basis size) /
endomorphism (external input). The middle bucket was introduced to
classify W1c sub-layer~3 (orbital-pair flexibility at NaH max\_n=2),
on the reading that the bonding combination $[\text{H 1s} \pm
\text{Na 3s}]_\pm$ is structurally a bimodule cross-shift (left+right
linear combination in the bimodule basis) that the FCI cannot realize
at max\_n=2 due to dimensional poverty alone.

\textbf{Source compilations / framing perspectives.} Bridge sprint's
three-bucket hypothesis (the refinement to test) vs.\ the
original two-bucket M-Z partition (the null hypothesis). Three
falsifiable predictions for NaH max\_n=3 were prespecified and frozen
into \texttt{debug/data/sprint\_w1c\_mz\_partition\_predictions.json}
before the test ran (curve-fit-audit discipline at inter-sprint scope):
\textbf{P1} (internal PES minimum at $R_\text{eq} \in [3.0, 4.5]$~bohr),
\textbf{P2} (binding $D_e \in [0.0375, 0.150]$~Ha within $2\times$
experimental $D_e \approx 0.075$~Ha), \textbf{P3} (multi-zeta differential
$\in [0.02, 0.30]$~Ha at $R_\text{eq}$).

\textbf{Convention difference and verdict.} The F1 max\_n=3 sprint
(\texttt{debug/sprint\_f1\_maxn3\_predictions\_test\_memo.md}, 2026-05-23)
returned 1 of 3 predictions passing: P1 FAIL ($R_\text{min} = 2.0$~bohr,
the smallest tested R, on the exact falsification criterion); P2 FAIL
(spurious-binding descent depth $D_e^\text{PES} = +0.7097$~Ha, $\sim 10\times$
experimental); P3 PASS (multi-zeta differential $0.2066$~Ha at $R = 2.0$,
consistent with max\_n=2's $0.21$~Ha at the same $R$). \textbf{CLEAN
NEGATIVE per the prespecified gate.} The proposed three-bucket
refinement is falsified; the two-bucket M-Z partition stands.

\textbf{Magnitude (structural finding).} The substantive new content
the prediction test did not anticipate is the structural success /
energetic failure split. At max\_n=3 under combined W1c+mz, the
dominant natural orbital IS a true bonding combination:
$\phi_\text{dom} = -0.698 \cdot \phi_\text{Na, 3s} - 0.687 \cdot
\phi_\text{H, 1s}$ with $50/50$ amplitude$^2$ split (the 2nd NO is the
antibonding combination). At max\_n=2 the dominant NO was 100\%
H-localized at every architecture. \emph{The basis enlargement DID
provide the orbital-mixing flexibility predicted by the bridge sprint.}
But the constructed bonding orbital has higher energy than the separated
configuration at every $R$; natural occupations remain $[1.0, 1.0]$
(open-shell singlet), NOT $[2.0, 0.0]$ (closed-shell bond). The
framework can BUILD the right orbital but cannot energetically PREFER
it. Sub-layer~3 reclassifies from the bridge sprint's ``basis-closable
cross-shift'' to ``module endomorphism (basis-irreducible at tested
scales)'' and joins sub-layer~2 (multi-zeta basis-shape) in the
endomorphism bucket.

\textbf{Class.} This is an architecture class~(iv) exposure
(framework-level structural taxonomy refinement), sibling to
\S\ref{sec:conv_multifocal_bimodule} (the parent M-Z partition entry).
The exposure is FALSIFICATION of a proposed sub-classification, not
discovery of a new convention mismatch. It is retained in the running
catalogue because the methodological lesson --- the structural
success / energetic failure split is a robust new diagnostic, with
the two-bucket partition surviving while the three-bucket refinement
evaporates --- sharpens where chemistry endomorphisms localize in the
M-Z framework.

\textbf{Same-day refinement: F2 + F3 sharpen sub-layer~3 further.}
F1 max\_n=3's verdict (sub-layer~3 as ``module endomorphism, basis-irreducible
at tested scales'') is structurally correct but was further sharpened by
two same-day same-arc follow-on sprints: F2 (cross-$V_{ne}$ kernel-shape
diagnostic, KERNEL-NOT-IT verdict) ruled out kernel-shape substitution as
a closure mechanism (multipole kernel bit-faithful to converged 3D
quadrature within $\sim 10^{-5}$~Ha, six orders below wall depth) and
discovered the framework's h1 in \texttt{composed\_qubit} is strictly
block-diagonal --- the cross-block off-diagonal h1 slot is
\emph{architecturally absent}, not just kernel-error-flawed. F3
(cross-block h1 architectural extension) CLOSED W1d (the architectural
absence named by F2) at the FCI level: naturals transform $[1.0, 1.0] \to
[1.9991, 0.0007]$ at NaH max\_n=2 (four-orders-of-magnitude advance in
dominant occupation signature). But F3 surfaced a new sub-layer W1e
(inner-region overattraction): the bonding pair lacks Pauli repulsion
against the [Ne] frozen core; PES still monotonically descending with
$D_e^\text{F3} = 4.37$~Ha = $58\times$ experimental. \textbf{Net refined
reading:} chemistry sub-walls split into THREE structurally distinct
sub-categories within the two-bucket M-Z partition (cross-shift /
endomorphism / \textbf{architectural-absence}), with architectural-absence
sub-walls conditionally closable via architectural extension (W1d via
\texttt{geovac/cross\_block\_h1.py}). Chemistry endomorphisms split
further into basis-shape data (sub-layer~2, mitigable via FrozenCore +
atomic-physics fits), kernel-energetics (F1's reading of sub-layer~3, now
sharpened to architectural-absence + Pauli-repulsion candidates), and
Pauli-class constraints (W1e, F4 target via bonding-orbital PK extension).
See \S\ref{sec:conv_w1c_f2_f3} for the F2+F3 follow-on entry.

\textbf{Sprint reference and closure.} W1c $\times$ M-Z bridge sprint
(\texttt{debug/sprint\_w1c\_mz\_partition\_analysis\_memo.md}); F1
max\_n=3 prediction test
(\texttt{debug/sprint\_f1\_maxn3\_predictions\_test\_memo.md}); Sprint
$\beta$-neg synthesis (F1-maturity, superseded for framing by the
full-arc memo)
(\texttt{debug/sprint\_beta\_neg\_synthesis\_memo.md}); F2
architectural-absence diagnostic
(\texttt{debug/sprint\_f2\_cross\_vne\_kernel\_memo.md}); F3 cross-block
h1 closure
(\texttt{debug/sprint\_f3\_cross\_block\_h1\_memo.md}); comprehensive
full-arc synthesis
(\texttt{debug/sprint\_w1c\_full\_arc\_synthesis\_memo.md}). Cross-references
Paper~19 \S\ref{sec:w1c_residual} (W1c-residual wall + F1-F2-F3 arc
extension), CLAUDE.md~\S1.7 (multi-focal-composition wall taxonomy),
CLAUDE.md~\S3 (dead-ends rows for the three-bucket hypothesis and
kernel-shape substitution as W1c-residual closure mechanism), and
\S\ref{sec:conv_multifocal_bimodule} (parent M-Z partition entry whose
three-bucket refinement is here falsified).

\subsubsection{Architectural-absence sub-category for chemistry sub-walls
(F2 diagnosis + F3 closure of W1d; W1e newly named) --- W1d CLOSED at FCI
level, W1e OPEN as F4 target}
\label{sec:conv_w1c_f2_f3}

\textbf{Observable.} Same-day same-arc follow-on to the three-bucket
falsification entry above (\S\ref{sec:conv_three_bucket_falsified}): two
sprints (F2 + F3) tightened the W1c-residual sub-layer hierarchy at NaH
max\_n=2 beyond F1 max\_n=3's verdict. F2 (cross-$V_{ne}$ kernel-shape
substitution diagnostic) executed a Step~1 algebraic gate before any
production extension: five matrix elements computed via converged 3D
quadrature vs the framework's default multipole expansion at
$L_\text{max}=4$. Largest converged differential
$2.0\times 10^{-5}$~Ha on the H 1s diagonal cross-$V_{ne}$ contribution
($2.7\times 10^{-4} = 0.027\%$ of bond-energy scale 0.075~Ha). F3
(cross-block h1 architectural extension) implemented the missing matrix
slot: new module \texttt{geovac/cross\_block\_h1.py} ($\sim 437$~lines, 18
new tests) computes off-diagonal h1 elements
$\langle\psi_a^A \mid T + \sum_C (-Z_C/|r-R_C|) \mid \psi_b^B\rangle$
between orbitals on different centers $A \neq B$ via 2D axial
Gauss-Legendre quadrature. \texttt{build\_composed\_hamiltonian} and
\texttt{build\_balanced\_hamiltonian} gain \texttt{cross\_block\_h1}
kwargs; bit-exact backward compat at default \texttt{cross\_block\_h1=False}.

\textbf{Source compilations / framing perspectives.} F2's diagnostic-first
gate logic ruled out the obvious kernel-error hypothesis at six orders of
magnitude below wall depth, ruling out kernel-shape substitution as a
W1c-residual closure mechanism. F3's algebraic Step~1 diagnostic predicted
strong bonding splitting (h1 lower-eigenvalue 3.46~Ha for hydrogenic and
4.20~Ha for multi-zeta, both $\gtrsim 70\times$ the 50~mHa decision-gate
threshold). F3 Steps~3+4 ran the production-wired FCI + mini-PES tests.

\textbf{Convention difference and verdict.} F2 verdict:
\textbf{KERNEL-NOT-IT}. The multipole expansion is bit-faithful where it
lives; the W1c-residual wall is NOT a kernel-truncation artifact within
the existing block-diagonal architecture. F2's substantive new content
(\textbf{architectural-absence finding, W1d named}): the framework's h1 in
\texttt{composed\_qubit} is strictly block-diagonal --- cross-block h1
matrix elements for orbitals on different centers are \emph{architecturally
absent}, not just kernel-error-flawed. F3 verdict: \textbf{PARTIAL-CLOSURE}
--- bonding-orbital construction CLOSED at the FCI level; inner-region
overattraction OPEN (\textbf{new sub-layer W1e}). Adding cross-block h1
transforms FCI naturals $[1.0000, 1.0000] \to [\mathbf{1.9991},
\mathbf{0.0007}]$ at NaH max\_n=2 R=3.5 --- four-orders-of-magnitude advance
in the dominant natural-orbital occupation signature, with the bonding
combination as a doubly-occupied ground state. Concurrently the h1 lowest
eigenvector character flips from H-localized at $-0.802$~Ha to bonding
combination Na/H = 0.51/0.49 at $-3.161$~Ha ($\sim 2.4$~Ha lower).

\textbf{Magnitude (substantive structural finding).} F3 mini-PES across
$R \in [2, 10]$~bohr is monotonically descending with $R_\text{min}=2.0$
(smallest tested) and well depth $D_e^\text{F3} = +4.37$~Ha vs experimental
NaH $D_e \approx 0.075$~Ha ($58\times$ too deep). Dominant NO is bonding
(50/50 Na/H) at every R --- the cross-block h1 construction is robust across
the full PES, not a localized feature. The new well-depth is WORSE than
W1c+mz alone (0.69~Ha): cross-block h1 enables bonding without supplying
opposing Pauli repulsion against the [Ne] frozen core. This is the W1e
mechanism: 2-electron FCI on NaH max\_n=2 has NO explicit frozen-core
electrons (the [Ne] core is a screening potential via W1c, not occupied
orbitals); no Pauli-exclusion mechanism prevents the bonding pair from
collapsing toward small R. Structural analog of the missing-Pauli-repulsion
mechanism that standard Phillips-Kleinman pseudopotentials supply, but the
framework's existing PK cross-center operates on partner-side valence
orbital diagonal --- not on the bonding-combination orbital that
cross-block h1 newly constructs.

\textbf{Class.} This is an architecture class~(iv) exposure
(framework-level structural taxonomy refinement) sibling to
\S\ref{sec:conv_three_bucket_falsified} (the parent F1-arc entry) and
\S\ref{sec:conv_multifocal_bimodule} (the grandparent M-Z partition
entry). The exposure CONFIRMS the F2 architectural-absence diagnosis at
the FCI level (F3) and surfaces W1e (inner-region overattraction) as a
new structural sub-layer to be addressed by F4 (bonding-orbital PK
extension, $\sim 1$~week). Refines the methodological lesson of
\S\ref{sec:conv_three_bucket_falsified}: chemistry sub-walls have a
richer sub-structure than the two-bucket M-Z partition alone captures
--- architectural-absence sub-walls are CONDITIONALLY CLOSABLE via
architectural extension (W1d demonstrated), and the architectural-extension
ladder (W1c $\to$ W1c-multi-zeta $\to$ W1d $\to$ W1e?) is the chemistry-side
analog of refinements that the QED side does not natively support
(QED endomorphisms strictly external-input via UV-completion physics).

\textbf{Sprint reference and closure.} F2 cross-$V_{ne}$ kernel-shape
diagnostic (\texttt{debug/sprint\_f2\_cross\_vne\_kernel\_memo.md}); F3
cross-block h1 architectural extension
(\texttt{debug/sprint\_f3\_cross\_block\_h1\_memo.md}); comprehensive
full-arc synthesis
(\texttt{debug/sprint\_w1c\_full\_arc\_synthesis\_memo.md}). Cross-references
Paper~19 \S\ref{sec:w1c_residual} (W1c-residual wall + F1-F2-F3 arc
extension with revised five-sub-layer hierarchy), Paper~17 \S6.10
(composed-architecture cross-reference, updated to F3 maturity),
Paper~32 \S VIII Sprint M-Z addendum (updated with W1d/W1e
architectural-extension reading), CLAUDE.md~\S2 (F1-F2-F3 comprehensive
sprint bullet), CLAUDE.md~\S3 (dead-ends rows for three-bucket hypothesis
and kernel-shape substitution), and \S\ref{sec:conv_three_bucket_falsified}
(parent entry).

\subsubsection{W1e three-sub-sub-mechanism refinement after F4--F6
($\beta$-update, 2026-05-23)}
\label{sec:conv_w1e_f4_f6_refinement}

\textbf{Observable.} Same-day same-arc follow-on to the F3 W1d-CLOSED
+ W1e-newly-named verdict above (\S\ref{sec:conv_w1c_f2_f3}): three
sprints (F4, F5, F6) tested all three candidate W1e closure mechanisms
F3 §6 had named, and refined the W1e mechanism characterization into
a structurally distinct multi-sub-mechanism reading. \textbf{F4}
(bonding-orbital Phillips--Kleinman extension on the F3-constructed
bonding orbital): RULED OUT at Step~1 algebraic diagnostic +
FCI rank-1 PK sensitivity injection. \textbf{F5} (explicit-core
Hartree treatment of the [Ne] core electrons via $\sum_c (J - K)$):
RULED OUT at Step~1 algebraic diagnostic. \textbf{F6} (valence
basis enlargement to NaH max\_n=4 with full F3 stack):
PARTIAL\_IMPROVEMENT only.

\textbf{Source compilations / framing perspectives.} The three
sprints each ran a Step~1 diagnostic before any production
implementation, per the diagnostic-before-engineering discipline
(CLAUDE.md §1.7, \texttt{feedback\_diagnostic\_before\_engineering.md}).
F4 Step~1 computed bonding-core overlap from the full F3 h1
diagonalization and ran an FCI sensitivity test injecting a rank-1
PK shift at varying magnitudes. F5 Step~1 computed the cross-block
2-body Coulomb correction $\sum_c (J - K)_{\text{bonding},c}$
between the F3 bonding orbital and the five Na [Ne] core orbitals
via Hartree mean-field plus Slater--Condon exchange. F6 executed
full production wiring at max\_n=4 (Q=120, M=60, FCI dim 3600;
${\sim}5$~min/single-point) across a 5-point PES.

\textbf{Convention difference and verdict.}

\emph{F4 verdict: RULED OUT.} Bonding orbital has substantial Na 2s
core overlap ($|S| = 17.5\%$, well above the $5\%$ strong-gate
threshold) but the Phillips--Kleinman barrier at strict absolute
reference $E_v = 0$ is only $+0.194$~Ha (22$\times$ smaller than
4.37~Ha wall depth). FCI sensitivity test injecting rank-1 PK shift
$\Delta H = c \cdot |\text{bond}\rangle\langle\text{bond}|$ shows
$0.05\%$ closure at the Step~1-predicted $c = +0.19$~Ha, and
saturates at $43\%$ closure even at $c = 2.0$~Ha (10$\times$
predicted). \textbf{W1e is NOT a single-particle Pauli-orthogonality
wall.}

\emph{F5 verdict: RULED OUT.} Hartree mean-field $J = +1.13$~Ha;
Slater--Condon exchange $K = +0.008$~Ha; predicted correction
$J - K = +1.12$~Ha. Right SIGN (repulsive, would lift inner-region
collapse) but only $25.7\%$ of wall depth. Linear extrapolation
predicts residual $D_e^\text{F5} \approx +3.25$~Ha = 43$\times$
experimental NaH $D_e \approx 0.075$~Ha. \textbf{W1e is NOT
addressable by Hartree-level core-bonding $J - K$ interaction
alone.}

\emph{F6 verdict: PARTIAL\_IMPROVEMENT.} At max\_n=4 with full
F3 stack, PES well depth $D_e^\text{PES, F6} = +3.929$~Ha at
$R_\text{min}=2.0$~bohr (smallest tested) vs F3 max\_n=2 baseline
$D_e^\text{F3} = +4.374$~Ha --- \textbf{$10.2\%$ closure}. The
2-point gate at R=3.5/R=10 (the standard F1/F2/F3 quick-gate) gives
$26.1\%$ apparent closure but \textbf{this is ARTIFACTUAL}: the
inner-region collapse continues from R=3.5 down to R=2.0 with another
$0.696$~Ha descent the 2-point gate misses. Bonding signature
preserved at every $R$ (dominant NO $\in [1.9971, 1.9992]$;
Na/H amplitude $\approx 50/50$). \textbf{W1e is NOT addressable by
basis enlargement at max\_n=4 alone.}

\textbf{Magnitude (substantive structural finding).} W1e decomposes
into three structurally INDEPENDENT sub-sub-mechanisms with
closure ceilings:
\begin{itemize}\setlength\itemsep{0pt}
\item \textbf{W1e-basis-truncation:} ${\sim}10.2\%$ closure ceiling
  at max\_n=4 PES level (F6 measured).
\item \textbf{W1e-Hartree-pressure:} ${\sim}25.7\%$ closure ceiling
  at max\_n=2 2-point level (F5 predicted; PES-level may be lower
  per F6 2-pt-vs-PES caveat).
\item \textbf{W1e-deep-correlation-residual:} ${\sim}65$--$90\%$ of
  the wall remains beyond mean-field + single-determinant +
  basis-enlargement architecture. Remaining closure candidates
  (Schmidt orthogonalization, fully-correlated [Ne] cores) are
  beyond-sprint-scale architectural investments.
\end{itemize}

\textbf{Methodological side finding: 2-point-vs-PES discrepancy.}
F6 reveals that the 2-point gate at $R = 3.5$~bohr vs $R = 10$~bohr
(the standard F1/F2/F3 quick-gate) \textbf{overestimates PES-derived
closure by 2.5$\times$ on monotonically-descending PES in this
regime}. The inner-region collapse extends below $R = 3.5$ down to
$R_\text{min} = 2.0$ with significant additional descent
(${\sim}0.7$~Ha at max\_n=4). Future quick-gate tests in this regime
should use a 3-point or 5-point mini-PES with $R_\text{min}$
unbounded below 3.5~bohr, not a 2-point differential at fixed
reference geometry. The F3-reported $D_e^\text{F3} = +4.37$~Ha is
itself a 2-pt value; the corresponding F3 PES well depth at
$R_\text{min} = 2.0$ is likely larger (untested but expected by
analogy to F6).

\textbf{Class.} This is an architecture class~(iv) exposure
(framework-level structural taxonomy refinement) sibling to
\S\ref{sec:conv_w1c_f2_f3} (the parent F3 entry) and
\S\ref{sec:conv_three_bucket_falsified} (the grandparent F1-arc
entry). The exposure refines the W1e mechanism characterization from
``missing Pauli repulsion'' (F3 working hypothesis) to ``deep
multi-determinant correlation requiring architectural extension
beyond mean-field + single-determinant + basis-enlargement.'' The
three candidate closure mechanisms F3 §6 ranked by priority are now
all empirically tested: priority~1 (F5 explicit-core Hartree) closes
at most $25.7\%$; priority~2 (F4 bonding-orbital PK) closes at most
$43\%$ in the FCI sensitivity test; priority~3 (F6 basis enlargement
to max\_n=4) closes only $10.2\%$ at PES level. The remaining
candidates not yet tested are beyond-sprint-scale architectural lifts
(Schmidt orthogonalization at basis-construction level; fully
correlated [Ne] cores requiring occupancy-restriction FCI
infrastructure not present in framework).

\textbf{Sprint reference and closure.} F4 bonding-orbital PK
STOP-at-Step-1
(\texttt{debug/sprint\_f4\_bonding\_pk\_memo.md}); F5 explicit-core
Hartree STOP-at-Step-1
(\texttt{debug/sprint\_f5\_explicit\_core\_memo.md}); F6 max\_n=4
NaH PARTIAL\_IMPROVEMENT
(\texttt{debug/sprint\_f6\_maxn4\_nah\_memo.md}); Sprint
$\beta$-update F4-F6-APPENDED-TO-F3-BASELINE synthesis
(\texttt{debug/sprint\_beta\_update\_f4f6\_memo.md}).
Cross-references Paper~19 \S\ref{sec:w1e_refinement_f4_f6}
(refined six-sub-layer hierarchy with the three W1e
sub-sub-mechanisms), Paper~17 \S6.10 (composed-architecture
cross-reference, updated to $\beta$-update F4--F6 maturity),
Paper~32 \S VIII Sprint M-Z addendum (extended with F4--F6 refinement),
CLAUDE.md~\S2 (F4--F6 $\beta$-update bullet appending to the
F3-maturity baseline), CLAUDE.md~\S3 (three new dead-ends rows
for the three ruled-out W1e closure mechanisms).

\subsubsection{Cross-domain instance of the multi-focal-composition wall (W1e chemistry, Schmidt + core correlation closeout, 2026-05-23)}
\label{sec:conv_w1e_cross_domain_wall}

\textbf{Observable.} Closeout of the W1e (inner-region overattraction)
chemistry-arc sprint-scale candidate space, with cross-domain pattern
recognition. Two Day-1 diagnostics tested the two beyond-sprint-scale
candidate
closure mechanisms named in \S\ref{sec:conv_w1e_f4_f6_refinement}
(Schmidt orthogonalization at basis-construction level; fully
correlated [Ne] cores), executing the diagnostic-before-engineering
rule at sprint-cycle scale. \textbf{Both ruled out.} Schmidt produces
a $+8.6$~mHa diagonal-element differential at NaH $R = 3.5$~bohr
(only $0.2\%$ of wall depth, $12\times$ below the $100$~mHa GO
threshold); literature order-of-magnitude estimates for [Ne]
core-correlation differential at NaH $R_\text{eq}$ converge across
three independent anchors to $0.001$--$0.003$~Ha
($14$--$1000\times$ below the $0.1$~Ha DEFER threshold). The W1e
wall is now characterized as deep multi-determinant correlation
beyond all sprint-scale closure mechanisms.

\textbf{Source compilations / framing perspectives.} The W1e wall
joins five previously-identified structural walls, all in physics
observables:
(1) Sprint H1 Yukawa (the framework's AC extension admits a Higgs
structurally but does not autonomously select $Y_F$; Marcolli--van
Suijlekom Yang-Mills-without-Higgs reading at the structural level);
(2) Sprint LS-8a $Z_2/\delta m$ (iterated Connes--Chamseddine
spectral action reproduces the UV-divergent integrand of two-loop
QED self-energy with correct prefactor, sign, and divergence
structure, but does not autonomously generate the renormalization
counterterms required for the finite $C_{2S} = +3.63$ extraction);
(3) HF-3 recoil (no cross-register two-body coordinate operator
$V_{ne}(\hat r_e, \hat R_n)$ in the framework; the Track-NI nuclear--
electronic composed architecture couples spin labels across
registers but not spatial coordinates);
(4) HF-4 Zemach (no magnetization-density operator at the proton
spatial coordinate; $r_Z$ must enter as Layer-2 calibration input);
(5) HF-5 multi-loop $a_e$ (same renormalization wall as LS-8a,
vertex sector instead of self-energy).
\textbf{The W1e wall is the first chemistry-domain instance of this
pattern.}

\textbf{Convention difference and verdict.} The unified structural
reading across the six instances: the framework determines the
\emph{structural skeleton} of each observable (selection rules,
transcendental classes, scaling laws, the build of the right basis
states, the correct projection-chain decomposition), but does
\emph{not} autonomously generate the calibration data that connects
the structural skeleton to specific numerical observable values
(Yukawa values, renormalization counterterms, cross-register
spatial-coordinate operators, magnetization radii, multi-loop QED
finite parts, second-row binding energetics). \emph{Calibration
data in both physics and chemistry domains is genuinely external
input within the framework's current scope.}

\textbf{Magnitude (chemistry-side closeout).} W1e closure-mechanism
candidate space at sprint scale (F4 + F5 + F6 + Schmidt + core
correlation): joint upper-bound closure $\sim 10$--$25\%$ of the
4.37~Ha wall depth at NaH $n_{\max} = 2$; experimental $D_e =
0.075$~Ha requires $\sim 98\%$ closure. Scales mismatch by
$\sim 50\times$. Remaining candidates are major architectural
commitments beyond sprint scale (self-consistent $Z_\text{eff}$
depending on cross-center electronic density; fully-correlated
all-electron treatment beyond the active-space architecture) or
acceptance of the structural-skeleton scope as the framework's
autonomous boundary.

\textbf{Class.} This is an architecture class~(iv) exposure
(framework-level structural taxonomy refinement) and \emph{the
first explicit cross-domain instance documented in the catalogue.}
The pattern's cross-domain transferability --- physics (five
observables) plus chemistry (W1e), structurally distinct
calibration-data classes hitting the same structural-skeleton-scope
boundary --- is the substantive new content this entry contributes:
not a new convention exposure on a single observable, but a
refinement of the framework's structural-skeleton scope statement
(CLAUDE.md~\S1.7) from a five-physics-observable pattern to a
six-observable cross-domain pattern.

\textbf{Sprint reference and closure.} Schmidt orthogonalization
Day-1 diagnostic STOP
(\texttt{debug/sprint\_w1e\_schmidt\_diagnostic\_memo.md});
[Ne] core-correlation literature order-of-magnitude estimate DEFER
(\texttt{debug/sprint\_w1e\_core\_correlation\_estimate\_memo.md}).
Cross-references Paper~19 \S\ref{sec:w1e_schmidt_core_correlation}
(W1e Schmidt + core correlation closeout), Paper~32 \S VIII Sprint
M-Z addendum (cross-domain wall extension), Paper~17 \S6.10
(composed-architecture cross-reference, updated to Schmidt + core
correlation closeout), and CLAUDE.md~\S1.7 multi-focal-composition
wall taxonomy.

\paragraph{Cross-pattern observation.}
The exposures partition into \emph{three} structural classes.
\emph{Class~(i) inter-compilation itemization} (entries
\S\ref{sec:conv_rZG_DHFS}, \S\ref{sec:conv_muH_VP},
\S\ref{sec:conv_karshenboim_recoil}, \S\ref{sec:conv_d_polarizability},
\S\ref{sec:conv_he_bethe_log},
\S\ref{sec:conv_d_polarizability_mud_lamb}):
six of the eight exposures
are between Friar--Payne~2005, Eides~2024, Krauth~2017, Karshenboim~2005,
Drake~1990 / Drake--Yan~1995, Pachucki--Yerokhin~2009/2010,
Antognini~2013, Carlson--Gorchtein--Vanderhaeghen~2014 / Carlson--Vanderhaeghen~2011 compilations on items at
sub-percent precision.  The pattern is consistent: most precision-physics
compilations agree on bottom-line totals but differ in itemization at
the $10^{-4}$ to $10^{-3}$ level.  This is exactly the precision range
the framework's projection-chain decomposition operates at, hence the
surface contact between the framework's diagnostic instrument and
the literature compilations.  \emph{Class~(ii) closed-form vs
operator-level convention} (entry \S\ref{sec:conv_friar_profile}):
one exposure is between Eides closed-form derivation (which assumes
RMS-radius input) and the framework's operator-level realization
(which calibrates to first-moment input), with the two paths
differing by a profile-dependent factor that the closed-form
derivation hides.  Class~(ii) is internal to the framework's
operator-level realization choices, not a literature convention
mismatch in the strict sense, but it is the same kind of finding:
sub-percent agreement on totals masks per-channel calibration choices
that surface only when the two paths are compared at operator level.
\emph{Class~(iii) framework-architecture limitation} (entry
\S\ref{sec:conv_alkali_hfs}): a new class introduced
2026-05-18 by the multi-track Roothaan Track-5 Li-7 HFS autopsy.
The single-zeta hydrogenic basis with empirical $Z_\text{eff}$ cannot
represent the PK-orthogonalized contact density $|\psi_{ns}(0)|^2$ at
the alkali valence; the cliff is $\sim 10\times$ for Li 2s and same
mechanism as the \S\ref{sec:autopsy_cs_hfs} Z\textgreater 20 cliff.
Class~(iii) is not a literature convention mismatch nor an operator-level
calibration choice but a property of the framework's graph-native
basis architecture itself.  Closing it requires extending the basis
(Hylleraas--Eckart for low-$Z$, BBB93/KTT multi-zeta or self-consistent
HF for high-$Z$), not the projection-chain dictionary.
All three classes vindicate the directive's class~(a) finding type:
multi-observable global fits using a single dictionary expose
itemization choices that single-observable analyses cannot see.

\textbf{Pattern broadening (V.D-prediction sprint, 2026-05-10).}  V.D.5
(D HFS polarizability sub-component) is HFS-class, consistent with the
original four entries.  V.D.6 (He $2{}^3P$ relativistic Bethe-log) is
\emph{Lamb-class precision spectroscopy}, NOT strictly HFS.  The
running pattern is therefore best read as ``precision-spectroscopy
observables with multi-component Layer-2 decomposition'' rather than
``HFS-only.''  All six entries remain in the same general category:
observables where the framework operates at sub-percent precision and
where multiple independent compilations exist with potentially different
itemizations.  The next falsifier-test for this pattern is to attempt
\S~V.D entries on (a) heavy-atom HFS at coarse experimental precision
(Cs, Rb), (b) molecular vibrational/rotational spectroscopy, or (c)
non-precision atomic spectra.  If 0 of those land class-(i), Pattern C
tightens to ``exclusive to high-precision atomic spectroscopy with
multi-loop QED $+$ nuclear-structure budgets.''

\paragraph{Protocol for adding new exposures.}
Each future sprint that surfaces a class~(a) finding adds a new
\S~V.D.$N$ subsection following the five-element template:
\emph{Observable~$|$~Source compilations~$|$~Convention difference~$|$
~Magnitude~$|$~Sprint reference $+$ Closure verdict.}  The exposure
is also added as a row to Table~\ref{tab:convention_exposures}.
Each entry should cross-reference the \S~\ref{sec:layer2} projection
chain whose operator-level decomposition surfaced the convention
difference, and the originating documentation in the catalogue
(\S\ref{sec:catalog}, \S\ref{sec:offprecision}, or \S\ref{sec:autopsies}).
The catalogue is a running register; expect it to grow as the
multi-observable focal-length decomposition program (CLAUDE.md
\S~1.8) advances.

\paragraph{Predictive refinement: where to look for the next
exposures.} % Source: synthesis pass §3.4 Pattern C predictions.
The four entries above all involve hyperfine-class or muonic-Lamb
observables.  The mechanism is structural: HFS observables decompose
into many sub-leading Layer-2 corrections at sub-percent precision
(leading Zemach, NLO recoil-mixing, finite-size, multi-loop QED,
polarizability, hadronic VP), each compiled differently across
Eides~2024, Karshenboim~2005, Krauth~2017, Pachucki--Yerokhin~2010,
and Friar--Payne~2005.  Sub-percent itemization differences
compound at exactly the precision the framework operates at.  This
suggests three pre-named candidate observables for the next
class~(a) findings, in decreasing prior probability:
\begin{enumerate}\setlength\itemsep{0pt}
\item \textbf{Deuteron polarizability convention.}  Pachucki--Yerokhin
2010 versus Friar--Payne~2005 itemization of the deuteron
polarizability contribution to D HFS.  Expected magnitude: tens of
ppm.  \textbf{Verdict (V.D-prediction sprint, 2026-05-10):
CONFIRMED $\to$ V.D.5 (\S\ref{sec:conv_d_polarizability}).}
\item \textbf{Positronium HFS annihilation channel itemization.}
Czarnecki--Melnikov--Yelkhovsky~2000 versus Karshenboim~2005
itemization of the $\alpha^5$ corrections to the Ps 1S annihilation
contribution.  Expected magnitude: $\sim 100$ ppm of LO HFS.
\textbf{Verdict (V.D-prediction sprint, 2026-05-10): INCONCLUSIVE} ---
the three relevant compilations (CMY 2000, Karshenboim 2005,
Adkins--Czarnecki 2014) operate at three different orders of the
expansion (m$\alpha^6$, review-aggregate, m$\alpha^7$) and are not
in convention conflict at the sprint scope; flagged for follow-up
PDF-level diagnostic (sprint-scale) on Karshenboim 2005 \S~6.2-6.5
vs CMY 2000 vs Adkins--Czarnecki 2014.
\item \textbf{Helium fine-structure $\alpha^3 (Z\alpha)^2$
itemization.}  Drake~1971 versus Pachucki--Yerokhin~2010
itemization of higher-order Breit--Pauli plus multi-electron
spin-spin / spin-other-orbit content.  Expected magnitude: tens of
MHz on GHz-scale intervals.  This is the only Lamb-class candidate
in the predicted set; the other two are HFS-class.
\textbf{Verdict (V.D-prediction sprint, 2026-05-10): CONFIRMED $\to$
V.D.6 (\S\ref{sec:conv_he_bethe_log})} --- Pachucki--Yerokhin~2009
(PRA \textbf{79}, 062516) explicitly resolved a 3$\sigma$ disagreement
with Drake on $\nu_{01}$ by reevaluating the relativistic Bethe-log
term, with magnitude $\sim$kHz on the 29~GHz observable
($\sim 30$~ppb on $\alpha$ determination).  This is the strongest of
the three V.D-prediction sprint candidates because the literature
explicitly documents both the disagreement and its resolution.
\end{enumerate}
The protocol classifies entries by predicted observable class:
class-(i) literature-itemization (V.D.1, V.D.2, V.D.3, V.D.5, V.D.6
of the existing catalogue) versus class-(ii) closed-form-vs-operator-level
(V.D.4 of the existing catalogue).  Future sprints adding entries
should tag which class the new entry falls under.

\textbf{Pattern C status (2026-05-10).}  After the V.D-prediction sprint:
6/6 \S~V.D entries are precision-spectroscopy class observables
(HFS, muonic-Lamb, or Lamb-class fine structure).  Pattern C is
\textbf{CONFIRMED at n=6 and STRENGTHENED}; the pattern broadens
slightly from ``HFS-only'' to ``precision-spectroscopy with
multi-component Layer-2 decomposition'' (V.D.6 is Lamb-class He
$2{}^3P$ fine structure, not strictly HFS).

\textbf{What would falsify Pattern C.}  If the next sprints produce
\S~V.D entries from heavy-atom HFS at coarse experimental precision
(Cs, Rb), molecular vibrational/rotational spectroscopy, or
non-precision atomic spectra, Pattern C is precision-spectroscopy-specific
rather than HFS-specific (a refinement, not a falsification).
A genuine falsifier requires \emph{0 of those land class-(i)}, in
which case Pattern C tightens to ``exclusive to high-precision atomic
spectroscopy with multi-loop QED $+$ nuclear-structure budgets''
rather than broadening further.  Conversely, a class-(i) finding on
a \emph{non}-precision-spectroscopy observable would falsify the
current pattern reading entirely.

% -------------------------------------------------------------------
\subsection{Falsified candidate: W3 spectral-zeta identification of
SM calibration data (sprint 2026-05-08, falsified same evening)}
\label{sec:calibration_candidates}
% -------------------------------------------------------------------

This subsection preserves the record of a single-session candidate
that was \emph{falsified} the same evening by three independent
follow-up tracks.  It is retained in the catalogue because the
methodological lesson -- how a $6$--$10\sigma$ signal under a
hand-curated prespecified basis can vanish under a mechanically
generated basis -- is a sharper instance of the curve-fit-audit
discipline (see \texttt{docs/curve\_fit\_audit\_memo.md}) than any
abstract description.

\paragraph{The original candidate.}
On 2026-05-08 evening, a conversational sprint produced candidate
spectral-zeta identifications of all four CKM Wolfenstein parameters
with master-Mellin-engine M1 and M2 ring elements:
$\lambda^2 \approx 1/\mathrm{Vol}(S^3) = 1/(2\pi^2)$ and
$\bar\rho \approx 1/(2\pi)$ in the M1 Hopf-base family;
$\bar\eta \approx (\ln 2)/2 = -\zeta'(0,1/2)$ and
$A^2 \approx \ln 2$ in the M2 chirality family.  The four-form fit
gave 7 hits at $<1\%$ in a hand-curated 30-form prespecified basis
(random expectation $0.7\pm 1.0$); a derived structural relation
$A^2 = 2\bar\eta$ survived PDG at $0.52\sigma$; a closed form
$\tan(\delta_{CP}) = \pi\cdot\ln(2)$ at $0.02\sigma$ provided a
striking single-number target.  The aggregate signal was reported as
approximately $6$--$10\sigma$ above chance for that basis size.  The
candidate motivated a draft addition to CLAUDE.md~\S 1.7 (WH7) and
prompted a forward plan in
\texttt{debug/w3\_forward\_plan\_memo.md} listing six tests.

\paragraph{Three-track falsification (same evening).}
The first three tests of the forward plan ran in parallel as
independent sub-agent dispatches.  All three returned negatives.

\begin{itemize}
\setlength\itemsep{2pt}
\item \textbf{PMNS recheck.}  The same 30-form prespecified basis
applied to the PMNS angles
$(\sin^2\theta_{12}, \sin^2\theta_{23}, \sin^2\theta_{13},
\delta_{CP}^{\rm PMNS})$ at PDG~2024 NH best-fit values returned
\emph{zero} matches at $<2\%$ across all four parameters.  Aggregate
signal: $-0.68\sigma$ below chance.  PMNS sits closer to tri-bimaximal
pure rationals $(1/3, 1/2, 0)$ than to any M1/M2 form.  The cleanest
single PMNS hit is $\theta_{12}\approx \arcsin(1/\sqrt{3})$ at
$0.11\,\sigma$ -- the TBM solar-angle prediction, not in the master
Mellin rings.  W3 universality across mixing matrices is falsified.
Files: \texttt{debug/w3\_pmns\_recheck.\{py,\_memo.md\}}.

\item \textbf{Lepton mass spectrum.}  The same methodology applied to
$\sim$30 charged-lepton mass observables (ratios, log spacings,
sqrt-mass ratios, Koide $K$, cone position) crossed with the 30-form
basis and 10 small-rational factor multipliers (9000 tests total)
returned only $2$ within-$1\sigma$ matches, both encoding the same
Koide $45^\circ$ cone fact via algebraically equivalent identities
($2/3 \times (\varphi^2 - \varphi) = 2/3$ and
$4 \times \pi/16 = \pi/4$).  Statistical signal: $+0.85\sigma$ above
chance, indistinguishable from random.  Hit-density gap CKM~vs~lepton:
$12\%$ vs $0.22\%$ ($50\times$).  W3 identification of mass spectra
is falsified.  Files:
\texttt{debug/w3\_lepton\_mass\_recheck.\{py,\_memo.md\}}.

\item \textbf{Mechanically-generated basis sweep.}  The hand-curated
30-form basis was replaced by a mechanically generated 21,448-form
basis: combinations of M1/M2/M3 seeds (Hopf-base measures, Hurwitz
shifts, Catalan~$G$, Dirichlet~$\beta$, $\zeta_R$ at integer points)
with small-integer rational prefactors and small algebraic factors
($\sqrt{2},\sqrt{3},\sqrt{5},\varphi$), plus a pure-rational control
class.  The full basis was frozen and saved to JSON before any test
run.  Re-running the four-Wolfenstein search returned aggregate
$z = -0.52$ (slightly \emph{below} random chance); per-parameter
z-scores all within $\pm 1\sigma$ of chance ($\lambda$ $+0.74$,
$A$ $-0.02$, $\bar\rho$ $-1.01$, $\bar\eta$ $+0.82$).  The original
M1$\leftrightarrow$real / M2$\leftrightarrow$CP-imaginary structural
reading does not survive: $\bar\eta$'s top mechanical match is in the
M1+algebraic class ($\pi^{3/2}/16$), not M2; $A$'s original
$\sqrt{\ln 2}$ does not appear in the top-20 sigma-ordered slice; and
$\lambda$'s top match is the pure rational $9/40 = 0.225$ exactly,
beating every transcendental form.  Selection-bias verdict: fully
attributable.  Files:
\texttt{debug/w3\_mechanical\_basis.\{py,\_memo.md\}}.
\end{itemize}

\paragraph{Three independent tests, three negatives.}
The W3 spectral-zeta candidate is falsified.  The original signal was
basis-selection bias.  WH7 is not promoted to CLAUDE.md~\S 1.7.

\paragraph{Methodological lesson.}
The original ``$5$--$10\sigma$ above chance'' calculation was
statistically correct \emph{for that basis size}.  The error was the
implicit assumption that the 30 hand-curated forms were drawn without
knowledge of the data; in fact, the basis was disproportionately near
the empirical parameter values.  Mechanical generation removes the
implicit selection, and when removed, the signal vanishes.  This is
the cleanest illustration the project has produced of why the
curve-fit-audit discipline (independent re-run with strictly
mechanically-generated basis) is necessary, not optional, for
empirical pattern claims that touch transcendental rings of any
size.

\paragraph{Surviving observations independent of W3.}
The charged-lepton Koide $45^\circ$ cone (1 arcsec) is independent of
the W3 sprint and stands.  One isolated match is preserved as a
record (no structural reading attached): the CKM Wolfenstein
parameter $\lambda = 0.225$ is exactly $9/40$ to $5$-digit PDG
precision, which in GeoVac vocabulary reads as $\lambda = 9\Delta$
with $\Delta^{-1} = g_3^{\rm Dirac} = 40$ (Paper~2 Dirac-mode
boundary count); whether this is structural or Diophantine
coincidence at PDG precision is not addressable from a single data
point and the framework does not commit to either reading.

\paragraph{Standing position on calibration data.}
GeoVac's structural-skeleton scope (Bertrand truncation, inner-factor
Mellin trivialization, H1/G3 verdicts, LS-8a renormalization gap)
is restored as the standing position on calibration data: parameter
values that are not autonomously generated by the framework are
genuinely external input.  The second-packing-axiom question for
calibration data remains open at the same level it was on
2026-05-07.

% -------------------------------------------------------------------
\subsection{Lorentzian transfer audit (Sprint L0, May 2026)}
\label{sec:lorentzian_transfer_audit}
% -------------------------------------------------------------------

This subsection is a structural partition table, not an empirical
match.  It is placed in \S~\ref{sec:matches} because it accumulates
sprint-by-sprint refinements parallel to \S~\ref{sec:offprecision}
(off-precision matches), \S~\ref{sec:autopsies} (Roothaan autopsies),
and \S~\ref{sec:convention_exposures} (literature convention
exposures).  The Sprint L0 audit
(\texttt{debug/lorentzian\_l0\_audit\_memo.md}, May~2026) classifies
this section's twenty-eight projections (\S~\ref{sec:layer2}) by their
transfer behavior under the Riemannian $(3, 0) \to$ Lorentzian $(3, 1)$
extension in the Bizi--Brouder--Besnard~\cite{bizi_brouder_besnard2018}
signature classification.  The full structural framing of the
signature partition is in Paper~31~\S~8~\cite{paper31}.  The principal
falsifier of the audit's structural-correspondence verdict on the
Wick-rotation projection (\S~\ref{sec:proj_wick_rotation}) is the
operator-system-level computation flagged in \S~\ref{sec:open} open
question 6.

The audit's trigger was the Nieuviarts~2025
twisted-spectral-triple morphism~\cite{nieuviarts2025a}, which provides
an explicit (Riemannian twisted spectral triple) $\to$
(pseudo-Riemannian spectral triple) morphism via a Krein-product
construction.  Whether the construction applies cleanly to
GeoVac $S^3 = \mathrm{SU}(2)$ (odd-dimension caveat) is the parallel
L2-Nieuviarts-scoping pass and is \emph{not} assumed in the audit.

\paragraph{Classification buckets.}
Four buckets are used (see Paper~31~\S~8.1 for the structural framing
in $\mathcal{A}/D$ partition language):

\begin{description}
\setlength\itemsep{2pt}
\item[\textsc{transfers freely (TF)}.]  Signature-blind: lives on the
$V$ or $D$ axis of \S~\ref{sec:dictionary} only, no Riemannian volume
form, no Mellin transform of $\mathrm{Tr}\,e^{-tD^2}$, no Latr\'emoli\`ere
propinquity.  Holds at $(3, 0)$ and $(3, 1)$ identically.
\item[\textsc{wick-map free (WMF)}.]  Reduces to the M1 sub-mechanism
of the master Mellin engine via the four-witness Wick-rotation
theorem (Hawking + Sewell + Bisognano-Wichmann + Unruh codified in
Sprint Unruh-pendant, May~2026; see Paper~32 \S~VIII Remark on the
Bisognano-Wichmann reading).  Lorentzian reading inherited at no
structural cost at the metric-functional level.
\item[\textsc{euclidean specific (ES)}.]  Requires Krein-space
machinery, Lorentzian-native heat-trace formalism, or Lorentzian
analog of Latr\'emoli\`ere propinquity.  Named blockers B1
(Lorentzian propinquity, $6$--$12$~months original NCG-mathematics
work, possibly shortened by Nieuviarts~2025 morphism), B2 (Krein
spectral triple at $(3, 1)$ per BBB~2018), B3 (Lorentzian-native
heat-trace formalism); see Paper~31~\S~8.2.
\item[\textsc{mixed (M)}.]  Parts transfer freely, parts require
extension.  Sub-mechanism-specific decomposition required.
\end{description}

\paragraph{The 28-projection table.}
Table~\ref{tab:lorentzian_transfer} gives the per-projection
classification with one-line mechanism.  The machine-readable form
with confidence ratings is at
\texttt{debug/data/lorentzian\_l0\_audit.json}.

\begin{table}[h]
\centering\footnotesize
\caption{Lorentzian transfer audit, Sprint L0.  Per-projection
classification of the twenty-eight entries of \S~\ref{sec:layer2}.
T.F.~=~\textsc{transfers freely}; W.M.F.~=~\textsc{wick-map free};
E.S.~=~\textsc{euclidean specific}; M.~=~\textsc{mixed}.  Confidence:
H high, MH medium-high, M medium.  Mechanism column gives one-line
rationale; for full discussion see \texttt{debug/lorentzian\_l0\_audit\_memo.md}.}
\label{tab:lorentzian_transfer}
\begin{tabular}{p{0.4cm} p{4.3cm} p{0.8cm} p{0.5cm} p{8cm}}
\toprule
\# & Projection (\S~III) & Class & Conf & Mechanism / Blocker \\
\midrule
1 & Fock conformal & E.S. & H & Riemannian sphere quotient; B2. \\
2 & Hopf bundle & E.S. & H & Riemannian fibration; Vol($S^2$) is M1 generator. \\
3 & Bargmann--Segal & T.F. & M & 1st-order complex Euler op; $\pi$-free certificate signature-blind. \\
4 & Stereographic / conformal & E.S. & H & Riemannian conformal map; B2. \\
5 & Sturmian reparameterization & T.F. & H & Algebraic re-labeling on $\mathcal{A}$. \\
6 & Spectral action (CC) & E.S. & H & Tr~$f(D^2/\Lambda^2)$ Euclidean by construction; B3. \\
7 & Camporesi--Higuchi spinor lift & E.S. & MH & Riemannian spin manifold; B2 Krein spinor at $(3,1)$. \\
8 & Wigner $3j$ angular coupling & T.F. & H & SU(2) rep theory; canonical universal sector. \\
9 & Wigner $D$ rotation & T.F. & H & Spatial rotations; SU(2). \\
10 & Wilson plaquette & M. & M & Combinatorial vocab T.F.; Marcolli-vS $=$ YM E.S. \\
11 & Vector-photon promotion & W.M.F. & M & $1/(4\pi)$ per loop $=$ M1 Hopf-base. \\
12 & Mol-frame hyperspherical & T.F. & MH & Angular separation on $\mathcal{A}$. \\
13 & Drake--Swainson asymptotic subtraction & T.F. & H & Central tier (Paper~31 TS-C); signature-blind. \\
14 & Rest-mass & T.F. & H & $D^2 \to D^2 + m^2\mathbb{1}$ central; signature-blind. \\
15 & Observation / temporal-window & W.M.F. & H & The Wick-engine: Matsubara $2\pi k/\beta$. \\
16 & Two-body Dirac / Breit retardation & M. & M & Angular T.F.; kinematic needs Lorentzian re-derivation. \\
17 & Charge density (Foldy/Friar) & T.F. & H & Layer-2 scalar input; signature-blind. \\
18 & Magnetization (Zemach) & T.F. & H & Layer-2 scalar input. \\
19 & Tensor multipole $Q_N$ & T.F. & H & Wigner-Eckart; pure SU(2). \\
20 & Phillips--Kleinman & T.F. & H & Pseudopotential projection on $\mathcal{A}$. \\
21 & Multipole / Gaunt termination & T.F. & H & 3j triangle inequality; pure SU(2). \\
22 & Bipolar / Drake combining & T.F. & H & Multi-electron angular; pure SU(2). \\
23 & Symmetry / Young tableau & T.F. & H & $S_N$ rep theory; signature-blind. \\
24 & Adiabatic / Born-Oppenheimer & T.F. & MH & Scale separation; relativistic BO well-defined. \\
25 & Coupled-channel / adiabatic curve & T.F. & M & ODE machinery; transcendental inherits from Fock. \\
26 & Gauge choice (Coulomb/Lorenz/Feynman) & W.M.F. & MH & Coulomb-gauge $1/(4\pi) = $ M1; Ward identity holds. \\
27 & Wick rotation / signature change & W.M.F. & H & THE BRIDGE; Krein-space period closure bit-exact but signature-blind --- Lorentzian ``literal identification'' withdrawn 2026-06 (compact-boost / $K^+$ descope; \S~\ref{sec:proj_wick_rotation}). \\
28 & Apparatus identity / state-side & T.F. & MH & State-side; PSLQ-disjoint from master Mellin engine. \\
\bottomrule
\end{tabular}
\end{table}

\paragraph{Headline distribution.}
Seventeen projections classify as \textsc{transfers freely}
($61\%$); four as \textsc{wick-map free} ($14\%$); five as
\textsc{euclidean specific} ($18\%$); two as \textsc{mixed} ($7\%$).
Three-quarters of the dictionary inherits a Lorentzian reading at no
further structural cost at the metric-functional level
(\textsc{transfers freely} + \textsc{wick-map free} $= 75\%$).  The
remaining $25\%$ (\textsc{euclidean specific} + \textsc{mixed}) sits
exactly at the load-bearing structural extensions that define future
Sprints L1 (operator-system literal-identification of Wick rotation,
$4$--$8$~weeks), L2 (BBB Krein lift, $3$--$6$~months), and L3
(Lorentzian propinquity, $6$--$12$~months, possibly skippable; see
Paper~31~\S~8.5).

\paragraph{Two predicted trivializations at $(3, 1)$.}
The audit produces one structural prediction worth flagging:\
under the BBB sign table at $(m, n) = (3, 1)$, the M3 sub-mechanism
of the master Mellin engine (vertex-parity Hurwitz / Dirichlet-$L$
content in Camporesi-Higuchi vertex sums, Paper~28~\S~QED-vertex)
becomes structurally empty on chirality-symmetric spectrum.  The
$\{J, \gamma_5\}$ anticommutation flip at $(3, 1)$ removes the
asymmetry that drives M3's half-integer-Hurwitz content.  This would
simplify Paper~32~\S~VIII's case-exhaustion theorem from three
mechanisms (M1, M2, M3) to two (M1, M2) on the Lorentzian side ---
not a contradiction with the Riemannian theorem, but a structural
sharpening that becomes a Sprint L2 consistency check
(cf.~Paper~31~\S~8.4 and Paper~32~\S~VIII.E).

\paragraph{Sprint L2-E:\ Krein-space period closure of the
Wick-rotation projection --- signature-blind (2026-05-17; the
``literal identification'' reading was withdrawn 2026-06, below).}
The Wick-rotation projection (\S~\ref{sec:proj_wick_rotation}) was
named at Sprint L0 as the BRIDGE projection from Riemannian to
Lorentzian content;\ the operator-system-level literal identification
on the Lorentzian Krein side was named there as the principal Sprint L1
target (Riemannian side) and Sprint L2-E target (Lorentzian side).
Sprint L2-E (\verb|geovac/modular_hamiltonian_lorentzian.py|,
Paper~32~\S~VIII.E.E
\cite{paper32}) closes the Lorentzian-side identification at finite
cutoff:\ on $\mathcal{K}_{n_{\max}, N_t} = \mathcal{H}_{\rm GV}^{n_{\max}}
\otimes \mathbb{C}^{N_t}$ at signature $(3, 1)$, with the Lorentzian
wedge $W_L = P_W^{\mathrm{spatial}} \otimes P_{t \geq 0}$ and the
BW choice $H_{\mathrm{local}} = K_L^{\alpha, W}/\beta$, the Krein-space
period-closure identities close bit-exactly:\ BW-$\alpha$
period closure, BW-$\gamma$ Tomita period closure, flow conjugacy
$\sigma_t^{L, \mathrm{TT}} = \sigma_{-t}^{L, \alpha}$, and six-witness
collapse all at residual $\leq 4 \times 10^{-16}$ across
$(n_{\max}, N_t) \in \{1, 2, 3\} \times \{1, 11, 21\}$.  At $N_t = 1$
the construction reduces bit-identically to Paper~42's Riemannian
closure~\cite{paper42} (the $N_t = 1$ slice), on the same compact KMS
$\beta = 2\pi$ circle.  The
structural finding that survives:\ Paper~42~\S~7.2 ``$H_{\mathrm{local}}
\neq D_W$'' open question O3 is \textbf{signature-INDEPENDENT at the
Riemannian limit} ($\|H_{\mathrm{local}} - D_L^W\|_F|_{N_t=1}$ equals
the Riemannian-side residual bit-exact at every tested $n_{\max}$);\
the spectral-action--vs--modular-Hamiltonian generator distinction is
a deeper structural feature of the framework, not a peculiarity of
the round-$S^3$ Riemannian truncation.  \emph{Status (2026-06):} the
``literal identification at the Krein operator-system level'' reading is
\textbf{withdrawn} --- the Paper~45 $K^+$-annihilation theorem
(2026-06-09) and the compact-boost closure (2026-06-19) show the
truncated BW boost $K = \mathrm{diag}(2m_j)$ has integer spectrum
($\mathrm{e}^{2\pi\mathrm{i}K} = I$ bit-exact), so the period closure is
the compact KMS $\beta = 2\pi$ circle and the Lorentzian signature is
metrically invisible at finite cutoff (Euclidean/convention; see the
status update in \S~\ref{sec:proj_wick_rotation}).  What survives is the
period-closure computation as a compact-side KMS statement.  No
finite-cutoff Lorentzian propinquity construction exists.

\paragraph{Sprint L2-D refinement of the M3 trivialization prediction
(2026-05-16).}
The L0 prediction above was tested directly by Sprint L2-D
(Paper~32~\S~VIII.D, computational driver
\verb|debug/data/l2_d_connes_axiom_audit_31.json|) and found to be
\emph{convention-dependent}.  Under the
$n_{\mathrm{Fock}}$-parity convention --- the Paper~28~\S~QED-vertex
convention that accesses Catalan $G$ and Dirichlet $\beta(4)$ via
quarter-integer Hurwitz shifts at $a = 3/4$ and $a = 5/4$ --- the
vertex-parity sum $D_{\mathrm{even}}(4) - D_{\mathrm{odd}}(4)$ on the
chirality-symmetric $(3, 1)$ truncation is bit-identical to the
Riemannian value at every $n_{\max} \in \{1, \ldots, 5\}$: the L0
prediction is \emph{falsified} under this reading.  Under the
\emph{chirality-pairing} convention (pair ``even'' with $\chi = +1$
Weyl and ``odd'' with $\chi = -1$ anti-Weyl), the chirality-symmetric
truncation gives the difference $= 0$ bit-exact, but tautologically
by symmetry alone; the L0 prediction is \emph{trivially confirmed}
under this reading but does not access Paper~28's M3 content.
\textbf{Structural refinement:}\ the BBB sign-flip
$\{J, \gamma_5\} = 0$ at $(m, n) = (4, 6)$ is a
spectral-triple-axiom statement, not a vertex-parity-sum statement.
M3's mechanism is sectional in $n_{\mathrm{Fock}}$-parity (the chirality-independent
label used in Paper~28's vertex-parity sum), not in spacetime
signature.  This is the corresponding informational annotation on the
M3 entries under \S~III.6 / \S~III.7 / \S~\ref{sec:proj_wick_rotation};
it is not a status flip on a Table~\ref{tab:lorentzian_transfer} row.
Paper~18~\S~III.7 records the corresponding sharpening of the master
Mellin engine domain partition.

\paragraph{Scope and limits.}
This is a first-pass classification by an analyst, not a theorem.
Confidence ratings flag medium-confidence rows that may shift to
either side under Sprint L1 / L2 actual computations.  In particular,
two MIXED rows (Wilson plaquette \#10, Breit retardation \#16) and
one TRANSFERS\_FREELY row at medium confidence (coupled-channel
\#25, which inherits an algebraic-implicit $\pi$ from Fock) may
move under finer audit.  The audit explicitly does not assume that
Nieuviarts~2025's morphism applies cleanly to GeoVac
$S^3 = \mathrm{SU}(2)$ (odd-dim caveat); that question is the
L2-Nieuviarts-scoping pass.  The Bekenstein-Hawking entropy sprint
(\texttt{debug/bekenstein\_hawking\_sprint\_scoping\_memo.md}) is
correctly DEFERRED and outside the L0/L1/L2/L3 sprint scope:\ BH
entropy is Euclidean spectral-action work that does not require
Lorentzian extension.

\paragraph{Cross-references.}
The structural framing of the Sig/Op partition is Paper~31~\S~8
\cite{paper31}.  The bridge projection (\S~\ref{sec:proj_wick_rotation})
already names the operator-system-level deepening as the principal
Sprint L1 falsifier;\ \S~\ref{sec:open} open question~6 carries the
forward-looking task list.  Paper~32~\S~VIII.E sharpens the master
Mellin engine's M1 sub-mechanism as the free-transfer engine via the
four-witness theorem.

% ===================================================================
\section{Dimensional Axiomatization}
\label{sec:axiomatization}
% ===================================================================

The projection dictionary of \S~\ref{sec:dictionary} treats each
projection as a graded map along three axes: variables added,
dimensions added, and transcendental class introduced.  Sprint TX-A
(May~2026) treated this graded structure as an axiomatic system and
tested three theorems analytically.  This section consolidates the
verdicts; full data are in
\texttt{debug/data/tx\_a\_projection\_table.json} and
\texttt{debug/data/tx\_a\_composition\_table.json}, with verification
driver \texttt{debug/tx\_a\_signature\_arithmetic.py}.

\begin{theorem}[Dimensional Consistency]
\label{thm:dimensional_consistency}
For any projection chain $S = P_n \circ \cdots \circ P_1$ applied to
the bare graph (Layer~1, dimensionless), the variable, dimension, and
transcendental signatures of $S$ equal the union of the signatures of
the constituent $P_i$:
\[
  \mathrm{sig}(P_n \circ \cdots \circ P_1)
   = \bigcup_{i=1}^{n} \mathrm{sig}(P_i).
\]
The catalogue rows of \S~\ref{sec:catalog} satisfy this consistency
check across all eight fully-tagged entries with non-trivial chains
(Lamb shift, $K = \pi(B + F - \Delta)$, Stefan-Boltzmann, He CI,
hydrogen $E_n$, scalar/Dirac spatial Casimir, Uehling VP).
\end{theorem}

The theorem is true by construction at the level of the dictionary
metadata: each projection's signature is fixed in
\S~\ref{sec:layer2}, so set/list union is a tautology.  Its
non-trivial content is the empirical check that the catalogue's
expected target signature matches the chain signature; this passes
for all eight tested rows.

\begin{theorem}[Weak Minimality]
\label{thm:minimality}
Two chains containing the same multiset of projections produce the
same signature, hence are observationally indistinguishable at the
level of dimensional types and transcendental classes.  They may
differ at the operator-system level when bundle-non-trivial
projections (Hopf, spectral action, Wilson) are composed with the
spinor lift; these inequivalences are not visible to signature
arithmetic and require explicit operator-system analysis (cf.\ the
WH1 spectral truncation rounds 1--3 documented in
\texttt{CLAUDE.md} \S~2).
\end{theorem}

The strong form of minimality (uniqueness up to commutation as
operators on the Hilbert space) fails for the Paper~2 $K$-chain:
\texttt{hopf} composed with \texttt{spinor\_lift} acts differently in
the two orderings (Hopf-then-spinor produces Hopf-$S^2$ with Dirac
fiber; spinor-then-Hopf gives a $\mathbb{Z}_2$-quotient of Dirac
$S^3$ with half-integer charges).  After the addition of
\texttt{breit\_retardation} (\S~\ref{sec:proj_breit_retardation}),
a second non-commuting pair appears:
\texttt{rest\_mass} composed with \texttt{breit\_retardation} is
ordered, with rest-mass first.  The non-commutation is semantic:
Breit retardation's variable $m_l/m_n$ is a refinement of the
rest-mass variable~$m$ (it is the ratio between two species' masses,
both of which are introduced by the rest-mass projection), so it is
not well-defined until rest-mass has acted.  The metadata-level
test correctly flags both
$(\texttt{hopf}, \texttt{spinor\_lift})$ and
$(\texttt{rest\_mass}, \texttt{breit\_retardation})$ as the
non-commuting pairs in the worked examples; the underlying
mechanisms differ
(bundle-non-trivial composition for the first, variable-refinement
dependence for the second).

\begin{theorem}[Composition Algebra]
\label{thm:composition_algebra}
The 16$\times$16 composition table classifies each ordered pair
$(P_1, P_2)$ as one of: \emph{commute} (admissible in both orders,
disjoint variables); \emph{one-way} (admissible in only one order,
constrained by Layer transitions or by semantic dependence between
variables); \emph{not composable}; \emph{idempotent} (Layer-1
self-composition); \emph{self-compose} (non-trivial diagonal);
\emph{variable conflict} (same variable name added by both);
\emph{inverse pair} (one is the inverse of the other).  Of
$16^2 = 256$ cells:
\begin{itemize}
\setlength\itemsep{0pt}
\item $\sim 196$ commute ($\sim 76.6\%$; updated from 15$\times$15
count of 182 by adding the 14 new commuting cells in P$_{16}$'s row
and column with projections that compose layer-cleanly under disjoint
variable axes);
\item $\sim 26$ one-way pairs (one-way cells flagged either by Layer
transitions or by semantic dependence between variables; updated by
$+4$ from the 15$\times$15 count of 22 to account for the four new
one-way cells flagged below in P$_{16}$'s row and column);
\item $6$ variable conflicts (energy-scale-conflating compositions
of Fock, Bargmann--Segal, and spectral action; the new P$_{16}$
introduces no further variable conflicts since its variable
$m_l/m_n$ is distinct from the existing variable set);
\item $4$ idempotent projections on the diagonal (Sturmian,
Wigner~$3j$, Drake--Swainson, and now two-body Dirac / Breit
retardation: applying the Breit Hamiltonian twice gives the same
two-body output);
\item $3$ inverse pairs (stereographic, Wigner~$3j$, Wigner~$D$;
unchanged by the new projection).
\end{itemize}
The previous version of this theorem flagged \texttt{drake\_swainson}
as the unique projection with input\_layer = output\_layer = 2.
After the addition of \texttt{breit\_retardation}
(\S~\ref{sec:proj_breit_retardation}), \emph{two} projections now
have input\_layer = output\_layer = 2: Drake--Swainson and Breit
retardation.  The structurally distinguishing fact is that they
operate on disjoint Layer-2 content (Drake--Swainson regularizes
divergent Bethe-log spectral sums; Breit retardation adds two-body
kinematic content to the framework's single-particle Lamb-shift
bracket), so they commute as a $(P_1, P_2)$ pair; their
non-uniqueness on the layer index is harmless for signature
arithmetic.

\textbf{New structural finding (the load-bearing T3 update from
this addition).}  The pair
$(P_{14}, P_{16}) = (\text{rest-mass projection},
\text{Breit retardation})$ does \emph{not} commute, but the
non-commutation is \emph{not} a Layer transition --- both projections
output to Layer~2 from input layers 1 and 2 respectively, and the
Layer indices alone would mark this as a one-way pair under the
existing classification.  The \emph{semantic} reason for one-way
ordering is sharper: the Breit projection's variable $m_l/m_n$ is
the rest-mass ratio between two species of particle in the bound
state, which is well-defined only after rest-mass projection has
assigned masses to each species.  In the reverse order
$P_{14} \circ P_{16}$ the variable $m_l/m_n$ would have to be
introduced before the masses themselves exist, which is semantically
inadmissible.  This is the first instance in the Paper~34 dictionary
of a one-way pair driven by a \emph{variable refinement} dependence
rather than by a Layer transition.  The signature-arithmetic test
correctly flags both the
$(P_{14}, P_{16})$ and the
$(P_{16}, P_{14})$ cells, and the test agrees with the structural
reading: Breit retardation must follow rest-mass projection in any
admissible chain.

\textit{Erratum (Track TS-C, 2026-05-04):} an earlier draft of this
section claimed that \texttt{spectral\_action} and
\texttt{observation\_window} were mutually one-way, encoding the
field-theoretic UV/IR ordering anomaly.  That claim was retracted
on direct inspection of the composition data: both projections
target output\_layer=2 from input\_layer=1 with disjoint variable
sets, and signature arithmetic correctly returns \texttt{commute}
in both directions.  The UV/IR anomaly ($\beta \to \infty$ does not
commute with $\Lambda \to \infty$) is real physics that the current
coarse layer-grading cannot see; refining the output\_layer index
to distinguish UV-cutoff, thermal, and flow Layer-2 outputs is the
natural next step.

\textit{Note on cell counts.}  The exact 16$\times$16 cell counts
are stated above as approximate ($\sim$) totals; the load-bearing
update is the 15$\times$15 $\to$ 16$\times$16 expansion plus the
qualitative shifts (new idempotent, new one-way pair driven by
variable refinement, no new variable conflicts).  Definitive
recomputation against the metadata in
\texttt{debug/data/tx\_a\_composition\_table.json} is deferred until
the next TX-A-style audit.

\textit{Pending T3 expansion (2026-05-09).}  With the addition of
the three intra-nuclear projections
(\S~\ref{sec:proj_charge_density},
\S~\ref{sec:proj_magnetization_density},
\S~\ref{sec:proj_tensor_multipole}), the T3 composition table needs
expansion from 16$\times$16 to 19$\times$19.  The qualitative
expectation: the three new projections all introduce intra-nuclear
$[L]$ (or $[L]^2$) variables disjoint from the existing length
axes, so they should commute layer-cleanly with most existing
projections; expected new entries are $\sim 50$ commute cells in
each new row and column, with possible variable-refinement one-way
pairs against rest-mass (the intra-nuclear distribution radii are
defined per nuclear species, so they require rest-mass projection
to assign mass to each species first --- structurally analogous to
the $(P_{14}, P_{16})$ one-way pair documented above).  No new
variable conflicts are expected since intra-nuclear $[L]$ scales
are categorically disjoint from atomic / internuclear $[L]$ scales.
Definitive cell-by-cell verification deferred to the next TX-A
audit.
\end{theorem}

\begin{observation}[axes are not independent]
\label{obs:axes_correlated}
Of the three axes (variable, dimension, transcendental), only the
transcendental axis is genuinely independent.  Cross-tabulating the
twenty-eight projections by axis-flag pattern, no projection adds a
dimension without also adding the variable that carries it (perfect
V/D correlation, 12/12); zero projections fall into the
$(\text{D-only})$ bucket.  The transcendental axis admits an
independence witness: vector-photon promotion (\S~\ref{sec:layer2},
projection~11) introduces $1/(4\pi)$ without adding any variable
or dimension.  Two projections (\texttt{spectral\_action},
\texttt{observation\_window}) move all three axes; these are the
``heavy'' projections, carrying the most physical content per step.
The newly added \texttt{breit\_retardation}
(\S~\ref{sec:proj_breit_retardation}) preserves the V/D correlation
(adds the variable $m_l/m_n$ together with the dimension energy)
and lives in the ring-preserving $\alpha^4 \cdot \mathbb{Q}$ tier on
the transcendental axis; it is therefore axis-pattern-equivalent to
the rest-mass projection at the V/D level, with the transcendental
class differing only by the leading $\alpha^4$ factor inherited from
its calibration sibling, the spectral action.
\end{observation}

This observation has structural force.  The variable axis and
dimension axis collapse to a single \emph{physical-content} axis
without information loss; the transcendental axis remains
independent.  The three-axis dictionary is therefore mildly
overcomplete on the V/D side but irreducible on the
transcendental side.  See Remark~\ref{rem:vd_correlation} in
\S~\ref{sec:dictionary} for the parallel statement at the
projection-row level.

% ===================================================================
\section{Falsifiable Prediction: Layer-2-Presence Bound on
Catalogue Residuals}
\label{sec:prediction}
% Source: May 2026 synthesis pass (debug/synthesis_pattern_review_memo.md
% §2 Pattern B falsification, §2.4 Layer-2-content axis, §2.5 replacement
% prediction).  The naive depth-vs-residual form was falsified on the
% catalogue itself; this section now states the replacement
% (Prediction 1') with the Sprint Calc-P refinement repositioned as
% the cleanest L2-count = 0 exact example.
% ===================================================================

The catalogue (\S~\ref{sec:catalog}, \S~\ref{sec:offprecision}) does
\emph{not} arrange itself by projection depth.  Re-tagging each
catalogue row by the number of Layer-2 inputs consumed in its chain
(Layer-2 inputs in the sense of \S~\ref{sec:layer2_boundary}:
external scalars such as $r_Z$, $g_p$, deuteron polarizability,
multi-loop QED coefficients) reveals a much sharper structural
ordering:
\begin{itemize}\setlength\itemsep{0pt}
\item \textbf{L2-count $= 0$ (framework-native chains).}  Residuals
run from machine-exact (Stefan--Boltzmann $\pi^2/90$ at depth~3,
hydrogen $1S$ static polarizability $\tfrac{9}{2}a_0^3$ at depth~3,
spatial $S^3$ Casimir $1/240$ at depth~2) to converged-basis
ceilings ($-0.534\%$ for the H Lamb shift one-loop closure,
$\approx 0.024\%$ for He NR variational).  The residual is bounded
by basis quality alone --- error code~B in the
\S~\ref{sec:offprecision} system --- when the chain consumes no
external Layer-2 input.
\item \textbf{L2-count $= 1$ (one external scalar).}  Residuals
tighten to tens of ppm or sub-ppm (H 21cm HF-4 with Zemach $r_Z$ at
$+18$~ppm; muonium 1S--2S under rest-mass projection at $-0.11$~ppm;
$\mu$H 1S Bohr--Fermi at $+2$~ppm).  The single Layer-2 input is
always a scalar parameter and its uncertainty propagates linearly.
\item \textbf{L2-count $= 2$+ (multi-focal-wall regime).}  Residuals
are dominated by the largest single Layer-2 contribution.  Mu HFS
at $+199$~ppm is the cleanest LS-8a-class isolation (no QCD nuclear
budget); positronium 1S--2S at $+64.75$~ppm framework-native is
closed by the $\alpha^4$ Breit Layer-2 input as the
\S~\ref{sec:proj_breit_retardation} sixteenth projection;
$\mu$H Lamb at $-0.10\%$ closes by Antognini~2013 multi-loop QED
plus deuteron polarizability plus recoil~NLO with the K\"all\'en--
Sabry two-loop VP dominating.
\end{itemize}

\begin{prediction}[Layer-2-presence bound on catalogue residuals]
\label{pred:l2_presence}
The residual error of a catalogue row is bounded above by the
larger of (i)~the framework's basis-quality residual at zero
Layer-2 inputs in the chain, and (ii)~the largest single Layer-2
input consumed by the chain:
\begin{equation*}
|\epsilon(\text{observable})| \;\leq\;
\max\!\bigl(\epsilon_\text{basis}(\text{framework}),
\;\max_i |L_2^{(i)}|\bigr).
\end{equation*}
The genuine ordering axis is observable-class (zero-, one-,
multi-Layer-2-input) crossed with the framework's basis-quality
limit at zero Layer-2 inputs.  Chained projection depth is a
downstream proxy: it correlates with opportunity for compounding
when the chain consumes Layer-2 inputs, but it fails in both
directions when the chain does not (deep zero-L2 chains can be
exact, shallow multi-L2 chains can have large residuals).
\end{prediction}

This is a falsifiable structural claim about the framework, not a
phenomenological fit.  The catalogue rows above are consistent
with it.  Tests for the prediction:
\begin{itemize}\setlength\itemsep{0pt}
\item Any new framework-native chain (L2-count $= 0$) producing a
rational-analytic answer should land at residual~$0$ at converged
basis (Sprint Calc-P generalization).
\item Any new chain with a single Layer-2 input should land at
residual bounded by the uncertainty on that input.
\item Any new chain with two or more Layer-2 inputs should have
its residual dominated by the largest single Layer-2 input,
not by the sum or quadrature.
\end{itemize}
A failure of the prediction would take one of three forms:
(a)~a catalogue row with L2-count $= 0$ and residual greatly
exceeding the framework's basis-quality limit (would indicate a
structural artefact not yet identified); (b)~a catalogue row with
L2-count $= 2+$ whose residual is smaller than the largest
single Layer-2 input by more than an order of magnitude (would
indicate an unexplained cancellation between Layer-2 inputs);
(c)~a statistical fit of the catalogue showing that depth alone
provides a tighter bound than Layer-2-presence.

\begin{remark}[Naive depth-vs-residual prediction, falsified]
\label{rem:depth_falsification}
A naive earlier form of this prediction asserted that residual
scales roughly linearly with chain depth, $\epsilon_k \sim k \cdot
\epsilon_1$ with $\epsilon_1 \sim 1\%$.  The May~2026 synthesis
pass falsified this on the catalogue itself.  At depth~3 alone the
catalogue contains the H~1S polarizability ($\tfrac{9}{2}\,a_0^3$
exact, residual $0$, Sprint Calc-P), the Stefan--Boltzmann
$\pi^2/90$ (Sprint TD Track~1, residual $0$), positronium 1S--2S
framework-native ($+64.75$~ppm), D HFS Bohr--Fermi strict
($+40$~ppm), and the He oscillator strength ($+3.4\%$).  At depth~4
the catalogue contains $\mu$H Bohr--Fermi ($+2$~ppm) and the H
Lamb shift one-loop closure ($-0.534\%$).  Residuals at every
depth $\geq 2$ are scattered across more than four orders of
magnitude.  Depth correlates with \emph{opportunity} for
compounding when Layer-2 inputs are consumed, but the residual is
controlled by the Layer-2-presence count and the framework's
basis-quality limit, not by depth itself.
\end{remark}

\begin{remark}[Sprint Calc-P, the cleanest L2-count $= 0$ exact
example]
\label{rem:depth_refinement}
The H~1S static dipole polarizability calculation (Sprint Calc-P,
\S~\ref{sec:catalog} three-projection chain row) returned
$\alpha_p(1S) = \tfrac{9}{2}\,a_0^3$ \emph{exactly} via a depth-3
chain (Fock $\circ$ Wigner $3j$ $\circ$ Sturmian), residual~$=~0$.
Under the Layer-2-presence bound of
Prediction~\ref{pred:l2_presence} this is the structurally clean
``L2-count $= 0$, rational analytic answer in finite-basis span''
class: there are no external Layer-2 inputs in the chain, the
analytic answer lies exactly within the framework's truncated
basis, so the residual sits at the basis-quality limit which is
$0$ for this observable.

The mechanism is sharp.  The Dalgarno--Lewis solution
$F|1s\rangle = -\tfrac{1}{2}\, z(r{+}2)|1s\rangle$ has $p$-wave
radial part $(r^2 + 2r) e^{-r}$ that lies \emph{exactly} in the
span of the first two Sturmian basis functions $\{S_2, S_3\}$ at
$\lambda = Z$.  Combined with the framework's symbolic-exact
matrix-element machinery (rational arithmetic via
\texttt{geovac/hypergeometric\_slater.py} and
\texttt{geovac/dirac\_matrix\_elements.py}), the projection chain
produces the analytic answer with no truncation residual.

This refines the L2-count $= 0$ class of
Prediction~\ref{pred:l2_presence} along the transcendental axis:
zero-Layer-2-input observables whose analytic answer is a pure
rational with a closed-form representation in the framework's
truncated basis land at residual $0$ at finite $n_\text{max}$;
zero-Layer-2-input observables whose analytic answer requires
irrational content (continuous integration over a
temporal/spectral parameter per Paper~35 Prediction~1, or limits
inaccessible at finite $n_\text{max}$) land at the framework's
basis-quality residual.  Polarizability is the first identified
example of the rational class; static linear-response observables
that admit Dalgarno--Lewis-type closed forms are candidates for
the same exact-at-finite-$n_\text{max}$ behaviour.  The dividing
line is empirically diagnosable from the transcendental
signature: rational-class observables have pure $\mathbb{Q}$
transcendental signature, irrational-class observables have at
least one transcendental factor.
\end{remark}

% ===================================================================
\section{Connection to Paper 18 (Transcendentals) and Paper 32 (Spectral Triple)}
\label{sec:connections}
% ===================================================================

This paper does not supersede Paper~18 or Paper~32; it extends
them with a complementary axis.

\paragraph{Versus Paper~18.}  Paper~18 classifies transcendental
content of GeoVac observables (intrinsic, calibration, embedding,
flow, composition).  Paper~34 classifies the projections that
\emph{produce} that content.  The two classifications are
duals: Observation~\ref{obs:duality} maps each Paper~18 tier to
specific Paper~34 projections.  Reading direction~A
(transcendental $\to$ projection) lets you predict which projection
must be active given an observed transcendental.  Reading
direction~B (projection $\to$ transcendental) lets you predict the
transcendental signature of a calculation before running it.  The
audit memo of 2026-05-02 documents what happens when this dictionary
is unavailable: PSLQ searches against an arbitrary basis surface
near-misses that look like identifications, because the right basis
(determined by the projection) was never specified.

\paragraph{Versus Paper~32.}  Paper~32 constructs the GeoVac
spectral triple $(\mathcal{A}_{\text{GV}}, \mathcal{H}_{\text{GV}},
D_{\text{GV}})$ and audits it against the Connes axioms.  In that
language, Layer~1 of this paper is the algebra
$\mathcal{A}_{\text{GV}}$ and the Hilbert space
$\mathcal{H}_{\text{GV}}$, while Layer~2 is the family of
spectral functions, gauge sectors, and embedding maps that take
$(\mathcal{A}, \mathcal{H}, D)$ to physical observables.
Paper~34 is therefore the projection-side reading of the spectral
triple synthesis.

\paragraph{Versus Paper~31.}  Paper~31's universal/Coulomb
partition is the same observation phrased on the algebra/Dirac side:
features that depend only on $\mathcal{A}$ (e.g.\ angular sparsity)
are universal, features that depend on $D_{\text{GV}}$ (e.g.\
calibration $\pi$, the Hopf bundle's projection-$\pi$) are
potential-specific.  Paper~34's claim is the same in projection
language: zero-projection matches and Wigner-$3j$-only matches
transfer between potentials, while spectral-action and Fock-conformal
matches do not.

% ===================================================================
\section{Open Questions}
\label{sec:open}
% ===================================================================

\begin{enumerate}
\setlength\itemsep{0pt}

\item \textbf{Closure of the projection list.}  Twenty-eight projections
are now documented.  The thirteenth (Drake--Swainson asymptotic
subtraction) was added in sprint LS-4 to close the 2P Bethe-log
structural floor of LS-3.  The fourteenth (rest-mass projection) and
fifteenth (observation / temporal-window projection) were added in
sprint KG / Paper~35 (May~2026) after verifying that the bare
relativistic Klein-Gordon spectrum on $S^3 \times \mathbb{R}$ is
$\pi$-free in $\mathbb{Q}[\sqrt{d_1}, \sqrt{d_2}, \ldots]$ and that
$\pi$ enters the spectrum at exactly one step --- temporal
compactification.  These two projections were previously bundled
implicitly with other parameter-projection rows; the split was
forced by the Casimir computation $E_\text{Cas}^{S^3} = 1/240$
(spatial only, exact rational, no $\pi$) versus the Stefan-Boltzmann
$\pi^2/90$ (high-$T$ limit on $S^3 \times S^1_\beta$, $\pi$ from
Matsubara).  The sixteenth (two-body Dirac / Breit retardation,
\S~\ref{sec:proj_breit_retardation}) was added in the
precision-catalogue sprint of 2026-05-08, motivated by the
positronium 1S--2S framework-native residual at $+64.75$ ppm (the
sixth and structurally cleanest instance of the
multi-focal-composition wall).  The seventeenth, eighteenth, and
nineteenth (Foldy/Friar charge density, Zemach magnetization-density,
tensor multipole; \S\S~\ref{sec:proj_charge_density},
\ref{sec:proj_magnetization_density},
\ref{sec:proj_tensor_multipole}) were added in the multi-track
Roothaan autopsy sprint of 2026-05-09, formalizing the intra-nuclear
Layer-2 input mechanisms that surfaced as operator-level
constructions during the H 21cm and D 1S HFS autopsies.  The
twentieth through twenty-fifth projections
(\S\S~\ref{sec:proj_phillips_kleinman}--%
\ref{sec:proj_coupled_channel}: Phillips--Kleinman / core--valence
orthogonality; multipole expansion / Gaunt termination; bipolar
harmonic / Drake combining; symmetry / Young tableau; adiabatic /
Born--Oppenheimer; coupled-channel / adiabatic curve) were added in
the Sprint~1 dictionary-completion sprint of 2026-05-15.  These six
projections are transcription of pre-existing framework structures
rather than new mechanisms: the framework had been using each of them
implicitly for years, but they were not previously named in the §III
dictionary.  The six together close the textbook-completeness gap
identified in the dictionary's structural audit and bring the
catalogue from nineteen to twenty-five.

The twenty-sixth through twenty-eighth projections
(\S\S~\ref{sec:proj_gauge_choice}--\ref{sec:proj_apparatus_identity}:
gauge choice / Coulomb-Lorenz-Feynman; Wick rotation / signature
change; apparatus identity / state-side reduction) were added in the
Sprint~2 axis-completeness sprint of 2026-05-15.  These three close
the V/D/transcendental axis-grid completeness identified by the
TX-A audit: gauge choice provides a second independence witness next
to vector-photon promotion in the transcendental-only corner; Wick
rotation promotes the metric-functional-level content of the
Bisognano--Wichmann reading (Sprint TD Track~4 + Sprint Unruh-pendant,
May~2026) from a §VIII open-question candidate to a named §III
projection slot; and apparatus identity opens the state-side
complement of the dictionary, which Sprint TD Track~5's PSLQ negative
proved is categorically distinct from the spectral-side master Mellin
engine.

Sprint~3 (2026-05-15) closed three additional scoping candidates as
ABSORBED into existing entries rather than promoted to §III slots:
(s1)~Tomita--Takesaki modular flow is absorbed into the
\S~\ref{sec:proj_observation} BW-reading footnote and
\S~\ref{sec:proj_wick_rotation}, with the framework-internal realization
named as the operator-system-level Lorentzian extension below
(Sprint~3 Track~10);
(s2)~loop expansion / $\alpha$-power-sorting is absorbed into
\S~\ref{sec:proj_spectral_action} (where it is now named explicitly in
the Structural reading paragraph), with the renormalization-counterterm
gap as the natural future-direction \emph{separate} projection candidate
(Sprint~3 Track~11);
(s3)~heat-kernel regularization / Schwinger proper-time is absorbed into
\S~\ref{sec:proj_spectral_action} (where it is now named explicitly with
the master Mellin engine framing) since proper-time integration and the
Connes--Chamseddine spectral action are the same operation under two
names; the Mellin transform itself is the meta-mechanism organizing
the spectral-side projection family, not a §III peer (Sprint~3 Track~12).
The Sprint~3 verdicts collectively confirm that the dictionary's
twenty-eight projections are structurally complete:\ the master Mellin
engine of Paper~32~\S~VIII / Paper~18~\S~III.7 is fully accounted for
through M1+M2+M3 at the existing §III entries, and the
operator-system-level Lorentzian extension and state-side enumeration
are correctly classified as open extensions rather than missing
projections.

Are there projections in the framework that are still not yet in
this list?  Two pre-Sprint-3 candidates remain on review:
(a)~the Mellin / heat-kernel evaluation at fractional~$s$
(currently treated as part of spectral-action;\ Sprint~3 Track~12
confirmed this candidate is structurally separate from the heat-kernel
absorption verdict and remains open);
(b)~the direction-resolved Hodge decomposition of vector-photon
edges (Paper~28 \S~vq\_sprint).
Plus two open extensions surfaced by Sprint~2:
(c)~the full operator-system-level Lorentzian extension that would
deepen Wick rotation (\S~\ref{sec:proj_wick_rotation}) from
structural-correspondence to literal-identification status (see
``Operator-system-level Lorentzian extension'' below; this replaces
the previously-listed candidate-Lorentz-boost-projection entry, which
has been subsumed by the §III.27 promotion);
(d)~state-side dictionary enumeration --- mutual information,
conditional entropy, fidelity, trace distance, relative entropy
$S(\rho\|\sigma)$, Wasserstein--Kantorovich distance on $P(\mathcal{H})$
--- candidates as future §III entries on the state-side complement
opened by \S~\ref{sec:proj_apparatus_identity}.
Decision pending on (a), (b); structural progress on (c), (d) named
in their respective entries below.

\item \textbf{Quantitative form of Prediction~\ref{pred:l2_presence}.}
The Layer-2-presence bound replaces the falsified naive
depth-vs-residual prediction (see
Remark~\ref{rem:depth_falsification}).  Outstanding work:
(i)~quantify the framework's basis-quality residual
$\epsilon_\text{basis}(\text{framework})$ for each major chain
class (one-loop QED Lamb-shift class, NR variational class,
graph-native CI class) so that the L2-count $= 0$ side of the
prediction has a per-class quantitative bound; (ii)~test whether
multi-Layer-2 chains are dominated by the largest single
Layer-2 input (the prediction's claim) versus quadrature-sum or
worst-case sum; (iii)~classify whether ``rational analytic answer
with finite-basis representation'' (Sprint Calc-P,
Remark~\ref{rem:depth_refinement}) is a sub-class of L2-count $=
0$ that always lands at residual $0$, distinct from
``irrational analytic content'' that lands at the
basis-quality limit.

The LS-4 Lamb shift datapoint adds historical context.  The LS-4
chained projection (Fock + spectral action + Sturmian +
Drake--Swainson) achieves $-0.39\%$ residual at the sweet-spot
$N=40$.  At converged $N$, the LS-4 result reverts to the LS-1
ceiling $-3.10\%$.  Under
Prediction~\ref{pred:l2_presence} this is the
$\epsilon_\text{basis}(\text{H one-loop QED Lamb})$ value and is
the framework's converged-basis quality limit at zero Layer-2
inputs; the $-0.39\%$ at $N=40$ is a transient finite-basis
cancellation, not a genuine reduction of the basis-quality
limit.  Prediction~\ref{pred:l2_presence} applies to
converged-basis residuals.

\item \textbf{$K = \pi(B + F - \Delta) \approx 1/\alpha$.}
The Paper~2 Observation is the most extreme multi-projection
coincidence in the catalogue: three structurally independent
projections (finite Casimir trace at $m=3$, infinite scalar
Dirichlet at $s = d_\text{max}$, Dirac mode degeneracy at the
spinor cutoff), all evaluated on the same $S^3$, summing to
$1/\alpha$ at $8.8 \times 10^{-8}$ relative residual.  Under
the (now falsified) naive depth-vs-residual prediction this was
the central anomaly: a three-projection match carrying
$10^{-7}$ residual versus $\sim 3\%$ expected.  Under the
Layer-2-presence bound of Prediction~\ref{pred:l2_presence} the
chain has L2-count $= 0$ (no external Layer-2 inputs are consumed),
so the residual should sit at the framework's basis-quality limit
for the chain class.  The framework's basis-quality limit for a
zero-Layer-2-input chain involving the K combination is not
characterized: $K$ is a single number, not a converging sequence
in $n_\text{max}$.  Either something structural we have not
identified is forcing the combination, or the match is a
coincidence below the framework's intrinsic resolution.  Working
hypothesis WH5 (Paper~2 in Observations status as of 2026-05-02)
reads it as a coincidence; this paper records the structural
anomaly without resolving it.

\textit{Track TS-C update (2026-05-04).} The ``depth model
under-counts'' reading was tested by examining whether the three
chains share a fourth implicit projection.  They do not: $F$ is
derived from a Dirichlet zeta on the $(n,l)$ lattice with no Hopf
invocation (Phase~4F $\alpha$-J); $\Delta^{-1}$ is a
Camporesi--Higuchi Dirac mode degeneracy with no Hopf invocation
(Phase~4H SM-D); and the Hopf-quotient graph at $n_{\max} = 3$ does
not encode $K$ through any independent spectral invariant beyond
a Casimir relabeling (Phase~4B $\alpha$-D).  The outer $\pi$
prefactor of $K$ \emph{is} identified with the Hopf-base measure
(Sprint~A $\alpha$-PI), but this is already counted as the Hopf
projection in the depth-3 chain.  Under-counting is therefore not
the mechanism.  Combined with the K-CC clean negative on a unified
Connes--Chamseddine spectral-action source, twelve mechanisms total
have been eliminated.  The remaining live readings are (ii) and
(iii) above; the standing interpretation is (iii), per WH5.

\item \textbf{Unit-introduction asymmetry between Fock and Bargmann.}
Both Fock and Bargmann--Segal projections introduce energy as the
dimension; the former adds $Z$ and produces $\pi$ via
$\mathrm{Vol}(S^3)$, the latter adds $\hbar\omega$ and produces
no $\pi$ at all.  The asymmetry is captured by the HO/Coulomb
rigidity theorems (Paper~24, Paper~31) at the operator-order level
(first-order complex versus second-order real).  Whether the
projection-dictionary axis can independently predict the asymmetry
without invoking operator order is open.

\item \textbf{Pure-axis projections.}  The current dictionary contains
no projection that moves only the dimension axis (e.g., a
unit-rescaling that introduces a length scale without a coordinate
variable).  Whether such projections are structurally absent from the
GeoVac framework, or merely absent from the current implementation,
is open.  Two candidate additions surfaced by the TX-A axiomatization
sprint: (a)~a gauge-choice projection (Coulomb vs Lorenz) that would
sit in the pure-transcendental bucket alongside vector-photon, (b)~a
spin-statistics projection that splits boson/fermion Matsubara
content out of observation/temporal-window.  Neither is forced by
current data.

\item \textbf{Operator-system-level Lorentzian extension (deepening
of \S~\ref{sec:proj_wick_rotation}).}

\medskip
\noindent\textit{Sprint L0 update (May 2026).}  The May 2026 NO-GO
scoping on full Lorentzian extension was reopened after a named
trigger (\texttt{debug/lorentzian\_literature\_update\_2026\_05\_16.md})
fired:\ the Nieuviarts~2025 twisted-spectral-triple morphism
(arXiv:2502.18105) provides a Riemannian $\to$ pseudo-Riemannian
construction via Krein-product twist that may significantly shorten
the blocker-B1 budget.  Sprint L0
(\texttt{debug/lorentzian\_l0\_audit\_memo.md}, May~2026) executed a
documentation-only pass:\ a 28-projection transfer audit classifying
this section's projections by their behavior under the
Riemannian $(3, 0) \to$ Lorentzian $(3, 1)$ extension.  Headline
distribution: $17$ TRANSFERS\_FREELY, $4$ WICK\_MAP\_FREE,
$5$ EUCLIDEAN\_SPECIFIC, $2$ MIXED.  Full per-projection table at
\S~\ref{sec:lorentzian_transfer_audit}; structural framing of the
Sig/Op partition at Paper~31~\S~8~\cite{paper31}.

\medskip
\noindent\textit{Sprint L1 closure (2026-05-16):\ Riemannian
operator-system identification verified.}  The principal Sprint L1
falsifier --- ``compute the modular Hamiltonian $K$ on
$\mathcal{T}_{n_{\max}}$ for a wedge-restricted Camporesi--Higuchi
state, verify $\sigma_{2\pi}(O) = O$ at finite $n_{\max}$'' ---
is now closed at the strongest level on the Riemannian side
(\texttt{debug/l1\_modular\_hamiltonian\_results\_memo.md},
Paper~32~\S~VIII.F).  Sprint L1 constructs $K$ via the
hemispheric-wedge BW-$\alpha$ geometric realization
($K = J_{\mathrm{polar}}$ with integer spectrum on the full-Dirac
basis), verifies the period closure $\sigma_{2\pi}(O) = O$ bit-exactly
at $n_{\max} \in \{2, 3, 4, 5\}$ for all four physical witnesses (BW,
HH, Sew, Unruh), with cross-witness collapse bit-exact (residuals
match across witnesses at machine precision; period closure is
independent of $\kappa_g$).  This lifts the four-witness Wick-rotation
theorem from ``structural correspondence'' to ``literal identification
at the operator-system level (Riemannian)'' at every finite cutoff.
(This Riemannian period-closure survives the 2026-06 descope of the
\emph{Lorentzian} ``literal identification at the Krein level'' reading
--- see the status update in \S~\ref{sec:proj_wick_rotation} --- because
$\sigma_{2\pi} = $ identity is precisely the compact KMS $\beta = 2\pi$
circle; only the Lorentzian-signature promotion was withdrawn.)
\S~\ref{sec:proj_wick_rotation} honest-scope paragraph updated;
\S~\ref{sec:offprecision} new ``Modular periodicity'' row added as
machinery-witness for the closure.

\medskip
\noindent\textit{Residual open scope (full Lorentzian extension).}\
The remaining open question now is the Lorentzian extension to
signature $(3, 1)$ via the
Bizi--Brouder--Besnard~2018~\cite{bizi_brouder_besnard2018}
Krein-space spectral triple --- this is the named Sprint L2 follow-up
(long-range scale per the Sprint L0 audit blocker-B1 estimate; the
Nieuviarts~2025 twisted-spectral-triple morphism is the named trigger
that may shorten this budget if it applies cleanly to $S^3 =
\mathrm{SU}(2)$).  Expected outcome (per Sprint L0 prediction):
M3 trivializes at $(3, 1)$ but M1 stays, so the four-witness theorem
extends to Lorentzian.  Sprint L1's Riemannian closure is the
prerequisite that confirms the operator-system identification is
sound before lifting signature.

% Source: Sprint Lorentz-Boost-Scoping (May 2026, Track B); refines
% synthesis pass §7 Hunt 4 REQUIRES-EXTENSION verdict with
% Bisognano-Wichmann reading.  See debug/lorentz_boost_scoping_memo.md.
% Sprint 2 update (2026-05-15): the metric-functional-level content of
% the Wick-rotation reading has been promoted to §III.27 as
% \S~\ref{sec:proj_wick_rotation}; this open-question entry recasts
% the remaining content as the operator-system-level deepening (full
% Lorentzian propinquity / modular Hamiltonian at finite n_max /
% literal-identification verifier for the BW reading) rather than as
% a separate candidate projection.
A natural candidate projection is a Lorentz boost relating two
observers at relative rapidity $v/c$ on the
$T_{S^3} \otimes T_{S^1_\beta}$ tensor-product spectral triple of
Sprint~TD Track~1: under the boost, $\lambda_x \to \lambda_x / \gamma$
and $\beta \to \beta\gamma$ with $\gamma = 1/\sqrt{1 - v^2/c^2}$.
The tagged signature would be: \emph{variable} $v/c$, \emph{dimension}
dimensionless, \emph{transcendental} algebraic over $\mathbb{Q}(v/c)$
via the relation $\gamma^2 (1 - (v/c)^2) = 1$.

\textbf{Verdict: REQUIRES-EXTENSION at the framework level, but
the M1 mechanism already does Lorentz-boost-class physics
(Bisognano--Wichmann reading).}  The GeoVac spectral triple is
Wick-rotated --- the Camporesi--Higuchi Dirac spectrum is on
Euclidean $S^3$, not on Minkowski space.  Lorentz boosts are
Minkowski-signature transformations and have no natural action on a
Riemannian spectral triple at the metric / curvature level: under
Wick rotation $t \to i\tau$ a Lorentz boost in the $(t, x)$ plane
becomes a rotation in the Euclidean $(\tau, x)$ plane that mixes
the spatial $S^3$ factor with the temporal $S^1_\beta$ factor; in
Track~1's tensor-product construction the two factors are
independent with disjoint operators, and a connection between them
would be a structural extension, not a projection refinement.

\textbf{However, the temperature mechanism IS already
Lorentz-boost physics.}  The Bisognano--Wichmann theorem~\cite{bisognano_wichmann1976}
identifies the modular flow of a vacuum state restricted to a
Rindler wedge with the unitary representation of Lorentz boosts
preserving the wedge.  Sewell~1980~\cite{sewell1980} extends this
to Schwarzschild geometry: the KMS state at $\beta = 8\pi M$ on the
Euclidean cigar IS the modular flow of the Lorentz boost (more
precisely, the Killing flow generated by the timelike Killing
vector outside the horizon).  Sprint~TD Track~4 reproduced
$T_H = 1/(8\pi M)$ on the Euclidean Schwarzschild cigar via the
master Mellin engine $M_1$ Hopf-base measure mechanism, with
$\beta = 8\pi M$ forced by regularity at $r = 2M$.  Under the
Bisognano--Wichmann reading, \textbf{the $2\pi$ in $\beta = 2\pi/\kappa_g$
(or in the Unruh formula $\beta = 2\pi/a$ for a Rindler observer at
acceleration $a$) IS the period of the boost orbit, and the
framework's $M_1$ Hopf-base measure signature is the same $2\pi$.}
Sprint~TD Track~4's reproduction of the Hawking temperature is
therefore not just a Wick-rotation observation; it is the
framework's instantiation of Lorentz-boost / modular-flow physics
at the level the master Mellin engine controls.

\textbf{What the framework does NOT do} is the metric / curvature
side at finite $v$: length contraction, Wigner rotation of angular
momentum under boost, the spinor-content side of how the
Camporesi--Higuchi spinors transform under boost, and the
continuous-spectrum side (Lorentzian Dirac has continuous spectrum,
breaking compact resolvent without temporal compactification).  A
clean Lorentz-boost projection acting on \emph{all} particles
simultaneously and relating two $T_{S^3} \otimes T_{S^1_\beta}$
constructions at different rapidities requires a Lorentzian
spectral triple at signature $(m, n) \in (\mathbb{Z}/8)^2 = (3, 1)$
in the Bizi--Brouder--Besnard~2018~\cite{bizi_brouder_besnard2018}
framework (currently $(3, 0)$), with Hilbert space replaced by Krein
space carrying fundamental symmetry $\mathcal{J}^2 = I$, and the
real structure $J$ extended via the
$(\mathbb{Z}/8) \times (\mathbb{Z}/8)$-valued sign table.

\textbf{The Lorentzian-NCG literature provides several published
prescriptions} (Strohmaier~2006~\cite{strohmaier2006};
Franco~2014~\cite{franco_eckstein2014} temporal Lorentzian
spectral triples; Bizi--Brouder--Besnard 2018 $(m,n)$
classification; van~den~Dungen~2016~\cite{van_den_dungen2016} Krein
fermionic action; Devastato--Lizzi--Martinetti~2018~\cite{devastato_lizzi_martinetti2018}
twisted-emergence) for what a Lorentzian extension would look like
at the spectral-triple level.  \emph{None} provides a Lorentzian
analog of the Latr\'emoli\`ere propinquity used in
Paper~38~\cite{paper38} for the WH1 PROVEN result; the propinquity
literature is purely Riemannian as of May~2026 (Latr\'emoli\`ere,
Toyota, Hekkelman--McDonald, Leimbach--van~Suijlekom).
Constructing a Lorentzian propinquity is original NCG-mathematics
work at long-range scale.

\textbf{Distinct from existing $\gamma$-extensions in the
catalogue.}  The Dirac-sector $\gamma = \sqrt{1 - (Z\alpha)^2}$ in
\S~\ref{sec:proj_spinor} is a \emph{bound-state radial Lorentz
factor} for a single-particle hydrogenic wavefunction in a central
potential; it is not a frame transformation between two observers.
Similarly, the rest-mass projection
(\S~\ref{sec:proj_restmass}) rescales a single particle's mass under
$m_e \to m_\text{red}$, and the Breit retardation projection
(\S~\ref{sec:proj_breit_retardation}) introduces the rest-mass
ratio $m_l/m_n$ as a two-body kinematic control parameter.  None of
these is a frame-level boost.  A Lorentz-boost projection would act
on \emph{all} particles in the system simultaneously and would
relate two $T_{S^3} \otimes T_{S^1_\beta}$ constructions at
different observer rapidities.

\textbf{Non-Lorentzian uniform-rescaling alternative.}  A
non-Lorentzian ``rapidity-like'' projection that uniformly rescales
all spatial focal lengths $\lambda_x$ by $1/\gamma$ and the
temporal compactification scale $\beta$ by $\gamma$ is admissible
in the current Riemannian framework: the rescaling preserves the
master Mellin engine ring assignments
(M1~Hopf-base measure, M2~Seeley--DeWitt heat-kernel coefficients,
M3~vertex-restricted Hurwitz Dirichlet content) and the Roothaan
multipole termination $L_\text{max} = 2\,l_\text{max}$ which is
pure-angular.  However, this uniform-rescaling alternative does
not surface a new structural feature beyond what the existing
\S~\ref{sec:proj_restmass} and \S~\ref{sec:proj_breit_retardation}
variable-refinement projections already encode.  In the V/D/P
tagging it would be V---, with no new dimension or transcendental
witness, structurally similar to the existing variable-refinement
projections.

\textbf{Recommendation: not adding now.}  The framework cannot host
a clean Lorentz-boost projection at the metric / curvature level
without a Minkowski-signature spectral triple (a substantial NCG
program, dominated by original Lorentzian-propinquity
construction).  The non-Lorentzian uniform-rescaling alternative is
admissible but reducible to existing variable-refinement
projections.  Sprint~TD Track~4's $T_H = 1/(8\pi M)$ result already
captures the Lorentz-boost / Bisognano--Wichmann content of the
$M_1$ mechanism, so the framework loses nothing strategically
important by deferring the full Lorentzian extension.  This open
question is flagged for future PI decision: \emph{if a Minkowski
spectral triple is built (a major structural extension), the
literal operator-system-level identification of the Wick-rotation
projection (\S~\ref{sec:proj_wick_rotation}) becomes possible at the
metric level; pre-Minkowski, the metric-functional-level
Bisognano--Wichmann reading of $\beta$-circle physics is what
\S~\ref{sec:proj_wick_rotation} already encodes, and the open
extension here is the literal-identification verifier (compute the
modular Hamiltonian $K$ on $\mathcal{T}_{n_\text{max}}$, verify
$\sigma_{i\cdot 2\pi} = $ identity at finite cutoff, take the
GH limit using Paper~38's lemmas).}  No near-term track; the
operator-system-level falsifier is named in
\S~\ref{sec:proj_wick_rotation}'s Honest scope.  Equivalently, this
extension is the framework-internal realization of Tomita--Takesaki
modular flow on the truncated wedge / static-exterior algebra of
$\mathcal{T}_{n_{\max}}$, with the modular Hamiltonian $K$ and modular
conjugation $J$ as the principal operator-system invariants;
Sprint~3 Track~10's absorption verdict
(2026-05-15) classifies Tomita--Takesaki as the algebraic vocabulary
of this extension rather than as an independent projection candidate.

\item \textbf{Candidate state-side dictionary enumeration.}
\S~\ref{sec:proj_apparatus_identity} (apparatus identity / state-side
reduction) is the first state-side entry in the dictionary; all
prior twenty-seven projections operate on the spectral data
$(\mathcal{A}, \mathcal{H}, D)$ while this one operates on density
matrices $\rho$ on $\mathcal{H}$.  Sprint TD Track~5's PSLQ negative
established that state-side correlation entropies are categorically
distinct from the spectral-side master Mellin engine, opening a
structurally distinct dictionary class.  Natural candidate
state-side projections for future enumeration:
\emph{mutual information} $I(A{:}B) = S(\rho_A) + S(\rho_B) - S(\rho_{AB})$,
\emph{conditional entropy} $S(A|B) = S(\rho_{AB}) - S(\rho_B)$,
\emph{relative entropy} $S(\rho \| \sigma) = \mathrm{Tr}\rho(\log\rho
- \log\sigma)$,
\emph{trace distance} $\tfrac{1}{2}\|\rho - \sigma\|_1$,
\emph{fidelity} $F(\rho, \sigma) = \mathrm{Tr}\sqrt{\sqrt{\rho}\sigma\sqrt{\rho}}$,
\emph{Wasserstein--Kantorovich distance} on the state space
$P(\mathcal{H})$ identified by Paper~38 as the natural metric
container for the propinquity limit.  These candidates have appeared
in the framework only as scaffolding inside specific calculations
(Paper~26 bipartite entanglement; Paper~38 Wasserstein--Kantorovich
identification of the propinquity limit) and have not been promoted
to dictionary entries.  Whether the state-side admits its own
internal sub-mechanism partition analogous to the spectral-side M1 /
M2 / M3 split (and if so, whether the four-mechanism structure ---
M1, M2, M3 spectral, plus a state-side fourth --- is the right
extension of the case-exhaustion theorem of Paper~32 \S~VIII) is the
deeper structural question.  No near-term track; this open question
is the natural Sprint~3 (or future-sprint) target if the
dictionary-completion arc continues.

\textbf{Falsifier.}  This open-question entry can be retired in
either direction.  (a)~If a Lorentzian propinquity is published
(monitor Latr\'emoli\`ere / Toyota / Hekkelman--McDonald arXiv
listings for ``Lorentzian / Krein-space propinquity''), or the
twisted-spectral-triple emergence approach
(Nieuviarts~2024--2025~\cite{nieuviarts2025a,nieuviarts2025b})
stabilizes as a standard prescription, the dominant cost shrinks
dramatically and the structural extension becomes tractable.
(b)~If a Sprint~TD-class GeoVac construction reproduces a
quantitative Lorentz-boost observable that is NOT already an $M_1$
$2\pi$-circle output (e.g., a Wigner-rotation matrix element on the
Camporesi--Higuchi spinor sector, or a length-contraction
matrix element relating two truncations at different $\beta$), the
framework already does more boost-class physics than this entry
recognizes and the entry needs sharpening.  Neither falsifier is
currently engaged.

\textbf{Track~D documentation closure (May~2026).}\ The Track~B
scoping pass~(\texttt{debug/lorentz\_boost\_scoping\_memo.md},
May~2026) recommended a sprint-scale documentation deliverable
(Candidate~A) on the Bisognano--Wichmann reading of Sprint~TD
Track~4.  Track~D~(\texttt{debug/bisognano\_wichmann\_track\_d\_memo.md},
May~2026) delivered Candidate~A:\ five-section structural analysis
documenting that the framework's M1 mechanism on the cigar's
$S^1_\tau$ factor reproduces $\beta = 8\pi M$, and that under the
verified Wick-rotation chain
(Hartle--Hawking~1976~\cite{hartle_hawking1976};
 Sewell~1982~\cite{sewell1982};
 Bisognano--Wichmann~1976~\cite{bisognano_wichmann1976}) this $2\pi$
IS the modular-flow / Lorentz-boost orbit period of a stationary
observer outside the Schwarzschild horizon.  Verdict:\ \emph{structural
correspondence}, not literal identification (the operator-level
statement $\sigma_{i\cdot 2\pi} = $ identity on the truncated metric
spectral triple $\mathcal{T}_{n_{\max}}$ is the principal Track~D
falsifier, follow-up at MEDIUM-HIGH priority).  Paper~32
\S~VIII gained a Bisognano--Wichmann-reading remark
(\texttt{rem:bisognano\_wichmann\_reading}); Paper~35 \S~VIII gained
a Bisognano--Wichmann subsection
(\texttt{subsec:bisognano\_wichmann\_reading});
\S~\ref{sec:proj_observation} (temporal-window projection) gained a
footnote naming the Bisognano--Wichmann reading; the
Hawking-temperature off-precision row in
\S~\ref{sec:offprecision} catalogues the Wick-rotated thermal
mechanism as a machinery-witness (error class C, $M$ external
input).  Two literature corrections inherited from Track~D's
verification pass:\ Sewell's publication year is 1982 (not 1980
as the Track~B scoping data had), and the Bisognano--Wichmann 1976
paper title is ``On the duality condition for quantum fields''
(``On the duality condition for a Hermitian scalar field'' is the
1975 paper title).  Track~D \emph{does not} change the
REQUIRES-EXTENSION verdict on the operator-system-level Lorentzian
extension itself;\ it closes the documentation loop on the
Bisognano--Wichmann reading without lifting projection
admissibility.

The Unruh pendant~(\texttt{debug/unruh\_pendant\_memo.md}, May~2026,
pendant follow-up to Track~D) added the flat-space face
$T_U = a/(2\pi)$ from the same M1 mechanism with $\beta_U = 2\pi/a$,
completing the Hawking--Sewell--BW--Unruh four-witness chain at the
framework's M1 sub-mechanism level (cf.~Paper~32~§VIII Unruh
extension to \texttt{rem:bisognano\_wichmann\_reading} and
Paper~35~\S~VIII \texttt{subsec:unruh\_pendant}).  The
REQUIRES-EXTENSION verdict on a literal Lorentz-boost projection is
unchanged, since the same operator-level falsifier
($\sigma_{i\cdot 2\pi} = $ identity on $\mathcal{T}_{n_{\max}}$ for
a wedge-restricted state) tests all four Wick-rotation faces
simultaneously --- the bridge runs through the same
KMS-Tomita-Takesaki algebra in every case.

\item \textbf{Born rule as calibration data --- CLOSED structural
position (added 2026-05-26).}  The question ``does GeoVac derive
the Born rule from substrate primitives?'' was attacked in sprint
\texttt{ahha\_born\_rule\_attempt} (2026-05-26, generative pass +
audit + literature check + explicit construction) and \emph{closed
in the negative as a structural position}, not left as an open
question.  The §III.28 apparatus identity / state-side reduction
projection produces Born-rule probabilities via the standard GNS
construction; the squared-modulus structure is inherited from the
complex inner product on $\mathcal{H}_{\mathrm{GV}}$ (taken as input
in the spectral-triple framework of Paper~32~\S~III), and the Born
rule follows from Hilbert-space + GNS via Gleason's theorem (1957).
No GeoVac-specific operation reduces Born rule to more primitive
structures.  Three attempted derivation routes (§III.28 + GNS;
Pythagorean orthogonality + apparatus preferred basis; master Mellin
engine $M_1$ modular trace) all return the same Born expression by
the same standard argument, with the Pythagorean route adding
decoherence-flavored content (preferred basis = centralizer of
$H_{\mathrm{local}}$ in $\mathcal{A}_W$) but not modifying the Born
weighting.  The framework's structural position is therefore:\ the
Born rule sits at the calibration-data tier (Paper~18~\S~IV.6
inner-factor input data; CLAUDE.md~\S~1.7 WH8; the
external-input three-class partition documented in memory file
\texttt{external\_input\_three\_class\_partition.md}), alongside the
Yukawa couplings, $\alpha$, the gauge-group lower-bound calibration,
and the multi-loop QED counterterms.  Probabilistic weighting of
measurement is external input to the structural skeleton, not
derivable from the packing axiom + Fock projection + spectral-triple
data.  This places GeoVac in the Wigner / calibration-tier camp on
the Born rule rather than the Deutsch--Wallace / Zurek /
Bub--Pitowsky derivation camps; in particular, GeoVac uses Gleason
1957 via standard Hilbert inheritance, does not improve on Gleason,
and does not claim to.  What the framework \emph{does} contribute to
the measurement-problem literature is the operator-algebraic
realization with Pythagorean orthogonality
$\langle H_{\mathrm{local}}, D_W^L\rangle_{\mathrm{HS}} = 0$
(Paper~43~\S~10.2, formal theorem with $1/\pi^2$ master Mellin
engine $M_1$ signature):\ a concrete realization of relational-QM-style
measurement structure where the modular (observer-defined) generator
and the spectral-action (substrate) generator decouple at the
Hilbert--Schmidt inner-product level.  This is novel as a theorem;
it is not novel as a Born rule derivation.  The honest contribution
of the framework on this question is therefore a positional
clarification:\ Born rule is calibration data in the same sense that
Yukawas are calibration data, structurally consistent with the
framework's broader scope statement.  The corresponding §III.28
catalogue paragraph documents this in the dictionary;\ this open-
questions entry records the closure.  No further derivation work on
Born rule is planned at the substrate level.

\emph{Quantitative converse tested (2026-07-01, WH8 Step-1 probe).}
The closure above is derivation-direction (no route from substrate
primitives to the Born expression).  The generation direction ---
could the skeleton's own native measure \emph{substitute} for Born?
--- is now also tested negative and pinned.  The graph's only native
measures are counting-type (rational, built from discrete
invariants); the probe shows (i)~the counting measure coincides with
Born exactly on, and only on, the \emph{flat locus} (states of
uniform amplitude over their support; for amplitudes
$\propto(1,2,3,4)$ on a 4-fold multiplet the deviation is the exact
rational $1/3$), (ii)~it diverges at total-variation distance $0.86$
from the He graph-native CI ground state ($n_{\max}=3$,
zero-parameter, exact Slater integrals; $90.8\%$ of the Born weight
sits on the $1s^2$ configuration, against $1/N_{\rm supp}$ uniform),
and (iii)~the flat locus is dynamically unstable under the graph's
own Hamiltonian (TV grows as $t^2$; measured small-$t$ exponent
$2.00$), so no state-independent skeleton rule can track Born
through dynamics.  Frozen falsifier
\texttt{tests/test\_wh8\_born\_probe.py}:\ a skeleton-side rule
reproducing Born statistics on non-flat physical states would break
these pins.  Registered internally as WH8 (CLAUDE.md~\S~1.7,
PI direction):\ the Born measure is the \emph{exchange constant} of
the observation projection (\S~III.28) --- the projection's
irreducible imported content, unique given the projection lattice
(Gleason) --- exactly as $\pi$ is the import of the Coulomb
projection and $2\pi/\beta$ of the temporal one.

\end{enumerate}

% ===================================================================
\section{Conclusion}
\label{sec:conclusion}
% ===================================================================

GeoVac is a two-layer framework.  Layer~1 is a bare combinatorial
graph: dimensionless integers, rational matrix elements, no
physics, no $\pi$.  Layer~2 is a finite family of named projections;
each adds specific physical variables, specific physical
dimensions, and specific transcendental content.  Every physical
quantity in any GeoVac calculation enters through exactly one
projection in Layer~2 and is fully traceable to it.

The dictionary in this paper makes that statement operational.
Twenty-eight projections are listed, each tagged by variables, dimension,
and transcendental signature.  An empirical matches catalogue uses
these tags to organise the framework's known machine-precision
agreements with standard physics values.  Sprint TX-A's
dimensional axiomatization (\S~\ref{sec:axiomatization}) verified
three structural theorems on this dictionary: dimensional consistency
holds across all worked examples, weak minimality holds at the
signature level (with the strong form failing on two non-commuting
pairs --- the Paper~2 \texttt{hopf}--\texttt{spinor\_lift} pair under
bundle-non-trivial composition, and the
\texttt{rest\_mass}--\texttt{breit\_retardation} pair under
variable-refinement dependence).  The 16$\times$16 composition table
expanded from the previous 15$\times$15 reflects the new sixteenth
projection (two-body Dirac / Breit retardation, added 2026-05-08);
under this expansion Drake--Swainson is no longer the unique
projection with input\_layer = output\_layer = 2, with Breit
retardation joining as the second.  Both projections continue to
account for all asymmetric (one-way) cells in the 16$\times$16
table.  The 2026-05-09 addition of three intra-nuclear projections
(charge-density / Foldy--Friar, magnetization-density / Zemach,
tensor multipole) extends the dictionary further to nineteen, and
sharpens the $L$--$E$ asymmetry of \S~\ref{sec:base_units}: $[L]$
admits six scalar instances plus one rank-2 tensor instance, while
$[E]$ admits four.  Sprint~1 (2026-05-15) added six more
projections that the framework had been using implicitly
(\S\S~\ref{sec:proj_phillips_kleinman}--\ref{sec:proj_coupled_channel}:
Phillips--Kleinman / core--valence orthogonality, multipole expansion
/ Gaunt termination, bipolar harmonic / Drake combining, symmetry /
Young tableau, adiabatic / Born--Oppenheimer, coupled-channel /
adiabatic curve), closing the textbook-completeness gap.  Sprint~2
(2026-05-15) added three more
(\S\S~\ref{sec:proj_gauge_choice}--\ref{sec:proj_apparatus_identity}:
gauge choice / Coulomb-Lorenz-Feynman; Wick rotation / signature
change; apparatus identity / state-side reduction), closing the
axis-grid completeness gap: gauge choice provides a second
independence witness in the transcendental-only corner of the
V/D/transcendental grid, Wick rotation promotes the
Bisognano--Wichmann reading of the M1 mechanism from a \S~VIII
open-question candidate to a named §III slot, and apparatus identity
opens the state-side complement of the dictionary (which Sprint TD
Track~5's PSLQ negative proved is categorically distinct from the
master Mellin engine).  These nine additions bring the dictionary
count from nineteen to twenty-eight.  Sprint~3 (2026-05-15) tested
three further candidates --- Tomita--Takesaki modular flow, loop
expansion / $\alpha$-power-sorting, heat-kernel regularization / Schwinger
proper-time --- and returned ABSORBED verdicts on all three:\ TT is
absorbed into \S\S~\ref{sec:proj_observation}+\ref{sec:proj_wick_rotation};
loop expansion is absorbed into \S\S~\ref{sec:proj_spectral_action}+%
\ref{sec:proj_vector_photon}+\ref{sec:proj_breit_retardation}+%
\ref{sec:proj_gauge_choice}; heat-kernel regularization is absorbed
into \S~\ref{sec:proj_spectral_action} where the master Mellin engine
framing is now made explicit in its Structural reading paragraph.  The
Sprint~3 scoping verdicts confirm structural completeness at twenty-eight:\
the dictionary does not require additional §III slots to cover the master
Mellin engine, the loop-order sorting of QED, or the Tomita--Takesaki
algebraic vocabulary of the Wick-rotation chain.  The new projections
add no new base units and preserve the V/D correlation property; T3
expansion to 28$\times$28 is flagged for the next TX-A audit.  A boundary subsection
(\S~\ref{sec:layer2_boundary}) names the structural distinction
between intrinsic graph-derived variables (which the framework
projects natively) and Layer-2 inputs (external classical fields,
multi-loop QED coefficients, lattice-QCD-derived radii); the same
day's external-field hunting verified that uniform Zeeman / Stark
fields are Layer-2, not new projections.  T3 expansion to
19$\times$19 is flagged for the next TX-A audit.  The
field-theoretic UV/IR ordering anomaly between \texttt{spectral\_action}
and \texttt{observation\_window} is real physics but is \emph{not}
encoded by the current coarse layer-grading (Track TS-C, 2026-05-04);
refining the output\_layer index to distinguish UV-cutoff, thermal,
and flow Layer-2 outputs is the natural next step.  Of the three axes, only the
transcendental axis is genuinely independent --- variable and
dimension are perfectly correlated.  A falsifiable structural
prediction (\S~\ref{sec:prediction}) emerges from the depth-grouped
catalogue: the corrected reading is the Layer-2-presence bound
articulated in \S\ref{sec:prediction} (the earlier ``error compounds
with projection depth'' conjecture was falsified by the
\S\ref{sec:autopsies} autopsy panel and replaced).  The corrected
prediction is consistent with current data, including the
208/208-confirmation panel of Paper~35 Prediction~1
(Sprint TX-B, May~2026), and refinable by future calculations.

This paper introduces no new computational result.  Its contribution
is a vocabulary that the framework has been needing.  With the
two-layer view in place, the question ``where does $\pi$ come from
in this observable?'' has a definite answer for every observable
GeoVac produces.  The same is true for $\alpha$, $Z$, $\hbar\omega$,
$\Lambda$, $\zeta(3)$, Catalan's $G$, and any other quantity that
enters from outside the bare graph.

% ===================================================================
\begin{thebibliography}{99}
% ===================================================================

\bibitem{paper0}
GeoVac Project, ``Angular Momentum as Information Geometry: Isotropic Packing and the Universal Angular Cross-Section,'' Paper~0 (2026).

\bibitem{paper2}
GeoVac Project, ``The Fine Structure Constant as a Spectral Coincidence on the Hopf Fibration,'' Paper~2 (2025; observation status, May~2026).

\bibitem{paper7}
GeoVac Project, ``The Dimensionless Vacuum: Recovering the Schr\"odinger Equation from Scale-Invariant Graph Topology,'' Paper~7 (2025).

\bibitem{paper14}
GeoVac Project, ``Structurally Sparse Qubit Hamiltonians from Spectral Graph Theory,'' Paper~14 (2025).

\bibitem{paper15}
GeoVac Project, ``The Level~4 Natural Geometry: Two-Center Two-Electron
Molecules in Molecule-Frame Hyperspherical Coordinates,'' Paper~15 (2025).

\bibitem{paper17}
GeoVac Project, ``Composed Natural Geometries for Core-Valence Diatomics: LiH and BeH$^+$ from Discrete Graph Topology,'' Paper~17 (2025).

\bibitem{paper18}
GeoVac Project, ``Spectral-Geometric Exchange Constants: Projecting Discrete Spectra onto Continuous Manifolds,'' Paper~18
(2026).

\bibitem{paper19}
GeoVac Project, ``Toward PK-Free Molecular Hamiltonians: Two-Center
Sturmian Integrals and Coupled Composition,'' Paper~19 (2026).

\bibitem{paper20}
GeoVac Project, ``Structurally Sparse Molecular Hamiltonians for Near-Term Quantum Simulation,''
Paper~20 (2026).

\bibitem{paper22}
GeoVac Project, ``Potential-Independent Angular Sparsity for Qubit Hamiltonians in Spherical Fermion Systems,''
Paper~22 (2026).

\bibitem{paper23}
GeoVac Project, ``Nuclear Shell Model Hamiltonians on the Hyperspherical Lattice,''
Paper~23 (2026).

\bibitem{paper24}
GeoVac Project, ``The Bargmann--Segal Lattice: $\pi$-Free Discretization of the Harmonic Oscillator on $S^5$,'' Paper~24
(2026).

\bibitem{paper25}
GeoVac Project, ``The Hopf Graph as a Lattice Gauge Structure,''
Paper~25 (2026).

\bibitem{paper27}
GeoVac Project, ``Entropy as a Projection Artifact: One-Body Operators are Entanglement-Inert on Sparse Lattices,'' Paper~27 (2026).

\bibitem{paper28}
GeoVac Project, ``QED on $S^3$: Spectral Geometry of Perturbative
Transcendentals,'' Paper~28 (2026).

\bibitem{paper30}
GeoVac Project, ``SU(2) Wilson Lattice Gauge Structure on the GeoVac
Hopf Graph,'' Paper~30 (2026).

\bibitem{paper31}
GeoVac Project, ``The Universal/Coulomb Partition of the GeoVac Spectral Triple,'' Paper~31 (2026).

\bibitem{paper32}
GeoVac Project, ``The GeoVac Spectral Triple: A Synthesis Paper Locking the Framework into the Marcolli--van~Suijlekom Gauge-Network Lineage,''
Paper~32 (2026).

\bibitem{paper33}
GeoVac Project, ``The 1+6+1 Partition of QED Selection Rules
on $S^3$,'' Paper~33 (2026).

\bibitem{paper35}
GeoVac Project, ``Time as Projection: The Klein-Gordon Spectrum on
$S^3 \times \mathbb{R}$ and the Algebraic / Observation Split,''
Paper~35 (2026).

\bibitem{ls1_memo}
GeoVac Project, ``LS-1 Sprint: Hydrogen Lamb Shift on the Dirac-S$^3$
Framework,'' \texttt{debug/ls1\_lamb\_shift\_memo.md} (May 2026).

\bibitem{ls2_memo}
GeoVac Project, ``LS-2 Sprint: Bethe Logarithm from GeoVac Spectral
Machinery,'' \texttt{debug/ls2\_bethe\_log\_memo.md} (May 2026).

\bibitem{ls3_memo}
GeoVac Project, ``LS-3 Sprint: Regularized Bethe Logarithm via
Acceleration Form,'' \texttt{debug/ls3\_bethe\_log\_regularized\_memo.md}
(May 2026).

\bibitem{ls4_memo}
GeoVac Project, ``LS-4 Sprint: Bethe Logarithm via Schwartz/Drake--Swainson
Regularization,'' \texttt{debug/ls4\_bethe\_log\_drake\_memo.md} (May 2026).

\bibitem{vp1_memo}
GeoVac Project, ``VP-1 Sprint: Vector-Photon Promotion on the
Graph-Native QED Pipeline,'' \texttt{debug/vp1\_vector\_graph\_native\_memo.md}
(May 2026).

\bibitem{vp2_memo}
GeoVac Project, ``VP-2 Sprint: Topology-Specific Projection Theory,''
\texttt{debug/vp2\_topology\_projections\_memo.md} (May 2026).

\bibitem{audit_memo}
GeoVac Project, ``Curve-Fit Audit (2026-05-02),''
\texttt{docs/curve\_fit\_audit\_memo.md} (May 2026).

\bibitem{fock1935}
V.~Fock, ``Zur Theorie des Wasserstoffatoms,''
\textit{Z.\ Phys.}\ \textbf{98}, 145 (1935).

\bibitem{bargmann1961}
V.~Bargmann, ``On a Hilbert space of analytic functions and an
associated integral transform, Part~I,'' \textit{Comm.\ Pure Appl.\
Math.}\ \textbf{14}, 187 (1961).

\bibitem{chamseddine_connes2010}
A.~H.~Chamseddine and A.~Connes, ``Noncommutative geometry as a
framework for unification of all fundamental interactions
including gravity,'' \textit{Fortsch.\ Phys.}\ \textbf{58}, 553--600
(2010).

\bibitem{marcolli_vs2014}
M.~Marcolli and W.~D.~van~Suijlekom, ``Gauge networks in
noncommutative geometry,'' \textit{J.\ Geom.\ Phys.}\ \textbf{75},
71 (2014); arXiv:1301.3480.

\bibitem{drake_swainson1990}
G.~W.~F.~Drake and R.~A.~Swainson, ``Bethe logarithms for hydrogenic
atoms,'' \textit{Phys.\ Rev.\ A}\ \textbf{41}, 1243 (1990).

\bibitem{camporesi_higuchi1996}
R.~Camporesi and A.~Higuchi, ``On the eigenfunctions of the Dirac
operator on spheres and real hyperbolic spaces,''
\textit{J.\ Geom.\ Phys.}\ \textbf{20}, 1 (1996).

% --- Lorentzian-NCG references (Sprint Lorentz-Boost-Scoping, May 2026; §VIII open question) ---

\bibitem{bisognano_wichmann1975}
J.~J.~Bisognano and E.~H.~Wichmann, ``On the duality condition for a
Hermitian scalar field,'' \textit{J.\ Math.\ Phys.}\ \textbf{16},
985--1007 (1975).

\bibitem{bisognano_wichmann1976}
J.~J.~Bisognano and E.~H.~Wichmann, ``On the duality condition for
quantum fields,'' \textit{J.\ Math.\ Phys.}\ \textbf{17},
303--321 (1976).
% Title and pagination corrected by Track D verification pass (May 2026):
% earlier version of this entry mistakenly used the 1975 paper's title
% (``On the duality condition for a Hermitian scalar field'') for the
% 1976 paper.  See debug/bisognano_wichmann_track_d_memo.md.

\bibitem{sewell1982}
G.~L.~Sewell, ``Quantum fields on manifolds: PCT and gravitationally
induced thermal states,'' \textit{Annals of Physics}\ \textbf{141},
201--224 (1982).
% Year corrected by Track D verification pass (May 2026): earlier
% version of this entry had 1980 (carried over from Track B scoping
% memo); the actual publication year is 1982.

\bibitem{sewell1980}
% Backward-compatibility alias for the same Sewell 1982 reference.
% Kept so that pre-Track-D citations \cite{sewell1980} still resolve.
G.~L.~Sewell, ``Quantum fields on manifolds: PCT and gravitationally
induced thermal states,'' \textit{Annals of Physics}\ \textbf{141},
201--224 (1982).

\bibitem{hartle_hawking1976}
J.~B.~Hartle and S.~W.~Hawking, ``Path-integral derivation of
black-hole radiance,'' \textit{Phys.\ Rev.\ D}\ \textbf{13},
2188--2203 (1976).

\bibitem{unruh1976}
W.~G.~Unruh, ``Notes on black-hole evaporation,''
\textit{Phys.\ Rev.\ D}\ \textbf{14}, 870--892 (1976).
% Added 2026-05-10 by the Unruh-pendant sprint
% (debug/unruh_pendant_memo.md), which extends Track D's
% Bisognano-Wichmann reading with the flat-space face $T_U = a/(2\pi)$
% via the same M1 mechanism on the Wick-rotated Rindler $\eta_E$-circle.

\bibitem{strohmaier2006}
A.~Strohmaier, ``On noncommutative and pseudo-Riemannian
geometry,'' \textit{J.\ Geom.\ Phys.}\ \textbf{56}, 175--195
(2006).  arXiv:math-ph/0110001.

\bibitem{bizi_brouder_besnard2018}
N.~Bizi, C.~Brouder, and F.~Besnard, ``Space and time dimensions of
algebras with applications to Lorentzian noncommutative geometry
and quantum electrodynamics,'' \textit{J.\ Math.\ Phys.}\ \textbf{59},
062303 (2018).  arXiv:1611.07062.

\bibitem{franco_eckstein2014}
N.~Franco, ``Temporal Lorentzian spectral triples,''
\textit{Reviews in Mathematical Physics}\ \textbf{26}, 1430007
(2014).  arXiv:1210.6575.

\bibitem{van_den_dungen2016}
K.~van~den~Dungen, ``Krein spectral triples and the fermionic
action,'' \textit{Math.\ Phys.\ Anal.\ Geom.}\ \textbf{19}, 4 (2016).

\bibitem{devastato_lizzi_martinetti2018}
A.~Devastato, S.~Farnsworth, F.~Lizzi, and P.~Martinetti, ``Lorentz
signature and twisted spectral triples,'' \textit{JHEP}\ \textbf{03}, 089 (2018).
arXiv:1710.04965.

\bibitem{nieuviarts2025a}
G.~Nieuviarts, ``Emergence of Lorentz symmetry from an
almost-commutative twisted spectral triple,'' arXiv:2502.18105
(2025).

\bibitem{nieuviarts2025b}
G.~Nieuviarts, ``Emergence of Time from a Twisted Spectral Triple
in Almost-Commutative Geometry,'' arXiv:2512.15450 (2025).

\bibitem{tener2025_bw_nonunitary}
J.~E.~Tener,
``The Bisognano-Wichmann property for non-unitary Wightman conformal
field theories,''
arXiv:2506.10625 (2025).

\bibitem{cavalletti_mondino2024}
F.~Cavalletti and A.~Mondino,
``Optimal transport in Lorentzian synthetic spaces, synthetic timelike
Ricci curvature lower bounds and applications,''
Cambridge J.\ Math.\ \textbf{12}, 417--534 (2024); arXiv:2004.08934.

\bibitem{paper38}
J.~Loutey, ``State-space Gromov--Hausdorff convergence of truncated Camporesi--Higuchi spectral triples on $S^3$,'' GeoVac project, May 2026.

\bibitem{paper42}
J.~Loutey, ``Tomita--Takesaki Modular Structure on Truncated SU(2)
Spectral Triples: Four-Witness Wick-Rotation Literal Identification
at Finite Cutoff,'' GeoVac project, May 2026.

\end{thebibliography}

\end{document}
